\documentclass[10pt,oneside]{book}

\ifdefined\pdfinfoomitdate\pdfinfoomitdate=1\fi
\ifdefined\pdftrailerid\pdftrailerid{}\fi
\ifdefined\pdfsuppressptexinfo\pdfsuppressptexinfo=15\fi

\usepackage[T1]{fontenc}
\usepackage[utf8]{inputenc}
\usepackage[spanish,es-nodecimaldot]{babel}
\usepackage{amsmath,amssymb,mathrsfs}
\usepackage{aurical}
\usepackage[all]{xy}
\usepackage[a4paper,margin=22mm]{geometry}
\usepackage[colorlinks=true,linkcolor=blue,hypertexnames=false,unicode=true,pdfencoding=auto]{hyperref}
\hypersetup{
  pdftitle={Elementos de geometría algebraica III, segunda parte — edición española},
  pdfauthor={Alexander Grothendieck y Jean Dieudonné},
  pdfsubject={Estudio cohomológico de los haces coherentes: secciones 6 y 7}
}

\renewcommand{\thesection}{\arabic{section}}
\renewcommand{\thesubsection}{\thesection.\arabic{subsection}}
\newcommand{\sh}[1]{\mathscr{#1}}
\newcommand{\shTor}{\sh{T}\textup{\kern-2.2pt{\Fontauri\slshape or}}}
\newcommand{\RR}{\mathrm R}
\newcommand{\HH}{\mathrm H}
\newcommand{\bb}[1]{\mathbf{#1}}
\newcommand{\cat}[1]{\mathtt{#1}}
\newcommand{\kres}{\kappa}
\newcommand{\vphi}{\phi}
\newcommand{\supertilde}{{\,\widetilde{\,}\,}}
\newcommand{\Spec}{\operatorname{Spec}}
\newcommand{\Proj}{\operatorname{Proj}}
\newcommand{\shProj}{\mathscr{P}\!roj}
\newcommand{\Hom}{\operatorname{Hom}}
\newcommand{\Tor}{\operatorname{Tor}}
\newcommand{\Ker}{\operatorname{Ker}}
\newcommand{\Coker}{\operatorname{Coker}}
\newcommand{\isoto}{\xrightarrow{\sim}}
\newcommand{\simto}{\xrightarrow{\sim}}
\newcommand{\shHom}{\mathscr{H}\!om}
\newcommand{\sref}[2][]{(#2)}
\newcommand{\oldpage}[2][]{\setcounter{footnote}{0}\marginpar{\scriptsize\textbf{#1}~|~#2}\ignorespaces}
\newenvironment{env}[1][]{\par\medskip\noindent\textbf{(#1)}\ }{\par}
\newenvironment{theorem}[1][]{\par\medskip\noindent\itshape Teorema (#1). ---\ }{\par}
\newenvironment{proposition}[1][]{\par\medskip\noindent\itshape Proposición (#1). ---\ }{\par}
\newenvironment{lemma}[1][]{\par\medskip\noindent\itshape Lema (#1). ---\ }{\par}
\newenvironment{corollary}[1][]{\par\medskip\noindent\itshape Corolario (#1). ---\ }{\par}
\newenvironment{remark}[1][]{\par\medskip\noindent\itshape Observación (#1). ---\ }{\par}
\newenvironment{remarks}[1][]{\par\medskip\noindent\itshape Observaciones (#1). ---\ }{\par}
\newenvironment{scholium}[1][]{\par\medskip\noindent\textit{Escolio (#1). ---\ }}{\par}
\newenvironment{definition}[1][]{\par\medskip\noindent\itshape Definición (#1). ---\ }{\par}
\newenvironment{proof}{\par\medskip\noindent\textit{Demostración. ---\ }}{\par}
\newenvironment{namedtheorem}[2][]{\par\medskip\noindent\itshape #2 (#1). ---\ }{\par}



\setlength{\emergencystretch}{6em}
\begin{document}
% BEGIN RIGHTS AND AI TRANSLATION NOTICE
\begingroup
\small
\noindent\textbf{Derechos, disponibilidad y autoría}

Los derechos sobre las obras originales y sobre las contribuciones protegidas a estas traducciones y transcripciones derivadas corresponden a sus respectivos titulares. En la medida en que el editor tenga derechos sobre la traducción, la transcripción, las correcciones y la composición producidas por IA, los dedica al dominio público bajo CC0 1.0 Universal (https://creativecommons.org/publicdomain/zero/1.0/deed.es); esto no afecta a los derechos sobre las obras originales ni concede permiso para usarlas. Las menciones de derechos de terceros y las autorizaciones ya concedidas siguen siendo aplicables.

Estas ediciones se ofrecen gratuitamente, con fines no comerciales y exclusivamente educativos, bajo la presunción de que los respectivos titulares de derechos no se opondrían a que estas traducciones y transcripciones estuvieran disponibles gratuitamente junto a los originales franceses en NUMDAM (www.numdam.org/). Esta presunción no significa que se haya obtenido autorización; el acceso gratuito a los originales no autoriza por sí solo su traducción, adaptación o republicación.

Si usted es titular de derechos sobre algún material presentado aquí y se opone a su disponibilidad, puede abrir una incidencia en GitHub (github.com/KokunoYumeto/ega-fga-es/issues). El responsable de la publicación retirará el material afectado a petición del titular de los derechos.

La traducción al español está producida por IA. Las correcciones, la edición y la composición también están producidas por IA. El responsable de la publicación no tradujo el texto. Se conservan la autoría y las atribuciones de las obras originales y de las contribuciones heredadas; no se atribuyen a la IA. Los modelos documentados se identifican en cada edición.
\par
\endgroup
\clearpage
% END RIGHTS AND AI TRANSLATION NOTICE

\begin{titlepage}
  \centering
  {\Large A. Grothendieck y J. Dieudonné\par}
  \vfill
  {\Huge Elementos de geometría algebraica III\par}
  \vspace{1.5ex}
  {\LARGE Segunda parte\par}
  \vfill
  {\large Traducción española\par}
  \vspace{1ex}
  {\small Texto original: A. Grothendieck y J. Dieudonné, Publications Mathématiques de l’IHÉS, volumen 17 (1963)\par}
\end{titlepage}

\chapter*{CAPÍTULO III (continuación)\\[1.2ex]
\Large ESTUDIO COHOMOLÓGICO DE LOS HACES COHERENTES}

\setcounter{section}{5}
\setcounter{subsection}{0}
\setcounter{equation}{0}
\setcounter{footnote}{0}
\section{Funtores «Tor» locales y globales; fórmula de Künneth}
\label{section:III.6}

\subsection{Introducción}
\label{subsection:III.6.1}

\oldpage[III]{137}

\begin{env}[6.1.1]
\label{III.6.1.1}
Sean $f:X\to Y$ un morfismo de preesquemas y $\mathscr F$ un
$\mathscr O_X$-módulo cuasicoherente. En el estudio de las «imágenes
directas superiores» $\mathrm R^n f_*(\mathscr F)$, se llega a considerar el
siguiente problema general: dado un morfismo de «cambio de base»
$g:Y'\to Y$, se pone
\[
  X'=X_{(Y')}=X\times_Y Y',\qquad
  \mathscr F'=\mathscr F\otimes_Y\mathscr O_{Y'},\qquad
  f'=f_{(Y')}:X'\to Y',
\]
y se pretende obtener información sobre las imágenes directas superiores
$\mathrm R^n f'_*(\mathscr F')$ (suponiendo conocidas las
$\mathrm R^n f_*(\mathscr F)$). Se ve fácilmente (cf., por ejemplo, (7.7.2))
que el problema se reduce a estudiar las variaciones de
$\mathrm R^n f_*(\mathscr F\otimes_{\mathscr O_Y}\mathscr G)$ para un
$\mathscr O_Y$-módulo cuasicoherente variable $\mathscr G$, es decir, el
funtor
\[
  \mathscr G\longmapsto
  \mathrm R^n f_*(\mathscr F\otimes_{\mathscr O_Y}\mathscr G).
\]
Si $\mathscr F$ es plano sobre $Y$, el funtor
$\mathscr G\mapsto\mathscr F\otimes_{\mathscr O_Y}\mathscr G$ es exacto
(0\textsubscript{I}, 6.7.4) y, por consiguiente, el funtor compuesto
$\mathscr G\mapsto\mathrm R^n f_*(\mathscr F\otimes_{\mathscr O_Y}\mathscr G)$
sigue siendo un funtor cohomológico. Pero esto ya no sucede en el caso
general; para poder aplicar los métodos cohomológicos, se ve uno llevado a
sustituir
$\mathscr G\mapsto\mathrm R^n f_*(\mathscr F\otimes_{\mathscr O_Y}\mathscr G)$
por otros funtores, que esta vez son siempre funtores cohomológicos. Estos
funtores, que generalizan los funtores «Tor» de la teoría de módulos, se
definen en los números 6.3 a 6.7; hay, por lo demás, dos generalizaciones de
este tipo, una «local» y la otra «global», relacionadas por sucesiones
espectrales que se discutirán en el número 6.7; como aplicación de estas
sucesiones espectrales, se obtiene en particular, bajo ciertas condiciones,
una «fórmula de Künneth» que expresa
\[
  \mathrm R^n(f_1\times f_2)_*(\mathscr F_1\otimes_Y\mathscr F_2)
\]
por medio de las imágenes directas superiores
$\mathrm R^p f_{1*}(\mathscr F_1)$ y $\mathrm R^q f_{2*}(\mathscr F_2)$.
Otras sucesiones espectrales (6.8) generalizan las sucesiones espectrales de
asociatividad del funtor «Tor» de módulos; por último, el problema del cambio
de base conduce también a sucesiones espectrales (6.9).
\end{env}

\begin{env}[6.1.2]
\label{III.6.1.2}
Además, se comprueba (6.10) que los funtores cohomológicos
$\mathscr G\mapsto\mathscr T_\bullet(\mathscr G)$ así definidos son
localmente (sobre $Y$) de la forma
\[
  \mathscr G\longmapsto
  \mathscr H_\bullet(\mathscr L_\bullet\otimes_Y\mathscr G),
\]
donde $\mathscr L_\bullet$ es un complejo de $\mathscr O_Y$-módulos
localmente libres (definido salvo homotopía) y $\mathscr H_\bullet$ denota la
homología. Conviene entonces dejar de lado la situación particular que dio
origen a $\mathscr T_\bullet$ y estudiar, de

\oldpage[III]{138}

manera general, los funtores de la forma precedente
$\mathscr G\mapsto\mathscr H_\bullet(\mathscr L_\bullet\otimes_Y\mathscr G)$
(imponiendo además, cuando proceda, hipótesis de finitud adecuadas sobre los
$\mathscr L_i$ o los $\mathscr H_i(\mathscr L_\bullet)$): este es el objeto
del § 7, cuya lectura es, en lo esencial, independiente del § 6. Las
propiedades más importantes de estos funtores conciernen a las propiedades de
exactitud de una componente $\mathscr T_i$ de $\mathscr T_\bullet$; se darán
diversos criterios que permiten establecer tales propiedades; como
aplicación, se obtendrán condiciones que permiten afirmar (con las notaciones
de (6.1.1)) que el funtor
$\mathscr G\mapsto\mathrm R^n f_*(\mathscr F\otimes_{\mathscr O_Y}\mathscr G)$
es exacto (lo que se expresará diciendo que $\mathscr F$ es
cohomológicamente plano sobre $Y$ en dimensión $n$). Otra propiedad
importante de las componentes $\mathscr T_i$ de
$\mathscr T_\bullet(\mathscr G)=
\mathscr H_\bullet(\mathscr L_\bullet\otimes_Y\mathscr G)$ es una propiedad
de semicontinuidad de la función
$y\mapsto\dim_{k(y)}(T_i(k(y)))$; cuando $\mathscr T_i$ es exacto, esta
propiedad es sustituida por una propiedad de continuidad, y el recíproco es
además cierto, por Grauert, cuando $Y$ es reducido (7.8.4).
\end{env}

\begin{env}[6.1.3]
\label{III.6.1.3}
En los §§ 6 y 7 hemos utilizado sistemáticamente la hipercohomología, tomando
en todas partes como argumentos complejos de haces en lugar de haces, aunque
la necesidad de este punto de vista solo se manifieste en capítulos
posteriores. El formalismo cohomológico desarrollado con este motivo se hará,
por lo demás, más transparente en el capítulo de este Tratado que estará
dedicado al desarrollo de un álgebra de funtores cohomológicos de haces
coherentes, incluido el formalismo de la dualidad. Pero esto requerirá
desarrollos que exceden el marco del presente capítulo.
\end{env}

\begin{env}[6.1.4]
\label{III.6.1.4}
Para abreviar, dados dos complejos $K^\bullet$, $K'^\bullet$ en una categoría
abeliana $C$, diremos que un morfismo de complejos
$f:K^\bullet\to K'^\bullet$ es una equivalencia homotópica si existe un
morfismo $g:K'^\bullet\to K^\bullet$ tal que los morfismos compuestos
$f\circ g$ y $g\circ f$ sean ambos homotópicos a la identidad (por abuso de
lenguaje, cuando exista una equivalencia homotópica de este tipo, diremos
asimismo que $K^\bullet$ y $K'^\bullet$ son homotópicamente equivalentes).
Cuando se puede definir la hipercohomología de un funtor covariante aditivo
$T$ de $C$ en una categoría abeliana $C'$, respecto de un complejo de $C$
(0, 11.4.3), es inmediato que una equivalencia homotópica
$K^\bullet\to K'^\bullet$ de complejos de $C$ define canónicamente un
isomorfismo
\[
  \mathrm R^\bullet T(K^\bullet)\xrightarrow{\sim}
  \mathrm R^\bullet T(K'^\bullet)
\]
en hipercohomología (loc. cit.).
\end{env}

\subsection{Hipercohomología de los complejos de módulos sobre un preesquema}
\label{subsection:III.6.2}

\begin{env}[6.2.1]
Sean $X$ un preesquema y
$\mathscr K^\bullet=(\mathscr K^i)_{i\in\mathbf Z}$ un complejo de
$\mathscr O_X$-módulos cuya diferencial es de grado $+1$.
Recordemos que, para todo morfismo $f:X\to Y$ de preesquemas, se han definido
(0, 12.4.1) los $\mathscr O_Y$-módulos de hipercohomología
$\mathscr H^n(f,\mathscr K^\bullet)$ (denotados también por
$\mathscr H_f^n(\mathscr K^\bullet)$ o
$\mathrm R^n f_*(\mathscr K^\bullet)$) para todo $n\in\mathbf Z$;
la hipercohomología $\mathscr H^\bullet(f,\mathscr K^\bullet)$ es el objeto de
convergencia de los dos funtores espectrales
$'\mathscr E(f,\mathscr K^\bullet)$ y
$''\mathscr E(f,\mathscr K^\bullet)$, cuyos términos $\mathscr E_2$ vienen
dados por
\begin{align}
  '\mathscr E_2^{pq}
    &=\mathscr H^p\bigl(\mathscr H^q(f,\mathscr K^\bullet)\bigr),
    \tag{6.2.1.1}\label{III.6.2.1.1}\\
  ''\mathscr E_2^{pq}
    &=\mathscr H^p\bigl(f,\mathscr H^q(\mathscr K^\bullet)\bigr)
      =\mathrm R^p f_*\bigl(\mathscr H^q(\mathscr K^\bullet)\bigr).
    \tag{6.2.1.2}\label{III.6.2.1.2}
\end{align}

\oldpage[III]{139}

donde $\mathscr H^q(f,\mathscr K^\bullet)$ es el complejo cuyo componente de
grado $i$ es
$\mathscr H^q(f,\mathscr K^i)=\mathrm R^q f_*(\mathscr K^i)$ (loc. cit.).
Recordemos también que, cuando $Y$ se reduce a un punto, se denota la
hipercohomología correspondiente por $\mathbf H^\bullet(X,\mathscr K^\bullet)$
(la cual está formada por módulos sobre $\Gamma(X,\mathscr O_X)$ independientes
del preesquema puntual $Y$ considerado); cuando $Y=X$ y $f=1_X$, se tiene
$\mathscr H^n(f,\mathscr K^\bullet)=\mathscr H^n(\mathscr K^\bullet)$
(cohomología del complejo $\mathscr K^\bullet$); cuando
$\mathscr K^i=0$, salvo para $i=i_0$, se tiene
\[
  \mathscr H^n(f,\mathscr K^\bullet)
    =\mathrm R^{n-i_0}f_*(\mathscr K^{i_0}).
\]
La sucesión espectral $'\mathscr E(f,\mathscr K^\bullet)$ es siempre
regular; las dos sucesiones espectrales son birregulares cuando
$\mathscr K^\bullet$ está acotado inferiormente (0, 12.4.1).

Toda equivalencia homotópica $h:\mathscr K^\bullet\to\mathscr K'^\bullet$ de complejos de
$\mathscr O_X$-módulos (6.1.4) da lugar a un isomorfismo
$\mathscr H^\bullet(f,\mathscr K^\bullet)\xrightarrow{\sim}
\mathscr H^\bullet(f,\mathscr K'^\bullet)$ en hipercohomología. Lo mismo ocurre
si se supone solamente que
$\mathscr H^\bullet(h):\mathscr H^\bullet(\mathscr K^\bullet)\to
\mathscr H^\bullet(\mathscr K'^\bullet)$ es un isomorfismo y que
$\mathscr K^\bullet$ y $\mathscr K'^\bullet$ están acotados inferiormente,
como se deduce inmediatamente de (0, 11.1.5), aplicado a la sucesión espectral
(6.2.1.2) y a su análoga para $\mathscr K'^\bullet$. Por último, para todo
cubrimiento abierto $\mathfrak U=(U_\alpha)$ de $X$, se ha definido también
(0, 12.4.5) la hipercohomología
$\mathbf H^\bullet(\mathfrak U,\mathscr K^\bullet)$ como la cohomología del
bicomplejo $C^\bullet(\mathfrak U,\mathscr K^\bullet)$ (cuyo componente de
índices $(i,j)$ es, por definición,
$C^i(\mathfrak U,\mathscr K^j)$); los
$\mathbf H^n(\mathfrak U,\mathscr K^\bullet)$ son también módulos sobre
$\Gamma(X,\mathscr O_X)$.
\end{env}

\begin{proposition}[6.2.2]
Sean $X$ un esquema y $\mathfrak U=(U_\alpha)$ un cubrimiento de $X$ por
abiertos afines. Para todo complejo $\mathscr K^\bullet$ de
$\mathscr O_X$-módulos cuasicoherentes, los módulos de hipercohomología
$\mathbf H^\bullet(X,\mathscr K^\bullet)$ y
$\mathbf H^\bullet(\mathfrak U,\mathscr K^\bullet)$ son canónicamente
isomorfos.
\end{proposition}

\begin{proof}
En efecto, toda intersección finita $V$ de abiertos del cubrimiento
$\mathfrak U$ es afín (I, 5.5.6), por lo que
$\mathrm H^q(V,\mathscr K^i)=0$ para todo $i$ y todo $q>0$ (1.3.1); la
proposición es, por tanto, un caso particular de (0, 12.4.7).
\end{proof}

\begin{proposition}[6.2.3]
Sea $f:X\to Y$ un morfismo cuasicompacto y separado de preesquemas. Para todo
complejo $\mathscr K^\bullet$ de $\mathscr O_X$-módulos cuasicoherentes, los
$\mathscr O_Y$-módulos $\mathscr H^n(f,\mathscr K^\bullet)$ son
cuasicoherentes.
\end{proposition}

\begin{proof}
Como los $\mathscr H^q(f,\mathscr K^i)=\mathrm R^q f_*(\mathscr K^i)$ son
$\mathscr O_Y$-módulos cuasicoherentes (1.4.10), lo mismo ocurre con
$'\mathscr E_2^{pq}$, que, según (6.2.1.1), es cociente del núcleo de un
homomorfismo de módulos cuasicoherentes por la imagen de un homomorfismo de
esta clase (I, 4.1.1). Por la misma razón, todos los
$\mathscr O_Y$-módulos $'\mathscr E_r^{pq}$,
$B_k('\mathscr E_2^{pq})$, $Z_k('\mathscr E_2^{pq})$ de la primera sucesión
espectral son cuasicoherentes. La regularidad de la sucesión espectral
$'\mathscr E(f,\mathscr K^\bullet)$ implica que
$Z_\infty('\mathscr E_2^{pq})$ es igual a uno de los
$Z_k('\mathscr E_2^{pq})$, luego es cuasicoherente, y lo mismo ocurre con
$B_\infty('\mathscr E_2^{pq})=\varinjlim_k B_k('\mathscr E_2^{pq})$
(0, 11.2.4 e I, 4.1.1); por tanto, los $'\mathscr E_\infty^{pq}$ son también
cuasicoherentes. Como la sucesión espectral precedente es regular, la
filtración de los $F^p(\mathscr H^n(f,\mathscr K^\bullet))$ es discreta y
exhaustiva; dicho de otro modo, el $\mathscr O_Y$-módulo
$\mathscr H^n(f,\mathscr K^\bullet)$ es la unión de una sucesión creciente
$(\mathscr G_k)_{k\geq0}$ de $\mathscr O_Y$-módulos tal que
$\mathscr G_0=0$ y cada $\mathscr G_k/\mathscr G_{k-1}$ es igual a uno de los
$\mathscr O_Y$-módulos $'\mathscr E_\infty^{pq}$ y, por tanto, es
cuasicoherente. Por inducción sobre $k$, se deduce que los $\mathscr G_k$ son
cuasicoherentes (1.4.17), y, como
$\mathscr H^n(f,\mathscr K^\bullet)=\varinjlim_k\mathscr G_k$, queda demostrada
la proposición (I, 4.1.1).
\end{proof}

\oldpage[III]{140}

\begin{corollary}[6.2.4]
Bajo las hipótesis de (6.2.3), para todo abierto afín $V$ de $Y$, el
homomorfismo canónico
\[
  \mathbf H^n(f^{-1}(V),\mathscr K^\bullet)
  \longrightarrow
  \Gamma(V,\mathscr H^n(f,\mathscr K^\bullet))
  \tag{6.2.4.1}\label{III.6.2.4.1}
\]
es biyectivo para todo $n\in\mathbf Z$.
\end{corollary}

\begin{proof}
La demostración es la misma que la de (1.4.11), utilizando (6.2.2),
reemplazando $\mathscr F$ por $\mathscr K^\bullet$,
$\mathscr K^\bullet$ por
$f_*(C^\bullet(\mathfrak U,\mathscr K^\bullet))$,
$\mathscr H^\bullet(\mathscr K^\bullet)$ por
$\mathscr H^\bullet(f,\mathscr K^\bullet)$, y observando que este último es un
$\mathscr O_Y$-módulo cuasicoherente por (6.2.3).
\end{proof}

\begin{proposition}[6.2.5]
Sean $Y$ un preesquema localmente noetheriano, $f:X\to Y$ un morfismo propio
y $\mathscr K^\bullet$ un complejo de $\mathscr O_X$-módulos tal que los
$\mathscr O_X$-módulos $\mathscr H^q(\mathscr K^\bullet)$ sean coherentes.
Entonces los $\mathscr O_Y$-módulos $\mathscr H^n(f,\mathscr K^\bullet)$ son
coherentes.
\end{proposition}

\begin{proof}
Como la cuestión es local sobre $Y$, podemos limitarnos al caso en que $Y$ es
noetheriano y afín; se trata entonces, en virtud de (6.2.4), de demostrar que
los $\mathbf H^n(X,\mathscr K^\bullet)$ son
$\Gamma(Y,\mathscr O_Y)$-módulos de tipo finito. Se tiene entonces
$\mathbf H^\bullet(X,\mathscr K^\bullet)=
\mathbf H^\bullet(\mathfrak U,\mathscr K^\bullet)$ (6.2.2), donde se puede
suponer que $\mathfrak U$ es finito, puesto que $X$ es cuasicompacto. Como las
cocadenas de cada complejo $C^\bullet(\mathfrak U,\mathscr K^j)$ son
alternadas por definición, existe un entero $r>0$ tal que
$C^i(\mathfrak U,\mathscr K^j)=0$ para $i<0$ y $i>r$; se concluye de ello
(0, 11.3.3) que las dos sucesiones espectrales del bicomplejo
$C^\bullet(\mathfrak U,\mathscr K^\bullet)$ son birregulares. Como las
intersecciones de los miembros de $\mathfrak U$ son abiertos afines
(I, 5.5.6), cada funtor
$\mathscr F\mapsto C^i(\mathfrak U,\mathscr F)$ es exacto en la categoría de
los $\mathscr O_X$-módulos cuasicoherentes; por consiguiente,
$\mathrm H_I^q(C^i(\mathfrak U,\mathscr K^\bullet))
=C^i(\mathfrak U,\mathscr H^q(\mathscr K^\bullet))$, y los términos
$E_2$ de la segunda sucesión espectral de
$C^\bullet(\mathfrak U,\mathscr K^\bullet)$ vienen dados (0, 11.3.2) por
\[
  ''E_2^{pq}
  =\mathrm H^p(C^\bullet(\mathfrak U,\mathscr H^q(\mathscr K^\bullet)))
  =\mathbf H^p(\mathfrak U,\mathscr H^q(\mathscr K^\bullet))
  =\mathrm H^p(X,\mathscr H^q(\mathscr K^\bullet))
\]
en virtud de (1.4.1); dado que $f$ es propio, estos son
$\Gamma(Y,\mathscr O_Y)$-módulos de tipo finito (3.2.1). Puesto que la sucesión
espectral $''E(C^\bullet(\mathfrak U,\mathscr K^\bullet))$ es birregular, se
deduce que los $\mathbf H^n(X,\mathscr K^\bullet)$ son
$\Gamma(Y,\mathscr O_Y)$-módulos de tipo finito (0, 11.1.8).

A fortiori, si $\mathscr K^\bullet$ es un complejo de
$\mathscr O_X$-módulos coherentes, los $\mathscr O_Y$-módulos
$\mathscr H^n(f,\mathscr K^\bullet)$ son coherentes bajo las hipótesis de
(6.2.5) relativas a $Y$ y $f$ (0\textsubscript{I}, 5.3.4).
\end{proof}

\begin{env}[6.2.6]
La hipercohomología $\mathscr H^\bullet(f,\mathscr K^\bullet)$ es un funtor
cohomológico en la categoría de los complejos de $\mathscr O_X$-módulos
acotados inferiormente (0, 12.4.4). Es un funtor cohomológico en la categoría
de todos los complejos de $\mathscr O_X$-módulos cuando el morfismo $f$ es
cuasicompacto y el espacio subyacente a $X$ es localmente noetheriano: en
efecto, de (G, II, 3.10.1) se deduce entonces que $f_*$ conmuta con los límites
inductivos (pues la cuestión es local sobre $Y$), y puede aplicarse
(0, 11.5.2).

Por último, si $f$ es separado, $\mathscr H^\bullet(f,\mathscr K^\bullet)$ es
un funtor cohomológico en la categoría de los complejos de
$\mathscr O_X$-módulos cuasicoherentes. Esto es inmediato cuando $Y$ es
afín, pues entonces $X$ es un esquema; por tanto, en virtud del isomorfismo
canónico (6.2.2), todo se reduce a ver que
$\mathscr K^\bullet\mapsto\mathbf H^\bullet(\mathfrak U,\mathscr K^\bullet)$
es un funtor cohomológico en la categoría de los complejos de
$\mathscr O_X$-módulos cuasicoherentes, lo cual es inmediato, ya que el
funtor $\mathscr K^\bullet\mapsto
C^\bullet(\mathfrak U,\mathscr K^\bullet)$ es exacto en esta categoría
(I, 1.3.7). En el caso general, para todo abierto afín $V$

\oldpage[III]{141}

de $Y$, $f^{-1}(V)$ es un esquema y, para aplicar lo que precede, basta
verificar que, para una sucesión exacta
\[
  0\longrightarrow\mathscr K'^\bullet
  \longrightarrow\mathscr K^\bullet
  \longrightarrow\mathscr K''^\bullet
  \longrightarrow0
\]
de complejos de $\mathscr O_X$-módulos cuasicoherentes, el homomorfismo
\[
  \partial:
  \mathscr H^n(f,\mathscr K''^\bullet)|V
  \longrightarrow
  \mathscr H^{n+1}(f,\mathscr K'^\bullet)|V
\]
no depende del cubrimiento abierto afín $\mathfrak U$ de $f^{-1}(V)$ utilizado
para definirlo. Pero esto se deduce de que, si $\mathfrak U'$ es un
cubrimiento abierto afín más fino que $\mathfrak U$, el diagrama
\[
  \xymatrix@C=34pt@R=11pt{
    \mathbf H^\bullet(\mathfrak U,\mathscr K^\bullet)
      \ar[rrd]^-{\sim}\ar[dd]^{\iota}
    && \\
    &&
    \mathbf H^\bullet(f^{-1}(V),\mathscr K^\bullet)
    \\
    \mathbf H^\bullet(\mathfrak U',\mathscr K^\bullet)
      \ar[rru]_-{\sim}
    &&
  }
\]
de isomorfismos canónicos es conmutativo, al igual que el diagrama
\[
  \xymatrix@C=46pt@R=25pt{
    \mathbf H^n(\mathfrak U,\mathscr K''^\bullet)
      \ar[r]^-{\partial}\ar[d]^{\iota}
    &
    \mathbf H^{n+1}(\mathfrak U,\mathscr K'^\bullet)
      \ar[d]^{\iota}
    \\
    \mathbf H^n(\mathfrak U',\mathscr K''^\bullet)
      \ar[r]_-{\partial}
    &
    \mathbf H^{n+1}(\mathfrak U',\mathscr K'^\bullet).
  }
\]

Cuando se cumple una de las condiciones precedentes y
$\mathscr K^\bullet\to\mathscr K'^\bullet$ es una equivalencia homotópica (6.1.4), el
isomorfismo correspondiente
$\mathscr H^\bullet(f,\mathscr K^\bullet)\xrightarrow{\sim}
\mathscr H^\bullet(f,\mathscr K'^\bullet)$ es entonces un isomorfismo de
$\partial$-funtores (0, 11.4.4).
\end{env}

\begin{env}[6.2.7]
Todo lo que precede se aplica naturalmente, sin más cambio que el de la
notación, a un complejo $\mathscr K_\bullet$ de $\mathscr O_X$-módulos
cuasicoherentes cuya diferencial es de grado $-1$; basta considerar el
complejo $\mathscr K^\bullet=(\mathscr K^i)$, donde
$\mathscr K^i=\mathscr K_{-i}$ para todo $i\in\mathbf Z$.
\end{env}

\subsection{Hipertor de dos complejos de módulos}
\label{subsection:III.6.3}

\begin{env}[6.3.1]
\label{III.6.3.1}
Sean $A$ un anillo conmutativo y $P_\bullet$, $Q_\bullet$ dos complejos de
$A$-módulos cuyos diferenciales son de grado $-1$; sea
$L_{\bullet\bullet}$ (resp. $M_{\bullet\bullet}$) una resolución proyectiva
de Cartan-Eilenberg de $P_\bullet$ (resp. $Q_\bullet$)
\sref[0]{0.11.6.1};
$L_{\bullet\bullet}\otimes_A M_{\bullet\bullet}$ es entonces (para la suma
de los primeros grados y la suma de los segundos grados) un bicomplejo (con
diferenciales de grado $-1$), cuya homología
$H_\bullet(L_{\bullet\bullet}\otimes_A M_{\bullet\bullet})$ no depende de las
resoluciones de Cartan-Eilenberg $L_{\bullet\bullet}$,
$M_{\bullet\bullet}$ elegidas y es, por definición, la hiperhomología del
bifuntor $P_\bullet\otimes_A Q_\bullet$ en $P_\bullet$ y $Q_\bullet$
\sref[0]{0.11.6.5}. Pondremos, por definición,
\[
  \mathbf{Tor}_n^A(P_\bullet,Q_\bullet)
  =H_n(L_{\bullet\bullet}\otimes_A M_{\bullet\bullet})
  \tag{6.3.1.1}\label{III.6.3.1.1}
\]

\oldpage[III]{142}

y diremos que este $A$-módulo es el hipertor de índice $n$ de los dos
complejos $P_\bullet$, $Q_\bullet$. Se sabe que, en la categoría de los
complejos de $A$-módulos acotados inferiormente, los
$\mathbf{Tor}_n^A(P_\bullet,Q_\bullet)$ forman un bifuntor homológico en
$P_\bullet$, $Q_\bullet$ \sref[0]{0.11.6.5}. Además:
\end{env}

\begin{proposition}[6.3.2]
\label{III.6.3.2}
El bifuntor $\mathbf{Tor}_\bullet^A(P_\bullet,Q_\bullet)$ es el término
límite común de dos bifuntores espectrales
$'E(P_\bullet,Q_\bullet)$, $''E(P_\bullet,Q_\bullet)$, cuyos términos $E_2$
son
\begin{align}
  'E^2_{pq}
    &=H_p\bigl(\operatorname{Tor}_q^A(P_\bullet,Q_\bullet)\bigr),
    \tag{6.3.2.1}\label{III.6.3.2.1}\\
  ''E^2_{pq}
    &=\bigoplus_{q'+q''=q}
      \operatorname{Tor}_p^A\bigl(H_{q'}(P_\bullet),H_{q''}(Q_\bullet)\bigr).
    \tag{6.3.2.2}\label{III.6.3.2.2}
\end{align}
donde, en \sref{III.6.3.2.1},
$\operatorname{Tor}_q^A(P_\bullet,Q_\bullet)$ designa el
bicomplejo formado por los $A$-módulos $\operatorname{Tor}_q^A(P_i,Q_j)$. La
sucesión espectral \sref{III.6.3.2.2} es siempre regular; si $P_\bullet$ y
$Q_\bullet$
están acotados inferiormente, o si $A$ es de dimensión cohomológica finita,
las dos sucesiones espectrales \sref{III.6.3.2.1} y
\sref{III.6.3.2.2} son birregulares.
\end{proposition}

\begin{proof}
Esto resulta de \sref[0]{0.11.6.5}, pues cuando $A$ es de dimensión
cohomológica
finita $n$, todo $A$-módulo admite una resolución proyectiva de longitud $n$
\emph{(M, VI, 2.1)}.
\end{proof}

\begin{corollary}[6.3.3]
\label{III.6.3.3}
Sean $P'_\bullet$, $Q'_\bullet$ dos complejos de $A$-módulos, y
$u:P_\bullet\to P'_\bullet$, $v:Q_\bullet\to Q'_\bullet$ dos homomorfismos
de complejos. Si los homomorfismos
$H_\bullet(u):H_\bullet(P_\bullet)\to H_\bullet(P'_\bullet)$,
$H_\bullet(v):H_\bullet(Q_\bullet)\to H_\bullet(Q'_\bullet)$ deducidos
respectivamente de $u$ y $v$ son biyectivos, entonces el homomorfismo
$\mathbf{Tor}_\bullet^A(P_\bullet,Q_\bullet)\to
\mathbf{Tor}_\bullet^A(P'_\bullet,Q'_\bullet)$ deducido de $u$ y $v$ es
biyectivo.
\end{corollary}

\begin{proof}
En efecto, el homomorfismo de sucesiones espectrales
$''E(P_\bullet,Q_\bullet)\to''E(P'_\bullet,Q'_\bullet)$ deducido de $u$ y $v$
es entonces un isomorfismo para los términos $E_2$, y la conclusión resulta
de que estas sucesiones son regulares en virtud de \sref{III.6.3.2} y
\sref[0]{0.11.1.5}.
\end{proof}

\begin{proposition}[6.3.4]
\label{III.6.3.4}
Sean $P_\bullet$, $Q_\bullet$ dos complejos de $A$-módulos acotados
inferiormente. Sea $L_{\bullet\bullet}$ (resp. $M_{\bullet\bullet}$) un
bicomplejo formado por $A$-módulos planos, tal que, para todo $i$,
$L_{i,\bullet}$ (resp. $M_{i,\bullet}$) sea una resolución de $P_i$
(resp. $Q_i$). Se tienen entonces isomorfismos canónicos
\[
  \mathbf{Tor}_\bullet^A(P_\bullet,Q_\bullet)
  \xrightarrow{\sim}H_\bullet(L_{\bullet\bullet}\otimes_AQ_\bullet)
  \xrightarrow{\sim}H_\bullet(P_\bullet\otimes_AM_{\bullet\bullet})
  \xrightarrow{\sim}H_\bullet(L_{\bullet\bullet}\otimes_AM_{\bullet\bullet})
  \tag{6.3.4.1}\label{III.6.3.4.1}
\]
\end{proposition}

\begin{proof}
Esto resulta de \sref[0]{0.11.6.5}, partes (ii) y (iii), y de la definición
de los
$A$-módulos planos.
\end{proof}

\begin{namedtheorem}[6.3.5]{Observaciones}
\label{III.6.3.5}
(i) Con las notaciones de \sref{III.6.3.1}, los bicomplejos
$L_{\bullet\bullet}\otimes_A M_{\bullet\bullet}$ y
$M_{\bullet\bullet}\otimes_A L_{\bullet\bullet}$ son canónicamente
isomorfos, de donde se obtiene un isomorfismo canónico
$\mathbf{Tor}_\bullet^A(P_\bullet,Q_\bullet)\xrightarrow{\sim}
\mathbf{Tor}_\bullet^A(Q_\bullet,P_\bullet)$.

(ii) Si $F$ y $G$ son dos $A$-módulos, y $P_\bullet$ y $Q_\bullet$ son los
complejos de $A$-módulos reducidos a $F$ y $G$, respectivamente, en grado 0 y
nulos en los demás grados, entonces dos resoluciones proyectivas $L_\bullet$,
$M_\bullet$ de $F$ y $G$, respectivamente, pueden considerarse como
resoluciones de Cartan-Eilenberg de $P_\bullet$ y $Q_\bullet$ completándolas
con ceros. Por consiguiente, en este caso se tiene
$\mathbf{Tor}_\bullet^A(P_\bullet,Q_\bullet)
=\operatorname{Tor}_\bullet^A(F,G)$.
\end{namedtheorem}

\begin{proposition}[6.3.6]
\label{III.6.3.6}
Sean $(P_\bullet^\lambda)$, $(Q_\bullet^\mu)$ dos sistemas inductivos
filtrantes de complejos de $A$-módulos; se tiene un isomorfismo canónico
\[
  \varinjlim_{\lambda,\mu}
  \mathbf{Tor}_\bullet^A(P_\bullet^\lambda,Q_\bullet^\mu)
  \xrightarrow{\sim}
  \mathbf{Tor}_\bullet^A
  \left(\varinjlim_\lambda P_\bullet^\lambda,
        \varinjlim_\mu Q_\bullet^\mu\right).
  \tag{6.3.6.1}\label{III.6.3.6.1}
\]
\end{proposition}

\begin{proof}
Pongamos $P_\bullet=\varinjlim_\lambda P_\bullet^\lambda$,
$Q_\bullet=\varinjlim_\mu Q_\bullet^\mu$; por funtorialidad, está claro que
los $\mathbf{Tor}_\bullet^A(P_\bullet^\lambda,Q_\bullet^\mu)$ forman un
sistema inductivo y que las aplicaciones
$\mathbf{Tor}_\bullet^A(P_\bullet^\lambda,Q_\bullet^\mu)\to
\mathbf{Tor}_\bullet^A(P_\bullet,Q_\bullet)$ deducidas

\oldpage[III]{143}

de las aplicaciones canónicas $P_\bullet^\lambda\to P_\bullet$,
$Q_\bullet^\mu\to Q_\bullet$, forman un sistema inductivo de homomorfismos,
de donde se obtiene un homomorfismo canónico \sref{III.6.3.6.1} y, más
generalmente,
un homomorfismo canónico
\[
  \varinjlim_{\lambda,\mu}''E(P_\bullet^\lambda,Q_\bullet^\mu)
  \longrightarrow ''E(P_\bullet,Q_\bullet)
\]
del cual \sref{III.6.3.6.1} es el homomorfismo de los términos límite.
Además, la sucesión espectral $''E(P_\bullet,Q_\bullet)$ es regular
\sref{III.6.3.2}, y lo mismo
ocurre con la sucesión espectral
$\varinjlim_{\lambda,\mu}''E(P_\bullet^\lambda,Q_\bullet^\mu)$, como resulta
de las definiciones \sref[0]{0.11.1.7} y de la demostración de
\sref[0]{0.11.3.3}; para demostrar que \sref{III.6.3.6.1} es biyectivo basta,
pues, por \sref[0]{0.11.1.5}, probar que el
homomorfismo
\[
  \varinjlim_{\lambda,\mu}''E(P_\bullet^\lambda,Q_\bullet^\mu)
  \longrightarrow ''E(P_\bullet,Q_\bullet)
  \tag{6.3.6.2}\label{III.6.3.6.2}
\]
es biyectivo para los términos $E^2$. Como el funtor $H_\bullet$ conmuta con
el límite inductivo de complejos de módulos, finalmente se reduce todo a
probar que, para dos sistemas inductivos filtrantes $(F^\lambda)$,
$(G^\mu)$ de $A$-módulos, el homomorfismo canónico
\[
  \varinjlim_{\lambda,\mu}\operatorname{Tor}_\bullet^A(F^\lambda,G^\mu)
  \longrightarrow
  \operatorname{Tor}_\bullet^A
  \left(\varinjlim_\lambda F^\lambda,
        \varinjlim_\mu G^\mu\right)
\]
es biyectivo. Para ello, consideremos para cada $F^\lambda$ la resolución
libre canónica
\[
  L_\bullet^\lambda:\quad
  \cdots\longrightarrow L_{i+1}^\lambda
  \longrightarrow L_i^\lambda\longrightarrow\cdots
  \longrightarrow L_1^\lambda\longrightarrow L_0^\lambda\longrightarrow0
\]
donde $L_0^\lambda$ es el $A$-módulo de las combinaciones lineales formales
de elementos de $F^\lambda$, y $L_{i+1}^\lambda$ es el $A$-módulo de las
combinaciones lineales formales de elementos de
$\operatorname{Ker}(L_i^\lambda\to L_{i-1}^\lambda)$; se comprueba
inmediatamente que los $L_\bullet^\lambda$ forman un sistema inductivo de
complejos, y si se pone
$F=\varinjlim_\lambda F^\lambda$, $L_i=\varinjlim_\lambda L_i^\lambda$, los
$L_i$ forman una resolución $L_\bullet$ de $F$, pues el funtor
$\varinjlim$ es exacto; además, los $L_i$, límites inductivos de
$A$-módulos libres, son planos \sref[0]{0.6.1.2}. Se considera de
igual modo, para cada $\mu$, la resolución libre canónica $M_\bullet^\mu$ de
$G^\mu$, y $M_\bullet=\varinjlim_\mu M_\bullet^\mu$ es una resolución plana
de $G=\varinjlim_\mu G^\mu$. Se tiene entonces
$\operatorname{Tor}_\bullet^A(\varinjlim_\lambda F^\lambda,
\varinjlim_\mu G^\mu)=H_\bullet(L_\bullet\otimes_A M_\bullet)$ en virtud de
\sref{III.6.3.5} y \sref{III.6.3.4}; pero
\[
  H_\bullet(L_\bullet\otimes_A M_\bullet)
  =\varinjlim_{\lambda,\mu}
   H_\bullet(L_\bullet^\lambda\otimes_A M_\bullet^\mu)
\]
puesto que $H_\bullet$ conmuta con los límites inductivos de complejos de
módulos; como $H_\bullet(L_\bullet^\lambda\otimes_A M_\bullet^\mu)
=\operatorname{Tor}_\bullet^A(F^\lambda,G^\mu)$, esto demuestra
\sref{III.6.3.6.1}.

Cuando se supone que existe $i_0$ tal que
$P_i^\lambda=Q_i^\mu=0$ para $i<i_0$, cualesquiera que sean $\lambda$ y
$\mu$, se demuestra del mismo modo que el homomorfismo canónico
\[
  \varinjlim_{\lambda,\mu}'E(P_\bullet^\lambda,Q_\bullet^\mu)
  \longrightarrow'E(P_\bullet,Q_\bullet)
  \tag{6.3.6.3}\label{III.6.3.6.3}
\]
es biyectivo.
\end{proof}

\begin{proposition}[6.3.7]
\label{III.6.3.7}
Supongamos que $P_\bullet$ y $Q_\bullet$ están acotados inferiormente. Si el
complejo $P_\bullet$ está formado por $A$-módulos planos, se tiene un
isomorfismo $A$-lineal canónico de $\partial$-funtores en $Q_\bullet$
\[
  \mathbf{Tor}_\bullet^A(P_\bullet,Q_\bullet)
  \xrightarrow{\sim}H_\bullet(P_\bullet\otimes_A Q_\bullet).
  \tag{6.3.7.1}\label{III.6.3.7.1}
\]
\end{proposition}

\begin{proof}
En efecto, la sucesión espectral \sref{III.6.3.2.1} es birregular y
degenerada, y la existencia del isomorfismo \sref{III.6.3.7.1} resulta de
\sref[0]{0.11.1.6}. Además,
calculando el hipertor a partir de una resolución proyectiva de
Cartan-Eilenberg de $P_\bullet$ como en \sref{III.6.3.4}, se ve
inmediatamente que el
isomorfismo así definido es un isomorfismo de $\partial$-funtores en
$Q_\bullet$.
\end{proof}
\begin{env}[6.3.8]
\label{III.6.3.8}
\oldpage[III]{144}
Sea $\rho:A\to A'$ un homomorfismo de anillos. Nos proponemos definir un
homomorfismo $A$-lineal funtorial de grado 0 asociado canónicamente a $\rho$:
\[
  \rho_{P_\bullet,Q_\bullet}:
  \mathbf{Tor}_\bullet^A(P_\bullet,Q_\bullet)
  \longrightarrow
  \mathbf{Tor}_\bullet^{A'}
  (P_\bullet\otimes_A A',Q_\bullet\otimes_A A').
  \tag{6.3.8.1}\label{III.6.3.8.1}
\]

Para ello, consideremos una resolución proyectiva de Cartan-Eilenberg
$L_{\bullet\bullet}$ de $P_\bullet$; consideremos, por otra parte, una
resolución proyectiva de Cartan-Eilenberg $L'_{\bullet\bullet}$ de
$P_\bullet\otimes_A A'$. Veremos que se puede definir un
homomorfismo $A'$-lineal de complejos
$L_{\bullet\bullet}\otimes_A A'\to L'_{\bullet\bullet}$, determinado salvo
homotopía. En efecto, la construcción de $L_{\bullet\bullet}$ queda
enteramente determinada cuando se elige (arbitrariamente), para cada $i$,
una resolución proyectiva $(X_{ij}^{\mathrm B})_{j\geqslant0}$ de
$\mathrm B_i(P_\bullet)$ y una resolución proyectiva
$(X_{ij}^{\mathrm H})_{j\geqslant0}$ de $\mathrm H_i(P_\bullet)$, que son
respectivamente iguales a $\mathrm B_i^{\mathrm I}(L_{\bullet\bullet})$ y
$\mathrm H_i^{\mathrm I}(L_{\bullet\bullet})$; se deduce sucesivamente
\[
  \mathrm Z_i^{\mathrm I}(L_{\bullet\bullet})
  =\mathrm H_i^{\mathrm I}(L_{\bullet\bullet})
   \oplus\mathrm B_i^{\mathrm I}(L_{\bullet\bullet}),
  \qquad
  L_{i,\bullet}=\mathrm Z_i^{\mathrm I}(L_{\bullet\bullet})
   \oplus\mathrm B_{i-1}^{\mathrm I}(L_{\bullet\bullet}).
\]
Ahora bien, $X_{i,\bullet}^{\mathrm B}\otimes_A A'$ ya no es, en general,
una resolución de $\mathrm B_i(P_\bullet)\otimes_A A'$, pero sigue siendo un
complejo formado por $A'$-módulos proyectivos, y por tanto existe un
homomorfismo $A'$-lineal
\[
  X_{i,\bullet}^{\mathrm B}\otimes_A A'
  \longrightarrow \mathrm B_i^{\mathrm I}(L'_{\bullet\bullet})
\]
compatible con las aumentaciones y determinado salvo homotopía
\emph{(M, V, 1.1)}. Se tiene asimismo un homomorfismo $A'$-lineal
\[
  X_{i,\bullet}^{\mathrm H}\otimes_A A'
  \longrightarrow \mathrm H_i^{\mathrm I}(L'_{\bullet\bullet})
\]
determinado salvo homotopía, de donde se deduce, mediante la construcción
recordada más arriba, un homomorfismo $A'$-lineal
$L_{i,\bullet}\otimes_A A'\to L'_{i,\bullet}$ para todo $i$; estos
homomorfismos (para $i\in\mathbf Z$) son compatibles con los diferenciales
$L_{i,\bullet}\to L_{i-1,\bullet}$ y con los análogos para
$L'_{\bullet\bullet}$, en virtud de la misma construcción, y constituyen por
tanto el homomorfismo $A'$-lineal
$L_{\bullet\bullet}\otimes_A A'\to L'_{\bullet\bullet}$ buscado.

Para definir \sref{III.6.3.8.1}, basta entonces considerar de igual modo una
resolución proyectiva de Cartan-Eilenberg $M_{\bullet\bullet}$ (resp.
$M'_{\bullet\bullet}$) de $Q_\bullet$ (resp.
$Q_\bullet\otimes_A A'$), y un homomorfismo $A'$-lineal
$M_{\bullet\bullet}\otimes_A A'\to M'_{\bullet\bullet}$. De estos
homomorfismos se deduce un homomorfismo $A'$-lineal
\[
  (L_{\bullet\bullet}\otimes_A A')\otimes_{A'}
  (M_{\bullet\bullet}\otimes_A A')
  \longrightarrow
  L'_{\bullet\bullet}\otimes_{A'}M'_{\bullet\bullet},
\]
y luego, por composición, un homomorfismo $A$-lineal de bicomplejos
\[
  L_{\bullet\bullet}\otimes_A M_{\bullet\bullet}
  \longrightarrow
  L'_{\bullet\bullet}\otimes_{A'}M'_{\bullet\bullet},
\]
y al pasar a la homología se obtiene \sref{III.6.3.8.1}, que está bien
definido
puesto que proviene de un morfismo de complejos determinado salvo homotopía.

Si $\rho':A'\to A''$ es un segundo homomorfismo de anillos, y
$\rho'':A\to A''$ es el homomorfismo compuesto $\rho'\circ\rho$, está claro
que
\[
  \rho''_{P_\bullet,Q_\bullet}
  =\rho'_{P'_\bullet,Q'_\bullet}\circ\rho_{P_\bullet,Q_\bullet},
  \qquad
  P'_\bullet=P_\bullet\otimes_A A',
  \quad Q'_\bullet=Q_\bullet\otimes_A A'.
\]
Observemos además que el morfismo de bicomplejos
$L_{\bullet\bullet}\otimes_A M_{\bullet\bullet}
\to L'_{\bullet\bullet}\otimes_{A'}M'_{\bullet\bullet}$ considerado
más arriba define morfismos funtoriales (en $P_\bullet$ y $Q_\bullet$) de
sucesiones espectrales
\[
  'E_{pq}^r(P_\bullet,Q_\bullet)
  \longrightarrow
  'E_{pq}^r(P_\bullet\otimes_A A',Q_\bullet\otimes_A A')
\]
y
\[
  ''E_{pq}^r(P_\bullet,Q_\bullet)
  \longrightarrow
  ''E_{pq}^r(P_\bullet\otimes_A A',Q_\bullet\otimes_A A'),
\]
independientes de las resoluciones de Cartan-Eilenberg consideradas y que
poseen también la propiedad de transitividad precedente.
\end{env}

\begin{proposition}[6.3.9]
\label{III.6.3.9}
Sea $\rho:A\to A'$ un homomorfismo de anillos tal que $A'$ sea un
$A$-módulo plano. Se tienen entonces isomorfismos canónicos funtoriales
\[
  \mathbf{Tor}_\bullet^{A'}
  (P_\bullet\otimes_A A',Q_\bullet\otimes_A A')
  \xrightarrow{\sim}
  \mathbf{Tor}_\bullet^A(P_\bullet,Q_\bullet)\otimes_A A'.
  \tag{6.3.9.1}\label{III.6.3.9.1}
\]
\[
  \left\{
  \begin{aligned}
  'E(P_\bullet\otimes_A A',Q_\bullet\otimes_A A')
    &\xrightarrow{\sim}'E(P_\bullet,Q_\bullet)\otimes_A A',\\
  ''E(P_\bullet\otimes_A A',Q_\bullet\otimes_A A')
    &\xrightarrow{\sim}''E(P_\bullet,Q_\bullet)\otimes_A A'.
  \end{aligned}
  \right.
  \tag{6.3.9.2}\label{III.6.3.9.2}
\]
\end{proposition}

\begin{proof}
\oldpage[III]{145}
En efecto, dada la exactitud del funtor $M\otimes_A A'$ en $M$,
$L_{\bullet\bullet}\otimes_A A'$ y $M_{\bullet\bullet}\otimes_A A'$ son
entonces resoluciones proyectivas de Cartan-Eilenberg de
$P_\bullet\otimes_A A'$ y $Q_\bullet\otimes_A A'$, respectivamente, de
donde resulta la conclusión.
\end{proof}

\begin{env}[6.3.10]
\label{III.6.3.10}
Sea $\rho:A\to A'$ un homomorfismo de anillos; para todo complejo
$P'_\bullet$ de $A'$-módulos, $P'_{\bullet[\rho]}$ es un complejo de
$A$-módulos; además, la aplicación identidad
$P'_{\bullet[\rho]}\to P'_\bullet$ puede considerarse como el compuesto de
las aplicaciones canónicas
\[
  P'_{\bullet[\rho]}\longrightarrow
  P'_{\bullet[\rho]}\otimes_A A'
  \xrightarrow{\mu}P'_\bullet,
\]
donde $\mu$ es el homomorfismo $A'$-lineal $\mu(x\otimes a')=a'x$. Si
$Q'_\bullet$
es un segundo complejo de $A'$-módulos, se tienen por tanto homomorfismos
canónicos funtoriales de grado 0
\[
  \mathbf{Tor}_\bullet^A(P'_{\bullet[\rho]},Q'_{\bullet[\rho]})
  \longrightarrow
  \mathbf{Tor}_\bullet^{A'}
  (P'_{\bullet[\rho]}\otimes_A A',Q'_{\bullet[\rho]}\otimes_A A')
  \longrightarrow
  \mathbf{Tor}_\bullet^{A'}(P'_\bullet,Q'_\bullet),
  \tag{6.3.10.1}\label{III.6.3.10.1}
\]
donde la primera flecha es el homomorfismo $A$-lineal definido en
\sref{III.6.3.8}, y la
segunda se deduce de los homomorfismos $A'$-lineales
$P'_{\bullet[\rho]}\otimes_A A'\to P'_\bullet$ y
$Q'_{\bullet[\rho]}\otimes_A A'\to Q'_\bullet$ por funtorialidad. Se tienen
homomorfismos análogos para las sucesiones espectrales de
\sref{III.6.3.2}, y
propiedades evidentes de transitividad, cuyo enunciado dejamos al cuidado del
lector.
\end{env}

\begin{proposition}[6.3.11]
\label{III.6.3.11}
Sea $\rho:A\to A'$ un homomorfismo de anillos que hace de $A'$ un
$A$-módulo plano. Para todo complejo $P'_\bullet$ de $A'$-módulos y todo
complejo $Q_\bullet$ de $A$-módulos acotado inferiormente, se tiene un
isomorfismo canónico funtorial
\[
  \mathbf{Tor}_\bullet^A(P'_{\bullet[\rho]},Q_\bullet)
  \xrightarrow{\sim}
  \mathbf{Tor}_\bullet^{A'}(P'_\bullet,Q_\bullet\otimes_A A').
  \tag{6.3.11.1}\label{III.6.3.11.1}
\]
\end{proposition}

\begin{proof}
En efecto, si $M_{\bullet\bullet}$ es una resolución proyectiva de
Cartan-Eilenberg de $Q_\bullet$, $M_{\bullet\bullet}\otimes_A A'$ es una
resolución proyectiva de Cartan-Eilenberg de $Q_\bullet\otimes_A A'$, y se
tiene, salvo isomorfismo canónico,
\[
  P'_{\bullet[\rho]}\otimes_A M_{\bullet\bullet}
  =P'_\bullet\otimes_{A'}(M_{\bullet\bullet}\otimes_A A');
\]
la conclusión resulta de \sref{III.6.3.4}.
\end{proof}

\begin{namedtheorem}[6.3.12]{Observación}
\label{III.6.3.12}
Sea $(A^\lambda)$ un sistema inductivo filtrante de anillos, y sean
$(P_\bullet^\lambda)$, $(Q_\bullet^\lambda)$ dos sistemas inductivos de
complejos de $(A^\lambda)$-módulos; se tiene entonces un isomorfismo canónico
que generaliza \sref{III.6.3.6.1}
\[
  \varinjlim_\lambda
  \mathbf{Tor}_\bullet^{A^\lambda}
  (P_\bullet^\lambda,Q_\bullet^\lambda)
  \xrightarrow{\sim}
  \mathbf{Tor}_\bullet^A(P_\bullet,Q_\bullet),
  \tag{6.3.12.1}\label{III.6.3.12.1}
\]
donde $A=\varinjlim A^\lambda$, $P_\bullet=\varinjlim P_\bullet^\lambda$,
$Q_\bullet=\varinjlim Q_\bullet^\lambda$. Una vez definidos, con ayuda de
\sref{III.6.3.10}, los homomorfismos
\[
  \mathbf{Tor}_\bullet^{A^\lambda}
  (P_\bullet^\lambda,Q_\bullet^\lambda)
  \longrightarrow
  \mathbf{Tor}_\bullet^{A^\mu}
  (P_\bullet^\mu,Q_\bullet^\mu)
\]
para $\lambda\leqslant\mu$, la demostración es la de \sref{III.6.3.6}.
\end{namedtheorem}

\begin{proposition}[6.3.13]
\label{III.6.3.13}
Sean $S$ un subconjunto multiplicativo de $A$, y $P_\bullet$ y $Q_\bullet$
dos complejos de $A$-módulos en los que las multiplicaciones por los
elementos de $S$ sean biyectivas, de modo que, si $A'=S^{-1}A$,
$P_\bullet$ y $Q_\bullet$ estén formados por $A'$-módulos. Entonces se tiene
un isomorfismo canónico
\[
  \mathbf{Tor}_\bullet^A(P_\bullet,Q_\bullet)
  \xrightarrow{\sim}
  \mathbf{Tor}_\bullet^{A'}(P_\bullet,Q_\bullet).
\]
\end{proposition}

\begin{proof}
En efecto, la hipótesis implica que los homomorfismos canónicos
$P_\bullet\to P_\bullet\otimes_A A'$,
$Q_\bullet\to Q_\bullet\otimes_A A'$ son biyectivos. Por otra parte, la
funtorialidad del hipertor muestra que la multiplicación por todo $s\in S$
es biyectiva en $\mathbf{Tor}_\bullet^A(P_\bullet,Q_\bullet)$ y, por
consiguiente,
\[
  \mathbf{Tor}_\bullet^A(P_\bullet,Q_\bullet)
  \longrightarrow
  \mathbf{Tor}_\bullet^A(P_\bullet,Q_\bullet)\otimes_A A'
\]

\oldpage[III]{146}

es también un homomorfismo biyectivo. Como $A'$ es un $A$-módulo plano, la
conclusión resulta de \sref{III.6.3.9}, y se tienen asimismo isomorfismos
canónicos
para las sucesiones espectrales.
\end{proof}

\subsection{Funtores hipertor locales de complejos de módulos cuasicoherentes: caso de los esquemas afines}
\label{subsection:III.6.4}
\label{III.6.4}

\begin{env}[6.4.1]
\label{III.6.4.1}
Sean $S$ un esquema afín de anillo $A$, y $X$ e $Y$ dos
$S$-esquemas afines de anillos $B$ y $C$, respectivamente; así,
$B$ y $C$ son $A$-álgebras. Todo complejo $\sh P_\bullet$ de
$\sh O_X$-módulos cuasicoherentes es de la forma
$\widetilde{P_\bullet}$ para un complejo $P_\bullet$ de
$B$-módulos, y análogamente todo complejo $\sh Q_\bullet$ de
$\sh O_Y$-módulos cuasicoherentes es de la forma
$\widetilde{Q_\bullet}$ para un complejo $Q_\bullet$ de
$C$-módulos (\sref[I]{I.1.3.7} y \sref[I]{I.1.3.8}).

Considerados como complejos de $A$-módulos, $P_\bullet$ y
$Q_\bullet$ tienen módulos hipertor
$\mathbf{Tor}_n^A(P_\bullet,Q_\bullet)$. Por bifuntorialidad, las
$A$-álgebras $B$ y $C$ actúan sobre este $A$-módulo; estas acciones
lo convierten en un $(B,C)$-bimódulo o, equivalentemente, en un
módulo sobre
\[
  B\otimes_AC=A(X\times_SY).
\]
Hemos definido así un $\sh O_{X\times_SY}$-módulo cuasicoherente
\[
  \shTor_n^{\sh O_S}(\sh P_\bullet,\sh Q_\bullet)
  =
  \bigl(\mathbf{Tor}_n^A(P_\bullet,Q_\bullet)\bigr)\supertilde.
  \tag{6.4.1.1}\label{III.6.4.1.1}
\]
Se denomina \emph{hipertor local de índice $n$} de
$\sh P_\bullet$ y $\sh Q_\bullet$, y también se denota por
$\shTor_n^S(\sh P_\bullet,\sh Q_\bullet)$.
\end{env}

\begin{lemma}[6.4.2]
\label{III.6.4.2}
Con las notaciones de \sref{III.6.4.1}, supongamos que
$A=R^{-1}A'$, donde $A'$ es un anillo y $R$ una parte
multiplicativa de $A'$. Sea $S'=\Spec(A')$. Entonces $X$ e $Y$
pueden considerarse como $S'$-preesquemas y
$X\times_{S'}Y=X\times_SY$
(\sref[I]{I.1.6.2} y \sref[I]{I.3.2.4}); además,
\[
  \shTor_\bullet^{S'}(\sh P_\bullet,\sh Q_\bullet)
  =
  \shTor_\bullet^S(\sh P_\bullet,\sh Q_\bullet).
\]
\end{lemma}

\begin{proof}
Esto se deduce de \sref{III.6.4.1.1} y \sref{III.6.3.13}.
\end{proof}

\begin{env}[6.4.3]
\label{III.6.4.3}
Con las notaciones e hipótesis de \sref{III.6.4.1}, sean
$\sh F=\widetilde F$ un $\sh O_X$-módulo cuasicoherente y
$\sh G=\widetilde G$ un $\sh O_Y$-módulo cuasicoherente.
Considerando $\sh F$ y $\sh G$ como complejos concentrados en
grado cero, denotamos por $\shTor_n^{\sh O_S}(\sh F,\sh G)$, o
por $\shTor_n^S(\sh F,\sh G)$, su hipertor local de índice $n$.
Por \sref{III.6.3.5}(ii),
\[
  \shTor_n^{\sh O_S}(\sh F,\sh G)
  =
  \bigl(\Tor_n^A(F,G)\bigr)\supertilde.
  \tag{6.4.3.1}\label{III.6.4.3.1}
\]

Volvamos ahora a dos complejos arbitrarios de módulos
cuasicoherentes $\sh P_\bullet$ y $\sh Q_\bullet$. Las fórmulas
\sref{III.6.4.1.1} y \sref{III.6.4.3.1}, junto con
\sref{III.6.3.2}, muestran que
$\shTor_\bullet^S(\sh P_\bullet,\sh Q_\bullet)$ es el objeto al
que convergen dos sucesiones espectrales
$'\sh E(\sh P_\bullet,\sh Q_\bullet)$ y
$''\sh E(\sh P_\bullet,\sh Q_\bullet)$, cuyos términos $E_2$ son
\[
  '\sh E_{pq}^2
  =
  \sh H_p\bigl(\shTor_q^S(\sh P_\bullet,\sh Q_\bullet)\bigr)
  \tag{6.4.3.2}\label{III.6.4.3.2}
\]
y
\[
  ''\sh E_{pq}^2
  =
  \bigoplus_{q'+q''=q}
  \shTor_p^S
  \bigl(\sh H_{q'}(\sh P_\bullet),
        \sh H_{q''}(\sh Q_\bullet)\bigr).
  \tag{6.4.3.3}\label{III.6.4.3.3}
\]
Aquí $\shTor_q^S(\sh P_\bullet,\sh Q_\bullet)$ en
\sref{III.6.4.3.2} denota el bicomplejo de
$\sh O_{X\times_SY}$-módulos cuasicoherentes
$\shTor_q^S(\sh P_i,\sh Q_j)$.
\end{env}

\begin{env}[6.4.4]
\label{III.6.4.4}
Sean
\[
  X^{(1)}=\Spec(B^{(1)}),
  \qquad
  Y^{(1)}=\Spec(C^{(1)})
\]
otros dos esquemas afines, donde $B^{(1)}$ y $C^{(1)}$ son
$A$-álgebras, y sean
\oldpage[III]{147}
\[
  u:X^{(1)}\longrightarrow X,
  \qquad
  v:Y^{(1)}\longrightarrow Y
\]
$S$-morfismos correspondientes a los $A$-homomorfismos
$\vphi:B\longrightarrow B^{(1)}$ y
$\psi:C\longrightarrow C^{(1)}$. Consideremos los complejos
\[
  u^*(\sh P_\bullet)
  =
  (P_\bullet\otimes_BB^{(1)})\supertilde,
  \qquad
  v^*(\sh Q_\bullet)
  =
  (Q_\bullet\otimes_CC^{(1)})\supertilde.
\]
Los $A$-homomorfismos canónicos
\[
  P_\bullet\longrightarrow P_\bullet\otimes_BB^{(1)},
  \qquad
  Q_\bullet\longrightarrow Q_\bullet\otimes_CC^{(1)}
\]
dan, por funtorialidad, un $A$-homomorfismo
\[
  \mathbf{Tor}_\bullet^A(P_\bullet,Q_\bullet)
  \longrightarrow
  \mathbf{Tor}_\bullet^A
  (P_\bullet\otimes_BB^{(1)},
   Q_\bullet\otimes_CC^{(1)}).
\]
De nuevo por funtorialidad, se trata en realidad de un homomorfismo
de $(B\otimes_AC)$-módulos. Por tanto, define un
$(u\times_Sv)$-morfismo
\[
  \theta:
  \shTor_\bullet^S(\sh P_\bullet,\sh Q_\bullet)
  \longrightarrow
  \shTor_\bullet^S(u^*\sh P_\bullet,v^*\sh Q_\bullet),
  \tag{6.4.4.1}\label{III.6.4.4.1}
\]
y, en consecuencia, un homomorfismo de
$\sh O_{X^{(1)}\times_SY^{(1)}}$-módulos
\[
  \theta^\sharp:
  (u\times_Sv)^*
  \shTor_\bullet^S(\sh P_\bullet,\sh Q_\bullet)
  \longrightarrow
  \shTor_\bullet^S(u^*\sh P_\bullet,v^*\sh Q_\bullet).
  \tag{6.4.4.2}\label{III.6.4.4.2}
\]
Este es manifiestamente un morfismo de bi-$\partial$-funtores en
las categorías de módulos cuasicoherentes acotados inferiormente.
El homomorfismo \sref{III.6.4.4.2} no es necesariamente biyectivo;
sin embargo:
\end{env}

\begin{lemma}[6.4.5]
\label{III.6.4.5}
Con las notaciones de \sref{III.6.4.4}, si $u$ y $v$ son
inmersiones abiertas, entonces \sref{III.6.4.4.2} es biyectivo.
\end{lemma}

\begin{proof}
Identifiquemos $X^{(1)}$ e $Y^{(1)}$ con abiertos de $X$ e $Y$,
respectivamente. Son uniones de abiertos de la forma $\mathrm D(f)$
y $\mathrm D(g)$, donde $f\in B$ y $g\in C$; los preesquemas
inducidos sobre $\mathrm D(\vphi(f))$ y $\mathrm D(f)$, y asimismo
sobre $\mathrm D(\psi(g))$ y $\mathrm D(g)$, son isomorfos.

Basta demostrar el lema cuando $X^{(1)}=\mathrm D(f)$ e
$Y^{(1)}=\mathrm D(g)$. En efecto, una vez sabido esto, en el caso
general basta mostrar que la restricción de $\theta^\sharp$ a cada
$\mathrm D(f)\times_S\mathrm D(g)$ es un isomorfismo. Si
\[
  u_1:\mathrm D(f)\longrightarrow X^{(1)},
  \qquad
  v_1:\mathrm D(g)\longrightarrow Y^{(1)}
\]
son las inclusiones canónicas, esta restricción es
$(u_1\times_Sv_1)^*(\theta^\sharp)$. De
\sref{III.6.4.4} y \sref[0]{0.4.4.8} se deduce inmediatamente que,
al componerla con el homomorfismo canónico
\[
  (u_1\times_Sv_1)^*
  \shTor_\bullet^S(u^*\sh P_\bullet,v^*\sh Q_\bullet)
  \longrightarrow
  \shTor_\bullet^S((u')^*\sh P_\bullet,(v')^*\sh Q_\bullet),
  \tag{6.4.5.1}\label{III.6.4.5.1}
\]
donde $u'=u\circ u_1$ y $v'=v\circ v_1$, se obtiene el
homomorfismo canónico
\[
  (u'\times_Sv')^*
  \shTor_\bullet^S(\sh P_\bullet,\sh Q_\bullet)
  \longrightarrow
  \shTor_\bullet^S((u')^*\sh P_\bullet,(v')^*\sh Q_\bullet).
  \tag{6.4.5.2}\label{III.6.4.5.2}
\]
Si \sref{III.6.4.5.1} y \sref{III.6.4.5.2} son isomorfismos,
también lo es $(u_1\times_Sv_1)^*(\theta^\sharp)$.

Supongamos, pues, que $X^{(1)}=\mathrm D(f)$ e
$Y^{(1)}=\mathrm D(g)$, de modo que $B^{(1)}=B_f$ y
$C^{(1)}=C_g$. Entonces $u^*(\sh P_\bullet)$ y
$v^*(\sh Q_\bullet)$ se identifican, respectivamente, con
$((P_\bullet)_f)\supertilde$ y $((Q_\bullet)_g)\supertilde$.
Además, $X^{(1)}\times_SY^{(1)}$ se identifica con el subesquema
abierto $\mathrm D(f\otimes g)$ de
\[
  X\times_SY=\Spec(B\otimes_AC)
\]
(\sref[II]{II.4.3.2.4}). Queda por demostrar que el homomorfismo
\[
  \bigl(\mathbf{Tor}_\bullet^A(P_\bullet,Q_\bullet)\bigr)_{f\otimes g}
  \longrightarrow
  \mathbf{Tor}_\bullet^A((P_\bullet)_f,(Q_\bullet)_g),
  \tag{6.4.5.3}\label{III.6.4.5.3}
\]
\oldpage[III]{148}
inducido por los homomorfismos canónicos
$P_\bullet\longrightarrow(P_\bullet)_f$ y
$Q_\bullet\longrightarrow(Q_\bullet)_g$, es biyectivo.

Por \sref[0]{0.1.6.1},
\[
  (P_\bullet)_f=\varinjlim_m P_\bullet^{(m)},
\]
donde cada $P_\bullet^{(m)}$ es el mismo complejo de $B$-módulos
que $P_\bullet$, y para $m\leq n$ la aplicación
$P_\bullet^{(m)}\longrightarrow P_\bullet^{(n)}$ es la
multiplicación por $f^{\,n-m}$. El enunciado análogo vale para
$Q_\bullet$, sustituyendo $f$ por $g$. El homomorfismo
\[
  \mathbf{Tor}_\bullet^A(P_\bullet^{(m)},Q_\bullet^{(m)})
  \longrightarrow
  \mathbf{Tor}_\bullet^A(P_\bullet^{(n)},Q_\bullet^{(n)})
\]
correspondiente a estas aplicaciones de transición es, por
definición, la multiplicación por $(f\otimes g)^{n-m}$. La
conclusión se deduce de \sref[0]{0.1.6.1}, aplicado al miembro
izquierdo de \sref{III.6.4.5.3}, y de \sref{III.6.3.6}.
\end{proof}

\begin{env}[6.4.6]
\label{III.6.4.6}
Con las notaciones de \sref{III.6.4.4}, se definen del mismo modo
homomorfismos canónicos de funtores espectrales
\[
  \left\{
  \begin{aligned}
  (u\times_Sv)^*\,{}'\sh E(\sh P_\bullet,\sh Q_\bullet)
  &\longrightarrow
  {}'\sh E(u^*\sh P_\bullet,v^*\sh Q_\bullet),\\
  (u\times_Sv)^*\,{}''\sh E(\sh P_\bullet,\sh Q_\bullet)
  &\longrightarrow
  {}''\sh E(u^*\sh P_\bullet,v^*\sh Q_\bullet).
  \end{aligned}
  \right.
  \tag{6.4.6.1}\label{III.6.4.6.1}
\]
El razonamiento de \sref{III.6.4.5} muestra que, cuando $u$ y $v$
son inmersiones abiertas, estos homomorfismos son biyectivos.
Teniendo en cuenta \sref{III.6.3.6.2} y \sref{III.6.3.6.3}, el
razonamiento demuestra que son isomorfismos en los términos $E^2$,
mientras que \sref{III.6.4.5} demuestra que son isomorfismos en los
objetos límite. La afirmación se deduce entonces de
\sref[0]{0.11.1.2} y \sref[0]{0.11.2.4}.
\end{env}

\subsection{Funtores hipertor locales de complejos de módulos cuasicoherentes: caso general}
\label{subsection:III.6.5}

\begin{env}[6.5.1]
\label{III.6.5.1}
Sea ahora $S$ un preesquema arbitrario, y sean $X$ e $Y$ dos
$S$-preesquemas. Sean $\sh P_\bullet$ un complejo de
$\sh O_X$-módulos cuasicoherentes y $\sh Q_\bullet$ un complejo
de $\sh O_Y$-módulos cuasicoherentes. Pongamos $Z=X\times_SY$.
Vamos a definir $\sh O_Z$-módulos cuasicoherentes
\[
  \shTor_n^S(\sh P_\bullet,\sh Q_\bullet),
\]
denominados \emph{hipertor locales} de $\sh P_\bullet$ y
$\sh Q_\bullet$, que se reducen a los de \sref{III.6.4} cuando
$S$, $X$ e $Y$ son afines.

Cuando $\sh P_\bullet$ y $\sh Q_\bullet$ están concentrados en
grado cero, con términos respectivos $\sh F$ y $\sh G$, escribimos
$\shTor_n^S(\sh F,\sh G)$ en lugar de
$\shTor_n^S(\sh P_\bullet,\sh Q_\bullet)$.
\end{env}

\begin{env}[6.5.2]
\label{III.6.5.2}
Supongamos primero que $S$ es afín, y sean $(X_\lambda)$ y
$(Y_\mu)$ recubrimientos de $X$ e $Y$ por abiertos afines.
Entonces
\[
  Z_{\lambda\mu}=X_\lambda\times_SY_\mu
\]
forma un recubrimiento abierto afín de $Z$. Pongamos
\[
  \sh P_{\lambda,\bullet}=\sh P_\bullet|X_\lambda,
  \qquad
  \sh Q_{\mu,\bullet}=\sh Q_\bullet|Y_\mu.
\]
Para cada par $(\lambda,\mu)$ tenemos así un
$\sh O_{Z_{\lambda\mu}}$-módulo cuasicoherente
\[
  \sh F_{\lambda\mu}
  =
  \shTor_n^S(\sh P_{\lambda,\bullet},\sh Q_{\mu,\bullet}).
\]
Debemos mostrar que los $\sh F_{\lambda\mu}$ satisfacen la
condición de pegado \sref[0]{0.3.3.1}. Basta comprobar que, para
todo abierto afín
\[
  U\subset X_\lambda\cap X_{\lambda'},
  \qquad
  V\subset Y_\mu\cap Y_{\mu'},
\]
las restricciones de $\sh F_{\lambda\mu}$ y
$\sh F_{\lambda'\mu'}$ a $U\times_SV$ son canónicamente
isomorfas. Esto se deduce inmediatamente de sus isomorfismos
canónicos con
\[
  \shTor_n^S(\sh P_\bullet|U,\sh Q_\bullet|V),
\]
dados por \sref{III.6.4.5}.

La definición y \sref{III.6.4.5} muestran también que el
$\sh O_Z$-módulo así obtenido es independiente, salvo isomorfismo
único compatible con las identificaciones locales, de los
recubrimientos abiertos elegidos.
\oldpage[III]{149}
Lo denotamos por
\[
  \shTor_n^S(\sh P_\bullet,\sh Q_\bullet).
\]
Finalmente, para todo abierto $U$ de $X$ y todo abierto $V$ de
$Y$, su restricción a $U\times_SV$ es canónicamente isomorfa a
\[
  \shTor_n^S(\sh P_\bullet|U,\sh Q_\bullet|V).
\]
\end{env}

\begin{env}[6.5.3]
\label{III.6.5.3}
Pasemos ahora al caso general, en que $S$ es arbitrario. Sea
$(S_\alpha)$ un recubrimiento de $S$ por abiertos afines, y sean
$X_\alpha$ e $Y_\alpha$ las imágenes inversas de $S_\alpha$ en
$X$ e $Y$. Debemos mostrar que los haces
\[
  \sh G_\alpha
  =
  \shTor_n^{S_\alpha}
  (\sh P_\bullet|X_\alpha,\sh Q_\bullet|Y_\alpha)
\]
satisfacen la condición de pegado.

Para todo abierto afín $T\subset S_\alpha\cap S_\beta$, sean
$U$ y $V$ sus imágenes inversas en $X$ e $Y$. Basta definir
isomorfismos canónicos entre las restricciones de $\sh G_\alpha$
y $\sh G_\beta$ a $U\times_SV$ y
\[
  \shTor_n^T(\sh P_\bullet|U,\sh Q_\bullet|V).
\]
Podemos restringirnos además al caso en que
$T=\mathrm D(f_\alpha)=\mathrm D(f_\beta)$, donde $f_\alpha$ y
$f_\beta$ son secciones de $\sh O_S$ sobre $S_\alpha$ y
$S_\beta$, respectivamente. Por \sref{III.6.4.2},
\[
  \shTor_n^T(\sh P_\bullet|U,\sh Q_\bullet|V)
\]
es canónicamente isomorfo tanto a
$\shTor_n^{S_\alpha}(\sh P_\bullet|U,\sh Q_\bullet|V)$ como a
$\shTor_n^{S_\beta}(\sh P_\bullet|U,\sh Q_\bullet|V)$.
Combinando estos isomorfismos con las identificaciones locales de
\sref{III.6.5.2}, queda definido el $\sh O_Z$-módulo
\[
  \shTor_n^S(\sh P_\bullet,\sh Q_\bullet).
\]
Para todo abierto $U$ de $X$ y todo abierto $V$ de $Y$, la
restricción de este módulo a $U\times_SV$ es canónicamente
isomorfa a
\[
  \shTor_n^S(\sh P_\bullet|U,\sh Q_\bullet|V).
\]

Esta construcción define manifiestamente, en las categorías de
complejos acotados inferiormente de módulos cuasicoherentes, un
bi-$\partial$-funtor
\[
  \shTor_\bullet^S(\sh P_\bullet,\sh Q_\bullet)
\]
con valores en los $\sh O_Z$-módulos. En efecto, la afirmación es
local en $X$, $Y$ y $S$, por \sref{III.6.4.5} y por el hecho de
que \sref{III.6.4.4.2} es un morfismo de bi-$\partial$-funtores.

Si $\sh P_\bullet$ y $\sh Q_\bullet$ están concentrados en grado
cero, con términos respectivos $\sh F$ y $\sh G$, entonces
\[
  \shTor_0^S(\sh F,\sh G)
  =
  \sh F\otimes_S\sh G,
\]
el producto tensorial externo definido en \sref[I]{I.9.1.2}; esto
se deduce de \sref{III.6.4.1.1} y \sref[I]{I.9.1.3}.
\end{env}

\begin{env}[6.5.4]
\label{III.6.5.4}
La construcción precedente y \sref{III.6.4.6} muestran que
$\shTor_\bullet^S(\sh P_\bullet,\sh Q_\bullet)$ es el objeto al
que convergen dos funtores espectrales
$'\sh E(\sh P_\bullet,\sh Q_\bullet)$ y
$''\sh E(\sh P_\bullet,\sh Q_\bullet)$, cuyos términos $E_2$ son
\[
  '\sh E_{pq}^2
  =
  \sh H_p\bigl(\shTor_q^S(\sh P_\bullet,\sh Q_\bullet)\bigr)
  \tag{6.5.4.1}\label{III.6.5.4.1}
\]
y
\[
  ''\sh E_{pq}^2
  =
  \bigoplus_{q'+q''=q}
  \shTor_p^S
  \bigl(\sh H_{q'}(\sh P_\bullet),
        \sh H_{q''}(\sh Q_\bullet)\bigr).
  \tag{6.5.4.2}\label{III.6.5.4.2}
\]
La sucesión espectral \sref{III.6.5.4.2} es siempre regular.
Ambas sucesiones espectrales son birregulares cuando
$\sh P_\bullet$ y $\sh Q_\bullet$ están acotados inferiormente.
A continuación se describe otro caso en que son birregulares.
\end{env}

\begin{env}[6.5.5]
\label{III.6.5.5}
Sea $T$ un espacio topológico. Se dice que un haz de anillos
$\sh A$ sobre $T$ tiene \emph{dimensión cohomológica a lo sumo
$n$} si el anillo $\sh A_t$ tiene dimensión cohomológica a lo
sumo $n$ para todo $t\in T$. En este caso se dice también que el
espacio anillado $(T,\sh A)$ tiene dimensión cohomológica a lo
sumo $n$. Un haz de anillos, o un espacio anillado, tiene
\emph{dimensión cohomológica finita} si esto se cumple para algún
entero $n$.

Si los $\sh A_t$ son anillos locales noetherianos conmutativos,
\oldpage[III]{150}
decir que tienen dimensión cohomológica a lo sumo $n$ significa
que son regulares y tienen dimensión de Krull a lo sumo $n$
(\sref[0]{0.17.3.1}). En la terminología de la teoría de la
dimensión que se introducirá en el capítulo IV, un preesquema
localmente noetheriano $T$ tiene dimensión cohomológica a lo sumo
$n$ exactamente cuando es regular (\sref[0]{0.4.1.4}) y tiene
dimensión a lo sumo $n$. Equivalentemente, para todo abierto afín
$U$ de $T$, el anillo $\Gamma(U,\sh O_T)$ tiene dimensión
cohomológica a lo sumo $n$ (\sref[0]{0.17.2.6}).

En consecuencia, si $S$ es localmente noetheriano y tiene
dimensión cohomológica finita, \sref{III.6.3.2} muestra que tanto
$'\sh E(\sh P_\bullet,\sh Q_\bullet)$ como
$''\sh E(\sh P_\bullet,\sh Q_\bullet)$ son birregulares.

Finalmente, bajo el isomorfismo canónico
$X\times_SY\xrightarrow{\sim}Y\times_SX$,
$\shTor_\bullet^S(\sh P_\bullet,\sh Q_\bullet)$ corresponde
canónicamente a
$\shTor_\bullet^S(\sh Q_\bullet,\sh P_\bullet)$.
\end{env}

\begin{proposition}[6.5.6]
\label{III.6.5.6}
Sea $(\sh P_{\alpha\bullet})$ un sistema inductivo filtrante de
complejos de $\sh O_X$-módulos cuasicoherentes. Existe un
isomorfismo canónico
\[
  \varinjlim_\alpha
  \shTor_\bullet^S(\sh P_{\alpha\bullet},\sh Q_\bullet)
  \xrightarrow{\sim}
  \shTor_\bullet^S
  \left(\varinjlim_\alpha\sh P_{\alpha\bullet},\sh Q_\bullet\right).
  \tag{6.5.6.1}\label{III.6.5.6.1}
\]
\end{proposition}

\begin{proof}
La cuestión es local en $S$, $X$ e $Y$. Por tanto, podemos
suponer que $S$, $X$ e $Y$ son afines, en cuyo caso la proposición
se reduce a \sref{III.6.3.6}.
\end{proof}

\begin{env}[6.5.7]
\label{III.6.5.7}
\emph{Observaciones.}

\textnormal{(i)} Consideremos en particular el caso $S=X=Y$, con
$\sh P_\bullet$ y $\sh Q_\bullet$ dos complejos de
$\sh O_S$-módulos cuasicoherentes. Entonces los
$\shTor_n^S(\sh P_\bullet,\sh Q_\bullet)$ son
$\sh O_S$-módulos cuasicoherentes. Además, para todo punto
$z\in S$, \sref{III.6.5.6} da un isomorfismo canónico
\[
  \bigl(\shTor_n^S(\sh P_\bullet,\sh Q_\bullet)\bigr)_z
  \xrightarrow{\sim}
  \mathbf{Tor}_n^{\sh O_z}
  \bigl((\sh P_\bullet)_z,(\sh Q_\bullet)_z\bigr),
  \tag{6.5.7.1}\label{III.6.5.7.1}
\]
porque la cuestión es local y, por \sref{III.6.4.1.1}, se reduce
al caso de módulos.

\textnormal{(ii)} La definición de hipertor puede generalizarse
a dos complejos $\sh P_\bullet,\sh Q_\bullet$ de
$\sh O_X$-módulos sobre un mismo espacio anillado $(X,\sh O_X)$.
Para cada abierto $U$ de $X$, pongamos
\[
  A(U)=\Gamma(U,\sh O_X),\qquad
  P_\bullet(U)=\Gamma(U,\sh P_\bullet),\qquad
  Q_\bullet(U)=\Gamma(U,\sh Q_\bullet).
\]
Los $A(U)$-módulos
\[
  \mathbf{Tor}_n^{A(U)}\bigl(P_\bullet(U),Q_\bullet(U)\bigr)
\]
forman entonces un prehaz sobre $X$; denotemos por
$\shTor_n^X(\sh P_\bullet,\sh Q_\bullet)$ el $\sh O_X$-módulo
asociado. Cuando $X$ es un preesquema, \sref{III.6.3.12} muestra
que este módulo es canónicamente isomorfo al hipertor definido
más arriba. No desarrollaremos más esta generalización.
\end{env}

\begin{proposition}[6.5.8]
\label{III.6.5.8}
Sean $X$ e $Y$ dos $S$-preesquemas, y sea $\sh F$
(respectivamente $\sh G$) un $\sh O_X$-módulo cuasicoherente
(respectivamente un $\sh O_Y$-módulo cuasicoherente). Si
$\sh F$ o $\sh G$ es plano sobre $S$, entonces
\[
  \shTor_n^S(\sh F,\sh G)=0\qquad(n\ne0).
\]
\end{proposition}

\begin{proof}
La cuestión es local en $X$ e $Y$, por lo que podemos suponer que
$X$, $Y$ y $S$ son afines, de anillos respectivos $B$, $C$ y
$A$, y escribir $\sh F=\widetilde M$ y $\sh G=\widetilde N$,
donde $M$ es un $B$-módulo y $N$ un $C$-módulo. Supongamos, por
ejemplo, que $\sh F$ es plano sobre $S$. Esto significa que, para
todo $s\in S$, $M_s$ es un $A_s$-módulo plano
(\sref[0]{0.6.7.1}); por consiguiente, $M$ es un $A$-módulo plano
(\sref[0]{0.6.3.3}). Así,
$\Tor_n^A(M,N)=0$ para $n>0$ y todo $C$-módulo $N$
(\sref[0]{0.6.1.1}), y la conclusión se deduce de
\sref{III.6.4.1.1}.
\end{proof}

\begin{corollary}[6.5.9]
\label{III.6.5.9}
Sean $X$ e $Y$ dos $S$-preesquemas, y sea $\sh P_\bullet$
(respectivamente $\sh Q_\bullet$) un complejo
\oldpage[III]{151}
acotado
inferiormente de $\sh O_X$-módulos cuasicoherentes
(respectivamente de $\sh O_Y$-módulos cuasicoherentes).
Supongamos que cada $\sh P_i$ es plano sobre $S$. Entonces existe
un isomorfismo canónico de $\partial$-funtores en $\sh Q_\bullet$
\[
  \shTor_\bullet^S(\sh P_\bullet,\sh Q_\bullet)
  \xrightarrow{\sim}
  \sh H_\bullet(\sh P_\bullet\otimes_S\sh Q_\bullet).
  \tag{6.5.9.1}\label{III.6.5.9.1}
\]
\end{corollary}

\begin{proof}
Cuando $S$, $X$ e $Y$ son afines, esto es precisamente
\sref{III.6.3.7}; el caso general se deduce mediante los
razonamientos de \sref{III.6.5.2} y \sref{III.6.5.3}.
\end{proof}

\begin{env}[6.5.10]
\label{III.6.5.10}
\emph{Corolario.}
Supongamos que $X$ es plano sobre $S$ (\sref[0]{0.6.7.1}), que
$\sh P_\bullet$ y $\sh Q_\bullet$ están acotados inferiormente,
y que cada $\sh P_i$ es un $\sh O_X$-módulo localmente libre, no
necesariamente de tipo finito. Entonces el homomorfismo
\sref{III.6.5.9.1} es biyectivo.

En efecto, se cumple la hipótesis de \sref{III.6.5.9}: la
planitud es, por definición, una propiedad puntual en $X$, y toda
suma directa de módulos planos es plana (\sref[0]{0.6.1.2}).
\end{env}

\begin{proposition}[6.5.11]
\label{III.6.5.11}
Sean $X'$ e $Y'$ dos $S$-preesquemas, y sean
$f:X\to X'$ y $g:Y\to Y'$ dos $S$-morfismos afines. Para un
complejo $\sh P_\bullet$ de $\sh O_X$-módulos cuasicoherentes y
un complejo $\sh Q_\bullet$ de $\sh O_Y$-módulos
cuasicoherentes, existe un isomorfismo canónico funtorial
\[
  (f\times_Sg)_*
  \bigl(\shTor_\bullet^S(\sh P_\bullet,\sh Q_\bullet)\bigr)
  \xrightarrow{\sim}
  \shTor_\bullet^S(f_*\sh P_\bullet,g_*\sh Q_\bullet).
  \tag{6.5.11.1}\label{III.6.5.11.1}
\]
\end{proposition}

\begin{proof}
Como $f$ y $g$ son afines, $f_*\sh P_\bullet$ y
$g_*\sh Q_\bullet$ son complejos de módulos cuasicoherentes
(\sref[II]{II.1.2.6}). Si $Z'=X'\times_SY'$, ambos miembros de
\sref{III.6.5.11.1} son $\sh O_{Z'}$-módulos cuasicoherentes por
\sref{III.6.5.1}. Nos reducimos inmediatamente al caso en que
$S$, $X'$ e $Y'$ son afines; entonces $X$ e $Y$ también son
afines, y la afirmación se deduce de \sref{III.6.4.1.1} y
\sref[I]{I.1.6.3}.
\end{proof}

\begin{env}[6.5.12]
\label{III.6.5.12}
\emph{Observación.}
Sean $X'$ e $Y'$ dos $S$-preesquemas, con $X'$ plano sobre $S$.
Sea $X$ un subpreesquema cerrado de $X'$, y sea
$i:X\to X'$ la inmersión canónica. Sea $\sh P_\bullet$
(respectivamente $\sh Q_\bullet$) un complejo acotado
inferiormente de $\sh O_X$-módulos cuasicoherentes
(respectivamente de $\sh O_{Y'}$-módulos cuasicoherentes).
Finalmente, sea $\sh L'_{\bullet\bullet}$ una resolución de
$i_*\sh P_\bullet$ formada por $\sh O_{X'}$-módulos localmente
libres, tal que cada punto de $X'$ tenga un entorno abierto afín
$U$ para el cual $\sh L'_{j,\bullet}|U$ sea una resolución libre de
$i_*\sh P_j|U$, para todo $j$. Existe entonces un isomorfismo
canónico
\[
  (i\times_S1)_*
  \bigl(\shTor_\bullet^S(\sh P_\bullet,\sh Q_\bullet)\bigr)
  \xrightarrow{\sim}
  \sh H_\bullet(\sh L'_{\bullet\bullet}\otimes_S\sh Q_\bullet).
  \tag{6.5.12.1}\label{III.6.5.12.1}
\]

Si $S$, $X'$ e $Y'$ son afines, y si
$\sh L'_{j,\bullet}$ es una resolución libre de $i_*\sh P_j$ para
todo $j$, \sref{III.6.5.11} nos reduce al caso $X'=X$, con $X$
plano sobre $S$, y basta aplicar \sref{III.6.3.4}. En general,
\sref{III.6.5.12.1} se define localmente, y queda comprobar que
estas definiciones locales dan un isomorfismo global.

Para ello, recordemos la definición del primer isomorfismo
\sref{III.6.3.4.1}. Proviene de un isomorfismo de sucesiones
espectrales (\sref[0]{0.11.6.5} y \sref[0]{0.11.5.3}), obtenido a
su vez de un morfismo de bicomplejos
\[
  L_{\bullet\bullet}\longrightarrow L''_{\bullet\bullet},
\]
donde $L_{\bullet\bullet}$ es la resolución dada de $P_\bullet$,
y $L''_{\bullet\bullet}$ es una resolución proyectiva de
$P_\bullet$ en la categoría de complejos acotados
\oldpage[III]{152}
inferiormente
de $A$-módulos (\sref[0]{0.11.5.2.2}). La afirmación se deduce del
hecho de que \sref{III.6.3.4.1} es independiente de la resolución
proyectiva elegida $L''_{\bullet\bullet}$, pues existe una equivalencia
homotópica entre cualesquiera dos resoluciones de este tipo
(\emph{M}, V, 1.2).
\end{env}

\begin{proposition}[6.5.13]
\label{III.6.5.13}
Sean $X$ e $Y$ dos $S$-preesquemas, y supongamos que se cumple
una de las condiciones siguientes:
\begin{enumerate}
  \item[(i)] $X$ y $Z=X\times_SY$ son localmente noetherianos, y
    $X$ es plano sobre $S$;
  \item[(ii)] $S$ y $X$ son localmente noetherianos, y $Y$ es de
    tipo finito sobre $S$.
\end{enumerate}
Sea $\sh P_\bullet$ (respectivamente $\sh Q_\bullet$) un complejo
acotado inferiormente de $\sh O_X$-módulos cuasicoherentes
(respectivamente de $\sh O_Y$-módulos cuasicoherentes).
Supongamos además que, para todo $n$,
$\sh H_n(\sh P_\bullet)$ (respectivamente
$\sh H_n(\sh Q_\bullet)$) es un $\sh O_X$-módulo
(respectivamente un $\sh O_Y$-módulo) de tipo finito. Entonces
\[
  \shTor_n^S(\sh P_\bullet,\sh Q_\bullet)
\]
es un $\sh O_Z$-módulo coherente para todo $n$.
\end{proposition}

\begin{proof}
Como $\sh P_\bullet$ y $\sh Q_\bullet$ están acotados
inferiormente, la sucesión espectral
$''\sh E(\sh P_\bullet,\sh Q_\bullet)$ es birregular por
\sref{III.6.5.4}. Tanto en el caso (i) como en el caso (ii), $Z$
es localmente noetheriano, de modo que \sref[0]{0.11.1.8} reduce
la afirmación a la coherencia de los términos
$''\sh E_{pq}^2$. La hipótesis de tipo finito sobre los módulos de
homología y la fórmula \sref{III.6.5.4.2} reducen la proposición
al caso particular en que $\sh P_\bullet$ y $\sh Q_\bullet$ están
concentrados en grado cero, es decir, al corolario siguiente.
\end{proof}

\begin{corollary}[6.5.14]
\label{III.6.5.14}
Supongamos que se cumple una de las condiciones \emph{(i)} y
\emph{(ii)} de \sref{III.6.5.13}, y sea $\sh F$
(respectivamente $\sh G$) un $\sh O_X$-módulo cuasicoherente
(respectivamente un $\sh O_Y$-módulo cuasicoherente) de tipo
finito. Entonces $\shTor_n^S(\sh F,\sh G)$ es un
$\sh O_Z$-módulo coherente para todo $n$.
\end{corollary}

\begin{proof}
La cuestión es local en $X$ e $Y$, por lo que podemos suponer que
$S$, $X$ e $Y$ son afines.

\textnormal{(i)} Bajo las hipótesis de (i), $S$, $X$ y $Z$ son
noetherianos. Existe, por tanto, una resolución localmente libre
$\sh L_\bullet$ de $\sh F$, formada por $\sh O_X$-módulos de
tipo finito (\sref[I]{I.1.3.7}). Como $X$ es plano sobre $S$,
\sref{III.6.3.4} da
\[
  \shTor_n^S(\sh F,\sh G)
  =
  \sh H_n(\sh L_\bullet\otimes_S\sh G).
\]
Los módulos $\sh L_i\otimes_S\sh G$ son
$\sh O_Z$-módulos cuasicoherentes de tipo finito
(\sref[I]{I.9.1.1}), luego son coherentes. Se sigue que
$\sh H_n(\sh L_\bullet\otimes_S\sh G)$ es coherente
(\sref[0]{0.5.3.4}).

\textnormal{(ii)} Supongamos ahora que se cumple la condición
(ii). Como $A(Y)$ es cociente de una $A(S)$-álgebra de polinomios
en un número finito de indeterminadas (\sref[I]{I.6.3.3}), $Y$ es
un subpreesquema cerrado de un $S$-esquema afín $Y'$ plano y de
tipo finito sobre $S$. El esquema $Y'$ es noetheriano
(\sref[I]{I.6.3.7}), por lo que existe una resolución localmente
libre $\sh M_\bullet$ de $j_*\sh G$ por
$\sh O_{Y'}$-módulos de tipo finito, donde $j:Y\to Y'$ es la
inmersión canónica. Por \sref{III.6.5.12},
$\shTor_n^S(\sh F,\sh G)$ es la imagen inversa por $1\times j$
del $\sh O_{Z'}$-módulo
\[
  \sh H_n(\sh F\otimes_S\sh M_\bullet),
  \qquad Z'=X\times_SY'.
\]
Como en (i), los $\sh F\otimes_S\sh M_i$ son
$\sh O_{Z'}$-módulos coherentes, y \sref[0]{0.5.3.4} da de nuevo
la conclusión.
\end{proof}

\begin{env}[6.5.15]
\label{III.6.5.15}
La teoría desarrollada arriba para dos complejos
$\sh P_\bullet,\sh Q_\bullet$ de haces cuasicoherentes sobre dos
$S$-preesquemas $X$ e $Y$ se extiende sin dificultad a una
familia finita arbitraria de $S$-preesquemas $X^{(i)}$
$(1\leq i\leq m)$, cada uno provisto de un complejo
$\sh P_\bullet^{(i)}$ de $\sh O_{X^{(i)}}$-módulos
cuasicoherentes. Si
\[
  Z=X^{(1)}\times_SX^{(2)}\times_S\cdots\times_SX^{(m)},
\]
se define así un $\sh O_Z$-módulo cuasicoherente
\[
  \shTor_q^S
  \bigl(\sh P_\bullet^{(1)},\sh P_\bullet^{(2)},\ldots,
        \sh P_\bullet^{(m)}\bigr).
\]
Dejamos al lector el desarrollo de la teoría en este caso general,
y registramos únicamente, para referencia posterior, el término
$E_2$ de la segunda sucesión espectral regular que converge a este
hipertor:

\oldpage[III]{153}

\[
  ''\sh E_{pq}^2
  =
  \bigoplus_{q_1+q_2+\cdots+q_m=q}
  \shTor_p^S
  \bigl(
    \sh H_{q_1}(\sh P_\bullet^{(1)}),\ldots,
    \sh H_{q_m}(\sh P_\bullet^{(m)})
  \bigr).
  \tag{6.5.15.1}\label{III.6.5.15.1}
\]
En \sref{III.6.8} estudiaremos las sucesiones espectrales de
asociatividad a las que dan lugar estos funtores hipertor de un
número finito arbitrario de complejos.
\end{env}

\subsection{Funtores de hipertor global para complejos de módulos
cuasicoherentes y sucesiones espectrales de Künneth: caso de base afín}
\label{subsection:III.6.6}

\begin{env}[6.6.1]
\label{III.6.6.1}
Consideremos un esquema afín $S=\operatorname{Spec}(A)$ y dos $S$-esquemas
cuasicompactos $X^{(i)}$ ($i=1,2$); sea $\mathscr P_\bullet^{(i)}$ un complejo
de $\mathscr O_{X^{(i)}}$-módulos cuasicoherentes, acotado inferiormente,
cuyo diferencial es de grado $-1$ ($i=1,2$). Consideremos, por otra parte, un
cubrimiento finito
$\mathfrak U^{(i)}=(U_\alpha^{(i)})$ de $X^{(i)}$ por abiertos afines; sea
\[
  X=X^{(1)}\times_S X^{(2)},
\]
que es un $S$-esquema cuasicompacto (I, 5.5.1 y 6.6.4), y sea
$\mathfrak U=\mathfrak U^{(1)}\times_S\mathfrak U^{(2)}$ el cubrimiento de
$X$ formado por los abiertos afines
$U_\alpha^{(1)}\times_S U_\beta^{(2)}$. Para todo par de enteros
$p\leqslant0$, $q\in\mathbf Z$, el grupo de las $(-p)$-cocadenas alternadas
\[
  C^{-p}(\mathfrak U^{(i)},\mathscr P_q^{(i)})
\]
del cubrimiento $\mathfrak U^{(i)}$ con coeficientes en el haz
$\mathscr P_q^{(i)}$ (G, II, 5.1) es un $A$-módulo; para $p>0$, pondremos
$C^p(\mathfrak U^{(i)},\mathscr P_q^{(i)})=0$; queda así definido un
bicomplejo
\[
  C^\bullet(\mathfrak U^{(i)},\mathscr P_\bullet^{(i)})
\]
de $A$-módulos, cuyos dos diferenciales son de grado $-1$. De las
definiciones ($0_{\mathrm I}$, 12.1.2) resulta que el $A$-módulo de homología
\[
  \mathrm H_n\bigl(C^\bullet
  (\mathfrak U^{(i)},\mathscr P_\bullet^{(i)})\bigr)
\]
de este bicomplejo (considerado, como de costumbre, como un complejo simple
para el grado total) no es sino el $A$-módulo de hipercohomología
\[
  \mathbf H^{-n}
  (\mathfrak U^{(i)},\mathscr P^{(i)\bullet}),
\]
donde $\mathscr P^{(i)\bullet}$ es el complejo con diferencial de grado $+1$
obtenido tomando $\mathscr P_{-q}^{(i)}$ como componente de grado $q$; por
abuso de notación, lo escribiremos
$\mathbf H^{-n}(\mathfrak U^{(i)},\mathscr P_\bullet^{(i)})$. Resulta entonces
de (6.2.2) que este $A$-módulo es canónicamente isomorfo al $A$-módulo de
hipercohomología
$\mathbf H^{-n}(X^{(i)},\mathscr P^{(i)\bullet})$, que escribiremos asimismo
$\mathbf H^{-n}(X^{(i)},\mathscr P_\bullet^{(i)})$; por tanto, no depende del
cubrimiento finito $\mathfrak U^{(i)}$ elegido.
\end{env}

\begin{env}[6.6.2]
\label{III.6.6.2}
A los dos bicomplejos de $A$-módulos
\[
  L_{\bullet\bullet}^{(i)}
  =C^\bullet(\mathfrak U^{(i)},\mathscr P_\bullet^{(i)})
  \qquad(i=1,2)
\]
y al bifuntor covariante que asocia a estos dos bicomplejos su producto
tensorial $L_{\bullet\bullet}^{(1)}\otimes_A L_{\bullet\bullet}^{(2)}$ les
aplicaremos la teoría general de la hiperhomología de funtores respecto de
bicomplejos ($0_{\mathrm I}$, 11.7.4). Como las cocadenas consideradas son
alternadas y los cubrimientos $\mathfrak U^{(i)}$ son finitos, observemos que
los módulos $C^{-p}(\mathfrak U^{(i)},\mathscr P_q^{(i)})$ solo son $\ne0$
para un número finito (independiente de $q$) de valores de $p$ y, en
particular, las dos graduaciones de cada $L_{\bullet\bullet}^{(i)}$ están
acotadas inferiormente. Denotaremos por
\[
  \mathbf{Tor}_n^A(L_{\bullet\bullet}^{(1)},L_{\bullet\bullet}^{(2)})
\]
o por
\[
  \mathbf{Tor}_n^S
  (\mathfrak U^{(1)},\mathfrak U^{(2)};
   \mathscr P_\bullet^{(1)},\mathscr P_\bullet^{(2)})
\]
el $n$-ésimo módulo de hiperhomología de
$L_{\bullet\bullet}^{(1)}\otimes_A L_{\bullet\bullet}^{(2)}$, al que
llamaremos el \emph{hipertor} de índice $n$ de
$\mathscr P_\bullet^{(1)}$ y $\mathscr P_\bullet^{(2)}$, relativo a los
cubrimientos $\mathfrak U^{(1)}$ y $\mathfrak U^{(2)}$. Cuando
$\mathscr P_\bullet^{(1)}$ y $\mathscr P_\bullet^{(2)}$ están reducidos a sus
términos de grado 0, $\mathscr F^{(1)}$ y $\mathscr F^{(2)}$, escribimos
\[
  \mathbf{Tor}_n^S
  (\mathfrak U^{(1)},\mathfrak U^{(2)};
   \mathscr F^{(1)},\mathscr F^{(2)})
\]
para su hipertor. Conforme a las convenciones generales, denotaremos por
\[
  \mathbf{Tor}_n^S
  (\mathfrak U^{(1)},\mathfrak U^{(2)};
   \mathscr P_\bullet^{(1)},\mathscr P_\bullet^{(2)})
\]
el bicomplejo cuya componente de índices $(j,k)$ es
\[
  \mathbf{Tor}_n^S
  (\mathfrak U^{(1)},\mathfrak U^{(2)};
   \mathscr P_j^{(1)},\mathscr P_k^{(2)}).
\]

Como $L_{\bullet\bullet}^{(i)}$ es un funtor exacto respecto de
$\mathscr P_\bullet^{(i)}$, puesto que las intersecciones de los miembros
\oldpage[III]{154}
de $\mathfrak U^{(i)}$ son afines (I, 5.5.6 y 1.3.11),
$\mathbf{Tor}_\bullet^S
(\mathfrak U^{(1)},\mathfrak U^{(2)};
 \mathscr P_\bullet^{(1)},\mathscr P_\bullet^{(2)})$ es un
bi-$\partial$-funtor covariante en $\mathscr P_\bullet^{(1)}$ y
$\mathscr P_\bullet^{(2)}$, con valores en la categoría de los $A$-módulos
($0_{\mathrm I}$, 11.7.3). Además, se sabe ($0_{\mathrm I}$, 11.7.2) que este
bifuntor es el término límite común de seis funtores espectrales birregulares,
que denotaremos por
\[
  {}^{(t)}\mathscr E^S
  (\mathfrak U^{(1)},\mathfrak U^{(2)};
   \mathscr P_\bullet^{(1)},\mathscr P_\bullet^{(2)})
\]
o por
\[
  {}^{(t)}\mathscr E
  (\mathfrak U^{(1)},\mathfrak U^{(2)};
   \mathscr P_\bullet^{(1)},\mathscr P_\bullet^{(2)}),
\]
donde $t$ ha de sustituirse por una de las letras $a$, $b$, $a'$, $b'$, $c$,
$d$, y cuyos términos $E_2$ son los siguientes:
\[
\begin{aligned}
  {}^{(a)}\mathscr E_{pq}^2
  &=\bigoplus_{q_1+q_2=q}
    \mathbf{Tor}_p^A
    \bigl(\mathrm H_{q_1}(L_{\bullet\bullet}^{(1)}),
          \mathrm H_{q_2}(L_{\bullet\bullet}^{(2)})\bigr),\\
  {}^{(b)}\mathscr E_{pq}^2
  &=\mathrm H_p\bigl(
    \mathbf{Tor}_q^{A,\mathrm{II}}
    (L_{\bullet\bullet}^{(1)},L_{\bullet\bullet}^{(2)})\bigr),\\
  {}^{(a')}\mathscr E_{pq}^2
  &=\bigoplus_{q_1+q_2=q}
    \mathbf{Tor}_p^A
    \bigl(\mathrm H_{q_1}^{\mathrm I}(L_{\bullet\bullet}^{(1)}),
          \mathrm H_{q_2}^{\mathrm I}(L_{\bullet\bullet}^{(2)})\bigr),\\
  {}^{(b')}\mathscr E_{pq}^2
  &=\mathrm H_p\bigl(
    \mathbf{Tor}_q^A
    (L_{\bullet\bullet}^{(1)},L_{\bullet\bullet}^{(2)})\bigr),\\
  {}^{(c)}\mathscr E_{pq}^2
  &=\bigoplus_{q_1+q_2=q}
    \mathbf{Tor}_p^A
    \bigl(\mathrm H_{q_1}^{\mathrm{II}}(L_{\bullet\bullet}^{(1)}),
          \mathrm H_{q_2}^{\mathrm{II}}(L_{\bullet\bullet}^{(2)})\bigr),\\
  {}^{(d)}\mathscr E_{pq}^2
  &=\mathrm H_p\bigl(
    \mathbf{Tor}_q^{A,\mathrm I}
    (L_{\bullet\bullet}^{(1)},L_{\bullet\bullet}^{(2)})\bigr),
\end{aligned}
\]
donde las notaciones concuerdan con las de la teoría general de la
hiperhomología. A continuación explicitaremos con mayor detalle estos términos
iniciales.
\end{env}

\begin{env}[6.6.3]
\label{III.6.6.3}
\emph{Sucesiones espectrales $(a)$ y $(a')$.} Hemos visto en (6.6.1) que el
módulo de homología $\mathrm H_n(L_{\bullet\bullet}^{(i)})$ del bicomplejo
$L_{\bullet\bullet}^{(i)}$ era igual a
$\mathbf H^{-n}(X^{(i)},\mathscr P_\bullet^{(i)})$; por tanto,
\[
  {}^{(a)}\mathscr E_{pq}^2
  =\bigoplus_{q_1+q_2=q}
   \mathbf{Tor}_p^A
   \bigl(\mathbf H^{-q_1}(X^{(1)},\mathscr P_\bullet^{(1)}),
         \mathbf H^{-q_2}(X^{(2)},\mathscr P_\bullet^{(2)})\bigr).
\]
Por definición, el complejo
$\mathrm H_n^{\mathrm I}(L_{\bullet\bullet}^{(i)})$ tiene por término de grado
$k$ el módulo de homología
$\mathrm H_n(C^\bullet(\mathfrak U^{(i)},\mathscr P_k^{(i)}))$, es decir,
por definición, el módulo de cohomología
$\mathbf H^{-n}(\mathfrak U^{(i)},\mathscr P_k^{(i)})$; se sabe por (1.4.1)
que este módulo es canónicamente isomorfo a
$\mathbf H^{-n}(X^{(i)},\mathscr P_k^{(i)})$; por tanto,
\[
  {}^{(a')}\mathscr E_{pq}^2
  =\bigoplus_{q_1+q_2=q}
   \mathbf{Tor}_p^A
   \bigl(\mathbf H^{-q_1}(X^{(1)},\mathscr P_\bullet^{(1)}),
         \mathbf H^{-q_2}(X^{(2)},\mathscr P_\bullet^{(2)})\bigr).
\]
\end{env}

\begin{env}[6.6.4]
\label{III.6.6.4}
\emph{Sucesiones espectrales $(b)$ y $(b')$.} Por definición,
$\mathbf{Tor}_q^{A,\mathrm{II}}
(L_{\bullet\bullet}^{(1)},L_{\bullet\bullet}^{(2)})$ es un bicomplejo cuya
componente de bigrado $(h,k)$ es el $A$-módulo
\[
  \mathbf{Tor}_q^A
  \bigl(C^{-h}(\mathfrak U^{(1)},\mathscr P_\bullet^{(1)}),
        C^{-k}(\mathfrak U^{(2)},\mathscr P_\bullet^{(2)})\bigr).
\]
Sea $\Phi^{(i)}$ el conjunto de índices de $\mathfrak U^{(i)}$; por
definición, el complejo de módulos
$C^r(\mathfrak U^{(i)},\mathscr P_\bullet^{(i)})$ ($r\geqslant0$) es suma
directa de los complejos
$\Gamma(U_\rho^{(i)},\mathscr P_\bullet^{(i)})$, donde $U_\rho^{(i)}$ es la
intersección de los $U_\xi^{(i)}$ para $\xi\in\rho$, y $\rho$ recorre
$\mathfrak P(\Phi^{(i)})$; por tanto, el $A$-módulo
\[
  \mathbf{Tor}_q^A
  \bigl(C^{-h}(\mathfrak U^{(1)},\mathscr P_\bullet^{(1)}),
        C^{-k}(\mathfrak U^{(2)},\mathscr P_\bullet^{(2)})\bigr)
\]
es suma directa de los $A$-módulos
\[
  \mathbf{Tor}_q^A
  \bigl(\Gamma(U_\sigma^{(1)},\mathscr P_\bullet^{(1)}),
        \Gamma(U_\tau^{(2)},\mathscr P_\bullet^{(2)})\bigr),
\]
donde $\sigma$ (respectivamente, $\tau$) recorre los elementos de
$\mathfrak P(\Phi^{(1)})$ (respectivamente, $\mathfrak P(\Phi^{(2)})$) tales
que $\operatorname{Card}(\sigma)=-(h+1)$ (respectivamente,
$\operatorname{Card}(\tau)=-(k+1)$). Como $X^{(1)}$ y $X^{(2)}$ son
esquemas, los $U_\rho^{(i)}$ son afines y, por tanto, se tiene (6.4.1.1)
\[
  \mathbf{Tor}_q^A
  \bigl(\Gamma(U_\sigma^{(1)},\mathscr P_\bullet^{(1)}),
        \Gamma(U_\tau^{(2)},\mathscr P_\bullet^{(2)})\bigr)
  =\Gamma\bigl(U_\sigma^{(1)}\times_S U_\tau^{(2)},
    \mathscr{T}\!or_q^S
    (\mathscr P_\bullet^{(1)},\mathscr P_\bullet^{(2)})\bigr).
\]
\oldpage[III]{155}
Se ve, pues, que ${}^{(b)}\mathscr E_{pq}^2$ es el $(-p)$-ésimo módulo de
cohomología del complejo
\[
  L^\bullet(\Phi^{(1)},\Phi^{(2)};\mathscr S)
\]
de cocadenas bialternadas sobre $\Phi^{(1)}$ y $\Phi^{(2)}$ con valores en el
sistema de coeficientes
\[
  \mathscr S:(\sigma,\tau)\leadsto
  \Gamma\bigl(U_\sigma^{(1)}\times_S U_\tau^{(2)},
    \mathscr{T}\!or_q^S
    (\mathscr P_\bullet^{(1)},\mathscr P_\bullet^{(2)})\bigr)
\]
($0_{\mathrm I}$, 11.8.4). Se sabe entonces ($0_{\mathrm I}$, 11.8.5 y
11.8.6) que la cohomología de este complejo es la misma que la del complejo
$C^\bullet(\Phi^{(1)},\Phi^{(2)};\mathscr S)$ de todas las cocadenas sobre
$\Phi^{(1)}$ y $\Phi^{(2)}$ con valores en $\mathscr S$, y también la del
complejo $P^\bullet(\Phi^{(1)},\Phi^{(2)};\mathscr S)$, cuyos elementos son
las combinaciones lineales de los
\[
  \lambda(\sigma,\tau)\in
  \Gamma\bigl(U_\sigma^{(1)}\times_S U_\tau^{(2)},
    \mathscr{T}\!or_q^S
    (\mathscr P_\bullet^{(1)},\mathscr P_\bullet^{(2)})\bigr),
\]
donde $\sigma=(\alpha_0,\ldots,\alpha_h)$ y
$\tau=(\beta_0,\ldots,\beta_h)$ son sucesiones con el mismo número de
elementos. Pero entonces se tiene
\[
  U_\sigma^{(1)}\times_S U_\tau^{(2)}
  =(U_{\alpha_0}^{(1)}\times_S U_{\beta_0}^{(2)})\cap\cdots\cap
    (U_{\alpha_h}^{(1)}\times_S U_{\beta_h}^{(2)})
  \qquad\text{(I, 3.2.7)}.
\]
Si se denota por $\mathfrak U$ el cubrimiento de
$Z=X^{(1)}\times_S X^{(2)}$ por los abiertos afines
$U_\alpha^{(1)}\times_S U_\beta^{(2)}$, se ve finalmente, teniendo en cuenta
que $X^{(1)}\times_S X^{(2)}$ es un esquema, que se tiene, en virtud de
(1.3.1),
\[
  {}^{(b)}\mathscr E_{pq}^2
  =\mathbf H^{-p}\bigl(X^{(1)}\times_S X^{(2)},
    \mathscr{T}\!or_q^S
    (\mathscr P_\bullet^{(1)},\mathscr P_\bullet^{(2)})\bigr).
\]

En segundo lugar,
$\mathbf{Tor}_q^A(L_{\bullet\bullet}^{(1)},L_{\bullet\bullet}^{(2)})$ es un
bicomplejo cuya componente de bigrado $(h,k)$ es la suma directa de los
$A$-módulos
\[
  \mathbf{Tor}_q^A
  \bigl(C^{-h_1}(\mathfrak U^{(1)},\mathscr P_{k_1}^{(1)}),
        C^{-h_2}(\mathfrak U^{(2)},\mathscr P_{k_2}^{(2)})\bigr)
\]
tales que $h_1+h_2=h$ y $k_1+k_2=k$; explicitando los módulos
$C^r(\mathfrak U^{(i)},\mathscr P_s^{(i)})$ como antes, se ve nuevamente que
esta componente es suma directa de los $A$-módulos
\[
  \Gamma\bigl(U_\sigma^{(1)}\times_S U_\tau^{(2)},
    \mathscr{T}\!or_q^S
    (\mathscr P_{k_1}^{(1)},\mathscr P_{k_2}^{(2)})\bigr),
\]
donde $k_1+k_2=k$, y $\sigma$ (respectivamente, $\tau$) recorre los elementos
de $\mathfrak P(\Phi^{(1)})$ (respectivamente,
$\mathfrak P(\Phi^{(2)})$) tales que
$\operatorname{Card}(\sigma)+\operatorname{Card}(\tau)=-h-2$. El término
${}^{(b')}\mathscr E_{pq}^2$ que estamos calculando es el $(-p)$-ésimo módulo
de cohomología de un bicomplejo $N^{\bullet\bullet}=(N^{hk})$, donde el
complejo simple $N^{\bullet k}$ es el complejo de cocadenas bialternadas sobre
$\Phi^{(1)}$ y $\Phi^{(2)}$, con valores en el sistema de coeficientes
\[
  \mathscr S_k:(\sigma,\tau)\leadsto
  \Gamma\left(U_\sigma^{(1)}\times_S U_\tau^{(2)},
    \bigoplus_{k_1+k_2=k}
    \mathscr{T}\!or_q^S
    (\mathscr P_{-k_1}^{(1)},\mathscr P_{-k_2}^{(2)})\right),
\]
y estos sistemas de coeficientes forman un complejo $\mathscr S^\bullet$ cuyo
diferencial proviene del diferencial del complejo simple asociado al bicomplejo
$\mathscr{T}\!or_q^S
(\mathscr P_\bullet^{(1)},\mathscr P_\bullet^{(2)})$. Se sabe que la
cohomología de $N^{\bullet\bullet}$ es la misma que la del bicomplejo
$C^\bullet(\Phi^{(1)},\Phi^{(2)};\mathscr S^\bullet)$
($0_{\mathrm I}$, 11.8.9), y también la misma que la del bicomplejo
$P^\bullet(\Phi^{(1)},\Phi^{(2)};\mathscr S^\bullet)$, donde los elementos de
bigrado $(h,k)$ son las combinaciones lineales de los
\[
  \lambda(\sigma,\tau)\in
  \Gamma\left(U_\sigma^{(1)}\times_S U_\tau^{(2)},
    \bigoplus_{k_1+k_2=k}
    \mathscr{T}\!or_q^S
    (\mathscr P_{-k_1}^{(1)},\mathscr P_{-k_2}^{(2)})\right),
\]
donde $\sigma=(\alpha_0,\ldots,\alpha_h)$ y
$\tau=(\beta_0,\ldots,\beta_h)$ son sucesiones con el mismo número de
elementos ($0_{\mathrm I}$, 11.8.10). Se ve entonces, como antes, que
${}^{(b')}\mathscr E_{pq}^2$ es el $(-p)$-ésimo módulo de cohomología del
bicomplejo $C^\bullet(\mathfrak U,\mathscr Q^\bullet)$, donde
$\mathscr Q^\bullet$ es el complejo simple asociado al bicomplejo
$\mathscr{T}\!or_q^S
(\mathscr P_\bullet^{(1)},\mathscr P_\bullet^{(2)})$ de
$\mathscr O_Z$-módulos. Con las convenciones adoptadas en (6.6.1), se tiene,
por tanto,
\[
  {}^{(b')}\mathscr E_{pq}^2
  =\mathbf H^{-p}\bigl(X^{(1)}\times_S X^{(2)},
    \mathscr{T}\!or_q^S
    (\mathscr P_\bullet^{(1)},\mathscr P_\bullet^{(2)})\bigr).
\]
\end{env}

\oldpage[III]{156}
\begin{env}[6.6.5]
\label{III.6.6.5}
\emph{Sucesiones espectrales $(c)$ y $(d)$.} Por definición, el complejo
$\mathrm H_n^{\mathrm{II}}(L_{\bullet\bullet}^{(i)})$ tiene por término de
grado $h$ el $A$-módulo
$C^{-h}(\mathfrak U^{(i)},\mathscr H_n(\mathscr P_\bullet^{(i)}))$, en virtud
de la exactitud del funtor $C^{-h}$. Por tanto, por definición del hipertor de
dos módulos relativo a dos cubrimientos (6.6.2),
\[
  {}^{(c)}\mathscr E_{pq}^2
  =\bigoplus_{q_1+q_2=q}
   \mathbf{Tor}_p^S
   \bigl(\mathfrak U^{(1)},\mathfrak U^{(2)};
         \mathscr H_{q_1}(\mathscr P_\bullet^{(1)}),
         \mathscr H_{q_2}(\mathscr P_\bullet^{(2)})\bigr).
\]
Por último, por definición,
$\mathbf{Tor}_q^{A,\mathrm I}
(L_{\bullet\bullet}^{(1)},L_{\bullet\bullet}^{(2)})$ es un bicomplejo cuya
componente de bigrado $(h,k)$ es el $A$-módulo
$\mathbf{Tor}_q^S
(\mathfrak U^{(1)},\mathfrak U^{(2)};
 \mathscr P_h^{(1)},\mathscr P_k^{(2)})$. Por tanto,
\[
  {}^{(d)}\mathscr E_{pq}^2
  =\mathrm H_p\bigl(
    \mathbf{Tor}_q^S
    (\mathfrak U^{(1)},\mathfrak U^{(2)};
     \mathscr P_\bullet^{(1)},\mathscr P_\bullet^{(2)})\bigr).
\]
\end{env}

\begin{env}[6.6.6]
\label{III.6.6.6}
La teoría de la hiperhomología de los funtores de bicomplejos
($0_{\mathrm I}$, 11.7.3) muestra, como en (6.3.4), que, para toda resolución
\emph{plana} de Cartan--Eilenberg $M_{\bullet\bullet}^{(i)}$ de
$L_{\bullet\bullet}^{(i)}$ (en la categoría de los complejos de módulos
acotados inferiormente) ($i=1,2$), se tienen isomorfismos canónicos de
bi-$\partial$-funtores
\begin{env}[6.6.6.1]
\label{III.6.6.6.1}
\[
\begin{aligned}
  \mathbf{Tor}_\bullet^S
  (\mathfrak U^{(1)},\mathfrak U^{(2)};
   \mathscr P_\bullet^{(1)},\mathscr P_\bullet^{(2)})
  &\xrightarrow{\sim}
   \mathrm H_\bullet
   (M_{\bullet\bullet}^{(1)}\otimes_A M_{\bullet\bullet}^{(2)})\\
  &\xrightarrow{\sim}
   \mathrm H_\bullet
   (M_{\bullet\bullet}^{(1)}\otimes_A L_{\bullet\bullet}^{(2)})
   \xrightarrow{\sim}
   \mathrm H_\bullet
   (L_{\bullet\bullet}^{(1)}\otimes_A M_{\bullet\bullet}^{(2)}).
\end{aligned}
\tag{6.6.6.1}
\]
\end{env}
\end{env}

\begin{env}[6.6.7]
\label{III.6.6.7}
Vamos a mostrar ahora que el hipertor global definido en (6.6.2), así como
las seis sucesiones espectrales correspondientes, no dependen de los
cubrimientos abiertos afines finitos $\mathfrak U^{(i)}$ que se utilizaron
para definirlos (salvo isomorfismos canónicos). Bastará para ello mostrar que,
si $\mathfrak V^{(i)}$ son otros dos cubrimientos de la misma naturaleza,
tales que $\mathfrak V^{(i)}$ sea \emph{más fino} que $\mathfrak U^{(i)}$
para $i=1,2$, entonces se tienen \emph{isomorfismos canónicos} de funtores
espectrales
\begin{env}[6.6.7, $t$]
\label{III.6.6.7.t}
\[
  {}^{(t)}\mathscr E
  (\mathfrak U^{(1)},\mathfrak U^{(2)};
   \mathscr P_\bullet^{(1)},\mathscr P_\bullet^{(2)})
  \xrightarrow{\sim}
  {}^{(t)}\mathscr E
  (\mathfrak V^{(1)},\mathfrak V^{(2)};
   \mathscr P_\bullet^{(1)},\mathscr P_\bullet^{(2)})
  \tag{6.6.7, $t$}\label{III.6.6.7.t.eq}
\]
\end{env}
donde $t$ se sustituye por $a$, $b$, $a'$, $b'$, $c$ o $d$.

Ahora bien, para $i=1,2$ se tienen homomorfismos de bicomplejos
\[
  C^\bullet(\mathfrak U^{(i)},\mathscr P_\bullet^{(i)})
  \longrightarrow
  C^\bullet(\mathfrak V^{(i)},\mathscr P_\bullet^{(i)})
\]
bien definidos salvo homotopía (G, II, 5.7.1); de ello resultan ya
homomorfismos (6.6.7, $t$) definidos canónicamente y compatibles con los
operadores de borde en los términos límite ($0_{\mathrm I}$, 11.3.2). Además,
el cálculo de los términos $E_2$ de las sucesiones espectrales $(a)$, $(b)$,
$(a')$, $(b')$ muestra que, para estas sucesiones espectrales, el homomorfismo
(6.6.7, $t$) es un isomorfismo sobre los términos $E_2$; como estas sucesiones
espectrales son birregulares, se ve que (6.6.7, $t$) es un isomorfismo para
estos funtores espectrales y, por consiguiente, induce un \emph{isomorfismo de
bi-$\partial$-funtores} en su término límite común
($0_{\mathrm I}$, 11.1.5).

En particular, para $\mathscr O_{X^{(i)}}$-módulos cuasicoherentes
$\mathscr F^{(i)}$ ($i=1,2$), el homomorfismo canónico
\[
  \mathbf{Tor}_\bullet^S
  (\mathfrak U^{(1)},\mathfrak U^{(2)};
   \mathscr F^{(1)},\mathscr F^{(2)})
  \longrightarrow
  \mathbf{Tor}_\bullet^S
  (\mathfrak V^{(1)},\mathfrak V^{(2)};
   \mathscr F^{(1)},\mathscr F^{(2)})
\]
es biyectivo; visto el cálculo de (6.6.5), se ve que (6.6.7, $t$) es también
un isomorfismo de los términos $E_2$ para $t=c$ y $t=d$. Se concluye como
antes que (6.6.7, $t$) es también un isomorfismo de sucesiones espectrales
para $t=c$ y $t=d$.
\oldpage[III]{157}
Se puede considerar que los isomorfismos (6.6.7, $t$) definen sistemas
inductivos de funtores espectrales sobre el conjunto filtrante de pares
$(\mathfrak U^{(1)},\mathfrak U^{(2)})$ de cubrimientos abiertos afines
finitos de $X^{(1)}$ y $X^{(2)}$. Denotaremos por
\[
  {}^{(t)}\mathscr E
  (X^{(1)},X^{(2)};
   \mathscr P_\bullet^{(1)},\mathscr P_\bullet^{(2)})
  \quad\text{o}\quad
  {}^{(t)}\mathscr E^S
  (X^{(1)},X^{(2)};
   \mathscr P_\bullet^{(1)},\mathscr P_\bullet^{(2)})
\]
el límite inductivo de este sistema, y por
$\mathbf{Tor}_\bullet^S
(X^{(1)},X^{(2)};\mathscr P_\bullet^{(1)},\mathscr P_\bullet^{(2)})$
el término límite de este funtor espectral, al que llamaremos el
\emph{hipertor global} de los dos complejos $\mathscr P_\bullet^{(1)}$ y
$\mathscr P_\bullet^{(2)}$; si $\mathscr P_\bullet^{(1)}$ y
$\mathscr P_\bullet^{(2)}$ están reducidos a sus términos de grado 0,
$\mathscr F^{(1)}$ y $\mathscr F^{(2)}$, escribiremos
\[
  \mathbf{Tor}_\bullet^S
  (X^{(1)},X^{(2)};\mathscr F^{(1)},\mathscr F^{(2)}),
\]
y, conforme a las convenciones generales,
\[
  \mathbf{Tor}_q^S
  (X^{(1)},X^{(2)};\mathscr P_\bullet^{(1)},\mathscr P_\bullet^{(2)})
\]
será, por tanto, el bicomplejo cuyas componentes son los módulos
\[
  \mathbf{Tor}_q^S
  (X^{(1)},X^{(2)};\mathscr P_h^{(1)},\mathscr P_k^{(2)}).
\]
\end{env}

\begin{env}[6.6.8]
\label{III.6.6.8}
Bajo las hipótesis de (6.6.1), consideremos ahora dos $S$-morfismos
$f_i:X^{(i)}\to Y^{(i)}$, donde
$Y^{(i)}=\operatorname{Spec}(B_i)$ es un $S$-esquema \emph{afín}, siendo
$B_i$ una $A$-álgebra ($i=1,2$); esto define, pues, un $A$-homomorfismo
$B_i\to\Gamma(X^{(i)},\mathscr O_{X^{(i)}})$ (I, 2.2.4) y, por consiguiente,
cada uno de los $L_{\bullet\bullet}^{(i)}$ definidos en (6.6.2) es un
bicomplejo de $B_i$-módulos; se concluye que
$L_{\bullet\bullet}^{(1)}\otimes_A L_{\bullet\bullet}^{(2)}$ es un
complejo cuádruple de $(B_1\otimes_A B_2)$-módulos, y sus seis funtores
espectrales de hiperhomología pueden considerarse, por tanto, con valores en
la categoría de las sucesiones espectrales de
$(B_1\otimes_A B_2)$-módulos. Si se pone
$Y=Y^{(1)}\times_S Y^{(2)}=\operatorname{Spec}(B_1\otimes_A B_2)$, se pueden
considerar los $\mathscr O_Y$-módulos cuasicoherentes asociados a estos
módulos (I, 1.3.4); denotaremos por
${}^{(t)}\mathscr E
(f_1,f_2;\mathscr P_\bullet^{(1)},\mathscr P_\bullet^{(2)})$
(para $t=a$, $b$, $a'$, $b'$, $c$ o $d$) las seis sucesiones espectrales
\[
  \bigl({}^{(t)}\mathscr E
  (X^{(1)},X^{(2)};
   \mathscr P_\bullet^{(1)},\mathscr P_\bullet^{(2)})\bigr)^{\sim}
\]
de $\mathscr O_Y$-módulos, y por
$\mathfrak{Tor}_\bullet^S
(f_1,f_2;\mathscr P_\bullet^{(1)},\mathscr P_\bullet^{(2)})$ su término
límite común
\[
  \bigl(\mathbf{Tor}_\bullet^S
  (X^{(1)},X^{(2)};
   \mathscr P_\bullet^{(1)},\mathscr P_\bullet^{(2)})\bigr)^{\sim}.
\]
Lo denotaremos por
$\mathfrak{Tor}_\bullet^S(f_1,f_2;\mathscr F^{(1)},\mathscr F^{(2)})$
cuando $\mathscr P_\bullet^{(i)}$ esté reducido a su término de grado 0,
$\mathscr F^{(i)}$ ($i=1,2$).
\end{env}

\subsection{Funtores de hipertor global de complejos de módulos
cuasicoherentes y sucesiones espectrales de Künneth: caso general}
\label{subsection:III.6.7}

\begin{env}[6.7.1]
\label{III.6.7.1}
Vamos a generalizar ahora las definiciones de (6.6.8) al caso en que $S$ es
un preesquema cualquiera, $X^{(i)}$ e $Y^{(i)}$ son $S$-preesquemas y
$f_i:X^{(i)}\to Y^{(i)}$ son morfismos \emph{separados y cuasicompactos}. Se
trata entonces, para todo par de complejos $\mathscr P_\bullet^{(i)}$ de
$\mathscr O_{X^{(i)}}$-módulos cuasicoherentes, acotados inferiormente
($i=1,2$), de definir para todo $n$ un $\mathscr O_Y$-módulo cuasicoherente
$\mathfrak{Tor}_n^S
(f_1,f_2;\mathscr P_\bullet^{(1)},\mathscr P_\bullet^{(2)})$, así como seis
funtores espectrales, que se reduzcan a las definiciones de (6.6.8) cuando
$S$, $Y^{(1)}$ e $Y^{(2)}$ son afines (se pone
$Y=Y^{(1)}\times_S Y^{(2)}$).

Supongamos primero que $S=\operatorname{Spec}(A)$ es \emph{afín}, pero que
$Y^{(1)}$ e $Y^{(2)}$ son cualesquiera; sea $W^{(i)}$ un abierto afín de
$Y^{(i)}$; entonces $f_i^{-1}(W^{(i)})$ es un $S$-esquema cuasicompacto y
$W=W^{(1)}\times_S W^{(2)}$ es un abierto afín de $Y$; sea
$f'_i:f_i^{-1}(W^{(i)})\to W^{(i)}$ la restricción de $f_i$, y sea
$\mathscr P'_\bullet{}^{(i)}$ la restricción
$\mathscr P_\bullet^{(i)}|f_i^{-1}(W^{(i)})$ ($i=1,2$). Se tienen entonces,
por (6.6.8), las sucesiones espectrales
${}^{(t)}\mathscr E
(f'_1,f'_2;\mathscr P'_\bullet{}^{(1)},\mathscr P'_\bullet{}^{(2)})$
de $(\mathscr O_Y|W)$-módulos cuasicoherentes, y hay que verificar que
satisfacen las condiciones de pegado ($0_{\mathrm I}$, 3.3.1).
Nos reducimos inmediatamente al caso en que
$Y^{(i)}=\operatorname{Spec}(B_i)$ es afín y $W^{(i)}=D(g_i)$, donde
$g_i\in B_i$, de modo que
$W=D(g_1\otimes g_2)$
\oldpage[III]{158}
en $Y=\operatorname{Spec}(B_1\otimes_A B_2)$ (II, 4.3.2.1); si
$X'^{(i)}=f_i^{-1}(W^{(i)})$, se trata de establecer un isomorfismo canónico
de funtores espectrales
\begin{env}[6.7.1.1]
\label{III.6.7.1.1}
\[
  {}^{(t)}\mathscr E
  (X'^{(1)},X'^{(2)};
   \mathscr P'_\bullet{}^{(1)},\mathscr P'_\bullet{}^{(2)})
  \xrightarrow{\sim}
  {}^{(t)}\mathscr E
  (X^{(1)},X^{(2)};
   \mathscr P_\bullet^{(1)},\mathscr P_\bullet^{(2)})
  \otimes_B B_g
  \tag{6.7.1.1}\label{III.6.7.1.1.eq}
\]
\end{env}

donde se ha puesto $B=B_1\otimes_A B_2$ y $g=g_1\otimes g_2$. Para ello,
partamos de cubrimientos abiertos afines finitos $\mathfrak U^{(i)}$ de
$X^{(i)}$ ($i=1,2$), y sea $\mathfrak U'^{(i)}$ la traza de
$\mathfrak U^{(i)}$ sobre $X'^{(i)}$, que sigue estando formada por abiertos
afines (I, 5.5.10); más precisamente, se tiene
\[
  C^\bullet(\mathfrak U'^{(i)},\mathscr P'_\bullet{}^{(i)})
  =C^\bullet(\mathfrak U^{(i)},\mathscr P_\bullet^{(i)})
   \otimes_{B_i}(B_i)_{g_i}.
\]
Si se pone
$L'_{\bullet\bullet}{}^{(i)}
=C^\bullet(\mathfrak U'^{(i)},\mathscr P'_\bullet{}^{(i)})$, se tiene, por
tanto,
\[
  L'_{\bullet\bullet}{}^{(1)}\otimes_A L'_{\bullet\bullet}{}^{(2)}
  =\bigl(L_{\bullet\bullet}^{(1)}\otimes_{B_1}(B_1)_{g_1}\bigr)
   \otimes_A
   \bigl(L_{\bullet\bullet}^{(2)}\otimes_{B_2}(B_2)_{g_2}\bigr);
\]
como $B_g=(B_1)_{g_1}\otimes_A(B_2)_{g_2}$ salvo isomorfismo canónico, se
tiene, también salvo isomorfismo canónico,
\[
  L'_{\bullet\bullet}{}^{(1)}\otimes_A L'_{\bullet\bullet}{}^{(2)}
  =(L_{\bullet\bullet}^{(1)}\otimes_A L_{\bullet\bullet}^{(2)})
    \otimes_B B_g.
\]
Si $M_{\bullet\bullet}^{(i)}$ es una resolución proyectiva de
Cartan--Eilenberg de $L_{\bullet\bullet}^{(i)}$, que puede suponerse formada
por $B_i$-módulos, del hecho de que $(B_i)_{g_i}$ es \emph{plano} sobre $B_i$
se deduce que
$M'_{\bullet\bullet}{}^{(i)}
=M_{\bullet\bullet}^{(i)}\otimes_{B_i}(B_i)_{g_i}$ es una resolución
proyectiva de Cartan--Eilenberg del bicomplejo
$L'_{\bullet\bullet}{}^{(i)}$; además, se tiene
\[
  M'_{\bullet\bullet}{}^{(1)}\otimes_A M'_{\bullet\bullet}{}^{(2)}
  =(M_{\bullet\bullet}^{(1)}\otimes_A M_{\bullet\bullet}^{(2)})
    \otimes_B B_g.
\]
El isomorfismo buscado (6.7.1.1) se deduce entonces inmediatamente de las
definiciones de la hiperhomología de un bicomplejo y de la exactitud del
funtor $G\otimes_B B_g$ respecto del $B$-módulo $G$.
\end{env}

\begin{env}[6.7.2]
\label{III.6.7.2}
Supongamos ahora que $S$ es cualquiera, y sean
$u_i:Y^{(i)}\to S$ los morfismos estructurales ($i=1,2$). Sea $(S_\alpha)$ un
cubrimiento abierto afín de $S$, pongamos
$Y_\alpha^{(i)}=u_i^{-1}(S_\alpha)$,
$X_\alpha^{(i)}=f_i^{-1}(Y_\alpha^{(i)})$, y sea
$f_{i\alpha}:X_\alpha^{(i)}\to Y_\alpha^{(i)}$ la restricción de $f_i$, que
es un morfismo \emph{separado y cuasicompacto}. Los
$Y_\alpha=Y_\alpha^{(1)}\times_{S_\alpha}Y_\alpha^{(2)}$ forman un
cubrimiento abierto de $Y$, y sobre cada $Y_\alpha$ quedan definidos, por
(6.7.1), los funtores espectrales
\[
  {}^{(t)}\mathscr E^{S_\alpha}
  \bigl(f_{1\alpha},f_{2\alpha};
    \mathscr P_\bullet^{(1)}|X_\alpha^{(1)},
    \mathscr P_\bullet^{(2)}|X_\alpha^{(2)}\bigr);
\]
hay que demostrar de nuevo que estos funtores verifican las condiciones de
pegado.

Nos reducimos inmediatamente a la situación siguiente:
$S=\operatorname{Spec}(A)$ es afín, $S'=D(h)$, donde $h\in A$, y
$u_i(Y^{(i)})\subset S'$; puede suponerse además que
$Y^{(i)}=\operatorname{Spec}(B_i)$ es afín; se trata de definir isomorfismos
canónicos
\begin{env}[6.7.2.1]
\label{III.6.7.2.1}
\[
  {}^{(t)}\mathscr E^S
  (X^{(1)},X^{(2)};
   \mathscr P_\bullet^{(1)},\mathscr P_\bullet^{(2)})
  \xrightarrow{\sim}
  {}^{(t)}\mathscr E^{S'}
  (X^{(1)},X^{(2)};
   \mathscr P_\bullet^{(1)},\mathscr P_\bullet^{(2)}).
  \tag{6.7.2.1}\label{III.6.7.2.1.eq}
\]
\end{env}
Ahora bien, con las notaciones de (6.6.2), los
$L_{\bullet\bullet}^{(i)}$ están formados por $A_h$-módulos, y se tiene
$L_{\bullet\bullet}^{(1)}\otimes_{A_h}L_{\bullet\bullet}^{(2)}
=L_{\bullet\bullet}^{(1)}\otimes_A L_{\bullet\bullet}^{(2)}$ salvo
isomorfismo canónico; como puede tomarse una resolución proyectiva de
Cartan--Eilenberg $M_{\bullet\bullet}^{(i)}$ de
$L_{\bullet\bullet}^{(i)}$ formada por $A_h$-módulos, esto proporciona
inmediatamente el isomorfismo canónico buscado.
\end{env}

Hemos demostrado, en resumen, el siguiente

\begin{theorem}[6.7.3]
\label{III.6.7.3}
Sean $S$ un preesquema, $f_i:X^{(i)}\to Y^{(i)}$ un $S$-morfismo separado y
cuasicompacto de $S$-preesquemas y $\mathscr P_\bullet^{(i)}$ un complejo
acotado inferiormente de $\mathscr O_{X^{(i)}}$-módulos cuasicoherentes
($i=1,2$); se pone $Y=Y^{(1)}\times_S Y^{(2)}$. Existe un
bi-$\partial$-funtor
  $\mathfrak{Tor}_\bullet^S
(f_1,f_2;\mathscr P_\bullet^{(1)},\mathscr P_\bullet^{(2)})$
con valores en la categoría de los $\mathscr O_Y$-módulos cuasicoherentes,
tal que, si $V^{(i)}$ es un abierto afín de $Y^{(i)}$ ($i=1,2$) y
$V=V^{(1)}\times_S V^{(2)}$, se tiene
\[
  \mathfrak{Tor}_\bullet^S
  (f_1,f_2;\mathscr P_\bullet^{(1)},\mathscr P_\bullet^{(2)})|V
  =\Bigl(\mathbf{Tor}_\bullet^S
    \bigl(f_1^{-1}(V^{(1)}),f_2^{-1}(V^{(2)});
      \mathscr P_\bullet^{(1)}|f_1^{-1}(V^{(1)}),
      \mathscr P_\bullet^{(2)}|f_2^{-1}(V^{(2)})\bigr)\Bigr)^{\sim}.
\]
\oldpage[III]{159}
Este bifuntor es el término límite de seis funtores espectrales
birregulares
\[
  {}^{(t)}\mathscr E
  (f_1,f_2;\mathscr P_\bullet^{(1)},\mathscr P_\bullet^{(2)})
  \qquad(t=a,b,a',b',c,d)
\]
cuyos términos $E_2$ vienen dados por
\[
\begin{aligned}
  {}^{(a)}\mathscr E_{pq}^2
  &=\bigoplus_{q_1+q_2=q}
    \mathscr{T}\!or_p^S
    \bigl(\mathscr H^{-q_1}(f_1,\mathscr P_\bullet^{(1)}),
          \mathscr H^{-q_2}(f_2,\mathscr P_\bullet^{(2)})\bigr),\\
  {}^{(b)}\mathscr E_{pq}^2
  &=\mathscr H^{-p}\bigl(f_1\times_S f_2,
    \mathfrak{Tor}_q^S
    (\mathscr P_\bullet^{(1)},\mathscr P_\bullet^{(2)})\bigr),\\
  {}^{(a')}\mathscr E_{pq}^2
  &=\bigoplus_{q_1+q_2=q}
    \mathfrak{Tor}_p^S
    \bigl(\mathscr H^{-q_1}(f_1,\mathscr P_\bullet^{(1)}),
          \mathscr H^{-q_2}(f_2,\mathscr P_\bullet^{(2)})\bigr),\\
  {}^{(b')}\mathscr E_{pq}^2
  &=\mathscr H^{-p}\bigl(f_1\times_S f_2,
    \mathscr{T}\!or_q^S
    (\mathscr P_\bullet^{(1)},\mathscr P_\bullet^{(2)})\bigr),\\
  {}^{(c)}\mathscr E_{pq}^2
  &=\bigoplus_{q_1+q_2=q}
    \mathfrak{Tor}_p^S
    \bigl(f_1,f_2;
      \mathscr H_{q_1}(\mathscr P_\bullet^{(1)}),
      \mathscr H_{q_2}(\mathscr P_\bullet^{(2)})\bigr),\\
  {}^{(d)}\mathscr E_{pq}^2
  &=\mathscr H_p\bigl(\mathfrak{Tor}_q^S
    (f_1,f_2;\mathscr P_\bullet^{(1)},\mathscr P_\bullet^{(2)})\bigr).
\end{aligned}
\]
Se dice que las sucesiones espectrales $(a)$ y $(b)$ son las
\emph{sucesiones espectrales de Künneth}.

Obsérvese que las sucesiones espectrales $(a)$ y $(a')$ (respectivamente,
$(b)$ y $(b')$) son idénticas cuando $\mathscr P_\bullet^{(1)}$ y
$\mathscr P_\bullet^{(2)}$ se reducen a sus términos de grado 0; en este caso,
las sucesiones $(c)$ y $(d)$ son degeneradas y, por tanto, no presentan
interés.
\end{theorem}

\begin{remark}[6.7.4]
\label{III.6.7.4}
Los hipertor globales que acabamos de definir comprenden como casos
particulares tanto los \emph{módulos de hipercohomología} definidos en
(6.2.1) como los \emph{hipertor locales} definidos en (6.5.3). Demostremos que
para todo morfismo $f:X\to Y$ \emph{cuasicompacto y separado} y todo complejo
$\mathscr P_\bullet$ de $\mathscr O_X$-módulos cuasicoherentes, acotado
inferiormente, se tiene un isomorfismo canónico de $\partial$-funtores en
$\mathscr P_\bullet$
\begin{env}[6.7.4.1]
\label{III.6.7.4.1}
\[
  \mathfrak{Tor}_n^Y(f,1_Y;\mathscr P_\bullet,\mathscr O_Y)
  \xrightarrow{\sim}
  \mathscr H^{-n}(f,\mathscr P_\bullet)
  \qquad\text{(para todo $n\in\mathbf Z$)}.
  \tag{6.7.4.1}\label{III.6.7.4.1.eq}
\]
\end{env}
En efecto, los métodos de pegado de (6.7.2) reducen inmediatamente al caso en
que $Y$ es afín; entonces, en virtud de (6.2.2), pueden calcularse los dos
miembros de (6.7.4.1) mediante un mismo cubrimiento finito $\mathfrak U$ de
$X$ por abiertos afines y, para el primer miembro, mediante el cubrimiento de
$Y$ formado por el propio $Y$; con las notaciones de (6.6.2), el bicomplejo
$L_{\bullet\bullet}^{(2)}$ se reduce entonces a su término de grados $(0,0)$,
igual a $A$, y la conclusión se deduce de ($0_{\mathrm I}$, 11.7.5). Para una
generalización de este resultado, véase (6.7.7); pero obsérvese que, cuando en
el primer miembro de (6.7.4.1) se reemplaza $\mathscr O_Y$ por un
$\mathscr O_Y$-módulo cuasicoherente \emph{cualquiera} $\mathscr F$, ya no
se tiene en general un isomorfismo con
$\mathscr H^{-n}(f,\mathscr P_\bullet\otimes_{\mathscr O_Y}\mathscr F)$,
aunque, en el cálculo precedente, el bicomplejo
$L_{\bullet\bullet}^{(1)}\otimes_A L_{\bullet\bullet}^{(2)}$ todavía se
identifica con el bicomplejo
$C^\bullet(\mathfrak U,\mathscr P_\bullet\otimes_Y\mathscr F)$.

Por otra parte, se tiene un isomorfismo canónico de
bi-$\partial$-funtores
\begin{env}[6.7.4.2]
\label{III.6.7.4.2}
\[
  \mathfrak{Tor}_\bullet^S
  (1_{X^{(1)}},1_{X^{(2)}};
   \mathscr P_\bullet^{(1)},\mathscr P_\bullet^{(2)})
  \xrightarrow{\sim}
  \mathscr{T}\!or_\bullet^S
  (\mathscr P_\bullet^{(1)},\mathscr P_\bullet^{(2)}).
  \tag{6.7.4.2}\label{III.6.7.4.2.eq}
\]
\end{env}
En efecto, por (6.7.1) y (6.7.2), nos reducimos de nuevo al caso en que $S$ y
los $X^{(i)}$ son afines; en el cálculo del primer miembro de (6.7.4.2), puede
tomarse entonces como
\oldpage[III]{160}
cubrimiento $\mathfrak U^{(i)}$ la familia reducida al
único elemento $X^{(i)}$, de modo que, con las notaciones de (6.6.2),
$L_{\bullet\bullet}^{(i)}$ se reduce a
$\Gamma(X^{(i)},\mathscr P_\bullet^{(i)})$ (considerado como un bicomplejo
cuyos términos de primer grado $\ne0$ son nulos), y la igualdad de los dos
miembros de (6.7.4.2) se deduce de (6.4.1.1) y (6.3.1).
\end{remark}

\begin{proposition}[6.7.5]
\label{III.6.7.5}
Sea $u:\mathscr P_\bullet^{(1)}\to\mathscr Q_\bullet^{(1)}$ un
homomorfismo de complejos acotados inferiormente de
$\mathscr O_{X^{(1)}}$-módulos cuasicoherentes tal que el homomorfismo
\[
  \mathscr H_\bullet(u):
  \mathscr H_\bullet(\mathscr P_\bullet^{(1)})
  \longrightarrow
  \mathscr H_\bullet(\mathscr Q_\bullet^{(1)})
\]
inducido por $u$ sea un isomorfismo. Entonces los homomorfismos
\[
  {}^{(t)}\mathscr E
  (f_1,f_2;\mathscr P_\bullet^{(1)},\mathscr P_\bullet^{(2)})
  \longrightarrow
  {}^{(t)}\mathscr E
  (f_1,f_2;\mathscr Q_\bullet^{(1)},\mathscr P_\bullet^{(2)})
\]
inducidos por $u$ son isomorfismos para $t=a$, $t=b$ y $t=c$.
\end{proposition}

\begin{proof}
La afirmación relativa a la sucesión espectral $(c)$ se deduce de que esta
sucesión es birregular y de que, por hipótesis, el homomorfismo considerado es
un isomorfismo sobre los términos $E_2$ ($0_{\mathrm I}$, 11.1.5). Esto
demuestra ya que
\[
  \mathfrak{Tor}_\bullet^S
  (f_1,f_2;\mathscr P_\bullet^{(1)},\mathscr P_\bullet^{(2)})
  \longrightarrow
  \mathfrak{Tor}_\bullet^S
  (f_1,f_2;\mathscr Q_\bullet^{(1)},\mathscr P_\bullet^{(2)})
\]
es un isomorfismo. Aplicando las relaciones (6.7.4.1) y (6.7.4.2), se ve
primero que los homomorfismos
\[
  \mathscr H^{-n}(f_1,\mathscr P_\bullet^{(1)})
  \longrightarrow
  \mathscr H^{-n}(f_1,\mathscr Q_\bullet^{(1)})
  \quad\text{y}\quad
  \mathscr{T}\!or_n^S
  (\mathscr P_\bullet^{(1)},\mathscr P_\bullet^{(2)})
  \longrightarrow
  \mathscr{T}\!or_n^S
  (\mathscr Q_\bullet^{(1)},\mathscr P_\bullet^{(2)})
\]
inducidos por $u$ son isomorfismos. La afirmación relativa a las sucesiones
$(a)$ y $(b)$ se deduce entonces de que estas sucesiones son birregulares
(6.7.3) y de que los homomorfismos considerados son biyectivos sobre los
términos $E_2$ ($0_{\mathrm I}$, 11.1.5).

Obsérvese, por otra parte, que, si
$u:\mathscr P_\bullet^{(1)}\to\mathscr Q_\bullet^{(1)}$ es una
\emph{equivalencia homotópica}, se deducen de ella isomorfismos canónicos
\[
  {}^{(t)}\mathscr E
  (f_1,f_2;\mathscr P_\bullet^{(1)},\mathscr P_\bullet^{(2)})
  \longrightarrow
  {}^{(t)}\mathscr E
  (f_1,f_2;\mathscr Q_\bullet^{(1)},\mathscr P_\bullet^{(2)})
\]
para las seis sucesiones espectrales. En efecto, si $S$ y los $Y^{(i)}$ son
afines, se deduce de $u$ una equivalencia homotópica de bicomplejos
\[
  C^\bullet(\mathfrak U^{(1)},\mathscr P_\bullet^{(1)})
  \longrightarrow
  C^\bullet(\mathfrak U^{(1)},\mathscr Q_\bullet^{(1)}),
\]
y la proposición se deduce de la teoría general de la hiperhomología
($0_{\mathrm I}$, 11.3.2); el paso al caso general se efectúa por pegado,
utilizando el hecho de que una equivalencia homotópica de complejos induce
una equivalencia homotópica entre resoluciones proyectivas de
Cartan--Eilenberg de dichos complejos (M, XVII, 1.2).
\end{proof}

\begin{proposition}[6.7.6]
\label{III.6.7.6}
Supongamos que el complejo $\mathscr P_\bullet^{(1)}$ o el complejo
$\mathscr P_\bullet^{(2)}$ está formado por módulos $S$-planos (estando ambos
complejos acotados inferiormente). Se tiene entonces un isomorfismo canónico
de bi-$\partial$-funtores
\begin{env}[6.7.6.1]
\label{III.6.7.6.1}
\[
  \mathfrak{Tor}_n^S
  (f_1,f_2;\mathscr P_\bullet^{(1)},\mathscr P_\bullet^{(2)})
  \xrightarrow{\sim}
  \mathscr H^{-n}\bigl(f_1\times_S f_2,
    \mathscr P_\bullet^{(1)}\otimes_S\mathscr P_\bullet^{(2)}\bigr).
  \tag{6.7.6.1}\label{III.6.7.6.1.eq}
\]
\end{env}
\end{proposition}

\begin{proof}
Supongamos primero que $S$, $Y^{(1)}$ e $Y^{(2)}$ son afines, de modo que
estamos en la situación de (6.6.2), cuyas notaciones conservamos. Supongamos,
por ejemplo, que $\mathscr P_\bullet^{(1)}$ está formado por módulos
$S$-planos, y calculemos el hipertor utilizando la observación (6.6.6): es,
por tanto, la homología de
$L_{\bullet\bullet}^{(1)}\otimes_A M_{\bullet\bullet\bullet}^{(2)}$, donde
$M_{\bullet\bullet\bullet}^{(2)}$ es una resolución proyectiva de
Cartan--Eilenberg de $L_{\bullet\bullet}^{(2)}$, en el sentido de
($0_{\mathrm I}$, 11.7.1). Por otra parte, los módulos
$L_{ij}^{(1)}$ son \emph{planos} sobre $A$ en virtud de la hipótesis
(I, 4.15.1); de ($0_{\mathrm I}$, 11.7.5) se deduce entonces un isomorfismo
canónico
\begin{env}[6.7.6.2]
\label{III.6.7.6.2}
\[
  \mathbf{Tor}_\bullet^S
  (\mathfrak U^{(1)},\mathfrak U^{(2)};
   \mathscr P_\bullet^{(1)},\mathscr P_\bullet^{(2)})
  \xrightarrow{\sim}
  \mathscr H_\bullet
  (L_{\bullet\bullet}^{(1)}\otimes_A L_{\bullet\bullet}^{(2)}).
  \tag{6.7.6.2}\label{III.6.7.6.2.eq}
\]
\end{env}
Por otra parte, se tiene un homomorfismo natural de bicomplejos de
$L_{\bullet\bullet}^{(1)}\otimes_A L_{\bullet\bullet}^{(2)}$ en
$C^\bullet(\mathfrak U,\mathscr Q_\bullet)$, donde $\mathfrak U$ es el
cubrimiento de $Z=X^{(1)}\times_S X^{(2)}$ por los abiertos afines
\oldpage[III]{161}
$U_\alpha^{(1)}\times_S U_\beta^{(2)}$ y
$\mathscr Q_\bullet
=\mathscr P_\bullet^{(1)}\otimes_S\mathscr P_\bullet^{(2)}$ (considerado como
complejo simple para el grado total); en efecto, la definición de este
homomorfismo se dio en sustancia durante el cálculo de la sucesión $(b')$ en
(6.6.4), para $q=0$; basta simplemente (conservando las notaciones de (6.6.4))
tener en cuenta que hay, por una parte, un homomorfismo natural del complejo
$N^{\bullet k}$ en el complejo
$C^\bullet(\Phi^{(1)},\Phi^{(2)};\mathscr S_k)$
($0_{\mathrm I}$, 11.8.5), por otra parte, un homomorfismo natural de este
último complejo en el complejo
$P^\bullet(\Phi^{(1)},\Phi^{(2)};\mathscr S_k)$
($0_{\mathrm I}$, 11.8.6), y, por último, un homomorfismo natural de este
último complejo de cocadenas en el subcomplejo de las cocadenas alternadas
($0_{\mathrm I}$, 11.8.7). Además, el homomorfismo de bicomplejos así definido
induce un isomorfismo en homología, como se vio en (6.6.4); componiendo con
(6.7.6.2), se obtiene, por tanto, un isomorfismo
\begin{env}[6.7.6.3]
\label{III.6.7.6.3}
\[
  \mathbf{Tor}_n^S
  (\mathfrak U^{(1)},\mathfrak U^{(2)};
   \mathscr P_\bullet^{(1)},\mathscr P_\bullet^{(2)})
  \xrightarrow{\sim}
  \mathbf H^{-n}
  (\mathfrak U,\mathscr P_\bullet^{(1)}\otimes_S\mathscr P_\bullet^{(2)}).
  \tag{6.7.6.3}\label{III.6.7.6.3.eq}
\]
\end{env}
Hay que demostrar a continuación que el isomorfismo así definido no depende
de los cubrimientos abiertos elegidos (el segundo miembro de (6.7.6.3) es
canónicamente isomorfo a
$\mathbf H^{-n}
(X^{(1)}\times_S X^{(2)},
 \mathscr P_\bullet^{(1)}\otimes_S\mathscr P_\bullet^{(2)})$ por (6.2.2));
esto se hace mediante (6.6.7), observando (con las notaciones de (6.6.7)) que
se tiene un diagrama conmutativo salvo homotopía
\[
\xymatrix@C=18pt@R=28pt{
  C^\bullet(\mathfrak U^{(1)},\mathscr P_\bullet^{(1)})
  \otimes_A
  C^\bullet(\mathfrak U^{(2)},\mathscr P_\bullet^{(2)})
  \ar[r]\ar[d]
  & C^\bullet(\mathfrak U,\mathscr Q_\bullet)\ar[d]
  \\
  C^\bullet(\mathfrak V^{(1)},\mathscr P_\bullet^{(1)})
  \otimes_A
  C^\bullet(\mathfrak V^{(2)},\mathscr P_\bullet^{(2)})
  \ar[r]
  & C^\bullet(\mathfrak V,\mathscr Q_\bullet)
}
\]
donde las flechas horizontales son los homomorfismos definidos anteriormente.
Por último, hay que pasar al caso general por pegado, lo que se hace sin
dificultad como en (6.7.1) y (6.7.2); dejamos los detalles al lector.
\end{proof}

\begin{proposition}[6.7.7]
\label{III.6.7.7}
Supongamos que $\mathscr P_\bullet^{(1)}$ y
$\mathscr P_\bullet^{(2)}$ están acotados inferiormente, y que todos los
módulos $\mathscr H^{-n}(f_1,\mathscr P_\bullet^{(1)})$ o todos los módulos
$\mathscr H^{-n}(f_2,\mathscr P_\bullet^{(2)})$ son $S$-planos. Se tiene
entonces un isomorfismo canónico de bi-$\partial$-funtores ($n$ recorre
$\mathbf Z$)
\begin{env}[6.7.7.1]
\label{III.6.7.7.1}
\[
  \mathfrak{Tor}_n^S
  (f_1,f_2;\mathscr P_\bullet^{(1)},\mathscr P_\bullet^{(2)})
  \xrightarrow{\sim}
  \bigoplus_{q_1+q_2=n}
  \mathscr H^{-q_1}(f_1,\mathscr P_\bullet^{(1)})
  \otimes_S
  \mathscr H^{-q_2}(f_2,\mathscr P_\bullet^{(2)}).
  \tag{6.7.7.1}\label{III.6.7.7.1.eq}
\]
\end{env}
\end{proposition}

\begin{proof}
En efecto, en vista de (6.5.8), la sucesión espectral $(a)$ de (6.7.3) es
degenerada, y la proposición se deduce inmediatamente de
($0_{\mathrm I}$, 11.1.6), pues esta sucesión es birregular (6.7.3).
\end{proof}

\begin{theorem}[6.7.8]
\label{III.6.7.8}
Supongamos que se cumplen las condiciones siguientes:
\begin{enumerate}
  \item[{$1^\circ$}]
    Los complejos $\mathscr P_\bullet^{(1)}$ y
    $\mathscr P_\bullet^{(2)}$ están acotados inferiormente.
  \item[{$2^\circ$}]
    El complejo $\mathscr P_\bullet^{(1)}$ o el complejo
    $\mathscr P_\bullet^{(2)}$ está formado por módulos $S$-planos.

  \item[{$3^\circ$}]
    \oldpage[III]{162}
    Todos los módulos $\mathscr H^{-n}(f_1,\mathscr P_\bullet^{(1)})$ o todos
    los módulos $\mathscr H^{-n}(f_2,\mathscr P_\bullet^{(2)})$ son
    $S$-planos.
\end{enumerate}
Se tiene entonces un isomorfismo canónico de bi-$\partial$-funtores
($n$ recorre $\mathbf Z$)
\begin{env}[6.7.8.1]
\label{III.6.7.8.1}
\[
  \mathscr H^n\bigl(f_1\times_S f_2,
    \mathscr P_\bullet^{(1)}\otimes_S\mathscr P_\bullet^{(2)}\bigr)
  \xrightarrow{\sim}
  \bigoplus_{n_1+n_2=n}
  \mathscr H^{n_1}(f_1,\mathscr P_\bullet^{(1)})
  \otimes_S
  \mathscr H^{n_2}(f_2,\mathscr P_\bullet^{(2)}).
  \tag{6.7.8.1}\label{III.6.7.8.1.eq}
\]
\end{env}
(«fórmula de Künneth»).
\end{theorem}

\begin{proof}
Esto se deduce de (6.7.6) y (6.7.7).

Cuando $S$, $Y^{(1)}$ e $Y^{(2)}$ son afines, el isomorfismo recíproco de
(6.7.8.1) se deduce (con las notaciones de (6.7.6)) del homomorfismo de
bicomplejos
\[
  C^\bullet(\mathfrak U^{(1)},\mathscr P_\bullet^{(1)})
  \otimes_A
  C^\bullet(\mathfrak U^{(2)},\mathscr P_\bullet^{(2)})
  \longrightarrow
  C^\bullet(\mathfrak U,\mathscr Q_\bullet)
\]
mediante el procedimiento definido en (G, I, 2.7), como se deduce de
(G, I, 5.5).
\end{proof}

\begin{proposition}[6.7.9]
\label{III.6.7.9}
Supongamos que se cumplen las tres condiciones siguientes:
\begin{enumerate}
  \item[{$1^\circ$}]
    $S$, $Y^{(1)}$ e $Y^{(2)}$ son localmente noetherianos, $f_1$ y $f_2$
    son propios, y $Y^{(1)}$ o $Y^{(2)}$ es de tipo finito sobre $S$.
  \item[{$2^\circ$}]
    $\mathscr P_\bullet^{(1)}$ y $\mathscr P_\bullet^{(2)}$ están acotados
    inferiormente.
  \item[{$3^\circ$}]
    Para todo $n\in\mathbf Z$,
    $\mathscr H_n(\mathscr P_\bullet^{(i)})$ es un módulo coherente
    ($i=1,2$).
\end{enumerate}
En estas condiciones,
$\mathfrak{Tor}_n^S
(f_1,f_2;\mathscr P_\bullet^{(1)},\mathscr P_\bullet^{(2)})$ es un
$\mathscr O_Y$-módulo coherente (con
$Y=Y^{(1)}\times_S Y^{(2)}$).
\end{proposition}

\begin{proof}
De (6.5.13) se deduce que los hipertor locales
$\mathscr{T}\!or_n^S(\mathscr P_\bullet^{(1)},\mathscr P_\bullet^{(2)})$ son
$\mathscr O_X$-módulos coherentes
($X=X^{(1)}\times_S X^{(2)}$ es localmente noetheriano, pues uno de los
$X^{(i)}$ es, por hipótesis, de tipo finito sobre $S$ (I, 6.3.4 y 6.3.8)).
Como $Y$ es localmente noetheriano y $f_1\times_S f_2$ es propio
(II, 5.4.2), de (6.2.5) se deduce que los términos
${}^{(b)}\mathscr E_{pq}^2$ de (6.7.3) son $\mathscr O_Y$-módulos coherentes.
Como todas las sucesiones espectrales de (6.7.3) son birregulares en virtud
de la hipótesis $2^\circ$, se concluye por ($0_{\mathrm I}$, 11.1.8).
\end{proof}

\begin{env}[6.7.10]
\label{III.6.7.10}
Sean ahora $Y'^{(i)}$ dos $S$-preesquemas,
$v_i:Y'^{(i)}\to Y^{(i)}$ dos $S$-morfismos ($i=1,2$) y
$v=v_1\times_S v_2$ su producto, que es un $S$-morfismo $Y'\to Y$, donde se
pone $Y'=Y'^{(1)}\times_S Y'^{(2)}$. Consideremos, por otra parte, para
$i=1,2$, un $S$-preesquema $X'^{(i)}$ y dos $S$-morfismos
$u_i:X'^{(i)}\to X^{(i)}$, $f'_i:X'^{(i)}\to Y'^{(i)}$, de modo que los
diagramas
\begin{env}[6.7.10.1]
\label{III.6.7.10.1}
\[
\xymatrix@C=36pt@R=28pt{
  X'^{(i)}\ar[r]^{u_i}\ar[d]_{f'_i}
  & X^{(i)}\ar[d]^{f_i}
  \\
  Y'^{(i)}\ar[r]_{v_i}
  & Y^{(i)}
}
\tag{6.7.10.1}\label{III.6.7.10.1.eq}
\]
\end{env}
sean conmutativos, siendo los morfismos $f'_i$ \emph{separados y
cuasicompactos}. Se tienen entonces homomorfismos canónicos
$\mathscr O_{Y'}$-lineales de funtores espectrales
\begin{env}[6.7.10.2]
\label{III.6.7.10.2}
\[
  v^*\bigl({}^{(t)}\mathscr E
  (f_1,f_2;\mathscr P_\bullet^{(1)},\mathscr P_\bullet^{(2)})\bigr)
  \longrightarrow
  {}^{(t)}\mathscr E
  (f'_1,f'_2;u_1^*(\mathscr P_\bullet^{(1)}),
                 u_2^*(\mathscr P_\bullet^{(2)}))
  \tag{6.7.10.2}\label{III.6.7.10.2.eq}
\]
\end{env}
para $t=a$, $a'$, $b$, $b'$, $c$, $d$. Para definirlos, supongamos primero
que $S=\operatorname{Spec}(A)$,
$Y^{(i)}=\operatorname{Spec}(B_i)$ e
$Y'^{(i)}=\operatorname{Spec}(B'_i)$ son afines; entonces $X^{(i)}$ y
$X'^{(i)}$ son esquemas cuasicompactos. Para calcular las sucesiones
espectrales
${}^{(t)}\mathscr E
(f_1,f_2;\mathscr P_\bullet^{(1)},\mathscr P_\bullet^{(2)})$, consideramos,
como en (6.6.1), cubrimientos finitos $\mathfrak U^{(i)}$ por abiertos afines
de $X^{(i)}$ ($i=1,2$); para calcular
${}^{(t)}\mathscr E
(f'_1,f'_2;u_1^*(\mathscr P_\bullet^{(1)}),
u_2^*(\mathscr P_\bullet^{(2)}))$, consideramos cubrimientos finitos
$\mathfrak U'^{(i)}$

\oldpage[III]{163}
de $X'^{(i)}$ por abiertos afines, \emph{más finos}, respectivamente, que los
cubrimientos $u_i^{-1}(\mathfrak U^{(i)})$ ($i=1,2$). Está claro que el
bicomplejo
$C^\bullet(\mathfrak U^{(i)},\mathscr P_\bullet^{(i)})
=L_{\bullet\bullet}^{(i)}$ puede considerarse canónicamente como un
subbicomplejo de
$C^\bullet(u_i^{-1}(\mathfrak U^{(i)}),u_i^*(\mathscr P_\bullet^{(i)}))$
($0_{\mathrm I}$, 4.4.3.2); además, eligiendo un morfismo simplicial
(G, II, 5.7) de $\mathfrak U'^{(i)}$ en
$u_i^{-1}(\mathfrak U^{(i)})$, se define un homomorfismo de bicomplejos
\[
  C^\bullet(u_i^{-1}(\mathfrak U^{(i)}),u_i^*(\mathscr P_\bullet^{(i)}))
  \longrightarrow
  C^\bullet(\mathfrak U'^{(i)},u_i^*(\mathscr P_\bullet^{(i)}))
\]
del cual se obtiene, por composición, un homomorfismo de bicomplejos
\[
  L_{\bullet\bullet}^{(i)}
  \longrightarrow
  L_{\bullet\bullet}^{\prime(i)}
  =C^\bullet(\mathfrak U'^{(i)},u_i^*(\mathscr P_\bullet^{(i)})).
\]
Además, este homomorfismo se sustituye por uno homotópico cuando se cambia el
morfismo simplicial (G, II, 5.7.1); así se obtiene un homomorfismo bien
definido de funtores espectrales:
\begin{env}[6.7.10.3]
\label{III.6.7.10.3}
\[
  {}^{(t)}\mathscr E
  (\mathfrak U^{(1)},\mathfrak U^{(2)};
   \mathscr P_\bullet^{(1)},\mathscr P_\bullet^{(2)})
  \longrightarrow
  {}^{(t)}\mathscr E
  (\mathfrak U'^{(1)},\mathfrak U'^{(2)};
   u_1^*(\mathscr P_\bullet^{(1)}),u_2^*(\mathscr P_\bullet^{(2)})).
  \tag{6.7.10.3}\label{III.6.7.10.3.eq}
\]
\end{env}
Se verifica inmediatamente que, si $\mathfrak V^{(i)}$ es un cubrimiento
afín finito de $X^{(i)}$ más fino que $\mathfrak U^{(i)}$, y
$\mathfrak V'^{(i)}$ un cubrimiento afín finito de $X'^{(i)}$, más fino que
$u_i^{-1}(\mathfrak V^{(i)})$ y que $\mathfrak U'^{(i)}$, el diagrama
\[
\xymatrix@C=34pt@R=30pt{
  C^\bullet(\mathfrak U^{(i)},\mathscr P_\bullet^{(i)})
  \ar[r]\ar[d]
  & C^\bullet(\mathfrak V^{(i)},\mathscr P_\bullet^{(i)})\ar[d]
  \\
  C^\bullet(\mathfrak U'^{(i)},u_i^*(\mathscr P_\bullet^{(i)}))
  \ar[r]
  & C^\bullet(\mathfrak V'^{(i)},u_i^*(\mathscr P_\bullet^{(i)}))
}
\]
es conmutativo, lo que implica que el homomorfismo (6.7.10.3) no depende
esencialmente de los cubrimientos $\mathfrak U^{(i)}$ y
$\mathfrak U'^{(i)}$ considerados. Por tanto, se ha definido en realidad un
homomorfismo de $A$-módulos
\begin{env}[6.7.10.4]
\label{III.6.7.10.4}
\[
  {}^{(t)}E
  (X^{(1)},X^{(2)};\mathscr P_\bullet^{(1)},\mathscr P_\bullet^{(2)})
  \longrightarrow
  {}^{(t)}E
  (X'^{(1)},X'^{(2)};
   u_1^*(\mathscr P_\bullet^{(1)}),u_2^*(\mathscr P_\bullet^{(2)})).
  \tag{6.7.10.4}\label{III.6.7.10.4.eq}
\]
\end{env}
pero, por la definición de los $u_i^*(\mathscr P_\bullet^{(i)})$ y en virtud
de la conmutatividad de (6.7.10.1), está claro que este homomorfismo es también
un homomorfismo de $(B_1\otimes_A B_2)$-módulos; como el segundo miembro de
(6.7.10.4) está formado por $(B'_1\otimes_A B'_2)$-módulos, se deduce
canónicamente de (6.7.10.4) un homomorfismo de
$(B'_1\otimes_A B'_2)$-módulos
\begin{env}[6.7.10.5]
\label{III.6.7.10.5}
\[
  {}^{(t)}E
  (X^{(1)},X^{(2)};\mathscr P_\bullet^{(1)},\mathscr P_\bullet^{(2)})
  \otimes_{B_1\otimes_A B_2}(B'_1\otimes_A B'_2)
  \longrightarrow
  {}^{(t)}E
  (X'^{(1)},X'^{(2)};
   u_1^*(\mathscr P_\bullet^{(1)}),u_2^*(\mathscr P_\bullet^{(2)}))
  \tag{6.7.10.5}\label{III.6.7.10.5.eq}
\]
\end{env}
lo que, teniendo en cuenta (I, 1.6.5), no es sino el homomorfismo buscado
(6.7.10.2) en el caso particular considerado.

Queda pasar al caso general siguiendo los procedimientos de pegado de (6.7.1)
y (6.7.2); el
segundo paso es inmediato; por lo que respecta al primero, se consideran, como
en (6.7.1), elementos $g_i\in B_i$ y sus imágenes $g'_i\in B'_i$, el producto
tensorial $g=g_1\otimes g_2$ en $B=B_1\otimes_A B_2$ y su imagen
$g'=g'_1\otimes g'_2$ en $B'=B'_1\otimes_A B'_2$, y todo se reduce a utilizar

\oldpage[III]{164}
el isomorfismo canónico
\[
  (M\otimes_B B')_{g'}
  \xrightarrow{\sim}
  M_g\otimes_{B_g}B'_{g'}
\]
($0_{\mathrm I}$, 1.5.4); dejamos los detalles al lector.
\end{env}

\begin{env}[6.7.11]
\label{III.6.7.11}
La teoría de los hipertor globales, desarrollada anteriormente para dos
$S$-morfismos $X^{(i)}\to Y^{(i)}$ y dos complejos
$\mathscr P_\bullet^{(i)}$ de módulos cuasicoherentes acotados inferiormente,
se extiende inmediatamente al caso general siguiente: se tiene un preesquema
$S$, una familia finita de $S$-preesquemas $Y^{(i)}$ ($i\in I$), una familia
finita de $S$-morfismos separados y cuasicompactos
$f_i:X^{(i)}\to Y^{(i)}$, y, para cada $i$, un complejo acotado inferiormente
de $\mathscr O_{X^{(i)}}$-módulos cuasicoherentes
$\mathscr P_\bullet^{(i)}$. Si $Y$ es el producto de los $S$-preesquemas
$Y^{(i)}$, se define entonces, para cada entero $n\in\mathbf Z$, un
$\mathscr O_Y$-módulo cuasicoherente
$\mathfrak{Tor}_n^S((f_i)_{i\in I};
(\mathscr P_\bullet^{(i)})_{i\in I})$;
estos módulos forman un $\partial$-funtor covariante en cada uno de los
complejos $\mathscr P_\bullet^{(i)}$; además, este funtor es el término límite
común de seis funtores espectrales
${}^{(t)}\mathscr E((f_i)_{i\in I};(\mathscr P_\bullet^{(i)})_{i\in I})$.
Dejamos al lector la tarea de repetir, para este caso general, las definiciones
y los razonamientos hechos anteriormente para $I=\{1,2\}$. Obsérvese
simplemente que, cuando $I$ se reduce a un solo elemento, se recupera la
hipercohomología $\mathscr H^\bullet(f,\mathscr P_\bullet)$ definida en
(6.2.7) (como ya se observó en (6.7.4)). Cuando $I$ es el intervalo
$1\leq i\leq m$ de $\mathbf N$, escribiremos
\phantomsection\label{III.6.7.11.2}
\[
  \mathfrak{Tor}_n^S
  (f_1,\ldots,f_m;
   \mathscr P_\bullet^{(1)},\ldots,\mathscr P_\bullet^{(m)})
  \quad\text{para}\quad
  \mathfrak{Tor}_n^S
  ((f_i)_{i\in I};(\mathscr P_\bullet^{(i)})_{i\in I}).
\]
\end{env}

\begin{proposition}[6.7.12]
\label{III.6.7.12}
Con las notaciones de (6.7.11), sea $J$ un subconjunto de $I$ tal que, para
$i\in I-J$, se tenga $X^{(i)}=Y^{(i)}=S$, que $f_i$ sea la identidad y que
$\mathscr P_\bullet^{(i)}$ sea el complejo reducido al término de grado $0$
igual a $\mathscr O_S$. Hay entonces un isomorfismo canónico de
$\partial$-funtores
\begin{env}[6.7.12.1]
\label{III.6.7.12.1}
\[
  \mathfrak{Tor}_\bullet^S
  ((f_i)_{i\in I};(\mathscr P_\bullet^{(i)})_{i\in I})
  \xrightarrow{\sim}
  \mathfrak{Tor}_\bullet^S
  ((f_i)_{i\in J};(\mathscr P_\bullet^{(i)})_{i\in J}).
  \tag{6.7.12.1}\label{III.6.7.12.1.eq}
\]
\end{env}
\end{proposition}

\begin{proof}
Podemos limitarnos a definir este isomorfismo cuando $S$ y los $Y^{(i)}$ son
afines, pues el pegado se efectúa como de costumbre. Para $i\in I-J$, puede
tomarse el cubrimiento $\mathfrak U^{(i)}$ formado por el único conjunto $S$,
y entonces
$L_{\bullet\bullet}^{(i)}=C^\bullet(\mathfrak U^{(i)},
\mathscr P_\bullet^{(i)})$ se reduce a su único término de grados $(0,0)$,
igual a $\Gamma(S,\mathscr O_S)=A(S)$; el isomorfismo (6.7.12.1) es entonces
evidente.
\end{proof}

\begin{remark}[6.7.13]
\label{III.6.7.13}
Con las notaciones de (6.7.3), consideremos el $S$-isomorfismo canónico
$Y^{(1)}\times_S Y^{(2)}\to Y^{(2)}\times_S Y^{(1)}$ (I, 3.3.5); entonces la
imagen por este isomorfismo de
$\mathfrak{Tor}_\bullet^S(f_1,f_2;
\mathscr P_\bullet^{(1)},\mathscr P_\bullet^{(2)})$ es
$\mathfrak{Tor}_\bullet^S(f_2,f_1;
\mathscr P_\bullet^{(2)},\mathscr P_\bullet^{(1)})$; como la cuestión es
local, nos reducimos al caso considerado en (6.6.2), y, si se designa por
$M_{\bullet\bullet\bullet}^{(i)}$ una resolución proyectiva de
Cartan--Eilenberg de $L_{\bullet\bullet}^{(i)}$ ($i=1,2$), el isomorfismo
considerado transforma
$M_{\bullet\bullet\bullet}^{(1)}\otimes_A
M_{\bullet\bullet\bullet}^{(2)}$ en
$M_{\bullet\bullet\bullet}^{(2)}\otimes_A
M_{\bullet\bullet\bullet}^{(1)}$, de donde se deduce la afirmación al
considerar la homología de los complejos simples asociados a estos
tricomplejos.
\end{remark}

\subsection{Sucesiones espectrales de asociatividad de los hipertor globales}
\label{subsection:III.6.8}
\label{III.6.8}

\begin{env}[6.8.1]
\label{III.6.8.1}
Mantengamos las hipótesis y notaciones de \sref{III.6.7.11}; en
particular, los complejos $\sh P_\bullet^{(i)}$ están acotados
inferiormente. Sea $(I_j)_{j\in J}$ una partición del conjunto de
índices $I$. Buscamos una relación de «asociatividad» entre
\[
  \mathfrak{Tor}_n^S
  \bigl(
    (f_i)_{i\in I};
    (\sh P_\bullet^{(i)})_{i\in I}
  \bigr)
\]
y los hipertor «parciales»
\[
  \mathcal T_{\bullet j}
  =
  \mathfrak{Tor}_\bullet^S
  \bigl(
    (f_i)_{i\in I_j};
    (\sh P_\bullet^{(i)})_{i\in I_j}
  \bigr).
\]
\oldpage[III]{165}

Para simplificar la notación, nos limitaremos al caso en que $I$
es el intervalo $1\leq i\leq m$ y la partición consta de los dos
intervalos
\[
  \{1,2,\ldots,r\}
  \qquad\hbox{y}\qquad
  \{r+1,\ldots,m\}.
\]
\end{env}

\begin{proposition}[6.8.2]
\label{III.6.8.2}
Existe un funtor espectral canónico birregular, denominado
\emph{funtor espectral de asociatividad}, que se denota por
\[
  {}^{(e)}\sh E^S
  \bigl(
    f_1,\ldots,f_m;
    \sh P_\bullet^{(1)},\ldots,\sh P_\bullet^{(m)}
  \bigr),
\]
o simplemente por ${}^{(e)}\sh E$ con los mismos argumentos, que
converge a
\[
  \mathfrak{Tor}_\bullet^S
  \bigl(
    f_1,\ldots,f_m;
    \sh P_\bullet^{(1)},\ldots,\sh P_\bullet^{(m)}
  \bigr)
\]
y cuyo término $E_2$ es
\[
\begin{split}
  {}^{(e)}\sh E_{pq}^2
  =
  \bigoplus_{q_1+q_2=q}
  \shTor_p^S
  \Bigl(
    &
    \mathfrak{Tor}_{q_1}^S
    \bigl(
      f_1,\ldots,f_r;
      \sh P_\bullet^{(1)},\ldots,\sh P_\bullet^{(r)}
    \bigr),
    \\
    &
    \mathfrak{Tor}_{q_2}^S
    \bigl(
      f_{r+1},\ldots,f_m;
      \sh P_\bullet^{(r+1)},\ldots,\sh P_\bullet^{(m)}
    \bigr)
  \Bigr).
\end{split}
\]
\end{proposition}

\begin{proof}
Aquí se identifica canónicamente $Y$ con
$Z^{(1)}\times_SZ^{(2)}$, donde
\[
  Z^{(1)}
  =
  Y^{(1)}\times_S\cdots\times_SY^{(r)}
  \quad\hbox{y}\quad
  Z^{(2)}
  =
  Y^{(r+1)}\times_S\cdots\times_SY^{(m)}.
\]
Basta tratar el caso en que $S$ y todos los $Y^{(i)}$ son afines.
El paso de este caso al caso general utiliza los métodos de
\sref{III.6.7.1} y \sref{III.6.7.2}; dejamos los detalles al
lector. Basta, pues, demostrar el siguiente enunciado afín.
\end{proof}

\begin{corollary}[6.8.3]
\label{III.6.8.3}
Sea $A$ un anillo, sea $S=\Spec(A)$ y sean $X^{(i)}$,
$1\leq i\leq m$, $S$-esquemas cuasicompactos. Para cada $i$, sea
$\sh P_\bullet^{(i)}$ un complejo acotado inferiormente de
$\sh O_{X^{(i)}}$-módulos cuasicoherentes. Existe un funtor
espectral canónico birregular que converge a
\[
  \mathbf{Tor}_\bullet^S
  \bigl(
    X^{(1)},\ldots,X^{(m)};
    \sh P_\bullet^{(1)},\ldots,\sh P_\bullet^{(m)}
  \bigr)
\]
y cuyo término $E_2$ es
\[
\begin{split}
  {}^{(e)}E_{pq}^2
  =
  \bigoplus_{q_1+q_2=q}
  \Tor_p^A
  \Bigl(
    &
    \mathbf{Tor}_{q_1}^S
    \bigl(
      X^{(1)},\ldots,X^{(r)};
      \sh P_\bullet^{(1)},\ldots,\sh P_\bullet^{(r)}
    \bigr),
    \\
    &
    \mathbf{Tor}_{q_2}^S
    \bigl(
      X^{(r+1)},\ldots,X^{(m)};
      \sh P_\bullet^{(r+1)},\ldots,\sh P_\bullet^{(m)}
    \bigr)
  \Bigr).
\end{split}
\]
\end{corollary}

\begin{proof}
Siguiendo la definición de \sref{III.6.6.2}, elijamos para cada
$i$ un recubrimiento abierto afín finito $\mathfrak U^{(i)}$ de
$X^{(i)}$, formemos el bicomplejo
\[
  L_{\bullet\bullet}^{(i)}
  =
  C^\bullet
  \bigl(
    \mathfrak U^{(i)},\sh P_\bullet^{(i)}
  \bigr),
\]
y elijamos una resolución proyectiva de Cartan--Eilenberg
$M_{\bullet\bullet\bullet}^{(i)}$ de este bicomplejo en el sentido
de \sref[0]{0.11.7.1}. El hipertor considerado es la homología del
producto tensorial
\[
  M_{\bullet\bullet\bullet}
  =
  \bigotimes_{i=1}^{m}
  M_{\bullet\bullet\bullet}^{(i)}.
\]
Consideremos $M_{\bullet\bullet\bullet}$ como un complejo simple
$N_\bullet$, producto tensorial de los dos complejos simples
\[
  N_\bullet'
  =
  \bigotimes_{i=1}^{r}
  M_{\bullet\bullet\bullet}^{(i)},
  \qquad
  N_\bullet''
  =
  \bigotimes_{i=r+1}^{m}
  M_{\bullet\bullet\bullet}^{(i)},
\]
donde $N_\bullet'$ y $N_\bullet''$ están graduados por las sumas
de los grados totales de los $M_{\bullet\bullet\bullet}^{(i)}$.

Los $A$-módulos que forman $N_\bullet'$ y $N_\bullet''$ son
proyectivos. Por tanto, \sref{III.6.5.9} da
\[
  H_\bullet(M_{\bullet\bullet\bullet})
  =
  \mathbf{Tor}_\bullet^A(N_\bullet',N_\bullet'').
\]
La sucesión espectral buscada es, pues, precisamente
\sref{III.6.3.2.2} aplicada a $N_\bullet'$ y $N_\bullet''$, con
la interpretación precedente de sus módulos de homología para los
dos productos parciales
\[
  X^{(1)}\times_S\cdots\times_SX^{(r)}
  \quad\hbox{y}\quad
  X^{(r+1)}\times_S\cdots\times_SX^{(m)}.
\]
Finalmente, las afirmaciones de regularidad se deducen de
\sref{III.6.3.2}. En efecto, como los
$\sh P_\bullet^{(i)}$ están acotados inferiormente, también lo
están los $M_{\bullet\bullet\bullet}^{(i)}$ y, por consiguiente,
$N_\bullet'$ y $N_\bullet''$.
\end{proof}

\oldpage[III]{166}


\subsection{Las sucesiones espectrales de cambio de base en los hipertor
globales.}
\label{subsection:III.6.9}
\label{III.6.9}

\begin{env}[6.9.1]
\label{III.6.9.1}
Se mantienen las hipótesis y notaciones de \sref{III.6.7.11} (y, en
particular, se supone que los $\mathscr P_\bullet^{(i)}$ están acotados
inferiormente). Consideremos un morfismo $g:S'\to S$ de preesquemas, y
pongamos
$Y'^{(i)}=Y_{(S')}^{(i)}$, $X'^{(i)}=X_{(S')}^{(i)}$ y
$\mathscr P_\bullet'^{(i)}=
\mathscr P_\bullet^{(i)}\otimes_{\mathscr O_S}\mathscr O_{S'}$,
de modo que $\mathscr P_\bullet'^{(i)}$ es un complejo de
$\mathscr O_{X'^{(i)}}$-módulos cuasicoherentes; sea
$f'_i=(f_i)_{(S')}:X'^{(i)}\to Y'^{(i)}$, que es un morfismo separado y
cuasicompacto (\sref[I]{I.5.5.1} y \sref{III.6.6.4}). Nos proponemos
estudiar las relaciones entre los $\mathscr O_{Y'}$-módulos cuasicoherentes
$\mathfrak{Tor}_n^{S'}((f'_i)_{i\in I};
(\mathscr P_\bullet'^{(i)})_{i\in I})$ y
$\mathfrak{Tor}_n^S((f_i)_{i\in I};
(\mathscr P_\bullet^{(i)})_{i\in I})
\otimes_{\mathscr O_S}\mathscr O_{S'}$, donde
$Y'=Y\times_S S'=Y_{(S')}$. Un caso particularmente sencillo es el
siguiente, que se reduce a \sref[I]{I.1.4.15} cuando $I$ se reduce a un solo
elemento y $\mathscr P_\bullet$ a un solo módulo:
\end{env}

\begin{proposition}[6.9.2]
\label{III.6.9.2}
Si el morfismo $g:S'\to S$ es plano, se tiene un isomorfismo canónico de
$\partial$-funtores (en los $\mathscr P_\bullet^{(i)}$):
\begin{env}[6.9.2.1]
\label{III.6.9.2.1}
\[
  \mathfrak{Tor}_\bullet^S
  ((f_i)_{i\in I};(\mathscr P_\bullet^{(i)})_{i\in I})
  \otimes_{\mathscr O_S}\mathscr O_{S'}
  \xrightarrow{\sim}
  \mathfrak{Tor}_\bullet^{S'}
  ((f'_i)_{i\in I};(\mathscr P_\bullet'^{(i)})_{i\in I}).
  \tag{6.9.2.1}\label{III.6.9.2.1.eq}
\]
\end{env}
\end{proposition}

\begin{proof}
Todavía se puede restringir al caso en que $S$, $S'$ y los $Y^{(i)}$ son
afines, efectuándose el pegado mediante los métodos de \sref{III.6.7.1} y
\sref{III.6.7.2}. Sean $S=\operatorname{Spec}(A)$,
$S'=\operatorname{Spec}(A')$, y tomemos para cada $i$ un cubrimiento abierto
afín $\mathfrak U^{(i)}$ de $X^{(i)}$; si
$u_i:X'^{(i)}\to X^{(i)}$ es la proyección canónica,
$u_i^{-1}(\mathfrak U^{(i)})$ es un cubrimiento abierto afín de
$X'^{(i)}$ (\sref[II]{II.1.5.5}), que denotaremos por
$\mathfrak U'^{(i)}$; está claro entonces que
\[
  C^\bullet(\mathfrak U'^{(i)},\mathscr P_\bullet'^{(i)})
  =C^\bullet(\mathfrak U^{(i)},\mathscr P_\bullet^{(i)})\otimes_A A',
\]
y la existencia del isomorfismo \eqref{III.6.9.2.1.eq} es inmediata, pues si
$M_{\bullet\bullet\bullet}^{(i)}$ es una resolución proyectiva de
Cartan-Eilenberg de
$L_{\bullet\bullet}^{(i)}=C^\bullet(\mathfrak U^{(i)},
\mathscr P_\bullet^{(i)})$ en el sentido de \sref[0]{0.11.7.1}, formada por
$A$-módulos, $M_{\bullet\bullet\bullet}^{(i)}\otimes_A A'$ es una resolución
proyectiva de Cartan-Eilenberg (en el mismo sentido) de
$L_{\bullet\bullet}^{(i)}\otimes_A A'$ formada por $A'$-módulos, en virtud
de la hipótesis de que $A'$ es un $A$-módulo plano; esta misma hipótesis
muestra además que
\[
  H_\bullet\left(
    \bigotimes_{i=1}^m
    (M_{\bullet\bullet\bullet}^{(i)}\otimes_A A')
  \right)
  =H_\bullet\left(
    \bigotimes_{i=1}^m M_{\bullet\bullet\bullet}^{(i)}
  \right)\otimes_A A'.
\]
\end{proof}

\begin{env}[6.9.2.2]
\label{III.6.9.2.2}
Obsérvese que, cuando $I$ se reduce al único elemento $1$, la fórmula
\eqref{III.6.9.2.1.eq} se deduce directamente de \sref{III.6.7.7}, aplicado
tomando $Y_2=X_2=S'$, $f_2=1_{S'}$, y el complejo
$\mathscr P_\bullet^{(2)}$ reducido a su término de grado $0$, igual a
$\mathscr O_{S'}$; se sabe entonces que la hipercohomología
$\mathscr H^n(1_{S'},\mathscr O_{S'})$ es nula para todo $n\ne0$ y se reduce
a $\mathscr O_{S'}$ para $n=0$ (\sref{III.6.2.1}).
\end{env}

En el caso general, introduciremos, en lugar de $\mathscr O_{S'}$, un
complejo $\mathscr Q'_\bullet$ de $\mathscr O_{S'}$-módulos cuasicoherentes
acotado inferiormente, de modo que, si para simplificar se toma
$I=\{1,2,\ldots,m\}$, se puede considerar el $\partial$-funtor
\[
  \mathfrak{Tor}_\bullet^S
  (f_1,\ldots,f_m,1_{S'};
   \mathscr P_\bullet^{(1)},\ldots,\mathscr P_\bullet^{(m)},
   \mathscr Q'_\bullet).
\]

\begin{proposition}[6.9.3]
\label{III.6.9.3}
Existen tres funtores espectrales canónicos birregulares, denotados por
\[
{}^{(t)}\mathscr E(f_1,\ldots,f_m;
\mathscr P_\bullet^{(1)},\ldots,\mathscr P_\bullet^{(m)},
\mathscr Q'_\bullet)
\qquad (t=e,\ f\text{ o }f'),
\]
que tienen por término límite común
\[
\mathfrak{Tor}_\bullet^S(f_1,\ldots,f_m,1_{S'};
\mathscr P_\bullet^{(1)},\ldots,\mathscr P_\bullet^{(m)},
\mathscr Q'_\bullet),
\]
y cuyos términos $E_2$ son respectivamente

\oldpage[III]{167}
\[
\begin{split}
  {}^{(e)}\mathscr E_{pq}^2
  =\bigoplus_{q'+q''=q}
  \mathscr{T}\!or_p^S\Bigl(
  &\mathfrak{Tor}_{q'}^S
  (f_1,\ldots,f_m;
   \mathscr P_\bullet^{(1)},\ldots,\mathscr P_\bullet^{(m)}),\\
  &\mathscr H_{q''}(\mathscr Q'_\bullet)
  \Bigr),
\end{split}
\]
\[
\begin{split}
  {}^{(f)}\mathscr E_{pq}^2
  =\bigoplus_{q_1+q_2+\cdots+q_{m+1}=q}
  \mathscr{T}\!or_p^{S'}\Bigl(
  &\mathfrak{Tor}_{q_1}^S
  (f_1,1_{S'};\mathscr P_\bullet^{(1)},\mathscr O_{S'}),\ldots,\\
  &\mathfrak{Tor}_{q_m}^S
  (f_m,1_{S'};\mathscr P_\bullet^{(m)},\mathscr O_{S'}),
  \mathscr H_{q_{m+1}}(\mathscr Q'_\bullet)
  \Bigr),
\end{split}
\]
\[
\begin{split}
  {}^{(f')}\mathscr E_{pq}^2
  =\bigoplus_{q_1+\cdots+q_m=q}
  \mathfrak{Tor}_p^{S'}\Bigl(
  &f'_1,\ldots,f'_m,1_{S'};\\
  &\mathscr{T}\!or_{q_1}^S
  (\mathscr P_\bullet^{(1)},\mathscr O_{S'}),\ldots,
  \mathscr{T}\!or_{q_m}^S
  (\mathscr P_\bullet^{(m)},\mathscr O_{S'}),
  \mathscr Q'_\bullet
  \Bigr).
\end{split}
\]
\end{proposition}

\begin{proof}
La sucesión $(e)$ no es otra que la sucesión de asociatividad de
\sref{III.6.8.2} para $r=m$. Para definir las otras dos sucesiones
espectrales, nos limitaremos una vez más al caso en que $S$, $S'$ y los
$Y^{(i)}$ son afines; el paso al caso general se hace mediante los métodos de
\sref{III.6.7.1} y \sref{III.6.7.2}, y se deja al lector. Por tanto, queda
por demostrar el siguiente enunciado afín.
\end{proof}

\begin{corollary}[6.9.4]
\label{III.6.9.4}
Sean $A$ un anillo, $A'$ una $A$-álgebra,
$S=\operatorname{Spec}(A)$, $S'=\operatorname{Spec}(A')$,
$X^{(i)}$ ($1\leq i\leq m$) unos $S$-esquemas cuasicompactos y, para cada
$i$, sea $X'^{(i)}=X_{(S')}^{(i)}$, que es un $S'$-esquema cuasicompacto.
Para cada $i$, sea $\mathscr P_\bullet^{(i)}$ un complejo de
$\mathscr O_{X^{(i)}}$-módulos cuasicoherentes; sea finalmente
$Q'_\bullet$ un complejo de $A'$-módulos, estando estos complejos acotados
inferiormente. Existen tres funtores espectrales birregulares en los
$\mathscr P_\bullet^{(i)}$ y en $Q'_\bullet$, que tienen por término límite
común
\[
  \mathbf{Tor}_\bullet^S
  (X^{(1)},\ldots,X^{(m)},S';
   \mathscr P_\bullet^{(1)},\ldots,\mathscr P_\bullet^{(m)},
   \widetilde{Q'_\bullet})
\]
y cuyos términos $E_2$ son respectivamente
\[
\begin{split}
  {}^{(e)}E_{pq}^2
  =\bigoplus_{q'+q''=q}
  \operatorname{Tor}_p^A\Bigl(
  &\mathbf{Tor}_{q'}^S
  (X^{(1)},\ldots,X^{(m)};
   \mathscr P_\bullet^{(1)},\ldots,\mathscr P_\bullet^{(m)}),\\
  &H_{q''}(Q'_\bullet)
  \Bigr),
\end{split}
\]
\[
\begin{split}
  {}^{(f)}E_{pq}^2
  =\bigoplus_{q_1+q_2+\cdots+q_{m+1}=q}
  \operatorname{Tor}_p^{A'}\Bigl(
  &\mathbf{Tor}_{q_1}^S
  (X^{(1)},S';\mathscr P_\bullet^{(1)},\mathscr O_{S'}),\ldots,\\
  &\mathbf{Tor}_{q_m}^S
  (X^{(m)},S';\mathscr P_\bullet^{(m)},\mathscr O_{S'}),
  H_{q_{m+1}}(Q'_\bullet)
  \Bigr),
\end{split}
\]
\[
\begin{split}
  {}^{(f')}E_{pq}^2
  =\bigoplus_{q_1+\cdots+q_m=q}
  \mathbf{Tor}_p^{S'}\Bigl(
  &X'^{(1)},\ldots,X'^{(m)},S';\\
  &\mathscr{T}\!or_{q_1}^S
  (\mathscr P_\bullet^{(1)},\mathscr O_{S'}),\ldots,
  \mathscr{T}\!or_{q_m}^S
  (\mathscr P_\bullet^{(m)},\mathscr O_{S'}),
  \widetilde{Q'_\bullet}
  \Bigr).
\end{split}
\]
\end{corollary}

\begin{proof}
No volveremos sobre el primero de estos funtores espectrales, que fue tratado
en \sref{III.6.8.3} y solo se incluye aquí como recordatorio. Para definir
los otros, consideremos para cada $i$ un cubrimiento abierto afín finito
$\mathfrak U^{(i)}$ de $X^{(i)}$ y, si
$u_i:X'^{(i)}\to X^{(i)}$ es la proyección canónica, el cubrimiento abierto
afín finito correspondiente
$\mathfrak U'^{(i)}=u_i^{-1}(\mathfrak U^{(i)})$.

En virtud de \sref{III.6.6.6},
$\mathbf{Tor}_\bullet^S(X^{(1)},\ldots,X^{(m)},S';
\mathscr P_\bullet^{(1)},\ldots,\mathscr P_\bullet^{(m)},
\widetilde{Q'_\bullet})$ se obtiene tomando, para $1\leq i\leq m$, una
resolución proyectiva de Cartan-Eilenberg
$M_{\bullet\bullet\bullet}^{(i)}$ de
$L_{\bullet\bullet}^{(i)}=C^\bullet(\mathfrak U^{(i)},
\mathscr P_\bullet^{(i)})$ (en el sentido de \sref[0]{0.11.7.1}),
considerando el tricomplejo
\[
  M_{\bullet\bullet\bullet}
  =M_{\bullet\bullet\bullet}^{(1)}\otimes_A
   M_{\bullet\bullet\bullet}^{(2)}\otimes_A\cdots\otimes_A
   M_{\bullet\bullet\bullet}^{(m)}\otimes_A Q'_\bullet
\]
(donde $Q'_\bullet$ se considera como un tricomplejo cuyos dos últimos
grados se reducen a $0$), y tomando su homología. Si se pone
$M_{\bullet\bullet\bullet}'^{(i)}=
M_{\bullet\bullet\bullet}^{(i)}\otimes_A A'$, se tiene (recordando que
$Q'_\bullet$ es un complejo de $A'$-módulos)
\[
  M_{\bullet\bullet\bullet}
  =M_{\bullet\bullet\bullet}'^{(1)}\otimes_{A'}
   M_{\bullet\bullet\bullet}'^{(2)}\otimes_{A'}\cdots\otimes_{A'}
   M_{\bullet\bullet\bullet}'^{(m)}\otimes_{A'} Q'_\bullet.
\]
Ahora bien, consideremos cada uno de los complejos
$M_{\bullet\bullet\bullet}'^{(i)}$ como un complejo simple (para su grado
total), y observemos que este complejo está formado por $A'$-módulos
proyectivos; de \sref{III.6.3.7} (extendido a un número arbitrario de
complejos) resulta que $H_\bullet(M_{\bullet\bullet\bullet})$ también es
igual a
\[
  \operatorname{Tor}_\bullet^{A'}
  (M_{\bullet\bullet\bullet}'^{(1)},\ldots,
   M_{\bullet\bullet\bullet}'^{(m)},Q'_\bullet);
\]
es, por tanto, por \sref{III.6.5.15}, el término límite de una sucesión
espectral que tiene las propiedades de regularidad requeridas (los tres
grados de $M_{\bullet\bullet\bullet}^{(i)}$ están acotados inferiormente
cuando $\mathscr P_\bullet^{(i)}$ está acotado inferiormente) y cuyo término
$E_2$ viene dado por
\[
  E_{pq}^2
  =\bigoplus_{q_1+\cdots+q_{m+1}=q}
  \operatorname{Tor}_p^{A'}
  (H_{q_1}(M_{\bullet\bullet\bullet}'^{(1)}),\ldots,
   H_{q_m}(M_{\bullet\bullet\bullet}'^{(m)}),H_{q_{m+1}}(Q'_\bullet)).
\]

\oldpage[III]{168}
Se tiene además
\[
  H_{q_i}(M_{\bullet\bullet\bullet}'^{(i)})
  =H_{q_i}(M_{\bullet\bullet\bullet}^{(i)}\otimes_A A')
  =\mathbf{Tor}_{q_i}^S
  (X^{(i)},S';\mathscr P_\bullet^{(i)},\mathscr O_{S'})
\]
en virtud de la definición de los hipertor globales, lo que da la sucesión
$(f)$ buscada. Por otra parte, se puede considerar
$M_{\bullet\bullet\bullet}'^{(i)}$ como un bicomplejo en el cual el primer
grado es la suma de los grados primero y segundo del tricomplejo
$M_{\bullet\bullet\bullet}^{(i)}$, y el segundo grado es el tercer grado de
este tricomplejo; como los módulos que forman los
$M_{\bullet\bullet\bullet}'^{(i)}$ son $A'$-módulos proyectivos, la teoría
general de la hiperhomología muestra que la homología del bicomplejo
\[
  M_{\bullet\bullet\bullet}'^{(1)}\otimes_{A'}
  M_{\bullet\bullet\bullet}'^{(2)}\otimes_{A'}\cdots\otimes_{A'}
  M_{\bullet\bullet\bullet}'^{(m)}\otimes_{A'} Q'_\bullet
\]
es canónicamente isomorfa a su hiperhomología \sref[0]{0.11.6.5}; por tanto,
es el término límite de una sucesión espectral cuyo término $E_2$ es igual a
\[
  E_{pq}^2
  =\bigoplus_{q_1+\cdots+q_{m+1}=q}
  \operatorname{Tor}_p^{A'}
  (H_{q_1}^{\mathrm{II}}(M_{\bullet\bullet\bullet}'^{(1)}),\ldots,
   H_{q_m}^{\mathrm{II}}(M_{\bullet\bullet\bullet}'^{(m)}),
   H_{q_{m+1}}^{\mathrm{II}}(Q'_\bullet)).
\]
Ahora bien, como el segundo grado de $Q'_\bullet$ se reduce a $0$, se tiene
$H_n^{\mathrm{II}}(Q'_\bullet)=0$ para $n\ne0$ y
$H_0^{\mathrm{II}}(Q'_\bullet)=Q'_\bullet$; la fórmula precedente también
se escribe
\[
  E_{pq}^2
  =\bigoplus_{q_1+\cdots+q_m=q}
  \operatorname{Tor}_p^{A'}
  (H_{q_1}^{\mathrm{II}}(M_{\bullet\bullet\bullet}'^{(1)}),\ldots,
   H_{q_m}^{\mathrm{II}}(M_{\bullet\bullet\bullet}'^{(m)}),Q'_\bullet).
\]
Además, se tiene
\[
  H_{q_i}^{\mathrm{II}}(M_{\bullet\bullet\bullet}'^{(i)})
  =H_{q_i}^{\mathrm{II}}(M_{\bullet\bullet\bullet}^{(i)}\otimes_A A')
  =\operatorname{Tor}_{q_i}^A(L_{\bullet\bullet}^{(i)},A')
\]
en virtud de \sref{III.6.3.4}; pero
$L_{-j,k}^{(i)}=C^j(\mathfrak U^{(i)},\mathscr P_k^{(i)})$, suma directa de
los $\Gamma(V,\mathscr P_k^{(i)})$, donde $V$ recorre las intersecciones
(afines) de $j+1$ conjuntos del cubrimiento $\mathfrak U^{(i)}$; si
$V'=u_i^{-1}(V)$, $V'$ es afín en $X'^{(i)}$ y de
\sref{III.6.4.1.1} resulta que se tiene
\[
  \Gamma\bigl(V',\mathscr{T}\!or_{q_i}^S
  (\mathscr P_k^{(i)},\mathscr O_{S'})\bigr)
  =\operatorname{Tor}_{q_i}^A(\Gamma(V,\mathscr P_k^{(i)}),A')
\]
de donde, para el bicomplejo
$H_{q_i}^{\mathrm{II}}(M_{\bullet\bullet\bullet}'^{(i)})$, se obtiene la
expresión
\[
  C^\bullet\bigl(\mathfrak U'^{(i)},
  \mathscr{T}\!or_{q_i}^S(\mathscr P_\bullet^{(i)},\mathscr O_{S'})\bigr),
\]
lo que da finalmente la expresión buscada para el término $E_2$ de la
sucesión $(f')$. El hecho de que esta sucesión sea birregular bajo las
condiciones indicadas se comprueba como de costumbre, teniendo en cuenta que,
si $\mathscr P_\bullet^{(i)}$ está acotado inferiormente, todos los grados de
$M_{\bullet\bullet\bullet}'^{(i)}$ están acotados inferiormente.
\end{proof}

\begin{remark}[6.9.5]
\label{III.6.9.5}
Se ve, como en \sref{III.6.7.6}, que sustituir los
$\mathscr P_\bullet^{(i)}$ y $\mathscr Q'_\bullet$ por complejos
respectivamente homotópicos a ellos no cambia las sucesiones $(e)$, $(f)$ y
$(f')$, salvo isomorfismo canónico. Además, para la sucesión $(f)$, unos
homomorfismos
$\mathscr P_\bullet^{(i)}\to\mathscr R_\bullet^{(i)}$,
$\mathscr Q'_\bullet\to\mathscr T'_\bullet$ de complejos que inducen
isomorfismos en homología
$\mathscr H_\bullet(\mathscr P_\bullet^{(i)})\xrightarrow{\sim}
\mathscr H_\bullet(\mathscr R_\bullet^{(i)})$,
$\mathscr H_\bullet(\mathscr Q'_\bullet)\xrightarrow{\sim}
\mathscr H_\bullet(\mathscr T'_\bullet)$ proporcionan un isomorfismo de
sucesiones espectrales
\[
\begin{split}
  {}^{(f)}\mathscr E
  (f_1,\ldots,f_m;
   \mathscr P_\bullet^{(1)},\ldots,\mathscr P_\bullet^{(m)},
   \mathscr Q'_\bullet)
  \xrightarrow{\sim}\;&
  {}^{(f)}\mathscr E
  (f_1,\ldots,f_m;\\
  &\mathscr R_\bullet^{(1)},\ldots,\mathscr R_\bullet^{(m)},
   \mathscr T'_\bullet);
\end{split}
\]
la demostración es la misma que la de \sref{III.6.7.6}, teniendo en cuenta
el resultado de esa sección y la regularidad de la sucesión $(f)$.
\end{remark}

\begin{corollary}[6.9.6]
\label{III.6.9.6}
Bajo las condiciones de \sref{III.6.9.1}, supongamos que:
\begin{enumerate}
  \item[$1^\circ$] Los complejos $\mathscr P_\bullet^{(i)}$ están formados
  por módulos planos sobre $S$, y los $\mathscr O_Y$-módulos
  \[
    \mathfrak{Tor}_n^S
    (f_1,\ldots,f_m;
     \mathscr P_\bullet^{(1)},\ldots,\mathscr P_\bullet^{(m)})
  \]
  son planos sobre $S$.
  \item[$2^\circ$] Los $\mathscr P_\bullet^{(i)}$ y
  $\mathscr Q'_\bullet$ están acotados inferiormente.
\end{enumerate}

\oldpage[III]{169}
Se tienen entonces, poniendo
$\mathscr P_\bullet'^{(i)}=
\mathscr P_\bullet^{(i)}\otimes_{\mathscr O_S}\mathscr O_{S'}$,
isomorfismos canónicos funtoriales
\begin{env}[6.9.6.1]
\label{III.6.9.6.1}
\[
\begin{split}
  &\mathfrak{Tor}_n^{S'}
  (f'_1,\ldots,f'_m,1_{S'};
   \mathscr P_\bullet'^{(1)},\ldots,
   \mathscr P_\bullet'^{(m)},\mathscr Q'_\bullet)
  \xrightarrow{\sim}\\
  &\qquad
  \bigoplus_{n'+n''=n}
  \mathfrak{Tor}_{n'}^S
  (f_1,\ldots,f_m;
   \mathscr P_\bullet^{(1)},\ldots,
   \mathscr P_\bullet^{(m)})
  \otimes_{\mathscr O_S}\mathscr H_{n''}(\mathscr Q'_\bullet).
\end{split}
\tag{6.9.6.1}\label{III.6.9.6.1.eq}
\]
\end{env}
En particular, para $\mathscr Q'_\bullet$ reducido a un solo término
$\mathscr F'$ de grado $0$, se tienen isomorfismos canónicos funtoriales
\begin{env}[6.9.6.2]
\label{III.6.9.6.2}
\[
\begin{split}
  &\mathfrak{Tor}_n^{S'}
  (f'_1,\ldots,f'_m,1_{S'};
   \mathscr P_\bullet'^{(1)},\ldots,
   \mathscr P_\bullet'^{(m)},\mathscr F')
  \xrightarrow{\sim}\\
  &\qquad
  \mathfrak{Tor}_n^S
  (f_1,\ldots,f_m;
   \mathscr P_\bullet^{(1)},\ldots,
   \mathscr P_\bullet^{(m)})
  \otimes_{\mathscr O_S}\mathscr F'
\end{split}
\tag{6.9.6.2}\label{III.6.9.6.2.eq}
\]
\end{env}
y, más particularmente, para $\mathscr F'=\mathscr O_{S'}$,
\begin{env}[6.9.6.3]
\label{III.6.9.6.3}
\[
\begin{split}
  &\mathfrak{Tor}_n^{S'}
  (f'_1,\ldots,f'_m;
   \mathscr P_\bullet'^{(1)},\ldots,
   \mathscr P_\bullet'^{(m)})
  \xrightarrow{\sim}\\
  &\qquad
  \mathfrak{Tor}_n^S
  (f_1,\ldots,f_m;
   \mathscr P_\bullet^{(1)},\ldots,
   \mathscr P_\bullet^{(m)})
  \otimes_{\mathscr O_S}\mathscr O_{S'}.
\end{split}
\tag{6.9.6.3}\label{III.6.9.6.3.eq}
\]
\end{env}
\end{corollary}

\begin{proof}
La hipótesis de planitud sobre los módulos que componen los
$\mathscr P_\bullet^{(i)}$ implica que los complejos
$\mathscr{T}\!or_{q_i}^S(\mathscr P_\bullet^{(i)},\mathscr O_{S'})$ son nulos
para $q_i\ne0$ por \sref{III.6.5.8}. La sucesión $(f')$ es, por tanto,
degenerada; la hipótesis $2^\circ$ implica además que es birregular
\sref{III.6.9.3}, de modo que el homomorfismo de borde
\[
\begin{split}
  &\mathfrak{Tor}_n^S
  (f_1,\ldots,f_m,1_{S'};
   \mathscr P_\bullet^{(1)},\ldots,
   \mathscr P_\bullet^{(m)},\mathscr Q'_\bullet)
  \longrightarrow {}^{(f')}\mathscr E_{n0}^2\\
  &\qquad=
  \mathfrak{Tor}_n^{S'}
  (f'_1,\ldots,f'_m,1_{S'};
   \mathscr P_\bullet'^{(1)},\ldots,
   \mathscr P_\bullet'^{(m)},\mathscr Q'_\bullet)
\end{split}
\tag{6.9.6.4}\label{III.6.9.6.4.eq}
\]
es biyectivo por \sref[0]{0.11.1.6}. La hipótesis de planitud sobre los
módulos
\[
  \mathfrak{Tor}_n^S
  (f_1,\ldots,f_m;
   \mathscr P_\bullet^{(1)},\ldots,\mathscr P_\bullet^{(m)})
\]
implica que ${}^{(e)}\mathscr E_{pq}^2=0$ para $p\ne0$ por
\sref{III.6.5.8}. La sucesión $(e)$ es, por tanto, también degenerada y, como
es birregular, el homomorfismo de borde
\[
\begin{split}
  &\mathfrak{Tor}_n^S
  (f_1,\ldots,f_m,1_{S'};
   \mathscr P_\bullet^{(1)},\ldots,
   \mathscr P_\bullet^{(m)},\mathscr Q'_\bullet)
  \longrightarrow {}^{(e)}\mathscr E_{0n}^2\\
  &\qquad=
  \bigoplus_{n'+n''=n}
  \mathfrak{Tor}_{n'}^S
  (f_1,\ldots,f_m;
   \mathscr P_\bullet^{(1)},\ldots,
   \mathscr P_\bullet^{(m)})
  \otimes_{\mathscr O_S}\mathscr H_{n''}(\mathscr Q'_\bullet)
\end{split}
\tag{6.9.6.5}\label{III.6.9.6.5.eq}
\]
es biyectivo; combinando los dos isomorfismos precedentes se obtiene el
isomorfismo \sref{III.6.9.6.1}. El isomorfismo \sref{III.6.9.6.2} se deduce
de este trivialmente, puesto que entonces se tiene
$\mathscr H_q(\mathscr Q'_\bullet)=0$ si $q\ne0$ y
$\mathscr H_0(\mathscr Q'_\bullet)=\mathscr F'$. Finalmente, el caso
$\mathscr F'=\mathscr O_{S'}$ en \sref{III.6.9.6.2}, teniendo en cuenta
\sref{III.6.7.12}, da el isomorfismo \sref{III.6.9.6.3}.
\end{proof}

\begin{corollary}[6.9.7]
\label{III.6.9.7}
Bajo las condiciones de \sref{III.6.9.1}, supongamos que $S$ y $S'$ son
afines, y que para cada $i$ se ha dado un entero $d_i$
($1\leq i\leq m$). Existe entonces un entero $N$ que solo depende de $S$, de
los $X^{(i)}$ y de los $d_i$, y que tiene la propiedad siguiente: para todo
entero $n_0$, se tienen los isomorfismos canónicos \sref{III.6.9.6.3} para
$n\leq n_0$ y para todo sistema de complejos
$\mathscr P_\bullet^{(i)}$ que satisfaga las condiciones siguientes:
\begin{enumerate}
  \item[$1^\circ$] $\mathscr P_k^{(i)}=0$ para $k<d_i$;
  \item[$2^\circ$] $\mathscr P_k^{(i)}$ es plano sobre $S$ para
  $k<n_0+N$;
  \item[$3^\circ$]
  $\mathfrak{Tor}_q^S(f_1,\ldots,f_m;
  \mathscr P_\bullet^{(1)},\ldots,\mathscr P_\bullet^{(m)})$ es plano sobre
  $S$ para $q<n_0+N$.
\end{enumerate}
\end{corollary}

\begin{proof}
Supongamos que $\mathscr P_k^{(i)}$ es plano sobre $S$ para $k<r$; entonces
$\mathscr{T}\!or_{q_i}^S(\mathscr P_k^{(i)},\mathscr O_{S'})=0$ para $k<r$
y $q_i\ne0$; calculemos
\[
\begin{split}
  \mathfrak{Tor}_p^{S'}\Bigl(
  &f'_1,\ldots,f'_m;
  \mathscr{T}\!or_{q_1}^S(\mathscr P_\bullet^{(1)},\mathscr O_{S'}),
  \ldots,\\
  &\mathscr{T}\!or_{q_m}^S(\mathscr P_\bullet^{(m)},\mathscr O_{S'})
  \Bigr).
\end{split}
\tag{6.9.7.1}\label{III.6.9.7.1.eq}
\]

\oldpage[III]{170}
por el método de \sref{III.6.6.2}, con ayuda de la imagen inversa de un
cubrimiento afín fijo $\mathfrak U^{(i)}$ de $X^{(i)}$
($1\leq i\leq m$) (independiente de $S'$ y de los
$\mathscr P_\bullet^{(i)}$); los términos $L_{jk}^{(i)}$ son nulos para
$j<-N_i$ (que solo depende de $\mathfrak U^{(i)}$); si alguno de los $q_i$
no es nulo, el complejo simple cuya homología de grado $p$ es
\eqref{III.6.9.7.1.eq} tiene términos nulos en todos los grados
\[
  <r+\sum_{j\ne i}d_j-\sum_{i=1}^m N_i,
\]
por lo que \eqref{III.6.9.7.1.eq} es nulo para $p<r-N$, designando por $N$
el mayor de los números
\[
  \sum_{i=1}^m N_i-\sum_{j\ne i}d_j.
\]
Se concluye que ${}^{(f')}\mathscr E_{pq}^2=0$ para $q\ne0$ y $p<r-N$;
como, por otra parte, ${}^{(f')}\mathscr E_{pq}^2=0$ para $q<0$, se ve que el
homomorfismo de borde \eqref{III.6.9.6.4.eq} es biyectivo para $n<r-N$
\emph{(M, XV, 5.6)} (para $\mathscr Q'_\bullet=\mathscr O_{S'}$). En segundo
lugar, si
\[
  \mathfrak{Tor}_q^S
  (f_1,\ldots,f_m;
   \mathscr P_\bullet^{(1)},\ldots,\mathscr P_\bullet^{(m)})
\]
es plano sobre $S$ para $q<r$, se tiene
${}^{(e)}\mathscr E_{pq}^2=0$ para $p\ne0$ y $q<r$; por otra parte,
${}^{(e)}\mathscr E_{pq}^2=0$ para $p<0$, por lo que el homomorfismo de borde
\eqref{III.6.9.6.5.eq} es biyectivo para $n<r$, lo que termina la
demostración.
\end{proof}

El caso de \sref{III.6.9.3} más importante en las aplicaciones es aquel en
que $m=1$, estando $\mathscr Q'_\bullet$ reducido a un solo término
$\mathscr F'$ de grado $0$; volveremos a enunciarlo en este caso con vistas
a referencias ulteriores\footnote{El caso tratado en \sref{III.6.9.8}, y en
particular las sucesiones espectrales \sref{III.6.9.8.3} y
\sref{III.6.9.8.4}, nos había sido señalado en 1957 por J.-P.~Serre.}:

\begin{proposition}[6.9.8]
\label{III.6.9.8}
Sean $S$ un preesquema, $g:S'\to S$ un morfismo,
$f:X\to Y$ un $S$-morfismo separado y cuasicompacto de $S$-preesquemas,
$\mathscr P_\bullet$ un complejo de $\mathscr O_X$-módulos cuasicoherentes
acotado inferiormente, y $\mathscr F'$ un $\mathscr O_{S'}$-módulo
cuasicoherente. Existen dos funtores espectrales birregulares en
$\mathscr P_\bullet$ y $\mathscr F'$, con valores en la categoría de los
$\mathscr O_{Y_{(S')}}$-módulos cuasicoherentes, que tienen el mismo término
límite $\mathfrak{Tor}_\bullet^S(f,1_{S'};
\mathscr P_\bullet,\mathscr F')$, y cuyos términos $E_2$ son
\begin{env}[6.9.8.1]
\label{III.6.9.8.1}
\[
  {}'\mathscr E_{pq}^2
  =\mathscr{T}\!or_p^S
  (\mathscr H^{-q}(f,\mathscr P_\bullet),\mathscr F')
  \tag{6.9.8.1}\label{III.6.9.8.1.eq}
\]
\end{env}
y
\begin{env}[6.9.8.2]
\label{III.6.9.8.2}
\[
  {}''\mathscr E_{pq}^2
  =\mathscr H^{-p}
  (f',\mathscr{T}\!or_q^S(\mathscr P_\bullet,\mathscr F')),
  \tag{6.9.8.2}\label{III.6.9.8.2.eq}
\]
\end{env}
donde $f'=f_{(S')}:X_{(S')}\to Y_{(S')}$.
\end{proposition}

\begin{proof}
Las sucesiones en cuestión también pueden obtenerse, no partiendo de
\sref{III.6.9.3}, sino de las sucesiones $(a)$ y $(b')$ de
\sref{III.6.7.3} para $X^{(1)}=X$, $Y^{(1)}=Y$,
$X^{(2)}=Y^{(2)}=S'$, $f_1=f$, $f_2=1_{S'}$. Cuando $S=S'=Y$, siendo $Y$
afín, se obtienen dos sucesiones espectrales cuyos términos $E_2$ son iguales
a
\[
  {}'\mathscr E_{pq}^2
  =\mathscr{T}\!or_p^Y
  (\mathscr H^{-q}(f,\mathscr P_\bullet),\mathscr F)
  \tag{6.9.8.3}\label{III.6.9.8.3}
\]
\[
  {}''\mathscr E_{pq}^2
  =\mathscr H^{-p}
  (f,\mathscr{T}\!or_q^Y(\mathscr P_\bullet,\mathscr F))
  \tag{6.9.8.4}\label{III.6.9.8.4}
\]
que convergen (en virtud de \sref{III.6.7.6}) a la hipercohomología
$\mathscr H^\bullet(f,\mathscr P_\bullet\otimes_Y\mathscr F)$ del funtor
$f_*$ con respecto al complejo
$\mathscr P_\bullet\otimes_Y\mathscr F$ de $\mathscr O_X$-módulos, para todo
$\mathscr O_Y$-módulo cuasicoherente y plano sobre $Y$, $\mathscr F$ (o para
todo $\mathscr O_Y$-módulo cuasicoherente $\mathscr F$ cuando
$\mathscr P_\bullet$ está formado por $\mathscr O_X$-módulos planos sobre
$Y$), las cuales son distintas de las de \sref{III.6.2.1}.
\end{proof}

\begin{corollary}[6.9.9]
\label{III.6.9.9}
Bajo las condiciones de \sref{III.6.9.8}, supongamos que el complejo
$\mathscr P_\bullet$ está acotado inferiormente, formado por módulos planos
sobre $S$, y que los $\mathscr O_Y$-módulos
$\mathscr H^n(f,\mathscr P_\bullet)$ son planos sobre $S$.

\oldpage[III]{171}
Se tienen entonces isomorfismos canónicos funtoriales
\begin{env}[6.9.9.1]
\label{III.6.9.9.1}
\[
  \mathfrak{Tor}_n^{S'}
  (f',1_{S'};\mathscr P'_\bullet,\mathscr F')
  \xrightarrow{\sim}
  \mathscr H^{-n}(f,\mathscr P_\bullet)
  \otimes_{\mathscr O_S}\mathscr F'
  \tag{6.9.9.1}\label{III.6.9.9.1.eq}
\]
\end{env}
donde
$\mathscr P'_\bullet=\mathscr P_\bullet\otimes_{\mathscr O_S}
\mathscr O_{S'}$; en particular, para $\mathscr F'=\mathscr O_{S'}$, se
tienen isomorfismos canónicos funtoriales
\begin{env}[6.9.9.2]
\label{III.6.9.9.2}
\[
  \mathscr H^n(f',\mathscr P'_\bullet)
  \xrightarrow{\sim}
  \mathscr H^n(f,\mathscr P_\bullet)
  \otimes_{\mathscr O_S}\mathscr O_{S'}.
  \tag{6.9.9.2}\label{III.6.9.9.2.eq}
\]
\end{env}
Este es el caso particular $m=1$ de \sref{III.6.9.6}. Más particularmente:
\end{corollary}

\begin{corollary}[6.9.10]
\label{III.6.9.10}
Sean $S$ un preesquema, $f:X\to Y$ un $S$-morfismo separado y cuasicompacto
de $S$-preesquemas, y $\mathscr P_\bullet$ un complejo acotado
inferiormente, formado por $\mathscr O_X$-módulos cuasicoherentes y planos
sobre $S$. Se supone además que los $\mathscr O_Y$-módulos
$\mathscr H^n(f,\mathscr P_\bullet)$ son planos sobre $S$. Para todo $s\in S$,
denotemos por $X_s$ y $Y_s$ las fibras
$X\otimes_S k(s)$, $Y\otimes_S k(s)$,
por $f_s:X_s\to Y_s$ el morfismo $f\times_S1$, y por
$\mathscr P_\bullet^s$ el complejo
$\mathscr P_\bullet\otimes_{\mathscr O_S}k(s)$ de
$\mathscr O_{X_s}$-módulos. Se tienen entonces isomorfismos canónicos
funtoriales
\begin{env}[6.9.10.1]
\label{III.6.9.10.1}
\[
  \mathscr H^n(f_s,\mathscr P_\bullet^s)
  \xrightarrow{\sim}
  \mathscr H^n(f,\mathscr P_\bullet)\otimes_{\mathscr O_S}k(s).
  \tag{6.9.10.1}\label{III.6.9.10.1.eq}
\]
\end{env}
\end{corollary}

Se tiene así, bajo hipótesis de planitud convenientes, un caso en el que la
formación de los funtores derivados $R^nf_*(\mathscr F)$ «conmuta con el paso
a las fibras», que volveremos a encontrar por otro método en el §~7.

\subsection{Estructura local de ciertos funtores cohomológicos}
\label{subsection:III.6.10}
\label{III.6.10}

\begin{proposition}[6.10.1]
\label{III.6.10.1}
Sea $S=\Spec(A)$ un esquema afín, y sea $Y^{(i)}$
$(1\leq i\leq n)$ una familia finita de $S$-esquemas afines,
planos sobre $S$. Para cada $i$, sea
\[
  f_i:X^{(i)}\longrightarrow Y^{(i)}
\]
un $S$-morfismo separado y cuasicompacto, y sea
$\sh P_\bullet^{(i)}$ un complejo acotado inferiormente de
$\sh O_{X^{(i)}}$-módulos cuasicoherentes. Sea $Y$ el producto
sobre $S$ de los esquemas $Y^{(i)}$.

Existe un complejo $\sh R_\bullet$ de $\sh O_Y$-módulos
cuasicoherentes, planos sobre $S$, con la propiedad siguiente.
Para todo $S$-esquema afín $S'$ y todo complejo acotado
inferiormente $\sh Q_\bullet'$ de $\sh O_{S'}$-módulos
cuasicoherentes, existe un isomorfismo
\begin{env}[6.10.1.1]
\label{III.6.10.1.1}
\[
  \mathbf{Tor}_\bullet^S
  \bigl(
    f_1,\ldots,f_n,1_{S'};
    \sh P_\bullet^{(1)},\ldots,\sh P_\bullet^{(n)},
    \sh Q_\bullet'
  \bigr)
  \xrightarrow{\sim}
  \sh H_\bullet
  \bigl(
    \sh R_\bullet\otimes_{\sh O_S}\sh Q_\bullet'
  \bigr),
  \tag{6.10.1.1}\label{III.6.10.1.1.eq}
\]
\end{env}
que es un isomorfismo de $\partial$-funtores en
$\sh Q_\bullet'$.

Además, para todo $S$-morfismo $u:S''\to S'$ de $S$-esquemas
afines, el diagrama
\begin{env}[6.10.1.2]
\label{III.6.10.1.2}
\[
\xymatrix@C=12pt@R=30pt{
  {\begin{gathered}
    \mathbf{Tor}_\bullet^S
    (f_1,\ldots,f_n,1_{S'};\\
    \sh P_\bullet^{(1)},\ldots,\sh P_\bullet^{(n)},
    \sh Q_\bullet')
  \end{gathered}}
  \ar[r]^-{\sim}\ar[d]
  &
  {\sh H_\bullet
    (\sh R_\bullet\otimes_{\sh O_S}\sh Q_\bullet')}
  \ar[d]
  \\
  {\begin{gathered}
    \mathbf{Tor}_\bullet^S
    (f_1,\ldots,f_n,1_{S''};\\
    \sh P_\bullet^{(1)},\ldots,\sh P_\bullet^{(n)},
    u^*(\sh Q_\bullet'))
  \end{gathered}}
  \ar[r]^-{\sim}
  &
  {\sh H_\bullet
    (\sh R_\bullet\otimes_{\sh O_S}u^*(\sh Q_\bullet'))}
}
\tag{6.10.1.2}\label{III.6.10.1.2.eq}
\]
\end{env}
es conmutativo. Aquí las flechas verticales son los
$(1_Y\times_Su)$-morfismos canónicos definidos en
\sref{III.6.7.10}.
\end{proposition}

\begin{proof}
\oldpage[III]{172}
Calculemos el hipertor mediante el método de \sref{III.6.6.2},
teniendo en cuenta la Observación \sref{III.6.6.6}. Con las
notaciones allí utilizadas, cada $L_{\bullet\bullet}^{(i)}$ es un
bicomplejo de $A_i$-módulos, donde $A_i$ denota el anillo de
$Y^{(i)}$. Tiene, por tanto, una resolución proyectiva de
Cartan--Eilenberg $M_{\bullet\bullet\bullet}^{(i)}$ (en el sentido
de \sref[0]{0.11.7.1}) por $A_i$-módulos proyectivos. Por
\sref{III.6.6.6}, el miembro izquierdo de
\eqref{III.6.10.1.1.eq} es canónicamente isomorfo a
\[
  \HH_\bullet
  \bigl(
    M_{\bullet\bullet\bullet}\otimes_AQ_\bullet'
  \bigr),
\]
donde
\[
  M_{\bullet\bullet\bullet}
  =
  M_{\bullet\bullet\bullet}^{(1)}
  \otimes_A
  M_{\bullet\bullet\bullet}^{(2)}
  \otimes_A\cdots\otimes_A
  M_{\bullet\bullet\bullet}^{(n)}
\]
y $\sh Q_\bullet'=\widetilde{Q_\bullet'}$.

Por hipótesis, los anillos $A_i$ son $A$-módulos planos. Por
tanto, los $M_{\bullet\bullet\bullet}^{(i)}$ son tricomplejos de
$A$-módulos planos \sref[0]{0.6.2.1}, y lo mismo ocurre con
$M_{\bullet\bullet\bullet}$. Además, si $B$ es el anillo de $Y$,
es decir, el producto tensorial sobre $A$ de los $A_i$, entonces
$M_{\bullet\bullet\bullet}$ es un tricomplejo de $B$-módulos.
En consecuencia, el complejo
\[
  \sh R_\bullet
  =
  \widetilde{M_{\bullet\bullet\bullet}},
\]
donde $M_{\bullet\bullet\bullet}$ se considera como complejo
simple, tiene las propiedades requeridas, como se deduce
inmediatamente de \sref{III.6.7.10}.
\end{proof}

\begin{corollary}[6.10.2]
\label{III.6.10.2}
En la Proposición \sref{III.6.10.1}, el complejo
$\sh R_\bullet$ puede tomarse acotado inferiormente. Si los
$\sh P_\bullet^{(i)}$ están acotados superiormente y los
$Y^{(i)}$ tienen dimensión cohomológica finita, entonces
$\sh R_\bullet$ puede tomarse acotado superiormente.
\end{corollary}

\begin{proof}
La primera afirmación se deduce de que los tres grados de cada
$M_{\bullet\bullet\bullet}^{(i)}$ están acotados inferiormente.
Por otra parte, si los anillos $A_i$ tienen dimensión
cohomológica finita, el tercer grado de cada
$M_{\bullet\bullet\bullet}^{(i)}$ solo toma un número finito de
valores, y por su construcción lo mismo ocurre con su primer grado
\sref{III.6.6.2}. Su segundo grado está acotado superiormente
exactamente cuando lo está el grado de $\sh P_\bullet^{(i)}$
\sref{III.6.6.2}, lo que demuestra la segunda afirmación.
\end{proof}

\begin{remarks}[6.10.3]
\label{III.6.10.3}
\textnormal{(i)}
Con las notaciones de \sref{III.6.10.1},
\[
  \sh H_\bullet
  (\sh R_\bullet\otimes_S\sh Q_\bullet')
  \simeq
  \mathbf{Tor}_\bullet(\sh R_\bullet,\sh Q_\bullet'),
\]
porque $\sh R_\bullet$ consta de $\sh O_Y$-módulos planos sobre
$S$ \sref{III.6.5.9}. Es, por tanto, según \sref{III.6.5.4}, el
objeto al que converge una sucesión espectral regular cuyo término
$E^2$ es
\begin{env}[6.10.3.1]
\label{III.6.10.3.1}
\[
  E^2_{pq}
  =
  \bigoplus_{q'+q''=q}
  \shTor_p^S
  \bigl(
    \sh H_{q'}(\sh R_\bullet),
    \sh H_{q''}(\sh Q_\bullet')
  \bigr).
  \tag{6.10.3.1}\label{III.6.10.3.1.eq}
\]
\end{env}
Esta es precisamente la sucesión espectral (e) de cambio de base
de \sref{III.6.9.3}.

\medskip
\textnormal{(ii)}
Sea $\sh R_\bullet'$ otro complejo de $\sh O_Y$-módulos
cuasicoherentes, planos sobre $S$, y sea
\[
  g:\sh R_\bullet'\longrightarrow\sh R_\bullet
\]
un homomorfismo de complejos para el cual
\[
  \sh H_\bullet(g):
  \sh H_\bullet(\sh R_\bullet')
  \longrightarrow
  \sh H_\bullet(\sh R_\bullet)
\]
es un isomorfismo. Por \sref{III.6.3.3} y
\sref{III.6.5.9}, $g$ induce un isomorfismo de
$\partial$-funtores en $\sh Q_\bullet'$,
\[
  \sh H_\bullet
  (\sh R_\bullet'\otimes_S\sh Q_\bullet')
  \xrightarrow{\sim}
  \sh H_\bullet
  (\sh R_\bullet\otimes_S\sh Q_\bullet'),
\]
tal que el diagrama
\[
\xymatrix@C=28pt@R=30pt{
  \sh H_\bullet
  (\sh R_\bullet'\otimes_S\sh Q_\bullet')
  \ar[r]^-{\sim}\ar[d]
  &
  \sh H_\bullet
  (\sh R_\bullet\otimes_S\sh Q_\bullet')
  \ar[d]
  \\
  \sh H_\bullet
  (\sh R_\bullet'\otimes_Su^*(\sh Q_\bullet'))
  \ar[r]^-{\sim}
  &
  \sh H_\bullet
  (\sh R_\bullet\otimes_Su^*(\sh Q_\bullet'))
}
\]
es conmutativo. Así, las propiedades de \sref{III.6.10.1} no
determinan por completo el complejo $\sh R_\bullet$.

\medskip
\oldpage[III]{173}
\textnormal{(iii)}
En la demostración de \sref{III.6.10.1}, los
$M_{\bullet\bullet\bullet}^{(i)}$ pueden elegirse formados por
$A_i$-módulos libres, como se deduce inmediatamente dualizando la
demostración de \sref[0]{0.11.5.2.1}. Entonces los
\[
  M_{\bullet\bullet\bullet}^{(i)}\otimes_{A_i}B
\]
constan de $B$-módulos libres; como
$M_{\bullet\bullet\bullet}$ es su producto tensorial sobre $B$,
se puede exigir además en \sref{III.6.10.1} que
$\sh R_\bullet$ esté asociado a un complejo de $B$-módulos
libres.

Además, por (M, XVII, 1.2), el tricomplejo
$M_{\bullet\bullet\bullet}$ depende funtorialmente de cada uno de
los bicomplejos $L_{\bullet\bullet}^{(i)}$ y, por tanto, de cada
$\sh P_\bullet^{(i)}$, una vez fijado un recubrimiento finito de
cada $X^{(i)}$. Aquí se entienden por «morfismos» de tricomplejos
las clases de homomorfismos módulo homotopía. Sustituir un
recubrimiento de $X^{(i)}$ por otro más fino proporciona, para
$L_{\bullet\bullet}^{(i)}$, homomorfismos definidos salvo
homotopía por \sref{III.6.6.8}; así, con la convención precedente
sobre los morfismos, $M_{\bullet\bullet\bullet}$ es un funtor en
cada $\sh P_\bullet^{(i)}$.

Precisaremos esta dependencia funtorial, en particular el
comportamiento de $\sh R_\bullet$ respecto de sucesiones exactas
cortas
\[
  0\longrightarrow\sh P_\bullet^{\prime(i)}
  \longrightarrow\sh P_\bullet^{(i)}
  \longrightarrow\sh P_\bullet^{\prime\prime(i)}
  \longrightarrow0,
\]
en el capítulo dedicado a un álgebra general de funtores
cohomológicos mencionado en \sref{III.6.1.3}.
\end{remarks}

\begin{scholium}[6.10.4]
\label{III.6.10.4}
El hecho de que $\sh R_\bullet$ conste de $\sh O_Y$-módulos
planos sobre $S$ implica inmediatamente que
\[
  \sh H_\bullet
  (\sh R_\bullet\otimes_S\sh Q_\bullet')
\]
es un funtor homológico en $\sh Q_\bullet'$; compárese el
razonamiento de \sref{III.7.7.1}. Como se señaló en
\sref{III.6.1.1}, esta propiedad motiva la introducción del
hipertor.

Pongamos
\[
  X=X^{(1)}\times_SX^{(2)}\times_S\cdots\times_SX^{(n)},
  \qquad
  f=f_1\times_Sf_2\times_S\cdots\times_Sf_n,
\]
\[
  \sh P_\bullet
  =
  \sh P_\bullet^{(1)}
  \otimes_S\sh P_\bullet^{(2)}
  \otimes_S\cdots\otimes_S\sh P_\bullet^{(n)},
\]
y
\[
  X'=X\times_SS',
  \qquad
  Y'=Y\times_SS',
  \qquad
  f'=f\times_S1_{S'}.
\]
Los problemas de cambio de base llevan a estudiar la
hipercohomología
\[
  \sh H^*_{f'}
  (\sh P_\bullet\otimes_S\sh N')
\]
como funtor del $\sh O_{S'}$-módulo cuasicoherente $\sh N'$, o,
equivalentemente, la hipercohomología
\[
  \sh H^*_f
  (\sh P_\bullet\otimes_S\sh N)
\]
como funtor del $\sh O_S$-módulo cuasicoherente $\sh N$. Cuando
los $\sh P_\bullet^{(i)}$, y por tanto $\sh P_\bullet$, son
planos sobre $S$, los resultados precedentes y
\sref{III.6.7.6} muestran que este es efectivamente un funtor
cohomológico en $\sh N$. Esto deja de ser cierto cuando no se
supone la planitud de los $\sh P_\bullet^{(i)}$, y el estudio de
\[
  \sh H^*_f(\sh P_\bullet\otimes_S\sh N)
\]
ya no puede abordarse mediante los métodos usuales del álgebra
homológica.

Utilizaremos principalmente el caso $n=1$, $Y=S$, en el que
$\sh P_\bullet$ consta de $\sh O_X$-módulos planos sobre $Y$.
En este caso se tiene el teorema siguiente.
\end{scholium}

\begin{theorem}[6.10.5]
\label{III.6.10.5}
Sea $Y=\Spec(A)$ un esquema afín noetheriano, sea
$f:X\to Y$ un morfismo propio y sea $\sh P_\bullet$ un complejo
acotado inferiormente de $\sh O_X$-módulos coherentes, planos
sobre $Y$. Existe un complejo acotado inferiormente
$\sh L_\bullet$ de $\sh O_Y$-módulos, cuyos términos son de la
forma
\[
  \sh L_i=\sh O_Y^{m_i},
\]
y un isomorfismo
\begin{env}[6.10.5.1]
\label{III.6.10.5.1}
\[
  \sh H^*
  (f,\sh P_\bullet\otimes_Y\sh Q_\bullet)
  \xrightarrow{\sim}
  \sh H_\bullet
  (\sh L_\bullet\otimes_Y\sh Q_\bullet)
  \tag{6.10.5.1}\label{III.6.10.5.1.eq}
\]
\end{env}
de $\partial$-funtores en el complejo acotado inferiormente
$\sh Q_\bullet$ de $\sh O_Y$-módulos cuasicoherentes.

Además, para todo morfismo $u:Y'\to Y$, pongamos
\[
  X'=X_{(Y')},
  \qquad
  f'=f_{(Y')},
  \qquad
  \sh P_\bullet'
  =
  \sh P_\bullet\otimes_Y\sh O_{Y'},
  \qquad
  \sh L_\bullet'=u^*(\sh L_\bullet).
\]
\oldpage[III]{174}
Entonces $\sh L_\bullet'$ es un complejo de
$\sh O_{Y'}$-módulos localmente libres de tipo finito, y existe
un isomorfismo
\begin{env}[6.10.5.2]
\label{III.6.10.5.2}
\[
  \sh H^*
  (f',\sh P_\bullet'\otimes_{Y'}\sh Q_\bullet')
  \xrightarrow{\sim}
  \sh H_\bullet
  (\sh L_\bullet'\otimes_{Y'}\sh Q_\bullet')
  \tag{6.10.5.2}\label{III.6.10.5.2.eq}
\]
\end{env}
de $\partial$-funtores en el complejo acotado inferiormente
$\sh Q_\bullet'$ de $\sh O_{Y'}$-módulos cuasicoherentes. Estos
isomorfismos pueden elegirse de modo que el diagrama
\begin{env}[6.10.5.3]
\label{III.6.10.5.3}
\[
\xymatrix@C=24pt@R=30pt{
  \sh H^*
  (f,\sh P_\bullet\otimes_Y\sh Q_\bullet)
  \ar[r]^-{\sim}\ar[d]
  &
  \sh H_\bullet
  (\sh L_\bullet\otimes_Y\sh Q_\bullet)
  \ar[d]
  \\
  \sh H^*
  (f',\sh P_\bullet'\otimes_{Y'}u^*(\sh Q_\bullet))
  \ar[r]^-{\sim}
  &
  \sh H_\bullet
  (\sh L_\bullet'\otimes_{Y'}u^*(\sh Q_\bullet))
}
\tag{6.10.5.3}\label{III.6.10.5.3.eq}
\]
\end{env}
sea conmutativo.
\end{theorem}

\begin{proof}
La Proposición \sref{III.6.10.1} da primero un complejo acotado
inferiormente $\sh R_\bullet$ \sref{III.6.10.2} de
$\sh O_Y$-módulos cuasicoherentes libres sobre $Y$
\sref{III.6.10.3}~(iii), junto con un isomorfismo
\begin{env}[6.10.5.4]
\label{III.6.10.5.4}
\[
  \sh H^*
  (f,\sh P_\bullet\otimes_Y\sh Q_\bullet)
  \xrightarrow{\sim}
  \sh H_\bullet
  (\sh R_\bullet\otimes_Y\sh Q_\bullet)
  \tag{6.10.5.4}\label{III.6.10.5.4.eq}
\]
\end{env}
de $\partial$-funtores en $\sh Q_\bullet$. Los términos de
$\sh R_\bullet$ no son necesariamente, \emph{a priori},
$\sh O_Y$-módulos de tipo finito.

Apliquemos \eqref{III.6.10.5.4.eq} cuando $\sh Q_\bullet$ es el
complejo concentrado en un grado cuyo único término es $\sh O_Y$.
Se sigue que $\sh H_\bullet(\sh R_\bullet)$ es isomorfo a
$\sh H^*(f,\sh P_\bullet)$ y, por tanto, consta de
$\sh O_Y$-módulos coherentes \sref{III.6.2.5}. Por
\sref[0]{0.11.9.2}, existe entonces un complejo acotado
inferiormente $\sh L_\bullet$ cuyos términos están asociados a
$A$-módulos libres de tipo finito, y un homomorfismo
\[
  \sh L_\bullet\longrightarrow\sh R_\bullet
\]
que induce un isomorfismo
\[
  \sh H_\bullet(\sh L_\bullet)
  \xrightarrow{\sim}
  \sh H_\bullet(\sh R_\bullet).
\]
La fórmula \eqref{III.6.10.5.1.eq} se deduce ahora de
\sref{III.6.10.3}~(ii).

Las demás afirmaciones de \sref{III.6.10.5} se deducen de
\sref{III.6.10.1} y \sref{III.6.10.3}~(ii) cuando $Y'$ es afín.
En el caso general, recubrámoslo por abiertos afines $(V_\alpha)$.
Si \eqref{III.6.10.5.2.eq} se forma sobre cada $V_\alpha$, sus
restricciones desde $V_\alpha$ y $V_\beta$ a un abierto afín
\[
  W\subset V_\alpha\cap V_\beta
\]
coinciden ambas con el isomorfismo formado sobre $W$. Esto se
deduce de la conmutatividad de \eqref{III.6.10.1.2.eq}, aplicada
a las inclusiones canónicas $W\to V_\alpha$ y $W\to V_\beta$, y
demuestra que los isomorfismos locales se pegan.
\end{proof}

\begin{remark}[6.10.6]
\label{III.6.10.6}
En los capítulos siguientes aplicaremos \sref{III.6.10.5}
principalmente cuando $\sh P_\bullet$ está concentrado en un grado
y su único término es un $\sh O_X$-módulo coherente $\sh F$ plano
sobre $Y$. Como
\[
  \sh H_n(\sh L_\bullet)
  =
  \sh H^{-n}(f,\sh F)
  =
  R^{-n}f_*(\sh F)
  \qquad\text{\sref{III.6.2.1}},
\]
los $\sh H_n(\sh L_\bullet)$ se anulan para $n>0$. Veremos más
adelante, en \sref{III.7.7.12}~(i), que $\sh L_\bullet$ puede
elegirse con términos solo en grados $\leq0$ y, por tanto, con un
número finito de términos, siempre que la exigencia de que sus
términos estén asociados a $A$-módulos libres de tipo finito se
sustituya por la de que sean localmente libres de tipo finito.
El complejo $\sh L_\bullet$ correspondiente a tal
$\sh O_X$-módulo $\sh F$ no parece tener ninguna propiedad
especial
\oldpage[III]{175}
aparte de la restricción precedente sobre sus grados. Cabe, pues,
preguntarse recíprocamente si, dado un complejo acotado
inferiormente $\sh L_\bullet$ cuyos términos son
$\sh O_Y$-módulos asociados a $A$-módulos proyectivos de tipo
finito y cuyos términos de grado $>0$ se anulan, existen un
$Y$-esquema $X$, proyectivo y plano sobre $Y$, y un
$\sh O_X$-módulo localmente libre $\sh F$ tales que exista un
isomorfismo
\[
  \sh H^*
  (f,\sh F\otimes_Y\sh Q_\bullet)
  \xrightarrow{\sim}
  \sh H_\bullet
  (\sh L_\bullet\otimes_Y\sh Q_\bullet)
\]
funtorial en $\sh Q_\bullet$. Un resultado de este tipo reduciría
por completo la teoría cohomológica de los módulos coherentes
planos sobre $Y$ en $Y$-esquemas propios a la teoría, «salvo
homotopía», de los complejos de $A$-módulos proyectivos de tipo
finito sobre un anillo noetheriano $A$.
\end{remark}

\section{Estudio del cambio de base en los funtores homológicos covariantes de módulos}
\label{section:III.7}

\subsection{Funtores de $A$-módulos}
\label{subsection:III.7.1}

\begin{env}[7.1.1]
\label{III.7.1.1}
Dado un anillo $A$, no necesariamente conmutativo, denotemos por
$\cat{Ab}_A$ la categoría de los $A$-módulos a izquierda, y
simplemente por $\cat{Ab}$ la categoría de los $\bb Z$-módulos,
que es la categoría de los grupos abelianos. Sea
\[
  T:\cat{Ab}_A\longrightarrow\cat{Ab}
\]
un funtor covariante aditivo, y sea $M$ un $(A,A)$-bimódulo.
Entonces $T(M)$ está dotado de manera natural de una estructura de
$A$-módulo a derecha.

En efecto, para cada $a\in A$, denotemos por $h_{a,M}$, o
simplemente por $h_a$, el endomorfismo
\[
  x\longmapsto xa
\]
del $A$-módulo a izquierda $M$. Por hipótesis, $T(h_a)$ es un
endomorfismo del $\bb Z$-módulo $T(M)$. Como $T$ es covariante y
aditivo, para $a,b\in A$ se tiene
\[
  T(h_{ab})
  =
  T(h_b\circ h_a)
  =
  T(h_b)\circ T(h_a)
\]
y
\[
  T(h_{a+b})
  =
  T(h_a+h_b)
  =
  T(h_a)+T(h_b).
\]
Por tanto,
\[
  (y,a)\longmapsto T(h_a)(y)
\]
define una multiplicación escalar de $A$-módulo a derecha sobre
$T(M)$. En particular, $T(A_s)$ es un $A$-módulo a derecha.
\end{env}

\begin{env}[7.1.2]
\label{III.7.1.2}
Supongamos que $A$ es conmutativo. De \sref{III.7.1.1} se deduce
que, para todo $A$-módulo $M$, el grupo $T(M)$ es naturalmente un
$A$-módulo. Además, si $u:M\to N$ es un homomorfismo
$A$-lineal, entonces, para todo $a\in A$,
\[
  u\circ h_{a,M}=h_{a,N}\circ u,
\]
y, por consiguiente,
\[
  T(u)\circ T(h_{a,M})
  =
  T(h_{a,N})\circ T(u).
\]
Así, $T(u):T(M)\to T(N)$ es $A$-lineal, y $T$ puede considerarse
como un funtor covariante aditivo de $\cat{Ab}_A$ en sí misma.

Más precisamente, esto define una equivalencia canónica entre la
categoría de los funtores covariantes aditivos
\[
  \cat{Ab}_A\longrightarrow\cat{Ab}
\]
y la categoría de los funtores covariantes $A$-lineales
\[
  T:\cat{Ab}_A\longrightarrow\cat{Ab}_A,
\]
es decir, de los funtores que satisfacen
\[
  T(h_{a,M})=h_{a,T(M)}
  \qquad(a\in A).
\]
El funtor de olvido
\[
  I:\cat{Ab}_A\longrightarrow\cat{Ab},
\]
que asocia a un $A$-módulo su $\bb Z$-módulo subyacente, es exacto
y fiel. Por tanto, los dos funtores correspondientes bajo la
equivalencia precedente tienen las mismas propiedades de exactitud.
\end{env}

\begin{env}[7.1.3]
\label{III.7.1.3}
Sigamos suponiendo que $A$ es conmutativo. Sea $B$ una
$A$-álgebra, no necesariamente conmutativa, y sea
\[
  \rho:A\longrightarrow B
\]
el homomorfismo de anillos que define la estructura de álgebra.
Este define un
\oldpage[III]{176}
funtor covariante aditivo
\[
  \rho_*:M\longmapsto M_{[\rho]}
\]
de la categoría $\cat{Ab}_B$ de los $B$-módulos a izquierda a
$\cat{Ab}_A$. Por composición se obtiene el funtor covariante
aditivo
\[
  T_{(B)}:
  \cat{Ab}_B
  \xrightarrow{\rho_*}
  \cat{Ab}_A
  \xrightarrow{T}
  \cat{Ab}.
\]
Por razones tipográficas, este funtor se denota también por
$T^{(B)}$ o por $T\otimes_AB$, y se dice que se obtiene de $T$ por
extensión de escalares de $A$ a $B$.

Si $B$ es conmutativa, $T_{(B)}$ puede considerarse, por supuesto,
como un funtor de $\cat{Ab}_B$ en sí misma
\sref{III.7.1.2}. Si $B$ es conmutativa y $C$ es una $B$-álgebra,
entonces
\[
  T_{(C)}=(T_{(B)})_{(C)}.
\]
La extensión de escalares es manifiestamente funtorial y aditiva
en $T$. Además, si $T$ conmuta con los límites inductivos o con las
sumas directas, lo mismo ocurre con $T_{(B)}$; y si $T$ es exacto
por la izquierda, exacto por la derecha o exacto,
respectivamente, lo mismo vale para $T_{(B)}$. En efecto,
$\rho_*$ es exacto y conmuta con los límites inductivos y las sumas
directas.
\end{env}

\begin{env}[7.1.4]
\label{III.7.1.4}
Supongamos de nuevo que $A$ es conmutativo, y sea
\[
  T:\cat{Ab}_A\longrightarrow\cat{Ab}_A
\]
un funtor covariante aditivo $A$-lineal que conmuta con los
límites inductivos. Para toda parte multiplicativa $S$ de $A$ y
todo $A$-módulo $M$, existe un isomorfismo canónico funtorial de
$A$-módulos
\begin{env}[7.1.4.1]
\label{III.7.1.4.1}
\[
  T(S^{-1}M)
  \xrightarrow{\sim}
  S^{-1}T(M).
  \tag{7.1.4.1}\label{III.7.1.4.1.eq}
\]
\end{env}

Supongamos primero que $S$ es el conjunto de las potencias $f^n$,
$n\geq0$, de un elemento $f\in A$. Entonces
\[
  M_f=\varinjlim M_n,
\]
donde $(M_n,\vphi_{nm})$ es el sistema inductivo de $A$-módulos
con $M_n=M$ y
\[
  \vphi_{nm}:z\longmapsto f^{\,n-m}z
  \qquad\text{\sref[0]{0.1.6.1}}.
\]
La fórmula \eqref{III.7.1.4.1.eq} se deduce en este caso de la
hipótesis sobre $T$. Para $S$ arbitraria, el módulo $S^{-1}M$ es
el límite inductivo de los $M_f$, para $f\in S$
\sref[0]{0.1.4.5}, y se aplica el mismo razonamiento.

La funtorialidad de \eqref{III.7.1.4.1.eq} muestra además que se
trata de un isomorfismo de $S^{-1}A$-módulos. Así, salvo
isomorfismo canónico, se puede escribir
\begin{env}[7.1.4.2]
\label{III.7.1.4.2}
\[
  T_{(S^{-1}A)}(S^{-1}M)
  =
  S^{-1}T(M)
  =
  T(S^{-1}M).
  \tag{7.1.4.2}\label{III.7.1.4.2.eq}
\]
\end{env}
Si $S=A-\mathfrak p$ es el complementario de un ideal primo
$\mathfrak p$ de $A$, se escribe $T_{\mathfrak p}$ en lugar de
$T_{(A_{\mathfrak p})}$.
\end{env}

\begin{proposition}[7.1.5]
\label{III.7.1.5}
Bajo las hipótesis de \sref{III.7.1.4}, supongamos que
$T_{\mathfrak m}$ es exacto por la izquierda, respectivamente
exacto por la derecha o exacto, para todo ideal maximal
$\mathfrak m$ de $A$. Entonces $T$ es exacto por la izquierda,
respectivamente exacto por la derecha o exacto.
\end{proposition}

\begin{proof}
Si dos submódulos $N$ y $P$ de un $A$-módulo $M$ satisfacen
\[
  N_{\mathfrak m}=P_{\mathfrak m}
\]
para todo ideal maximal $\mathfrak m$ de $A$, entonces $N=P$
(Bourbaki, \emph{Algèbre commutative}, cap.~II, \S~3,
n.º~3, teo.~1).
\end{proof}

\subsection{Caracterizaciones del funtor producto tensorial}
\label{subsection:III.7.2}

\begin{env}[7.2.1]
\label{III.7.2.1}
Sea $A$ un anillo, no necesariamente conmutativo; sea $M$
(respectivamente $N$) un $A$-módulo a izquierda
(respectivamente a derecha), y sea $P$ un $\bb Z$-módulo.
Recordemos que dar un $\bb Z$-homomorfismo
\[
  v:N\otimes_A M\longrightarrow P
\]
equivale a dar una aplicación $\bb Z$-bilineal
\[
  u:N\times M\longrightarrow P
\]
tal que
\[
  u(ta,x)=u(t,ax)
  \qquad(a\in A,\ t\in N,\ x\in M),
\]
estando relacionadas ambas aplicaciones por
\[
  v(t\otimes x)=u(t,x).
\]
Por otra parte, dar $u$ equivale a dar un
\oldpage[III]{177}
$\bb Z$-homomorfismo
\[
  x\longmapsto f_x
\]
de $M$ en $\Hom_{\bb Z}(N,P)$ tal que
\[
  f_{ax}(t)=f_x(ta)
  \qquad(a\in A,\ t\in N,\ x\in M),
\]
estando relacionadas ambas aplicaciones por $u(t,x)=f_x(t)$.
\end{env}

\begin{env}[7.2.2]
\label{III.7.2.2}
Sea
\[
  T:\cat{Ab}_A\longrightarrow\cat{Ab}
\]
un funtor covariante aditivo. Para todo $A$-módulo a izquierda
$M$ vamos a definir un homomorfismo canónico de $\bb Z$-módulos,
funtorial en $M$,
\begin{env}[7.2.2.1]
\label{III.7.2.2.1}
\[
  t_M:T(A_s)\otimes_A M\longrightarrow T(M).
  \tag{7.2.2.1}\label{III.7.2.2.1.eq}
\]
\end{env}
Por \sref{III.7.2.1}, basta definir un $\bb Z$-homomorfismo
\[
  x\longmapsto t'_M(x)
\]
de $M$ en $\Hom_{\bb Z}(T(A_s),T(M))$ tal que
\[
  t'_M(ax)(y)=t'_M(x)(ya)
  \qquad(a\in A,\ x\in M,\ y\in T(A_s)).
\]
Para ello, obsérvese que $\Hom_{\bb Z}(T(A_s),T(M))$ es
canónicamente un $A$-módulo a izquierda, en virtud de la
estructura de $A$-módulo a derecha de $T(A_s)$, con
multiplicación escalar definida por
\[
  (a\cdot w)(y)=w(ya)
\]
para $a\in A$, $w\in\Hom_{\bb Z}(T(A_s),T(M))$ e
$y\in T(A_s)$. Definamos $t'_M$ como el compuesto de los dos
homomorfismos canónicos
\[
  M
  \xrightarrow{\sim}
  \Hom_A(A_s,M)
  \xrightarrow{T}
  \Hom_{\bb Z}(T(A_s),T(M)).
\]

La segunda flecha envía $w$ a $T(w)$; la primera es el
isomorfismo canónico de $A$-módulos que envía $x$ a la aplicación
$\theta_x:A_s\to M$ definida por
\[
  \theta_x(\xi)=\xi x
  \qquad(\xi\in A,\ x\in M).
\]
Como $\theta_{ax}=\theta_x\circ h_a$, se tiene
\[
  T(\theta_{ax})
  =
  T(\theta_x\circ h_a)
  =
  T(\theta_x)\circ T(h_a),
\]
y, por consiguiente, para $y\in T(A_s)$,
\[
  T(\theta_{ax})(y)
  =
  T(\theta_x)\bigl(T(h_a)(y)\bigr)
  =
  T(\theta_x)(ya),
\]
por la definición de la estructura de $A$-módulo a derecha sobre
$T(A_s)$. Esto prueba la existencia de $t_M$.

Es inmediato que este homomorfismo es funtorial en $M$: para todo
homomorfismo $w:M\to M'$ de $A$-módulos a izquierda, el diagrama
\begin{env}[7.2.2.2]
\label{III.7.2.2.2}
\[
\xymatrix{
  T(A_s)\otimes_A M
    \ar[r]^-{t_M}\ar[d]_{1\otimes w}
  & T(M)\ar[d]^{T(w)}
  \\
  T(A_s)\otimes_A M'
    \ar[r]_-{t_{M'}}
  & T(M')
}
\tag{7.2.2.2}\label{III.7.2.2.2.eq}
\]
\end{env}
es conmutativo. La funtorialidad de
\eqref{III.7.2.2.1.eq} muestra asimismo que, cuando $A$ es
conmutativo, $t_M$ es un homomorfismo de $A$-módulos
\sref{III.7.1.2}.
\end{env}

\begin{env}[7.2.3]
\label{III.7.2.3}
Cuando $A$ es conmutativo, se puede definir más generalmente un
homomorfismo canónico de $A$-módulos
\begin{env}[7.2.3.1]
\label{III.7.2.3.1}
\[
  T(N)\otimes_A M
  \longrightarrow
  T(N\otimes_A M).
  \tag{7.2.3.1}\label{III.7.2.3.1.eq}
\]
\end{env}
para todo $A$-módulo $N$: en la construcción de
\sref{III.7.2.2}, sustitúyase $\theta_x$ por el homomorfismo de
$A$-módulos
\[
  N\longrightarrow N\otimes_A M,
  \qquad
  y\longmapsto y\otimes x.
\]
Es inmediato que este homomorfismo es funtorial tanto en $M$ como
en $N$.

\oldpage[III]{178}
En particular, si $B$ es una $A$-álgebra, no necesariamente
conmutativa, existe un homomorfismo, funtorial en $M$,
\begin{env}[7.2.3.2]
\label{III.7.2.3.2}
\[
  (T(M))_{(B)}
  =
  T(M)\otimes_A B
  \longrightarrow
  T(M\otimes_A B)
  =
  T_{(B)}(M_{(B)}).
  \tag{7.2.3.2}\label{III.7.2.3.2.eq}
\]
\end{env}
Por la funtorialidad de \eqref{III.7.2.3.1.eq} en $M$, este es un
homomorfismo de $B$-módulos.

Se tiene además un diagrama conmutativo
\begin{env}[7.2.3.3]
\label{III.7.2.3.3}
\[
\xymatrix{
  T(A)\otimes_A M
    \ar[r]^-{t_M}\ar[d]
  & T(M)\ar[d]
  \\
  T_{(B)}(B_s)\otimes_B M_{(B)}
    \ar[r]_-{t_{M_{(B)}}}
  & T_{(B)}(M_{(B)})
}
\tag{7.2.3.3}\label{III.7.2.3.3.eq}
\]
\end{env}
La flecha vertical derecha es el compuesto
\[
  T(M)
  \longrightarrow
  T(M)\otimes_A B
  \longrightarrow
  T(M\otimes_A B)
  =
  T_{(B)}(M_{(B)})
\]
del homomorfismo canónico y \eqref{III.7.2.3.2.eq}. La flecha
vertical izquierda de \eqref{III.7.2.3.3.eq} es el homomorfismo
\[
  T(A)\otimes_A M
  \longrightarrow
  T_{(B)}(B_s)\otimes_B(B\otimes_A M)
  =
  T_{(B)}(B_s)\otimes_A M,
\]
donde
\[
  T(A)\longrightarrow T_{(B)}(B_s)=T(B)
\]
es $T(\rho)$, considerando $\rho$ como homomorfismo de
$A$-módulos $A\to B_{[\rho]}$.
\end{env}

\begin{lemma}[7.2.4]
\label{III.7.2.4}
Si $T$ es un funtor covariante aditivo de $\cat{Ab}_A$ en
$\cat{Ab}$ que conmuta con las sumas directas, entonces el
homomorfismo canónico $t_L$ de \eqref{III.7.2.2.1.eq} es un
isomorfismo para todo $A$-módulo libre $L$.
\end{lemma}

\begin{proof}
Escribamos
\[
  L=\bigoplus_{\alpha\in I}L_\alpha,
\]
donde cada $L_\alpha$ es isomorfo a $A_s$. La definición de $t_M$
en \sref{III.7.2.2} muestra que
\[
  t_L=\bigoplus_{\alpha\in I}t_{L_\alpha},
\]
porque
\[
  T:
  \Hom_A(A_s,L)
  \longrightarrow
  \Hom_{\bb Z}(T(A_s),T(L))
\]
es la suma directa de las aplicaciones $\bb Z$-lineales
\[
  T_\alpha:
  \Hom_A(A_s,L_\alpha)
  \longrightarrow
  \Hom_{\bb Z}(T(A_s),T(L_\alpha)),
\]
por la hipótesis sobre $T$. Solo queda, por tanto, demostrar el
lema para $L=A_s$; pero entonces $t_L$ es precisamente el
isomorfismo canónico
\[
  T(A_s)\otimes_A A_s
  \xrightarrow{\sim}
  T(A_s),
\]
válido para todo $A$-módulo a derecha.
\end{proof}

\begin{proposition}[7.2.5]
\label{III.7.2.5}
Sea $T$ un funtor covariante aditivo de $\cat{Ab}_A$ en
$\cat{Ab}$ que conmuta con las sumas directas. Las condiciones
siguientes son equivalentes:
\begin{enumerate}
  \item[(a)] $T$ es exacto por la derecha;
  \item[(b)] el homomorfismo canónico $t_M$ de
    \eqref{III.7.2.2.1.eq} es un isomorfismo para todo $A$-módulo
    a izquierda $M$;
  \item[($\mathrm{b}'$)] $T$ es semiexacto y $t_M$ es
    sobreyectivo para todo $A$-módulo a izquierda $M$;
  \item[(c)] $T$ es isomorfo a un funtor de $M$ de la forma
    $N\otimes_A M$, donde $N$ es un $A$-módulo a derecha.
\end{enumerate}
\end{proposition}

\begin{proof}
Es claro que (b) implica (c), y que (c) implica (a). Demostremos
que (a) implica (b).
\oldpage[III]{179}
Pongamos
\[
  T'(M)=T(A_s)\otimes_A M
\]
para todo $A$-módulo a izquierda $M$. Existe una sucesión exacta
\[
  L'\longrightarrow L\longrightarrow M\longrightarrow0,
\]
donde $L$ y $L'$ son $A$-módulos libres a izquierda. Como $T$ y
$T'$ son exactos por la derecha, existe un diagrama conmutativo
\[
\xymatrix{
  T'(L')\ar[r]\ar[d]^{t_{L'}}
  & T'(L)\ar[r]\ar[d]^{t_L}
  & T'(M)\ar[r]\ar[d]^{t_M}
  & 0
  \\
  T(L')\ar[r]
  & T(L)\ar[r]
  & T(M)\ar[r]
  & 0
}
\]
cuyas filas son exactas. Por \sref{III.7.2.4}, $t_L$ y
$t_{L'}$ son isomorfismos; por el lema de los cinco, también lo es
$t_M$.

Es claro que (b) implica ($\mathrm{b}'$). Para mostrar que
($\mathrm{b}'$) implica (a), basta demostrar el lema siguiente.
\end{proof}

\begin{lemma}[7.2.5.1]
\label{III.7.2.5.1}
Sean $\cat K$ y $\cat K'$ categorías abelianas, sean $F$ y $G$
funtores covariantes aditivos de $\cat K$ en $\cat K'$, y sea
\[
  f:F\longrightarrow G
\]
una transformación natural \emph{(T, I, 1.2)} tal que
\[
  f_E:F(E)\longrightarrow G(E)
\]
es un epimorfismo para todo objeto $E$ de $\cat K$. Si $F$ es
exacto por la derecha y $G$ es semiexacto, entonces $G$ es exacto
por la derecha.
\end{lemma}

\begin{proof}
Basta mostrar que, para todo epimorfismo $v:E'\to E$ en $\cat K$,
\[
  G(v):G(E')\longrightarrow G(E)
\]
es un epimorfismo. Existe un diagrama conmutativo
\[
\xymatrix{
  F(E')\ar[r]^{F(v)}\ar[d]_{f_{E'}}
  & F(E)\ar[d]^{f_E}
  \\
  G(E')\ar[r]_{G(v)}
  & G(E)
}
\]
Aquí $F(v)$, $f_{E'}$ y $f_E$ son epimorfismos; por consiguiente,
$G(v)$ también es un epimorfismo.
\end{proof}

\begin{remark}[7.2.6]
\label{III.7.2.6}
Para todo $A$-módulo a derecha $N$, pongamos
\[
  T_N(M)=N\otimes_A M
\]
para todo $A$-módulo a izquierda $M$. Así, $T_N$ es un funtor
covariante aditivo de $\cat{Ab}_A$ en $\cat{Ab}$, es exacto por
la derecha y conmuta con las sumas directas. Si se identifica
canónicamente $T_N(A_s)$ con $N$, se comprueba inmediatamente que
el homomorfismo correspondiente \eqref{III.7.2.2.1.eq} es la
identidad. Se sigue que el $A$-módulo a derecha $N$ de
\sref{III.7.2.5}(c) está determinado salvo isomorfismo único y es
canónicamente isomorfo a $T(A_s)$.

Equivalentemente, las correspondencias
\[
  T\longmapsto T(A_s)
  \qquad\text{y}\qquad
  N\longmapsto T_N
\]
definen una equivalencia \emph{(T, I, 1.2)} entre la categoría de
los $A$-módulos a derecha y la categoría de los funtores
covariantes aditivos
\[
  \cat{Ab}_A\longrightarrow\cat{Ab}
\]
que son exactos por la derecha y conmutan con las sumas directas.
\end{remark}

\begin{proposition}[7.2.7]
\label{III.7.2.7}
Sea $A$ un anillo artiniano a izquierda cuyo cociente por su
radical $\mathfrak m$ es un cuerpo $k$, y sea $T$ un funtor
covariante aditivo de $\cat{Ab}_A$ en $\cat{Ab}$ que conmuta con
las sumas directas. Las condiciones de \sref{III.7.2.5} son
también equivalentes a la condición siguiente:
\begin{enumerate}
  \item[(d)] $T$ es semiexacto y el homomorfismo
    \[
      T(\varepsilon):T(A_s)\longrightarrow T(k)
    \]
    inducido por el homomorfismo canónico
    $\varepsilon:A_s\to k$ es sobreyectivo.
\end{enumerate}
\end{proposition}

\begin{proof}
\oldpage[III]{180}
Es claro que la condición ($\mathrm{b}'$) de
\sref{III.7.2.5} implica (d). Demostremos que (d) implica
($\mathrm{b}'$). Existe un entero $n$ tal que
$\mathfrak m^n=0$. Para todo $A$-módulo $M$, pongamos
\[
  M_h=\mathfrak m^hM.
\]
Demostraremos, por inducción descendente sobre $h$, que $t_{M_h}$
es sobreyectivo. La afirmación es clara para $h=n$. Si $h<n$,
existe una sucesión exacta
\[
  0\longrightarrow M_{h+1}
  \longrightarrow M_h
  \longrightarrow M_h/M_{h+1}
  \longrightarrow0,
\]
y la hipótesis de inducción dice que $t_{M_{h+1}}$ es
sobreyectivo. Además, $M_h/M_{h+1}$ es anulado por $\mathfrak m$
y, por tanto, es un $(A/\mathfrak m)$-módulo; en otras palabras,
es una suma directa de $A$-módulos isomorfos a $k$. Como $T$
conmuta con las sumas directas, para demostrar que
$t_{M_h/M_{h+1}}$ es sobreyectivo basta demostrar que $t_k$ lo es.
El diagrama
\[
\xymatrix{
  T(A_s)\otimes_A A_s
    \ar[r]^-{t_{A_s}}\ar[d]_{1\otimes\varepsilon}
  & T(A_s)\ar[d]^{T(\varepsilon)}
  \\
  T(A_s)\otimes_A k
    \ar[r]_-{t_k}
  & T(k)
}
\]
es conmutativo. La condición (d), junto con
\sref{III.7.2.4}, muestra entonces que $t_k$ es sobreyectivo.

Para terminar la demostración, queda mostrar que si
\[
  0\longrightarrow M'
  \longrightarrow M
  \xrightarrow{v}M''
  \longrightarrow0
\]
es una sucesión exacta de $A$-módulos para la que $t_{M'}$ y
$t_{M''}$ son sobreyectivos, entonces $t_M$ es sobreyectivo.
Existe un diagrama conmutativo
\[
\xymatrix{
  T'(M')\ar[r]\ar[d]_{t_{M'}}
  & T'(M)\ar[r]\ar[d]_{t_M}
  & T'(M'')\ar[r]\ar[d]_{t_{M''}}
  & 0\ar[d]
  \\
  T(M')\ar[r]
  & T(M)\ar[r]_{T(v)}
  & T(M'')\ar[r]
  & \Coker(T(v))
}
\]
cuyas dos filas son exactas, por la hipótesis de que $T$ es
semiexacto. La hipótesis de inducción dice que $t_{M'}$ y
$t_{M''}$ son epimorfismos, mientras que la última flecha vertical
es un monomorfismo. El lema de los cinco \emph{(M, I, 1.1)}
muestra entonces que $t_M$ es un epimorfismo.
\end{proof}

\subsection{Criterios de exactitud de los funtores homológicos de módulos}
\label{subsection:III.7.3}

\begin{proposition}[7.3.1]
\label{III.7.3.1}
Sea $A$ un anillo, no necesariamente conmutativo, y sea
$T_\bullet$ un funtor homológico covariante
\emph{(T, II, 2.1)} de $\cat{Ab}_A$ en $\cat{Ab}$ que
conmuta con las sumas directas. Sea $p$ un entero tal que
$T_p$ y $T_{p-1}$ estén definidos. Las condiciones siguientes son
equivalentes:
\begin{enumerate}
  \item[(a)] $T_p$ es exacto por la derecha;
  \item[(b)] $T_{p-1}$ es exacto por la izquierda;

\oldpage[III]{181}

  \item[(c)] para todo $A$-módulo a izquierda $M$, el homomorfismo
    canónico y funtorial \eqref{III.7.2.2.1.eq}
    \[
      T_p(A_s)\otimes_A M
      \longrightarrow
      T_p(M)
      \tag{7.3.1.1}\label{III.7.3.1.1.eq}
    \]
    es un isomorfismo;
  \item[(d)] para todo $A$-módulo a izquierda $M$, el homomorfismo
    \eqref{III.7.3.1.1.eq} es un epimorfismo;
  \item[(e)] $T_p$ es isomorfo a un funtor de la forma
    \[
      M\longmapsto N\otimes_A M,
    \]
    donde $N$ es un $A$-módulo a derecha.
\end{enumerate}
Si, además, se verifican las condiciones de \sref{III.7.2.7} sobre
$A$ y $\mathfrak m$, las condiciones precedentes son también
equivalentes a
\begin{enumerate}
  \item[(f)] el homomorfismo canónico
    \[
      T_p(\varepsilon):T_p(A_s)\longrightarrow T_p(k)
    \]
    es un epimorfismo.
\end{enumerate}
\end{proposition}

\begin{proof}
Por definición de funtor homológico, $T_i$ es semiexacto para todo
$i$ para el cual esté definido; además, para toda sucesión exacta
\[
  0\longrightarrow M'
  \xrightarrow{u}M
  \xrightarrow{v}M''
  \longrightarrow0,
\]
se tiene
\[
  \Ker\bigl(T_{i-1}(u)\bigr)
  =
  \Coker\bigl(T_i(v)\bigr).
\]
Por tanto, es claro que (a) y (b) son equivalentes, y las demás
afirmaciones resultan inmediatamente de \sref{III.7.2.5} y
\sref{III.7.2.7}.
\end{proof}

\begin{corollary}[7.3.2]
\label{III.7.3.2}
Sea $A$ un anillo conmutativo. Con las notaciones de
\sref{III.7.3.1}, supongamos que $T_p$ es exacto por la derecha. Si
$f\in A$ no pertenece al anulador de ningún elemento no nulo de un
$A$-módulo $M$, entonces tampoco pertenece al anulador de ningún
elemento no nulo de $T_{p-1}(M)$. En particular, si $A$ es un
dominio entero, el $A$-módulo $T_{p-1}(A)$ es sin torsión.
\end{corollary}

\begin{proof}
Si $h_f$ designa la homotecia $x\mapsto fx$ de $M$, la hipótesis
significa que $h_f$ es inyectiva. La condición (b) de
\sref{III.7.3.1} muestra entonces que $T_{p-1}(h_f)$ también es
inyectiva.
\end{proof}

\begin{proposition}[7.3.3]
\label{III.7.3.3}
Sea $A$ un anillo, y sea $T_\bullet$ un funtor homológico
covariante de $\cat{Ab}_A$ en $\cat{Ab}$ que conmuta con las
sumas directas. Sea $p$ un entero tal que $T_{p-1}$, $T_p$ y
$T_{p+1}$ estén definidos. Las condiciones siguientes son
equivalentes:
\begin{enumerate}
  \item[(a)] $T_p$ es exacto;
  \item[(b)] $T_{p+1}$ y $T_p$ son exactos por la derecha;
  \item[(c)] $T_p$ y $T_{p-1}$ son exactos por la izquierda;
  \item[(d)] $T_{p+1}$ es exacto por la derecha y $T_{p-1}$ es
    exacto por la izquierda;
  \item[(e)] para todo $A$-módulo $M$, los homomorfismos canónicos
    \[
      T_i(A_s)\otimes_A M
      \longrightarrow
      T_i(M)
      \tag{7.3.3.1}\label{III.7.3.3.1.eq}
    \]
    son isomorfismos para $i=p$ e $i=p+1$;
  \item[($\mathrm{e}'$)] para todo $A$-módulo $M$, los
    homomorfismos \eqref{III.7.3.3.1.eq} son epimorfismos para
    $i=p$ e $i=p+1$;
  \item[(f)] para todo $A$-módulo $M$, el homomorfismo
    \eqref{III.7.3.3.1.eq} es un isomorfismo para $i=p$, y
    $T_p(A_s)$ es plano como $A$-módulo a derecha;
  \item[($\mathrm{f}'$)] para todo $A$-módulo $M$, el
    homomorfismo \eqref{III.7.3.3.1.eq} es un epimorfismo para
    $i=p$, y $T_p(A_s)$ es plano como $A$-módulo a derecha.
\end{enumerate}
\end{proposition}

\begin{proof}
La equivalencia de (a), (b), (c) y (d) resulta de la equivalencia
de las condiciones (a) y (b) de \sref{III.7.3.1}. La equivalencia
de (b), (e) y ($\mathrm{e}'$) resulta de la equivalencia de (a),
(c) y (d) en \sref{III.7.3.1}. Finalmente, decir que $T_p(A_s)$
es plano significa que el funtor
\[
  M\longmapsto T_p(A_s)\otimes_A M
\]
es exacto por la izquierda; la equivalencia de (a), (f) y
($\mathrm{f}'$) resulta de nuevo de la equivalencia de (a), (c) y
(d) en \sref{III.7.3.1}.
\end{proof}

\oldpage[III]{182}

\begin{corollary}[7.3.4]
\label{III.7.3.4}
Supongamos que $A$ es conmutativo, que $T_p$ es exacto y que
$T_p(A)$ es un $A$-módulo de presentación finita. Entonces la
función
\[
  x\longmapsto
  \operatorname{rank}_{\kres(x)}
  \bigl(T_p(\kres(x))\bigr)
\]
es localmente constante en $X=\Spec(A)$ y, por tanto, es constante
si $\Spec(A)$ es conexo.
\end{corollary}

\begin{proof}
Por la condición (f) de \sref{III.7.3.3}, $T_p(A)$ es un
$A$-módulo plano. Por consiguiente, es proyectivo de tipo finito, de
modo que $\widetilde{T_p(A)}$ es un $\sh O_X$-módulo localmente
libre (Bourbaki, \emph{Alg\`ebre commutative}, chap.~II, \S~5,
no.~2, th.~1). Además,
\[
  T_p(\kres(x))
  =
  T_p(A)\otimes_A\kres(x)
\]
por la condición (e) de \sref{III.7.3.3}, y se sabe que el rango en
$x$ del $A$-módulo $T_p(A)$ es localmente constante
\emph{(loc.\ cit.)}. Esto prueba el corolario.
\end{proof}

\begin{proposition}[7.3.5]
\label{III.7.3.5}
Supongamos que $A$ es un anillo artiniano a izquierda cuyo cociente
por su radical $\mathfrak m$ es un cuerpo $k$. Las condiciones de
\sref{III.7.3.3} son también equivalentes a cada una de las
siguientes:
\begin{enumerate}
  \item[(g)] el homomorfismo canónico
    \[
      T_i(\varepsilon):T_i(A_s)\longrightarrow T_i(k)
    \]
    es un epimorfismo para $i=p$ e $i=p+1$;
  \item[(h)] $T_p(\varepsilon)$ es un epimorfismo y
    $T_p(A_s)$ es plano como $A$-módulo a derecha o,
    equivalentemente, es un $A$-módulo libre a derecha (Bourbaki,
    \emph{Alg\`ebre commutative}, chap.~II, \S~3, no.~2,
    cor.~2 de la prop.~5).
\end{enumerate}
Supongamos además que $A$ es conmutativo y que el $A$-módulo
$T_p(k)$ tiene longitud finita $d$. Las condiciones precedentes son
entonces también equivalentes a cada una de las siguientes:
\begin{enumerate}
  \item[(i)] para todo $A$-módulo $M$ de longitud finita,
    \[
      \operatorname{length}\bigl(T_p(M)\bigr)
      =
      d\,\operatorname{length}(M);
      \tag{7.3.5.1}\label{III.7.3.5.1.eq}
    \]
  \item[(j)]
    \[
      \operatorname{length}\bigl(T_p(A)\bigr)
      =
      d\,\operatorname{length}(A).
      \tag{7.3.5.2}\label{III.7.3.5.2.eq}
    \]
\end{enumerate}
\end{proposition}

La equivalencia de (g) y (h) con las condiciones de
\sref{III.7.3.3} resulta inmediatamente de \sref{III.7.2.7}. Para
las demás afirmaciones utilizaremos el lema siguiente.

\begin{lemma}[7.3.5.3]
\label{III.7.3.5.3}
Sean $\cat{K}$ y $\cat{K}'$ dos categorías abelianas, y sea
\[
  F:\cat{K}\longrightarrow\cat{K}'
\]
un funtor covariante aditivo. Supongamos que $F$ es semiexacto y que
$F(S)$ tiene longitud finita en $\cat{K}'$ para todo objeto simple
$S$ de $\cat{K}$. Entonces $F(E)$ tiene longitud finita en
$\cat{K}'$ para todo objeto $E$ de longitud finita en $\cat{K}$.
Para toda sucesión exacta
\[
  0\longrightarrow E'
  \xrightarrow{u}E
  \xrightarrow{v}E''
  \longrightarrow0
\]
de objetos de longitud finita de $\cat{K}$, se tiene
\[
  \operatorname{length}F(E)
  \leq
  \operatorname{length}F(E')
  +
  \operatorname{length}F(E''),
  \tag{7.3.5.4}\label{III.7.3.5.4.eq}
\]
y se da la igualdad en \eqref{III.7.3.5.4.eq} si y solo si
\[
  0\longrightarrow F(E')
  \longrightarrow F(E)
  \longrightarrow F(E'')
  \longrightarrow0
\]
es exacta.
\end{lemma}

\begin{proof}
La sucesión
\[
  F(E')\xrightarrow{F(u)}F(E)\xrightarrow{F(v)}F(E'')
\]
es exacta por hipótesis. Si $F(E')$ y $F(E'')$ tienen longitud
finita, también la tienen $\Im(F(u))$ y $\Im(F(v))$. Puesto que
$\Ker(F(v))=\Im(F(u))$, el objeto $F(E)$ tiene longitud finita, y
\[
  \begin{split}
  \operatorname{length}F(E)
  &=
  \operatorname{length}\Im(F(u))
  +
  \operatorname{length}\Im(F(v))
  \\
  &\leq
  \operatorname{length}F(E')
  +
  \operatorname{length}F(E'').
  \end{split}
  \tag{7.3.5.5}\label{III.7.3.5.5.eq}
\]

\oldpage[III]{183}

Una inducción sobre la longitud de $E$ prueba la primera
afirmación. Además, la igualdad en \eqref{III.7.3.5.5.eq} solo
puede darse cuando
\[
  \operatorname{length}\Im(F(u))
  =
  \operatorname{length}F(E')
\]
y
\[
  \operatorname{length}\Im(F(v))
  =
  \operatorname{length}F(E'').
\]
La primera igualdad equivale a
$\operatorname{length}\Ker(F(u))=0$, y por tanto a $\Ker(F(u))=0$;
la segunda equivale a $\operatorname{length}\Coker(F(v))=0$, y por
tanto a $\Coker(F(v))=0$.
\end{proof}

Sea ahora $M$ un $A$-módulo de longitud finita, con $A$
conmutativo. Los cocientes de una serie de Jordan--Hölder de $M$
son necesariamente isomorfos al $A$-módulo $k$. Una inducción
sobre la longitud de $M$, usando \eqref{III.7.3.5.4.eq}, da por
tanto
\[
  \operatorname{length}T_p(M)
  \leq
  d\,\operatorname{length}(M).
  \tag{7.3.5.6}\label{III.7.3.5.6.eq}
\]
Se sigue también de \sref{III.7.3.5.3} que, si $T_p$ es exacto,
entonces se da la igualdad \eqref{III.7.3.5.1.eq}. Así, la condición
(a) de \sref{III.7.3.3} implica (i), y (i) implica claramente (j).
Queda por probar el lema siguiente.

\begin{lemma}[7.3.5.7]
\label{III.7.3.5.7}
La igualdad
\[
  \operatorname{length}T_p(A)
  =
  d\,\operatorname{length}A
\]
implica que $T_p(\varepsilon)$ es un epimorfismo y que $T_p(A)$ es
un $A$-módulo plano.
\end{lemma}

\begin{proof}
Partiendo de la sucesión exacta
\[
  0\longrightarrow\mathfrak m
  \longrightarrow A
  \longrightarrow k
  \longrightarrow0,
\]
las fórmulas \eqref{III.7.3.5.4.eq} y
\eqref{III.7.3.5.6.eq} dan
\[
  \begin{split}
  \operatorname{length}T_p(A)
  &\leq
  \operatorname{length}T_p(\mathfrak m)
  +
  \operatorname{length}T_p(k)
  \\
  &\leq
  d\bigl(
    \operatorname{length}\mathfrak m
    +
    \operatorname{length}k
  \bigr)
  =
  d\,\operatorname{length}A.
  \end{split}
\]
Por \sref{III.7.3.5.3}, la igualdad solo puede darse si
\[
  0\longrightarrow T_p(\mathfrak m)
  \longrightarrow T_p(A)
  \longrightarrow T_p(k)
  \longrightarrow0
  \tag{7.3.5.8}\label{III.7.3.5.8.eq}
\]
es exacta. Por \sref{III.7.2.7} y \sref{III.7.2.5}, $T_p$ es
isomorfo a un funtor
\[
  M\longmapsto N\otimes_A M.
\]
La exactitud de \eqref{III.7.3.5.8.eq}, junto con la sucesión exacta
de $\Tor$, muestra que
\[
  \Tor^A_1(N,k)=0.
\]
Se sigue que $N=T_p(A)$ es un $A$-módulo plano
\emph{(0, 10.1.3)}.
\end{proof}

\begin{lemma}[7.3.6]
\label{III.7.3.6}
Sea $A$ un anillo, y sea $T_\bullet$ un funtor homológico
covariante de $\cat{Ab}_A$ en $\cat{Ab}$ que conmuta con las
sumas directas. Supongamos que $T_p$ y $T_{p+1}$ están definidos y
que $T_p$ es exacto por la izquierda. Entonces $T_{p+1}$ es exacto
si y solo si $T_{p+1}(A_s)$ es plano como $A$-módulo a derecha.
\end{lemma}

\begin{proof}
Por \sref{III.7.3.1}, el homomorfismo canónico
\[
  T_{p+1}(A_s)\otimes_A M
  \longrightarrow
  T_{p+1}(M)
\]
es un isomorfismo de funtores. Basta ahora aplicar la definición de
$A$-módulo plano.
\end{proof}

\begin{proposition}[7.3.7]
\label{III.7.3.7}
Sea $A$ un anillo, y sea $T_\bullet$ un funtor homológico
covariante de $\cat{Ab}_A$ en $\cat{Ab}$ que conmuta con las
sumas directas. Supongamos que existe un $i_0$ tal que $T_i$ es
exacto para $i\leq i_0$. Entonces, para todo entero $p>i_0$, las
condiciones siguientes son equivalentes:
\begin{enumerate}
  \item[(a)] $T_q$ es exacto para $q\leq p$;
  \item[(b)] $T_q(A_s)$ es plano como $A$-módulo a derecha para
    $q\leq p$;
  \item[(c)] para todo $A$-módulo $M$, el homomorfismo canónico
    \[
      T_q(A_s)\otimes_A M
      \longrightarrow
      T_q(M)
    \]
    es sobreyectivo para $q\leq p+1$.
\end{enumerate}
\end{proposition}

\begin{proof}
La equivalencia de (a) y (b) resulta por inducción sobre $q$ de
\sref{III.7.3.6}, pues $T_{i_0}$ es exacto por hipótesis. La
equivalencia de (a) y (c) resulta de la equivalencia de las
condiciones (a) y ($\mathrm{e}'$) en \sref{III.7.3.3}.
\end{proof}

\oldpage[III]{184}

\begin{env}[7.3.8]
\label{III.7.3.8}
Sea $A$ un anillo conmutativo, sea $B$ una $A$-álgebra, no
necesariamente conmutativa, y sea $T_\bullet$ un funtor homológico
covariante de $\cat{Ab}_A$ en $\cat{Ab}$. Se sigue de las
definiciones de \sref{III.7.1.3} que el funtor de $\cat{Ab}_B$ en
$\cat{Ab}$ obtenido por extensión de escalares de $A$ a $B$,
denotado por
\[
  T_\bullet^{(B)}=\bigl(T_i^{(B)}\bigr),
\]
es también un funtor homológico.
\end{env}

\begin{corollary}[7.3.9]
\label{III.7.3.9}
Supongamos que $T_\bullet$ satisface las condiciones generales de
\sref{III.7.3.7} y conmuta con los límites inductivos. Supongamos
además que $A$ es un dominio entero y que cada $T_n(A)$ es un
$A$-módulo de presentación finita. Entonces, para todo entero $N$,
existe un $f\in A-\{0\}$ tal que
\[
  T_p^{(A_f)}:\cat{Ab}_{A_f}\longrightarrow\cat{Ab}
\]
es exacto para $p\leq N$.
\end{corollary}

\begin{proof}
Por hipótesis, $T_i$ es exacto para $i\leq i_0$, de modo que
$T_i(A)$ es plano para estos valores de $i$. Por la condición (b) de
\sref{III.7.3.7}, basta elegir $f$ de modo que
\[
  T_p^{(A_f)}(A_f)=T_p(A_f)
\]
sea un $A_f$-módulo libre para $i_0<p\leq N$. Ahora bien,
\[
  T_p(A_f)=\bigl(T_p(A)\bigr)_f,
\]
pues $T_p$ conmuta con los límites inductivos \sref{III.7.1.4}. Si
$x$ es el punto genérico de $\Spec(A)$, entonces
$\bigl(T_p(A)\bigr)_x$ es un espacio vectorial de dimensión finita
sobre el cuerpo de fracciones de $A$. Como cada $T_p(A)$ es de
presentación finita, existe un $f$ con la propiedad deseada
(Bourbaki, \emph{Alg\`ebre commutative}, chap.~II, \S~5,
no.~1, cor.\ de la prop.~2).

Si solo hay un número finito de índices $i$ tales que $T_i\neq0$,
se puede elegir $f\in A-\{0\}$ de modo que todos los
$T_p^{(A_f)}$ sean exactos.
\end{proof}

\begin{corollary}[7.3.10]
\label{III.7.3.10}
Supongamos que $T_\bullet$ satisface las condiciones generales de
\sref{III.7.3.7} y conmuta con los límites inductivos, que $A$ es
conmutativo y noetheriano, y que cada $T_n(A)$ es un $A$-módulo de
tipo finito. Entonces, para todo entero $N$, existe un abierto denso
$U$ de $\Spec(A)$ tal que, para todo entero $p\leq N$, la función
\[
  x\longmapsto
  \operatorname{rank}_{\kres(x)}
  \bigl(T_p(\kres(x))\bigr)
\]
es constante en $U$.
\end{corollary}

\begin{proof}
Sea $\mathfrak p$ un ideal primo minimal de $A$. El anillo
$B=A/\mathfrak p$ es un dominio entero, y $\Spec(B)$ se identifica
con una componente irreducible del espacio topológico $\Spec(A)$.
Probaremos por inducción sobre $p\leq N$ que existe un
$f_p\in B-\{0\}$ tal que, si $B'=B_{f_p}$, entonces
$T_i^{(B')}$ es exacto y $T_i(B')$ es un $B'$-módulo de tipo finito
para $i\leq p$.

La afirmación es cierta para $p\leq i_0$ por hipótesis, tomando
$f_p=1$, de modo que $B'=B=A/\mathfrak p$. En efecto, $T_p$ es
entonces exacto, y $T_p(B)$ es isomorfo a
\[
  T_p(A)/T_p(\mathfrak p);
\]
por tanto, es un $A$-módulo y, \emph{a fortiori}, un $B$-módulo de
tipo finito.

Para el paso de inducción, sea $f_p$ la imagen canónica en $B$ de un
elemento $g_p\in A$, y pongamos $A'=A_{g_p}$. Entonces
\[
  B'=A'/\mathfrak p',
\]
donde $\mathfrak p'$ es el ideal primo minimal
$\mathfrak pA_{g_p}$ de $A'$. Puesto que
\[
  T_i(A_{g_p})=\bigl(T_i(A)\bigr)_{g_p},
\]
los $T_i(A')$ son $A'$-módulos de tipo finito. Así,
$T_\bullet^{(A')}$ satisface las mismas hipótesis que $T_\bullet$,
con $i_0$ sustituido por $p$.

Podemos, por tanto, reducirnos al caso $A'=A$, en el que $T_p$ es
exacto. La sucesión exacta
\[
  0\longrightarrow\mathfrak p
  \longrightarrow A
  \longrightarrow A/\mathfrak p
  \longrightarrow0
\]
da entonces una sucesión exacta
\[
  T_{p+1}(A)
  \longrightarrow
  T_{p+1}(A/\mathfrak p)
  \xrightarrow{\partial}
  T_p(\mathfrak p)
  \longrightarrow
  T_p(A).
\]
Como $T_p$ es exacto, la última flecha es inyectiva. Por
consiguiente, $T_{p+1}(A/\mathfrak p)$ es un cociente de
$T_{p+1}(A)$ y es, por tanto, de tipo finito.

La demostración de \sref{III.7.3.9} solo ha utilizado la hipótesis
de tipo finito sobre $T_i(A)$ para $i\leq N$. Podemos, por tanto,
aplicar ese corolario al dominio entero $B$, al funtor
$T_\bullet^{(B)}$ y a $N=p+1$, lo que completa la inducción. Existe,
en consecuencia, un $f_N\in B-\{0\}$ tal que $\Spec(B_{f_N})$ es un
abierto $V$, denso en $\Spec(B)$, que no interseca ninguna otra
componente irreducible de $\Spec(A)$.

\oldpage[III]{185}

Si se ha probado la afirmación para $B_{f_N}$, existe un abierto
denso $W$ de $V$ en el que las funciones del enunciado son
constantes, pues
\[
  A_x=(B_{f_N})_x
  \qquad(x\in W).
\]
Repitiendo el mismo razonamiento para cada componente irreducible de
$\Spec(A)$ se prueba el corolario. Podemos, por tanto, reducirnos al
caso en que $A$ es un dominio entero. El razonamiento de
\sref{III.7.3.9} da entonces un $f\in A-\{0\}$ tal que los
$T_p(A_f)$ son $A_f$-módulos libres de tipo finito para $p\leq N$.
La conclusión resulta de \sref{III.7.3.4}.
\end{proof}

\begin{proposition}[7.3.11]
\label{III.7.3.11}
Sea $A$ un anillo local conmutativo de cuerpo residual $k$, y sea
$T_\bullet$ un funtor homológico covariante de $\cat{Ab}_A$ en
$\cat{Ab}$ que conmuta con las sumas directas. Supongamos que
existe un $i_0$ tal que $T_i$ es exacto para $i\leq i_0$, y que cada
$T_n(A)$ es un $A$-módulo de presentación finita. Las condiciones
equivalentes (a), (b) y (c) de \sref{III.7.3.7} implican las dos
condiciones siguientes, y son equivalentes a ellas cuando $A$ es
además reducido:
\begin{enumerate}
  \item[(d)] para todo $x\in\Spec(A)$,
    \[
      \operatorname{rank}_{\kres(x)}
      T_q(\kres(x))
      =
      \operatorname{rank}_k T_q(k)
      \qquad(q\leq p);
    \]
  \item[($\mathrm{d}'$)] para todo punto genérico $x_j$ de una
    componente irreducible de $\Spec(A)$,
    \[
      \operatorname{rank}_{\kres(x_j)}
      T_q(\kres(x_j))
      =
      \operatorname{rank}_k T_q(k)
      \qquad(q\leq p).
    \]
\end{enumerate}
\end{proposition}

\begin{proof}
Puesto que $T_q(A)$ es un $A$-módulo de presentación finita, la
condición (b) de \sref{III.7.3.7} equivale a decir que $T_q(A)$ es
libre para $q\leq p$ (Bourbaki, \emph{Alg\`ebre commutative},
chap.~II, \S~3, no.~2, cor.~2 de la prop.~5). La condición (c)
implica
\[
  T_q(\kres(x))
  =
  T_q(A)\otimes_A\kres(x)
  \qquad(q\leq p).
\]
Así, las condiciones equivalentes de \sref{III.7.3.7} implican (d),
y es inmediato que (d) implica ($\mathrm{d}'$).

Queda por mostrar que ($\mathrm{d}'$) implica (a) cuando $A$ es
reducido. Razonamos por inducción sobre $q\leq p$, pues $T_q$ es
exacto para $q\leq i_0$. Supongamos que $T_k$ es exacto para
$k\leq q<p$. Mostraremos que $T_{q+1}(A)$ es un $A$-módulo libre.
Por la hipótesis de inducción, el homomorfismo canónico
\[
  T_{q+1}(A)\otimes_A M
  \longrightarrow
  T_{q+1}(M)
\]
es un isomorfismo para todo $A$-módulo $M$, por la condición (c) de
\sref{III.7.3.7} y \sref{III.7.3.3}. Aplicando esto a
$M=\kres(x_j)$ y $M=k$, la condición ($\mathrm{d}'$) da
\[
  \operatorname{rank}_{\kres(x_j)}
  \bigl(T_{q+1}(A)\otimes_A\kres(x_j)\bigr)
  =
  \operatorname{rank}_k T_{q+1}(k)
\]
para todo $j$. Se sigue que $T_{q+1}(A)$ es libre (Bourbaki,
\emph{Alg\`ebre commutative}, chap.~II, \S~3, no.~2, prop.~7), lo
que completa la demostración.
\end{proof}

Los resultados precedentes se mejorarán considerablemente para los
funtores homológicos del tipo particular que estudiaremos en
\sref{III.7.4}: los criterios de exactitud así obtenidos solo hacen
intervenir uno de los $T_p$.

\subsection{\texorpdfstring{Criterios de exactitud para los funtores
$H_\bullet(P_\bullet\otimes_A M)$}{Criterios de exactitud para los funtores H.(P. tensor M)}}
\label{subsection:III.7.4}
\label{III.7.4}

\begin{env}[7.4.1]
\label{III.7.4.1}
Sean $A$ un anillo (no necesariamente conmutativo) y $P_\bullet$ un complejo
de $A$-módulos planos a derecha. Como el funtor
$M\mapsto P_k\otimes_A M$ es entonces exacto en $\cat{Ab}_A$ para todo $k$,
el $\partial$-funtor
\[
  T_\bullet(M)=H_\bullet(P_\bullet\otimes_A M)
  \tag{7.4.1.1}\label{III.7.4.1.1.eq}
\]
\oldpage[III]{186}
es un funtor homológico de $\cat{Ab}_A$ en $\cat{Ab}$,
evidentemente $A$-lineal cuando $A$ es conmutativo
\sref{III.7.1.2}, y conmuta con los límites inductivos.

Si $A$ es conmutativo, entonces, para toda $A$-álgebra $B$, el funtor
homológico $T_\bullet^{(B)}$ \sref{III.7.3.8} viene dado, por definición, por
\[
  T_\bullet^{(B)}(N)
  =
  H_\bullet(P_\bullet\otimes_A N_{[\rho]})
  \tag{7.4.1.2}\label{III.7.4.1.2.eq}
\]
donde $\rho:A\to B$ es el homomorfismo que define la estructura de álgebra de
$B$; como también se puede escribir
\[
  P_\bullet\otimes_A N_{[\rho]}
  =P_\bullet\otimes_A(B\otimes_B N)_{[\rho]}
  =(P_\bullet\otimes_A B)\otimes_B N,
\]
se ve que
\[
  T_\bullet^{(B)}(N)=H_\bullet(P'_\bullet\otimes_B N)
  \tag{7.4.1.3}\label{III.7.4.1.3.eq}
\]
para todo $B$-módulo $N$, siendo $P'_\bullet$ el complejo
$P_\bullet\otimes_A B$ de $B$-módulos planos
\sref[0]{0.6.2.1}.
\end{env}

\begin{proposition}[7.4.2]
\label{III.7.4.2}
Bajo las condiciones generales de \sref{III.7.4.1}, y para un entero dado
$p\in\bb Z$, las propiedades siguientes son equivalentes:
\begin{enumerate}
  \item[(a)] $T_p$ es exacto por la izquierda (o, lo que es lo mismo,
    $T_{p+1}$ es exacto por la derecha);
  \item[(b)]
    \[
      Z'_p(P_\bullet)
      =
      \Coker(P_{p+1}\longrightarrow P_p)
    \]
    es un $A$-módulo plano a derecha;
  \item[(c)] existe un complejo $P'_\bullet$ de $A$-módulos planos a derecha
    tal que la diferencial
    \[
      d'_{p+1}:P'_{p+1}\longrightarrow P'_p
    \]
    sea nula, y un isomorfismo de funtores homológicos de
    $H_\bullet(P_\bullet\otimes_A M)$ sobre
    $H_\bullet(P'_\bullet\otimes_A M)$.
\end{enumerate}
\end{proposition}

\begin{proof}
Por definición, se tiene una sucesión exacta funtorial en $M$
\[
  0\longrightarrow T_p(M)
  \longrightarrow Z'_p(P_\bullet\otimes M)
  \longrightarrow P_{p-1}\otimes M
\]
donde
\[
  Z'_p(P_\bullet\otimes M)
  =
  \Coker(P_{p+1}\otimes M\longrightarrow P_p\otimes M)
  =
  Z'_p(P_\bullet)\otimes M
\]
en virtud de la exactitud por la derecha del producto tensorial. Para todo
homomorfismo $f:M\to N$, se tiene, por tanto, un diagrama conmutativo
\[
\xymatrix@C=30pt@R=38pt{
  0\ar[r]
  &T_p(M)\ar[r]\ar[d]_u
  &Z'_p(P_\bullet)\otimes M\ar[r]\ar[d]_v
  &P_{p-1}\otimes M\ar[d]_w
  \\
  0\ar[r]
  &T_p(N)\ar[r]
  &Z'_p(P_\bullet)\otimes N\ar[r]
  &P_{p-1}\otimes N
}
\tag{7.4.2.1}\label{III.7.4.2.1.eq}
\]
cuyas filas son exactas. Si $f$ es un monomorfismo, lo mismo ocurre con $w$,
puesto que $P_{p-1}$ es plano; si $T_p$ es exacto por la izquierda, $u$ es
también un monomorfismo; se concluye que $v$ es un monomorfismo, lo cual
implica que $Z'_p(P_\bullet)$ es plano. Recíprocamente, si es así, $v$ es un
monomorfismo para todo monomorfismo $f:M\to N$, luego el diagrama
\eqref{III.7.4.2.1.eq} muestra que $u$ es un monomorfismo, y, por consiguiente,
$T_p$ (que ya es semiexacto) es exacto por la izquierda. Así, (a) y (b) son
equivalentes. Es inmediato que (c) implica (a), pues si
$d'_{p+1}:P'_{p+1}\to P'_p$ es nula, y
$0\to M'\to M\to M''\to0$ es una sucesión exacta de $A$-módulos, el
homomorfismo de borde en la sucesión exacta
\[
  H_{p+1}(P'_\bullet\otimes M'')
  \xrightarrow{\partial}H_p(P'_\bullet\otimes M')
  \longrightarrow H_p(P'_\bullet\otimes M)
\]
\oldpage[III]{187}
es nulo por definición \emph{(M, IV, 1)}, luego $T_p$ es exacto por la
izquierda. Demostremos recíprocamente que (b) implica (c). Si
\[
  Z_{p+1}(P_\bullet)=\Ker(P_{p+1}\longrightarrow P_p),
\]
se tiene una sucesión exacta
\[
  0\longrightarrow Z_{p+1}(P_\bullet)\longrightarrow P_{p+1}
  \longrightarrow Z'_p(P_\bullet)\longrightarrow0
\]
en la cual $P_{p+1}$ y $Z'_p(P_\bullet)$ son planos, luego
$Z_{p+1}(P_\bullet)$ es plano \sref[0]{0.6.1.2}. Tomaremos
\[
  P'_i=P_i\quad\text{para }i\ne p\text{ y }i\ne p+1,
  \qquad
  P'_p=Z'_p(P_\bullet),
  \qquad
  P'_{p+1}=Z_{p+1}(P_\bullet);
\]
como diferencial $d'_i:P'_i\to P'_{i-1}$, se tomará la del complejo
$P_\bullet$ para $i\ne p$ e $i\ne p+1$, $0$ para $i=p+1$ y, para $i=p$,
el homomorfismo $Z'_p(P_\bullet)\to P_{p-1}$ deducido de $d_p$ por paso al
cociente. Como los $P_i$ son planos, se tiene
\[
  Z'_i(P_\bullet\otimes M)=Z'_i(P_\bullet)\otimes M,
  \qquad
  Z_i(P_\bullet\otimes M)=Z_i(P_\bullet)\otimes M
\]
y
\[
  B_i(P_\bullet\otimes M)=B_i(P_\bullet)\otimes M
\]
(poniendo $B_i(P_\bullet)=\Im(P_{i+1}\to P_i)$); se deducen inmediatamente,
para todo $M$, isomorfismos funtoriales
\[
  H_i(P_\bullet\otimes M)\xrightarrow{\sim}H_i(P'_\bullet\otimes M)
\]
para todo $i$, y la comprobación de que se trata de un isomorfismo de
$\partial$-funtores resulta sin dificultad de la definición de $\partial$
\emph{(M, IV, 1)}.
\end{proof}

Obsérvese que las condiciones de \sref{III.7.4.2} implican también que
$B_p(P_\bullet)$ es plano, pues se tiene una sucesión exacta
\[
  0\longrightarrow B_p(P_\bullet)\longrightarrow P_p
  \longrightarrow Z'_p(P_\bullet)\longrightarrow0,
\]
en la cual $P_p$ y $Z'_p(P_\bullet)$ son planos
\sref[0]{0.6.1.2}.

\begin{corollary}[7.4.3]
\label{III.7.4.3}
Supongamos que $A$ es un anillo noetheriano regular de dimensión $1$ (es
decir, un producto de anillos de Dedekind \sref[0]{0.17.1.3} y
\sref[0]{0.17.3.7}, por ejemplo un anillo de ideales principales). Entonces,
para que $T_p$ sea exacto por la izquierda, es necesario y suficiente que
$T_p(A)$ sea un $A$-módulo plano. Para que $T_p$ sea exacto, es necesario y
suficiente que $T_p(A)$ y $T_{p-1}(A)$ sean $A$-módulos planos.
\end{corollary}

\begin{proof}
Recordemos que, para un módulo $M$ sobre un anillo de Dedekind, es equivalente
decir que $M$ es plano o que es libre de torsión \sref[0]{0.6.3.3} y
\sref[0]{0.6.3.4}; bajo las hipótesis de \sref{III.7.4.3}, todo submódulo de un
$A$-módulo plano es, por tanto, plano.

La segunda aserción de \sref{III.7.4.3} resulta de la primera, puesto que decir
que $T_p$ es exacto significa que $T_p$ y $T_{p-1}$ son exactos por la
izquierda. Para demostrar la primera aserción, observemos que se tiene una
sucesión exacta
\[
  0\longrightarrow H_p(P_\bullet)\longrightarrow Z'_p(P_\bullet)
  \longrightarrow B_{p-1}(P_\bullet)\longrightarrow0
\]
en la cual $B_{p-1}(P_\bullet)$ es un $A$-módulo plano, por ser un submódulo
del $A$-módulo plano $P_{p-1}$. Por tanto, es equivalente decir que
$H_p(P_\bullet)$ es plano o que $Z'_p(P_\bullet)$ es plano
\sref[0]{0.6.1.2}.
\end{proof}

Las aplicaciones más importantes de \sref{III.7.4.2} son las siguientes:

\begin{proposition}[7.4.4]
\label{III.7.4.4}
Sean $A$ un anillo noetheriano y $P_\bullet$ un complejo de $A$-módulos
planos; se supone, o bien que los $P_i$ son de tipo finito, o bien que los
$H_i(P_\bullet)$ son $A$-módulos de tipo finito y que existe $i_0$ tal que
$H_i(P_\bullet)=0$ para $i<i_0$. Sea $T$ el funtor homológico definido por
\eqref{III.7.4.1.1.eq}. Entonces, el conjunto $U$ de los
$y\in\Spec(A)$ tales que $(T_p)_y$ \sref{III.7.1.4} sea exacto por la derecha
(resp. exacto por la izquierda, exacto) es abierto en $\Spec(A)$.
\end{proposition}

\begin{proof}
En la segunda hipótesis sobre $P_\bullet$, se puede reemplazar $P_\bullet$ por
un complejo $P'_\bullet$ de
\oldpage[III]{188}
$A$-módulos libres de tipo finito para el cual el funtor
$H_\bullet(P'_\bullet\otimes_A M)$ es isomorfo (como $\partial$-funtor) a
$T_\bullet(M)$ \sref[0]{0.11.9.3}. Por tanto, siempre se puede reducir a la
primera hipótesis, y, en este caso, los $Z'_i(P_\bullet)$ son de tipo finito;
además, basta limitarse a demostrar las aserciones relativas a la exactitud por
la izquierda (cf. \sref{III.7.4.2}, (a)). Dicho esto, sea $x\in U$; como el
funtor $M\mapsto M_x$ es exacto, se tiene
\[
  (Z'_p(P_\bullet))_x=Z'_p((P_\bullet)_x),
\]
y (teniendo en cuenta \eqref{III.7.4.1.3.eq}) la hipótesis implica, en virtud
de \sref{III.7.4.2}, (b), que $(Z'_p(P_\bullet))_x$ es un $A_x$-módulo plano,
luego libre, pues es de tipo finito y $A_x$ es un anillo local noetheriano
(Bourbaki, \emph{Algèbre commutative}, cap.~II, \S~3,
n\textsuperscript{o}~2, cor.~2 de la prop.~5). Se deduce que existe $f\in A$
tal que $(Z'_p(P_\bullet))_f$ es libre sobre $A_f$ (Bourbaki,
\emph{Algèbre commutative}, cap.~II, \S~5, n\textsuperscript{o}~1, cor.~de la
prop.~2), y, \emph{a fortiori}, $(Z'_p(P_\bullet))_y$ es libre sobre $A_y$ para
todo $y\in D(f)$, lo cual termina la demostración en virtud de
\sref{III.7.4.2}, (b).
\end{proof}

\begin{corollary}[7.4.5]
\label{III.7.4.5}
Bajo las hipótesis de \sref{III.7.4.4}, supongamos además que $A$ es un
dominio entero.
Entonces, el conjunto $U$ de los $x\in\Spec(A)$ tales que $(T_p)_x$ sea
exacto es abierto no vacío.
\end{corollary}

\begin{proof}
En efecto, basta demostrar que $(T_p)_x$ es exacto para el punto genérico $x$
de $\Spec(A)$, lo cual es inmediato, pues $A_x$ es un cuerpo, y, por tanto,
todo funtor aditivo sobre $\cat{Ab}_{A_x}$ es exacto.
\end{proof}

\begin{proposition}[7.4.6]
\label{III.7.4.6}
Bajo las hipótesis generales de \sref{III.7.4.4}, las condiciones (a), (b) y
(c) de \sref{III.7.4.2} son también equivalentes a:
\begin{enumerate}
  \item[(d)] existe un $A$-módulo $Q$ y un isomorfismo funtorial
    \[
      T_p(M)\xrightarrow{\sim}\Hom_A(Q,M).
      \tag{7.4.6.1}\label{III.7.4.6.1.eq}
    \]
    Además, el $A$-módulo $Q$ queda determinado, salvo isomorfismo único, por
    esta propiedad, y es de tipo finito.
\end{enumerate}
\end{proposition}

\begin{proof}
La unicidad de $Q$ es un caso particular de la unicidad de un objeto
representante de un funtor representable \sref[0]{0.8.1.5}. Es claro que el
segundo miembro de \eqref{III.7.4.6.1.eq} es exacto por la izquierda.
Recíprocamente, para demostrar la existencia de $Q$, cuando $T_p$ es exacto
por la izquierda, primero se puede, como en \sref{III.7.4.4}, reducir al caso
en que los $P_i$ son planos y de tipo finito, luego (puesto que $A$ es
noetheriano) proyectivos de tipo finito (Bourbaki,
\emph{Algèbre commutative}, cap.~II, \S~5, n\textsuperscript{o}~2, cor.~del
teor.~1). El dual $\check P_i$ de $P_i$ es entonces también un $A$-módulo
proyectivo de tipo finito, $P_i$ es canónicamente isomorfo al dual de
$\check P_i$, y el homomorfismo canónico
\[
  P_i\otimes_A M\longrightarrow\Hom_A(\check P_i,M)
\]
es biyectivo (Bourbaki, \emph{Algèbre}, cap.~II, 3.\textsuperscript{a} ed.,
\S~4, n\textsuperscript{o}~2, prop.~2). Por otra parte, se sabe por
\sref{III.7.4.2}, (c), que se puede suponer que
$d_{p+1}:P_{p+1}\to P_p$ es nula, de modo que se tiene una sucesión exacta
\[
  0\longrightarrow T_p(M)\xrightarrow{u}P_p\otimes M
  \xrightarrow{v}P_{p-1}\otimes M
\]
donde $v=d_p\otimes1$. Pongamos entonces $Q'=\Ker(d_p)$, de modo que se tiene
la sucesión exacta
\[
  0\longrightarrow Q'\xrightarrow{w}P_p\xrightarrow{d_p}P_{p-1},
\]
de donde, por transposición, la sucesión exacta
\[
  \check P_{p-1}\xrightarrow{{}^t d_p}\check P_p
  \xrightarrow{{}^t w}\check Q'\longrightarrow0.
\]
Vamos a ver que
\[
  Q=\check Q'=\Coker({}^t d_p)
\]
responde a la cuestión. En efecto, se tiene la sucesión exacta
\[
  0\longrightarrow\Hom_A(Q,M)
  \longrightarrow\Hom_A(\check P_p,M)
  \xrightarrow{v'}\Hom_A(\check P_{p-1},M)
\]
\oldpage[III]{189}
donde $v'=\Hom({}^t d_p,1)$; al identificar canónicamente
$P_i\otimes M$ con $\Hom(\check P_i,M)$, $v'$ se identifica, por tanto, con
$v=d_p\otimes1$, y se obtiene así el isomorfismo funtorial
\[
  T_p(M)\xrightarrow{\sim}\Hom_A(Q,M)
\]
buscado. Además, $Q$, por ser un cociente de $\check P_p$, es de tipo finito.
\end{proof}

\begin{proposition}[7.4.7]
\label{III.7.4.7}
Supongamos verificadas las condiciones generales de \sref{III.7.4.4}. Entonces,
para todo $A$-módulo $M$ de tipo finito:
\begin{enumerate}
  \item[(i)] los $T_i(M)$ son $A$-módulos de tipo finito;
  \item[(ii)] para todo ideal $\mathfrak m$ de $A$, el homomorfismo canónico
    \[
      \widehat{T_i(M)}
      \longrightarrow
      \varprojlim_n T_i\bigl(M\otimes_A(A/\mathfrak m^{n+1})\bigr)
      \tag{7.4.7.1}\label{III.7.4.7.1.eq}
    \]
    (donde el primer miembro es el completado separado de $T_i(M)$ para la
    topología $\mathfrak m$-ádica) es biyectivo.
\end{enumerate}
\end{proposition}

\begin{proof}
Como en \sref{III.7.4.4}, se reduce primero al caso en que los $P_i$ son de
tipo finito; como $A$ es noetheriano, los submódulos de los
$P_i\otimes_A M$ son de tipo finito, de donde se deduce trivialmente la
aserción (i). En cuanto a la aserción (ii), resulta más generalmente del lema
siguiente:
\end{proof}

\begin{lemma}[7.4.7.2]
\label{III.7.4.7.2}
Sean $A$ un anillo noetheriano y $u:E\to F$ un homomorfismo de $A$-módulos de
tipo finito. Para todo $A$-módulo $M$ de tipo finito, pongamos
\[
  K(M)=\Ker(u\otimes1_M),
  \qquad
  C(M)=\Coker(u\otimes1_M);
\]
entonces, los homomorfismos canónicos
\[
  \widehat{K(M)}\longrightarrow\varprojlim_n K(M_n),
  \qquad
  \widehat{C(M)}\longrightarrow\varprojlim_n C(M_n)
  \tag{7.4.7.3}\label{III.7.4.7.3.eq}
\]
(donde se ha puesto
$M_n=M\otimes_A(A/\mathfrak m^{n+1})=M/\mathfrak m^{n+1}M$) son biyectivos
para todo ideal $\mathfrak m$ de $A$.
\end{lemma}

\begin{proof}
Como $E\otimes M$ y $F\otimes M$ son de tipo finito, y el funtor
$M\mapsto\widehat M$ es exacto en la categoría de los $A$-módulos de tipo
finito \sref[0]{0.7.3.3}, $\widehat{K(M)}$ y $\widehat{C(M)}$ son,
respectivamente, el núcleo y el conúcleo de
\[
  \widehat{(u\otimes1)}:
  \widehat{(E\otimes M)}\longrightarrow\widehat{(F\otimes M)}.
\]
La exactitud por la izquierda del funtor $\varprojlim$ muestra, por tanto, que
$\widehat{K(M)}=\varprojlim_n K(M_n)$; por otra parte, la exactitud por la
derecha del producto tensorial prueba que
$C(M_n)=C(M)\otimes_A(A/\mathfrak m^{n+1})$, luego
$\widehat{C(M)}=\varprojlim_n C(M_n)$ por definición.
\end{proof}

\begin{remark}[7.4.8]
\label{III.7.4.8}
Teniendo en cuenta \sref{III.6.10.5} y \sref{III.6.10.6}, se ve que, bajo una
hipótesis suplementaria de planitud, \sref{III.7.4.7} vuelve a dar el hecho de
que (4.1.7.1) es un isomorfismo, es decir, lo esencial del «primer teorema de
comparación» para los morfismos propios; además, el enunciado se aplica ya no
solo a un $\sh{O}_X$-módulo coherente, sino a un complejo de tales módulos.
Sería interesante obtener un enunciado que comprendiese a la vez
\sref{III.7.4.7} y (4.1.7.1) como casos particulares. Obsérvese que, cuando los
$P_i$ son de tipo finito, la demostración de \sref{III.7.4.7} no utiliza el
hecho de que sean módulos planos; convendría examinar si la conclusión de
\sref{III.7.4.7} sigue siendo válida cuando no se supone que los $P_i$ sean
planos ni de tipo finito, sino que se supone que los $H_i(P_\bullet)$ son de
tipo finito para todo $i$ y nulos para $i<i_0$. ¿Es posible entonces reemplazar
$P_\bullet$ por un complejo $P'_\bullet$ de $A$-módulos de tipo finito tal que
los funtores $H_\bullet(P_\bullet\otimes M)$ y
$H_\bullet(P'_\bullet\otimes M)$ (que ya no son homológicos) sigan siendo
isomorfos?
\end{remark}

\oldpage[III]{190}
\subsection{Caso de los anillos locales noetherianos}
\label{subsection:III.7.5}

\begin{env}[7.5.1]
\label{III.7.5.1}
Sean $A$ un anillo local noetheriano y $\mathfrak m$ su ideal
maximal. Para todo $A$-módulo $M$, designemos por $\widehat M$ su
completado separado para la topología $\mathfrak m$-ádica, el cual es
isomorfo a
\[
  \varprojlim_n
  \bigl(M\otimes_A(A/\mathfrak m^{n+1})\bigr)
  =
  \varprojlim_n
  \bigl(M/\mathfrak m^{n+1}M\bigr).
\]
Sea $T$ un funtor covariante aditivo de $\cat{Ab}_A$ en
$\cat{Ab}$. Los homomorfismos canónicos
\sref{III.7.2.3.1}
\[
  T(M)\otimes_A(A/\mathfrak m^{n+1})
  \longrightarrow
  T\bigl(M\otimes_A(A/\mathfrak m^{n+1})\bigr)
\]
forman evidentemente un sistema proyectivo de $A$-homomorfismos. Por
paso al límite dan, por tanto, un $\widehat A$-homomorfismo funtorial
en $M$
\[
  \widehat{T(M)}
  \longrightarrow
  \varprojlim_n T(M_n),
  \tag{7.5.1.1}\label{III.7.5.1.1.eq}
\]
donde se ha puesto
\[
  M_n=M\otimes_A(A/\mathfrak m^{n+1}),
  \qquad
  A_n=A/\mathfrak m^{n+1}.
\]
\end{env}

\begin{proposition}[7.5.2]
\label{III.7.5.2}
Sean $A$ un anillo local noetheriano de ideal maximal $\mathfrak m$,
$k=A/\mathfrak m$ su cuerpo residual y $T$ un funtor covariante
aditivo de $\cat{Ab}_A$ en $\cat{Ab}$, semiexacto y que conmuta con
los límites inductivos. Supongamos además que, para todo $A$-módulo
de tipo finito $M$, $T(M)$ es un $A$-módulo de tipo finito y que el
homomorfismo canónico \eqref{III.7.5.1.1.eq} es un isomorfismo. En
estas condiciones, las propiedades siguientes son equivalentes:
\begin{enumerate}
  \item[(a)] $T$ es exacto por la derecha;
  \item[(b)] para todo $n$, el funtor $N\mapsto T(N)$ es exacto por
    la derecha en la categoría de los $A_n$-módulos de tipo finito;
    equivalentemente, $T$ es exacto por la derecha en la categoría
    de los $A$-módulos de longitud finita;
  \item[(c)] el homomorfismo canónico
    \[
      T(\varepsilon):T(A)\longrightarrow T(k)
    \]
    es sobreyectivo;
  \item[(d)] para todo $n$ suficientemente grande, el homomorfismo
    canónico
    \[
      T(A_n)\longrightarrow T(k)
    \]
    es sobreyectivo.
\end{enumerate}
\end{proposition}

\begin{proof}
Es claro que (a) implica (b). Demostremos el recíproco. Sea
$u:M\to N$ un epimorfismo de $A$-módulos; debemos probar que $T(u)$
es un epimorfismo. Como $T$ conmuta con los límites inductivos y los
límites inductivos filtrantes son exactos en la categoría de los
módulos, basta considerar el caso en que $M$ y $N$ son de tipo
finito.

Los módulos $T(M)$ y $T(N)$ son entonces de tipo finito. Como $A$ es
noetheriano y local, basta probar que
\[
  \widehat{T(u)}:
  \widehat{T(M)}
  \longrightarrow
  \widehat{T(N)}
\]
es sobreyectivo \sref[0]{0.7.3.5} y \sref[0]{0.6.4.1}. Por hipótesis,
$\widehat{T(M)}$ y $\widehat{T(N)}$ son, respectivamente,
\[
  \varprojlim_n T(M_n)
  \quad\hbox{y}\quad
  \varprojlim_n T(N_n).
\]
Así, $\widehat{T(u)}$ es el límite inverso del sistema proyectivo de
homomorfismos
\[
  T(u\otimes 1_{A_n}):T(M_n)\longrightarrow T(N_n).
\]
La condición (b) significa precisamente que estos homomorfismos son
sobreyectivos. Además, $T(M_n)$ es un $A_n$-módulo de tipo finito y
$A_n$ es artiniano. Se sigue que $\widehat{T(u)}$ es sobreyectivo
\sref[0]{0.13.1.2} y \sref[0]{0.13.2.2}.

Es claro que (a) implica (c), y (c) implica (d). Finalmente,
\sref{III.7.2.7} muestra que (b) y (d) son equivalentes, pues $T$ es
semiexacto en $\cat{Ab}_{A_n}$. Esto termina la demostración.
\end{proof}

\begin{corollary}[7.5.3]
\label{III.7.5.3}
Bajo las hipótesis generales de \sref{III.7.5.2}, si $T(k)=0$,
entonces $T(M)=0$ para todo $A$-módulo $M$.
\end{corollary}

\begin{proof}
\oldpage[III]{191}
Como $k$ es el único $A$-módulo simple, de
\eqref{III.7.3.5.4.eq} se deduce que $T(E)=0$ para todo $A$-módulo
de longitud finita $E$. Si $M$ es de tipo finito, entonces
\[
  \widehat{T(M)}
  \simeq
  \varprojlim_n T(M_n).
\]
Cada $M_n$ es de longitud finita, de modo que
$\widehat{T(M)}=0$. Por hipótesis, $T(M)$ es de tipo finito y se
inyecta en $\widehat{T(M)}$ \sref[0]{0.7.3.5}; por consiguiente,
$T(M)=0$. Finalmente, para un $A$-módulo arbitrario $M$, el módulo
$T(M)$ es el límite inductivo de los $T(N_\alpha)$, donde los
$N_\alpha$ recorren los submódulos de tipo finito de $M$. Esto
termina la demostración.
\end{proof}

\begin{proposition}[7.5.4]
\label{III.7.5.4}
Sean $A$ un anillo local noetheriano de ideal maximal $\mathfrak m$,
$k=A/\mathfrak m$ su cuerpo residual y $T_\bullet$ un funtor
homológico de $\cat{Ab}_A$ en $\cat{Ab}$ que conmuta con los límites
inductivos. Supongamos además que, para todo $i$ y todo $A$-módulo
de tipo finito $M$, $T_i(M)$ es de tipo finito y el homomorfismo
canónico
\[
  \widehat{T_i(M)}
  \longrightarrow
  \varprojlim_n T_i(M_n)
\]
es biyectivo. Para un entero dado $p$, las condiciones siguientes
son equivalentes:
\begin{enumerate}
  \item[(a)] $T_p$ es exacto;
  \item[(b)] $T_p$ es exacto por la derecha y $T_p(A)$ es un
    $A$-módulo libre;
  \item[(c)] los homomorfismos canónicos
    \[
      T_{p+1}(A)\longrightarrow T_{p+1}(k),
      \qquad
      T_p(A)\longrightarrow T_p(k)
    \]
    son sobreyectivos;
  \item[(d)] para todo $n$, los homomorfismos canónicos
    \[
      T_{p+1}(A_n)\longrightarrow T_{p+1}(k),
      \qquad
      T_p(A_n)\longrightarrow T_p(k)
    \]
    son sobreyectivos;
  \item[(e)] para todo $n$, el funtor $N\mapsto T_p(N)$ es exacto
    en la categoría de los $A_n$-módulos de tipo finito.
\end{enumerate}
\end{proposition}

\begin{proof}
Por \sref{III.7.3.3}, la condición (a) equivale a decir que
$T_{p+1}$ y $T_p$ son exactos por la derecha. Como $T_\bullet$ es
un funtor homológico en $\cat{Ab}_{A_n}$, el mismo razonamiento que
en \sref{III.7.3.1} muestra que (e) equivale a decir que $T_p$ y
$T_{p+1}$ son exactos por la derecha en la categoría de los
$A_n$-módulos de tipo finito. Se deduce de \sref{III.7.5.2} que (a)
y (e) son equivalentes; la misma proposición da también la
equivalencia de (a), (c) y (d).

Finalmente, todo módulo plano de tipo finito sobre el anillo local
$A$ es libre (Bourbaki, \emph{Algèbre commutative}, cap.~II,
\S~3, n\textsuperscript{o}~2, cor.~2 de la prop.~5). La equivalencia
de (a) y (b) resulta entonces de \sref{III.7.3.1} y
\sref{III.7.3.3}.
\end{proof}

\begin{corollary}[7.5.5]
\label{III.7.5.5}
Supongamos verificadas las hipótesis generales de
\sref{III.7.5.4}.
\begin{enumerate}
  \item[(i)] Si $T_p(k)=0$, entonces $T_p=0$, el funtor $T_{p+1}$
    es exacto por la derecha y $T_{p-1}$ es exacto por la izquierda.
  \item[(ii)] Si $T_{p-1}(k)=T_{p+1}(k)=0$, entonces $T_p$ es
    exacto, el homomorfismo canónico
    \[
      T_p(A)\otimes_A M
      \longrightarrow
      T_p(M)
    \]
    es biyectivo y $T_p(A)$ es un $A$-módulo libre.
\end{enumerate}
\end{corollary}

\begin{proof}
La aserción (i) se deduce inmediatamente de \sref{III.7.5.3}, pues
$T_p$ es semiexacto; la última aserción resulta de la definición de
funtor homológico. Teniendo en cuenta \sref{III.7.3.3}, (i) da de
inmediato las dos primeras aserciones de (ii), y el hecho de que
$T_p(A)$ sea libre resulta de \sref{III.7.5.4}.
\end{proof}

\begin{corollary}[7.5.6]
\label{III.7.5.6}
Supongamos verificadas las hipótesis generales de
\sref{III.7.5.4}, y supongamos además que $A$ es un anillo de
valuación discreta.
\begin{enumerate}
  \item[(i)] El funtor $T_p$ es exacto por la derecha si y solo si
    $T_{p-1}(A)$ es un $A$-módulo libre.
  \item[(ii)] El funtor $T_p$ es exacto si y solo si $T_p(A)$ y
    $T_{p-1}(A)$ son $A$-módulos libres.
\end{enumerate}
\end{corollary}

\oldpage[III]{192}
\begin{proof}
Es claro que (ii) implica (i) (cf.\ \sref{III.7.3.3}). Para probar
(i), sea $f$ un generador del ideal maximal de $A$, es decir, un
uniformizador de $A$. Un $A$-módulo de tipo finito $M$ es libre,
equivalentemente plano, si y solo si la homotecia
\[
  h_f:M\longrightarrow M,\qquad x\longmapsto fx,
\]
es inyectiva, pues en la situación presente esto equivale a decir
que $M$ es sin torsión \sref[0]{0.6.3.4}. La sucesión exacta
\[
  0
  \longrightarrow A
  \xrightarrow{h_f}
  A
  \longrightarrow k
  \longrightarrow 0
\]
da la sucesión exacta de homología
\[
  T_p(A)
  \longrightarrow T_p(k)
  \longrightarrow T_{p-1}(A)
  \xrightarrow{h_f}
  T_{p-1}(A).
\]
Así, $T_{p-1}(A)$ es libre si y solo si
$T_p(A)\to T_p(k)$ es sobreyectivo, y la conclusión resulta de
\sref{III.7.5.2}.
\end{proof}

\begin{remark}[7.5.7]
\label{III.7.5.7}
En virtud de \sref{III.7.4.7}, las hipótesis generales de
\sref{III.7.4.4} implican que el funtor homológico $T_\bullet$
definido por \eqref{III.7.4.1.1.eq} satisface las hipótesis generales
de \sref{III.7.5.4}. En este caso, \sref{III.7.5.6} está, por tanto,
contenido en \sref{III.7.4.3}.
\end{remark}

\subsection{Descenso de las propiedades de exactitud. Teorema de
semicontinuidad y criterio de exactitud de Grauert}
\label{subsection:III.7.6}
\label{III.7.6}

\begin{proposition}[7.6.1]
\label{III.7.6.1}
Bajo las hipótesis de \sref{III.7.4.1}, sea $B$ una $A$-álgebra
conmutativa. Si $T_p$ es exacto por la derecha (respectivamente,
exacto por la izquierda, exacto), lo mismo ocurre con $T_p^{(B)}$;
el recíproco es cierto cuando $B$ es un $A$-módulo fielmente plano.
\end{proposition}

\begin{proof}
La primera aserción es un caso particular de una aserción trivial de
\sref{III.7.1.3}. Recíprocamente, supongamos primero que $B$ es un
$A$-módulo plano. Entonces, para todo $A$-módulo $M$,
\[
  H_\bullet\bigl(P_\bullet\otimes_A(M\otimes_A B)\bigr)
  =
  H_\bullet(P_\bullet\otimes_A M)\otimes_A B,
\]
lo cual también puede escribirse, para todo $p$,
\[
  T_p(M)\otimes_A B
  \simeq
  T_p^{(B)}(M_{(B)})
  \tag{7.6.1.1}\label{III.7.6.1.1}
\]
salvo isomorfismo canónico. Supongamos que $T_p^{(B)}$ es exacto por
la derecha (respectivamente, exacto por la izquierda, exacto). Como
$M\mapsto M_{(B)}$ es un funtor exacto, el miembro izquierdo de
\eqref{III.7.6.1.1} es un funtor exacto por la derecha
(respectivamente, exacto por la izquierda, exacto) en $M$. Si ahora
$B$ es fielmente plano sobre $A$, se sigue que $T_p$ tiene la misma
propiedad de exactitud \sref[0]{0.6.4.1}.
\end{proof}

\long\def\egaIIISubsectionSevenSixEleven{%
\begin{remark}[7.6.11]
\label{III.7.6.11}
Cabe preguntarse si \sref{III.7.6.9}[(i)] puede reforzarse mediante la
desigualdad
\[
  \operatorname{rank}_{\kres(x)}T_p(\kres(x))
  \geq
  \operatorname{rank}_{\kres(x)}
  \bigl(T_p(A)\otimes_A\kres(x)\bigr)
  \tag{7.6.11.1}\label{III.7.6.11.1}
\]
para todo $x\in\Spec(A)$. Esta desigualdad se verifica cuando $A$ es
un anillo de valuación discreta y $x$ es su ideal maximal
\sref{III.7.6.5}. Limitémonos ahora al caso en que $A$ es un anillo
local noetheriano, de ideal maximal $\mathfrak m$ y cuerpo residual
$k$. Las condiciones siguientes son equivalentes:
\begin{enumerate}
  \item[a)] Para todo complejo $P_\bullet$ de $A$-módulos planos de
    tipo finito,
    \[
      \operatorname{rank}_k T_p(k)
      \geq
      \operatorname{rank}_k\bigl(T_p(A)\otimes_A k\bigr)
      \qquad\text{para todo entero }p.
      \tag{7.6.11.2}\label{III.7.6.11.2}
    \]
  \item[b)] Para todo $A$-módulo $M$ de tipo finito,
    \[
      \operatorname{rank}_k(M\otimes_A k)
      \geq
      \operatorname{rank}_k(\check M\otimes_A k).
      \tag{7.6.11.3}\label{III.7.6.11.3}
    \]
  \item[c)] Para todo $A$-módulo $N$ de tipo finito,
    \[
      \operatorname{rank}_k\Tor_1^A(N,k)
      \geq
      \operatorname{rank}_k\Tor_2^A(N,k).
      \tag{7.6.11.4}\label{III.7.6.11.4}
    \]
\end{enumerate}
Equivalentemente, por cambio de dimensión (M, V, 7.2), puede
exigirse, para todo $i\geq1$,
\[
  \operatorname{rank}_k\Tor_i^A(N,k)
  \geq
  \operatorname{rank}_k\Tor_{i+1}^A(N,k).
  \tag{7.6.11.5}\label{III.7.6.11.5}
\]

Demos brevemente algunas indicaciones sobre la demostración. Para ver
que \textnormal{a)} implica \textnormal{b)}, considérese una sucesión
exacta
\[
  L_1\xrightarrow{d}L_0\longrightarrow M\longrightarrow0,
\]
donde $L_0$ y $L_1$ son libres de tipo finito, y aplíquese
\textnormal{a)} al complejo
\[
  P_1\xrightarrow{{}^t d}P_0,
  \qquad
  P_0=\check L_1,\quad P_1=\check L_0,
\]
siendo nulos todos sus demás términos. Entonces
\[
  T_1(A)=\check M
\]
y
\[
  T_1(k)
  =
  \Hom_A(M,k)
  =
  \Hom_k(M/\mathfrak mM,k).
\]
Así, $T_1(k)$ es el dual del espacio vectorial $M\otimes_A k$ y tiene
el mismo rango. Para probar que \textnormal{b)} implica
\textnormal{c)}, establezcamos primero el lema siguiente.

\begin{lemma}[7.6.11.6]
\label{III.7.6.11.6}
Dado un complejo de $A$-módulos planos
\[
  \cdots\longrightarrow0\longrightarrow
  P_1\xrightarrow{d}P_0
  \longrightarrow0\longrightarrow\cdots,
\]
se tiene una sucesión exacta
\[
  0\longrightarrow\Tor_2^A(Z'_0,k)
  \longrightarrow T_1(A)\otimes_A k
  \longrightarrow T_1(k)
  \longrightarrow\Tor_1^A(Z'_0,k)
  \longrightarrow0.
  \tag{7.6.11.7}\label{III.7.6.11.7}
\]
\end{lemma}

\oldpage[III]{197}
\begin{proof}
De la sucesión exacta
\[
  0\longrightarrow Z_1\longrightarrow P_1
  \longrightarrow B_0\longrightarrow0
\]
se obtiene
\[
  0\longrightarrow\Tor_1^A(B_0,k)
  \longrightarrow Z_1\otimes_A k
  \longrightarrow P_1\otimes_A k
  \xrightarrow{u}B_0\otimes_A k
  \longrightarrow0.
\]
La sucesión exacta
\[
  0\longrightarrow B_0\longrightarrow Z_0
  \longrightarrow Z'_0\longrightarrow0
\]
y la planitud de $Z_0=P_0$ dan
\[
  \Tor_1^A(B_0,k)=\Tor_2^A(Z'_0,k).
\]
Por definición, $Z_1=T_1(A)$. Además,
$T_1(k)=\Ker(d\otimes1)$, y $d\otimes1$ se factoriza como
\[
  P_1\otimes_A k
  \xrightarrow{u}B_0\otimes_A k
  \xrightarrow{v}Z_0\otimes_A k
  =
  P_0\otimes_A k.
\]
Por tanto, $T_1(k)=u^{-1}(R)$, donde $R=\Ker v$. Como $u$ es
sobreyectivo, $R=u(T_1(k))$; finalmente,
$R=\Tor_1^A(Z'_0,k)$ por definición de $v$. Esto demuestra
\eqref{III.7.6.11.7}.
\end{proof}

Para deducir \textnormal{c)} de \textnormal{b)}, tómese una sucesión
exacta
\[
  L_1\xrightarrow{d}L_0\longrightarrow N\longrightarrow0
\]
con $L_0,L_1$ libres de tipo finito, y considérese el funtor $T$
asociado al complejo formado por $L_1$ y $L_0$. Como $L_0$ y $L_1$
son libres, se identifican con sus biduales. Si
\[
  M=\Coker({}^t d),
\]
entonces
\[
  T_1(A)=\Ker d=\check M.
\]
Además,
\[
  M\otimes_A k=\Coker({}^t d\otimes1_k)
\]
tiene el mismo rango sobre $k$ que $\Ker(d\otimes1_k)$. La hipótesis
\textnormal{b)} da, por tanto,
\[
  \operatorname{rank}_k\bigl(T_1(A)\otimes_A k\bigr)
  \leq
  \operatorname{rank}_kT_1(k).
\]
Como $Z'_0=N$, la desigualdad \eqref{III.7.6.11.4} se deduce de la
sucesión exacta \eqref{III.7.6.11.7}.

Finalmente, para probar que \textnormal{c)} implica
\textnormal{a)}, aplíquese \sref{III.7.6.11.6} al complejo de dos
términos
\[
  P_p\longrightarrow P_{p-1}.
\]
La hipótesis \textnormal{c)}, aplicada a $Z'_{p-1}$, da
\[
  \operatorname{rank}_k R\geq\operatorname{rank}_k S,
\]
donde
\[
  R=\Ker(P_p\otimes_A k\longrightarrow P_{p-1}\otimes_A k),
  \qquad
  S=Z_p\otimes_A k.
\]
Factorícese $d_{p+1}:P_{p+1}\to P_p$ como
\[
  P_{p+1}\xrightarrow{v}Z_p\xrightarrow{j}P_p.
\]
Entonces
\[
  \Im(d_{p+1}\otimes1)
  =
  (j\otimes1)\bigl(\Im(v\otimes1)\bigr).
\]
Como
\[
  T_p(A)\otimes_A k
  =
  (Z_p/B_p)\otimes_A k
  =
  (Z_p\otimes_A k)/\Im(v\otimes1)
\]
y
\[
  T_p(k)=R/\Im(d_{p+1}\otimes1),
\]
se obtiene la desigualdad \eqref{III.7.6.11.2}.

Supongamos ahora que el anillo local $A$ es regular de dimensión
$n$. Se sabe [17] que el $A$-módulo $\Tor_i^A(k,k)$ es isomorfo a la
potencia exterior
\[
  \bigwedge\nolimits^i(\mathfrak m/\mathfrak m^2).
\]
Así, \eqref{III.7.6.11.4} no se verifica para $N=k$ en cuanto
$n\geq4$. Por otra parte, supongamos que $A$ es un anillo local
íntegro tal que todo $A$-módulo reflexivo de tipo finito es libre,
como ocurre cuando $A$ es un anillo regular de dimensión $2$.
Entonces se verifica \eqref{III.7.6.11.3}. En efecto, el dual
$\check M$ de un $A$-módulo $M$ de tipo finito es reflexivo, luego
libre, y por consiguiente
\[
  \operatorname{rank}_k(\check M\otimes_A k)
  =
  \operatorname{rank}_K(\check M)
  =
  \operatorname{rank}_K(M),
\]
donde $K$ es el cuerpo de fracciones de $A$. Toda base de
$M\otimes_A k$ sobre $k$ está formada por las imágenes de un sistema
de generadores de $M$ (Bourbaki, \emph{Alg\`ebre commutative},
cap.~II, \S~3, n\textsuperscript{o}~2, cor.~2 de la prop.~4). Por
tanto,
\[
  \operatorname{rank}_K(M)
  \leq
  \operatorname{rank}_k(M\otimes_A k),
\]
lo cual demuestra la aserción.
\end{remark}
}

\begin{proposition}[7.6.2]
\label{III.7.6.2}
Bajo las hipótesis de \sref{III.7.4.1}, supongamos además que $A$ es
un anillo noetheriano reducido y que los $P_i$ son $A$-módulos de
tipo finito. Para que $T_p$ sea exacto por la derecha
(respectivamente, exacto por la izquierda, exacto), es necesario y
suficiente que $T_p^{(B)}$ tenga la misma propiedad para toda
$A$-álgebra $B$ que sea un anillo de valuación discreta.
\end{proposition}

\begin{proof}
Por \sref{III.7.3.1} y \sref{III.7.3.3}, podemos limitarnos a la
exactitud por la derecha, y solo hay que demostrar la suficiencia.
Por \sref{III.7.4.2}, basta probar que $Z'_{p-1}(P_\bullet)$ es un
$A$-módulo plano. Como $P_{p-1}$ es de tipo finito,
$Z'_{p-1}(P_\bullet)$ también es de tipo finito; el criterio de
\sref[0]{0.10.2.8} muestra entonces que basta con que
\[
  Z'_{p-1}(P_\bullet)\otimes_A B
\]
sea un $B$-módulo plano para toda $A$-álgebra $B$ que sea un anillo
de valuación discreta. Como $P_\bullet$ es un complejo de
$A$-módulos planos, se tiene
\[
  Z'_{p-1}(P_\bullet)\otimes_A B
  =
  Z'_{p-1}(P_\bullet\otimes_A B).
\]

\oldpage[III]{193}
El complejo $P_\bullet\otimes_A B$ está formado por $B$-módulos
planos \sref[0]{0.6.2.1}, y para todo $B$-módulo $N$ se tiene
\[
  H_\bullet(P_\bullet\otimes_A N)
  =
  H_\bullet\bigl((P_\bullet\otimes_A B)\otimes_B N\bigr),
\]
y, por tanto,
\[
  T_p^{(B)}(N)
  =
  H_p\bigl((P_\bullet\otimes_A B)\otimes_B N\bigr).
\]
Aplicando \sref{III.7.4.2} a $T_p^{(B)}$, vemos que la hipótesis de
que $T_p^{(B)}$ es exacto por la derecha equivale a que
$Z'_{p-1}(P_\bullet\otimes_A B)$ sea un $B$-módulo plano.

El criterio precedente nos lleva a examinar más de cerca el caso de
los anillos de valuación discreta.
\end{proof}

\begin{proposition}[7.6.3]
\label{III.7.6.3}
Bajo las hipótesis de \sref{III.7.4.1}, supongamos que $A$ es un
anillo noetheriano regular de dimensión $1$ (es decir, que $A$ es
noetheriano y, para todo $x\in\Spec(A)$, el anillo $A_x$ es un cuerpo
o un anillo de valuación discreta). Entonces, para todo entero $p$ y
todo $A$-módulo $M$, se tiene una sucesión exacta canónica, funtorial
en $M$,
\[
  0
  \longrightarrow T_p(A)\otimes_A M
  \xrightarrow{t_M} T_p(M)
  \longrightarrow
  \Tor_1^A\bigl(T_{p-1}(A),M\bigr)
  \longrightarrow 0.
  \tag{7.6.3.1}\label{III.7.6.3.1}
\]
\end{proposition}

\begin{proof}
En lo que sigue, para simplificar la notación omitiremos el complejo
$P_\bullet$ de las notaciones homológicas usuales
$H_p(P_\bullet)$, $B_p(P_\bullet)$, $Z_p(P_\bullet)$ y
$Z'_p(P_\bullet)$. Se tienen las tres sucesiones exactas
\begin{gather*}
  0\longrightarrow H_p\longrightarrow Z'_p
  \longrightarrow B_{p-1}\longrightarrow 0,\\
  0\longrightarrow B_{p-1}\longrightarrow Z_{p-1}
  \longrightarrow H_{p-1}\longrightarrow 0,\\
  0\longrightarrow Z_{p-1}\longrightarrow P_{p-1}
  \longrightarrow B_{p-2}\longrightarrow 0.
\end{gather*}
Como $P_{p-1}$ y $P_{p-2}$ son planos, también lo son sus
submódulos respectivos $B_{p-1}$, $Z_{p-1}$ y $B_{p-2}$: para todo
$x\in\Spec(A)$, los $A_x$-módulos planos y los $A_x$-módulos sin
torsión son los mismos. Tensorizando por $M$ se obtienen, por tanto,
las sucesiones exactas
\[
  0=\Tor_1^A(B_{p-1},M)
  \longrightarrow H_p\otimes_A M
  \longrightarrow Z'_p\otimes_A M
  \xrightarrow{u}B_{p-1}\otimes_A M
  \longrightarrow 0,
  \tag{7.6.3.2}\label{III.7.6.3.2}
\]
\[
  0=\Tor_1^A(Z_{p-1},M)
  \longrightarrow \Tor_1^A(H_{p-1},M)
  \longrightarrow B_{p-1}\otimes_A M
  \xrightarrow{v} Z_{p-1}\otimes_A M,
  \tag{7.6.3.3}\label{III.7.6.3.3}
\]
y
\[
  0=\Tor_1^A(B_{p-2},M)
  \longrightarrow Z_{p-1}\otimes_A M
  \xrightarrow{w}P_{p-1}\otimes_A M.
  \tag{7.6.3.4}\label{III.7.6.3.4}
\]

Por definición,
\[
  T_p(M)=
  \Ker(d_p\otimes 1)/\Im(d_{p+1}\otimes 1).
\]
Es, por tanto, el núcleo del homomorfismo
\[
  (P_p\otimes_A M)/\Im(d_{p+1}\otimes 1)
  \longrightarrow P_{p-1}\otimes_A M
\]
inducido por $d_p\otimes 1$. Por la definición
$Z'_p=P_p/B_p$, este homomorfismo también puede escribirse
$Z'_p\otimes_A M\to P_{p-1}\otimes_A M$; es el compuesto
\[
  Z'_p\otimes_A M
  \xrightarrow{u}B_{p-1}\otimes_A M
  \xrightarrow{v}Z_{p-1}\otimes_A M
  \xrightarrow{w}P_{p-1}\otimes_A M.
\]
Como $w$ es inyectivo por \eqref{III.7.6.3.4}, se tiene una sucesión
exacta
\[
  0\longrightarrow\Ker u\longrightarrow T_p(M)
  \longrightarrow\Ker v\longrightarrow 0.
\]
En vista de \eqref{III.7.6.3.2} y \eqref{III.7.6.3.3}, y puesto que
$H_p=T_p(A)$ por definición, esta es precisamente
\eqref{III.7.6.3.1}.
\end{proof}

\begin{remarks}[7.6.4]
\label{III.7.6.4}
\textnormal{(i)}
El grupo $H_\bullet(P_\bullet\otimes_A M)$ es la homología del
bicomplejo $P_\bullet\otimes_A M$, donde $M$ se considera como un
complejo concentrado en grado $0$. Es, por consiguiente,
\sref{III.6.3.6}\sref{III.6.3.2} el objeto al que converge la sucesión
espectral regular cuyo término $E_2$ es
\[
  E^2_{pq}
  =
  \Tor_p^A\bigl(H_q(P_\bullet),M\bigr)
  =
  \Tor_p^A\bigl(T_q(A),M\bigr).
\]

\oldpage[III]{194}
La hipótesis sobre $A$ implica que $\Tor_p^A(E,F)=0$ para $p\geq 2$
y para $A$-módulos cualesquiera $E,F$ \sref[0]{0.17.2.2}. Se sabe
(M, XV) que esto implica la exactitud de
\[
  0\longrightarrow E^2_{0,q}
  \longrightarrow H_q(P_\bullet\otimes_A M)
  \longrightarrow E^2_{1,q-1}
  \longrightarrow 0,
\]
que es precisamente \eqref{III.7.6.3.1}.

\medskip
\textnormal{(ii)}
Teniendo en cuenta \sref{III.7.3.1}, la sucesión exacta
\eqref{III.7.6.3.1} recupera \sref{III.7.4.3} como caso particular.
\end{remarks}

\begin{corollary}[7.6.5]
\label{III.7.6.5}
Bajo las hipótesis de \sref{III.7.4.1}, supongamos que $A$ es un
anillo de valuación discreta, con cuerpo de fracciones $K$ y cuerpo
residual $k$, y que los $T_i(A)$ son $A$-módulos de tipo finito.
Entonces
\[
  \operatorname{rank}_k T_p(k)
  \geq
  \operatorname{rank}_k\bigl(T_p(A)\otimes_A k\bigr)
  \geq
  \operatorname{rank}_A T_p(A)
  =
  \operatorname{rank}_K T_p(K).
  \tag{7.6.5.1}\label{III.7.6.5.1}
\]
Además, los dos términos extremos de esta desigualdad son iguales si
y solo si $T_p$ es exacto o, equivalentemente, si y solo si $T_p(A)$
y $T_{p-1}(A)$ son $A$-módulos libres.
\end{corollary}

\begin{proof}
Tómese $M=k$ en la sucesión exacta \eqref{III.7.6.3.1}. Como los
términos son espacios vectoriales sobre $k$, se obtiene
\[
  \operatorname{rank}_k T_p(k)
  =
  \operatorname{rank}_k\bigl(T_p(A)\otimes_A k\bigr)
  +
  \operatorname{rank}_k
  \Tor_1^A\bigl(T_{p-1}(A),k\bigr).
\]
Por otra parte, como $T_p(A)$ es un módulo de tipo finito sobre el
anillo local íntegro $A$, se tiene (Bourbaki,
\emph{Alg\`ebre commutative}, cap.~II, \S~3,
n\textsuperscript{o}~2, cor.~1 de la prop.~4)
\[
  \operatorname{rank}_k\bigl(T_p(A)\otimes_A k\bigr)
  \geq
  \operatorname{rank}_A T_p(A)
  =
  \operatorname{rank}_K\bigl(T_p(A)\otimes_A K\bigr),
  \tag{7.6.5.2}\label{III.7.6.5.2}
\]
y los dos miembros de \eqref{III.7.6.5.2} son iguales si y solo si
$T_p(A)$ es un $A$-módulo libre (\emph{loc.\ cit.}, prop.~7). Como
$K$ es un $A$-módulo plano, también se tiene, por definición,
\[
  T_p(A)\otimes_A K
  =
  H_p(P_\bullet)\otimes_A K
  =
  H_p(P_\bullet\otimes_A K)
  =
  T_p(K).
\]
Esto demuestra \eqref{III.7.6.5.1}. También muestra que la igualdad
solo puede darse si, en primer lugar, $T_p(A)$ es libre y, en segundo
lugar,
\[
  \Tor_1^A\bigl(T_{p-1}(A),k\bigr)=0,
\]
siendo esta última condición equivalente a que $T_{p-1}(A)$ sea
libre \sref[0]{0.10.1.3}. Finalmente, para los módulos de tipo finito
sobre $A$, la planitud y la libertad son equivalentes (Bourbaki,
\emph{Alg\`ebre commutative}, cap.~II, \S~3,
n\textsuperscript{o}~2, cor.~2 de la prop.~5), y el resultado se
deduce de \sref{III.7.4.3}.
\end{proof}

\begin{env}[7.6.6]
\label{III.7.6.6}
Manteniéndose las hipótesis de \sref{III.7.4.1}, para todo
$x\in\Spec(A)$ pongamos
\[
  d_p(x)=d_p^T(x)
  =
  \operatorname{rank}_{\kres(x)}T_p(\kres(x)).
  \tag{7.6.6.1}\label{III.7.6.6.1}
\]
\end{env}

\begin{lemma}[7.6.7]
\label{III.7.6.7}
Sea $\varphi:A\to A'$ un homomorfismo de anillos, y sea
\[
  f={}^a\!\varphi:\Spec(A')\longrightarrow\Spec(A)
\]
la aplicación correspondiente \sref[I]{I.1.2.1}. Si
$T'_\bullet=T_\bullet^{(A')}$ \sref{III.7.1.3}, entonces
\[
  d_p^{T'}=d_p^T\circ f.
  \tag{7.6.7.1}\label{III.7.6.7.1}
\]
\end{lemma}

\oldpage[III]{195}
\begin{proof}
Para $x'\in\Spec(A')$, pongamos $x=f(x')$. Entonces
\begin{align*}
  H_\bullet(P_\bullet\otimes_A\kres(x'))
  &=
  H_\bullet\bigl(
    (P_\bullet\otimes_A\kres(x))
    \otimes_{\kres(x)}\kres(x')
  \bigr)\\
  &=
  H_\bullet(P_\bullet\otimes_A\kres(x))
  \otimes_{\kres(x)}\kres(x'),
\end{align*}
porque $\kres(x')$ es plano sobre $\kres(x)$. Esto demuestra
\eqref{III.7.6.7.1}.
\end{proof}

\begin{lemma}[7.6.8]
\label{III.7.6.8}
Si $A$ es noetheriano y $P_\bullet$ está formado por $A$-módulos de
tipo finito, la función $x\mapsto d_p^T(x)$ sobre $\Spec(A)$ es
constructible.
\end{lemma}

\begin{proof}
Para toda parte cerrada irreducible $Y$ de $X=\Spec(A)$, debemos
probar que existe un abierto no vacío $U$ de $Y$ en el cual $d_p$ es
constante \sref[0]{0.9.2.2}. Escribamos
$Y=\Spec(A/\mathfrak a)$, donde $\mathfrak a$ es un ideal de $A$ tal
que $A/\mathfrak a$ es reducido. Por \sref{III.7.6.7}, podemos
reducirnos al caso en que $Y=X$ y $A$ es un anillo noetheriano
íntegro; la aserción resulta entonces de \sref{III.7.4.5}.
\end{proof}

\begin{theorem}[7.6.9]
\label{III.7.6.9}
Sean $A$ un anillo noetheriano, $P_\bullet$ un complejo de
$A$-módulos planos de tipo finito y
\[
  T_\bullet(M)=H_\bullet(P_\bullet\otimes_A M)
\]
el funtor homológico definido por $P_\bullet$. Para
$x\in\Spec(A)$, pongamos
\[
  d_p(x)=\operatorname{rank}_{\kres(x)}T_p(\kres(x)).
\]
Entonces:
\begin{enumerate}
  \item[(i)] La función $d_p$ es constructible y semicontinua
    superiormente sobre $\Spec(A)$.
  \item[(ii)] Si $T_p$ es exacto, $d_p$ es continua (y, por tanto,
    localmente constante) sobre $\Spec(A)$. El recíproco es cierto
    cuando $A$ es reducido.
\end{enumerate}
\end{theorem}

\begin{proof}
\textnormal{(i)}
La constructibilidad se demostró en \sref{III.7.6.8}. Para probar la
semicontinuidad superior, basta, por tanto, \sref[0]{0.9.2.4},
demostrar que, si $x'\neq x$ es una generalización de $x$ en
$\Spec(A)$, entonces $d_p(x')\leq d_p(x)$. Existe un anillo de
valuación discreta $B$ y un morfismo
\[
  f:\Spec(B)\longrightarrow\Spec(A)
\]
tales que, si $a$ es el punto cerrado de $\Spec(B)$ y $b$ su punto
genérico, entonces $f(a)=x$ y $f(b)=x'$
\sref[II]{II.7.1.9}. Por \eqref{III.7.6.7.1}, nos reducimos a probar
$d_p(a)\geq d_p(b)$ sobre $\Spec(B)$, que es precisamente
\eqref{III.7.6.5.1}.\footnote{El principio de demostrar
\textnormal{(i)} por reducción al caso de un anillo de valuación
discreta nos fue comunicado oralmente por Hironaka.}

\medskip
\textnormal{(ii)}
La primera aserción ya se demostró en \sref{III.7.3.4}. Para el
recíproco, utilícese el criterio valuativo \sref{III.7.6.2}. En vista
de \eqref{III.7.6.7.1}, nos reducimos al caso en que $A$ es un anillo
de valuación discreta. Como $\Spec(A)$ tiene entonces solamente dos
puntos, la hipótesis de que $d_p$ es constante implica que $T_p$ es
exacto, por \sref{III.7.6.5}.
\end{proof}

\begin{corollary}[7.6.10]
\label{III.7.6.10}
Sean $A$ un anillo noetheriano, $\mathfrak p_i$ $(1\leq i\leq r)$
sus ideales primos minimales, y $k_i$ el cuerpo residual de
$A_{\mathfrak p_i}$.
\begin{enumerate}
  \item[(i)] Para todo $x\in\Spec(A)$, existe un índice $i$ tal que
    \[
      d_p(x)\geq
      \operatorname{rank}_{k_i}T_p(k_i).
      \tag{7.6.10.1}\label{III.7.6.10.1}
    \]
\end{enumerate}
En particular, si $A$ es un anillo íntegro y $K$ es su cuerpo de
fracciones, entonces
\[
  d_p(x)\geq\operatorname{rank}_K T_p(K)
  \tag{7.6.10.2}\label{III.7.6.10.2}
\]
para todo $x\in\Spec(A)$.

\oldpage[III]{196}
\begin{enumerate}
  \item[(ii)] Supongamos además que $A$ es local y reducido, y sea
    $k$ su cuerpo residual. Entonces $T_p$ es exacto si y solo si
    \[
      \operatorname{rank}_k T_p(k)
      =
      \operatorname{rank}_{k_i}T_p(k_i)
      \qquad(1\leq i\leq r).
      \tag{7.6.10.3}\label{III.7.6.10.3}
    \]
\end{enumerate}
\end{corollary}

\begin{proof}
La aserción \textnormal{(i)} es inmediata: todo entorno de $x$
contiene uno de los $\mathfrak p_i$, y basta aplicar la definición de
semicontinuidad. Si $A$ es local, el único entorno del ideal maximal
$\mathfrak m$ en $\Spec(A)$ es todo $\Spec(A)$; por tanto,
\[
  d_p(x)\leq\operatorname{rank}_k T_p(k)
\]
para todo $x\in\Spec(A)$. Junto con \textnormal{(i)}, esto muestra que
\eqref{III.7.6.10.3} implica que $d_p$ es constante sobre
$\Spec(A)$ y, por tanto, que $T_p$ es exacto por
\sref{III.7.6.9}[(ii)]. El recíproco se deduce inmediatamente del
mismo resultado.
\end{proof}

\egaIIISubsectionSevenSixEleven

\subsection{Aplicación a los morfismos propios: I. La propiedad de cambio de base}
\label{subsection:III.7.7}
\label{III.7.7}

Los tres números que siguen son, en lo esencial, traducciones al lenguaje de
los morfismos de preesquemas de los resultados de los números precedentes.

\begin{env}[7.7.1]
\label{III.7.7.1}
Sea $f:X\to Y$ un morfismo cuasicompacto y separado de preesquemas, y sea
$\mathscr P_\bullet$ un complejo de $\mathscr O_X$-módulos cuasicoherentes cuyo
operador de derivación es de grado $-1$; supongamos además que los
$\mathscr O_X$-módulos $\mathscr P_i$ son planos sobre $Y$
\sref[0]{0.6.7.1}.

\oldpage[III]{198}
Consideraremos el $\partial$-funtor
$\mathscr M\mapsto\mathscr T_\bullet(\mathscr M)$ (también denotado por
$\mathscr T_\bullet(\mathscr P_\bullet,\mathscr M)$) de la categoría de los
$\mathscr O_Y$-módulos cuasicoherentes en la categoría de los
$\mathscr O_Y$-módulos cuasicoherentes, en virtud de \sref{III.6.2.3}, definido
por
\[
  \mathscr T_n(\mathscr P_\bullet,\mathscr M)
  =\mathscr T_n(\mathscr M)
  =\mathscr H^{-n}
    \bigl(f,\mathscr P^\bullet\otimes_{\mathscr O_Y}\mathscr M\bigr)
  \qquad\text{para }n\in\mathbf Z,
  \tag{7.7.1.1}\label{III.7.7.1.1}
\]
donde $\mathscr P^\bullet$ es el complejo cuyo término de grado $j$ es
$\mathscr P_{-j}$, de modo que su operador de derivación es entonces de grado
$+1$. El funtor $\mathscr T_\bullet$ así definido es un \emph{funtor
homológico} en $\mathscr M$ \sref{III.6.2.6}.
\end{env}

\begin{env}[7.7.2]
\label{III.7.7.2}
Sea $g:Y'\to Y$ un morfismo, y pongamos
\[
  X'=X_{(Y')}=X\times_Y Y',
  \qquad
  f'=f_{(Y')}:X'\longrightarrow Y',
\]
que es un morfismo cuasicompacto y separado; pongamos además
\[
  \mathscr P'_\bullet
  =\mathscr P_\bullet\otimes_{\mathscr O_Y}\mathscr O_{Y'}.
\]
Este es un complejo de $\mathscr O_{X'}$-módulos cuasicoherentes que son planos
sobre $Y'$, en virtud de \sref[I]{I.9.1.12} y \sref[0]{0.6.2.1}. Pondremos,
con las mismas convenciones sobre los grados,
\[
  \mathscr T^{Y'}_\bullet(\mathscr M')
  =\mathscr H^\bullet
    \bigl(f',\mathscr P'^\bullet\otimes_{\mathscr O_{Y'}}\mathscr M'\bigr)
  =\mathscr H^\bullet
    \bigl(f',\mathscr P^\bullet\otimes_{\mathscr O_Y}\mathscr M'\bigr)
  \tag{7.7.2.1}\label{III.7.7.2.1}
\]
que es un funtor homológico en el $\mathscr O_{Y'}$-módulo cuasicoherente
$\mathscr M'$. Cuando $Y'$ sea un esquema afín de anillo $A'$, escribiremos
$\mathscr T^{A'}_\bullet$ en lugar de $\mathscr T^{Y'}_\bullet$; para todo
$A'$-módulo $M'$ se tiene entonces
\[
  \mathscr T^{A'}_\bullet(\widetilde{M'})
  =\bigl(\Gamma(Y',\mathscr T^{Y'}_\bullet(\widetilde{M'}))\bigr)^\sim.
\]
Pondremos
\[
  T^{A'}_\bullet(M')
  =\Gamma(Y',\mathscr T^{Y'}_\bullet(\widetilde{M'})),
\]
que es un funtor homológico de $A'$-módulos con valores en la categoría de los
$A'$-módulos.

Obsérvese que, si $Y=\operatorname{Spec}(A)$ es también afín, el funtor de
$A'$-módulos $T^{A'}_\bullet$ coincide con el funtor obtenido por extensión de
escalares de $A$ a $A'$ a partir del funtor homológico de $A$-módulos
$T^A_\bullet$ \sref{III.7.1.3}. En efecto, sea $g:Y'\to Y$ el morfismo
correspondiente al homomorfismo de anillos $A\to A'$, y sea $g':X'\to X$ el
morfismo correspondiente, que es afín \sref[II]{II.1.6.2}; si $\mathfrak U$ es
un recubrimiento abierto afín de $X$, entonces
$\mathfrak U'=g'^{-1}(\mathfrak U)$ es un recubrimiento abierto afín de $X'$.
En virtud de \sref{III.6.2.2}, todo se reduce a ver que
\[
  C^\bullet
    \bigl(\mathfrak U,
      \mathscr P_\bullet\otimes_{\mathscr O_Y}g_*(\mathscr M')\bigr)
  =C^\bullet
    \bigl(\mathfrak U',
      \mathscr P_\bullet\otimes_{\mathscr O_Y}\mathscr M'\bigr),
\]
y, finalmente, que para todo abierto afín $U$ de $X$, poniendo
$U'=g'^{-1}(U)$, se tiene
\[
  \Gamma
    \bigl(U,\mathscr P_\bullet\otimes_{\mathscr O_Y}g_*(\mathscr M')\bigr)
  =\Gamma
    \bigl(U',\mathscr P_\bullet\otimes_{\mathscr O_Y}\mathscr M'\bigr),
\]
lo cual es trivial \sref[I]{I.1.3}, \sref[I]{I.3.2}.

En particular, si $U$ es un abierto de $Y$, para todo
$\mathscr O_Y$-módulo cuasicoherente $\mathscr M$ se tiene
\[
  \mathscr T^U_\bullet(\mathscr M|U)
  =\bigl(\mathscr T_\bullet(\mathscr M)\bigr)|U.
  \tag{7.7.2.2}\label{III.7.7.2.2}
\]
\end{env}

\begin{env}[7.7.3]
\label{III.7.7.3}
Para todo $\mathscr O_Y$-módulo cuasicoherente $\mathscr M$ existe un
homomorfismo canónico, funtorial en $\mathscr M$,
\[
  \mathscr T_p(\mathscr O_Y)\otimes_{\mathscr O_Y}\mathscr M
  \longrightarrow
  \mathscr T_p(\mathscr M).
  \tag{7.7.3.1}\label{III.7.7.3.1}
\]
En efecto, si $Y$ es afín, este homomorfismo se definió en
\sref{III.7.2.2}; esta definición se extiende sin dificultad al caso general,
observando que, si $U$ y $V$ son dos abiertos afines de $Y$ tales que
$V\subset U$, el diagrama

\oldpage[III]{199}
\[
\xymatrix@C=18pt@R=32pt{
  {\begin{gathered}
    \bigl(\mathscr T_p(\mathscr O_Y)
      \otimes_{\mathscr O_Y}\mathscr M\bigr)|U\\
    =\mathscr T^U_p(\mathscr O_Y|U)
      \otimes_{\mathscr O_Y|U}(\mathscr M|U)
  \end{gathered}}
  \ar[r]\ar[d]
  &
  {\begin{gathered}
    \mathscr T^U_p(\mathscr M|U)\\
    =\bigl(\mathscr T_p(\mathscr M)\bigr)|U
  \end{gathered}}
  \ar[d]
  \\
  {\begin{gathered}
    \bigl(\mathscr T_p(\mathscr O_Y)
      \otimes_{\mathscr O_Y}\mathscr M\bigr)|V\\
    =\mathscr T^V_p(\mathscr O_Y|V)
      \otimes_{\mathscr O_Y|V}(\mathscr M|V)
  \end{gathered}}
  \ar[r]
  &
  {\begin{gathered}
    \mathscr T^V_p(\mathscr M|V)\\
    =\bigl(\mathscr T_p(\mathscr M)\bigr)|V
  \end{gathered}}
}
\]
es conmutativo por \sref{III.7.2.3.3}.

Para todo morfismo $g:Y'\to Y$ existe un homomorfismo canónico
\[
  \mathscr T_p(\mathscr O_Y)\otimes_{\mathscr O_Y}\mathscr O_{Y'}
  \longrightarrow
  \mathscr T^{Y'}_p(\mathscr O_{Y'}),
  \tag{7.7.3.2}\label{III.7.7.3.2}
\]
que no es sino el caso particular de \sref{III.6.7.11.2}, para los objetos
límite, en el caso en que $S=Y$, $v_1=f$, $v_2=1_Y$ y
$\mathscr P^{(2)}_\bullet$ se reduce al único término $\mathscr M$ de grado
$0$.

Cuando $Y=\operatorname{Spec}(A)$ e $Y'=\operatorname{Spec}(A')$ son afines,
\eqref{III.7.7.3.2} no es sino el homomorfismo de haces correspondiente al
homomorfismo canónico de $A'$-módulos definido en \sref{III.7.2.2},
\[
  T^A_p(A)\otimes_A A'
  \longrightarrow
  T^{A'}_p(A')=T^A_p(A'),
\]
como resulta fácilmente de \sref{III.6.7.11}, pues en el caso considerado se
puede tomar $\mathfrak U'^{(i)}=u_i^{-1}(\mathfrak U^{(i)})$.
\end{env}

\begin{env}[7.7.4]
\label{III.7.7.4}
Cuando $f$ es un morfismo \emph{propio}, $Y=\operatorname{Spec}(A)$ es un
esquema afín noetheriano y $\mathscr P_\bullet$ es un complejo acotado
inferiormente de $\mathscr O_X$-módulos coherentes y planos sobre $Y$, vimos en
\sref{III.6.10.5} que se puede escribir, salvo isomorfismo,
\[
  \mathscr T_p(\mathscr M)
  =\mathscr H_p
    \bigl(\mathscr L_\bullet\otimes_{\mathscr O_Y}\mathscr M\bigr),
  \qquad
  \mathscr L_\bullet=\widetilde{L_\bullet},
\]
donde $L_\bullet$ es un complejo acotado inferiormente de $A$-módulos libres
de tipo finito; por tanto, el funtor $\mathscr T_\bullet$ es del tipo estudiado
en detalle en \sref{III.7.4} y \sref{III.7.6}. Traduciremos ahora los resultados
de este estudio.
\end{env}

\begin{theorem}[7.7.5]
\label{III.7.7.5}
Sean $Y$ un preesquema localmente noetheriano, $(U_\alpha)$ un recubrimiento
de $Y$ formado por abiertos afines, $f:X\to Y$ un morfismo propio y
$\mathscr P_\bullet$ un complejo acotado inferiormente de
$\mathscr O_X$-módulos coherentes y planos sobre $Y$. El funtor homológico
$\mathscr T_\bullet(\mathscr M)$ definido por \eqref{III.7.7.1.1} posee
entonces las propiedades siguientes:

\medskip
\noindent
\textnormal{I) (La propiedad de semicontinuidad).}\footnote{Un caso particular
de este teorema se encuentra ya en la nota [3] de Chow--Igusa. La propiedad de
semicontinuidad fue descubierta, en el marco de los espacios analíticos y bajo
hipótesis bastante particulares, por Kodaira--Spencer (``On the variations of
almost-complex structures'', \emph{Algebraic Geometry and Topology, A Symposium
in honor of S. Lefschetz}, Princeton Series n\textsuperscript{o} 12,
p.~139--150, Princeton, 1957), y la versión general fue demostrada por Grauert
[5].}
\emph{La función}
\[
  y\longmapsto
  d_p(y)
  =
  \operatorname{rank}_{\kres(y)}
    \mathscr T_p^{\kres(y)}(\kres(y))
  \tag{7.7.5.1}\label{III.7.7.5.1}
\]
\emph{es semicontinua superiormente.}

\oldpage[III]{200}
\medskip
\noindent
\textnormal{II) (La propiedad de cambio de base).}
\emph{Para un entero $p$ dado, las condiciones siguientes son equivalentes:}
\begin{enumerate}
  \item[a)] $\mathscr T_p$ es exacto por la derecha.

  \item[a$'$)] $\mathscr T_p(\mathscr M)$ es isomorfo a un funtor de la forma
    $\mathscr N\otimes_{\mathscr O_Y}\mathscr M$, siendo $\mathscr N$
    necesariamente isomorfo a
    \[
      \mathscr T_p(\mathscr O_Y)
      =
      \mathscr H^{-p}(f,\mathscr P^\bullet).
    \]

  \item[a$''$)] El homomorfismo canónico funtorial
    \eqref{III.7.7.3.1} es un isomorfismo.

  \item[b)] $\mathscr T_{p-1}$ es exacto por la izquierda.

  \item[b$'$)] Existe un $\mathscr O_Y$-módulo $\mathscr Q$ (necesariamente
    coherente y determinado salvo isomorfismo único) y un isomorfismo de
    funtores
    \[
      \mathscr T_{p-1}(\mathscr M)
      \xrightarrow{\sim}
      \mathscr{H}\!om_{\mathscr O_Y}(\mathscr Q,\mathscr M).
      \tag{7.7.5.2}\label{III.7.7.5.2}
    \]

  \item[c)] Denotando por $A_\alpha$ el anillo del abierto afín $U_\alpha$,
    para todo índice $\alpha$ el funtor de $A_\alpha$-módulos
    $T_p^{A_\alpha}$ es exacto por la derecha.

  \item[d)] Para todo morfismo $g:Y'\to Y$, el homomorfismo canónico
    \[
      \mathscr T_p(\mathscr O_Y)
      \otimes_{\mathscr O_Y}\mathscr O_{Y'}
      \longrightarrow
      \mathscr T_p^{Y'}(\mathscr O_{Y'})
      \tag{7.7.5.3}\label{III.7.7.5.3}
    \]
    es un isomorfismo.
\end{enumerate}
\end{theorem}

\begin{proof}
La propiedad de semicontinuidad es local sobre $Y$ y resulta, por tanto, de
\sref{III.7.7.4} y \sref{III.7.6.9}. Es claro que a$''$) implica a$'$) y que
a$'$) implica a). La equivalencia de a), a$''$), b) y b$'$) se demostró en
\sref{III.7.3.1} y \sref{III.7.4.6}, habida cuenta de
\sref{III.7.7.4}, cuando $Y$ es afín.

Para pasar al caso general, probemos primero que a) equivale a c), lo cual
probará que la propiedad a) es local sobre $Y$; la demostración se aplica
igualmente para probar el carácter local de a$''$) y b). Como es claro que c)
implica a), todo se reduce a probar el recíproco. Basta evidentemente demostrar
que, para todo abierto afín $U$ de $Y$ y toda sucesión exacta
\[
  0\longrightarrow\mathscr F'
  \longrightarrow\mathscr F
  \longrightarrow\mathscr F''
  \longrightarrow0
\]
de $(\mathscr O_Y|U)$-módulos cuasicoherentes, existe una sucesión exacta
\[
  0\longrightarrow\mathscr G'
  \longrightarrow\mathscr G
  \longrightarrow\mathscr G''
  \longrightarrow0
\]
de $\mathscr O_Y$-módulos cuasicoherentes tal que
\[
  \mathscr F'=\mathscr G'|U,
  \qquad
  \mathscr F=\mathscr G|U,
  \qquad
  \mathscr F''=\mathscr G''|U.
\]
Ahora bien, esto resulta inmediatamente de la hipótesis de que $Y$ es
localmente noetheriano y de \sref[I]{I.9.4.2}: en efecto, se prolonga
$\mathscr F$ a un $\mathscr O_Y$-módulo cuasicoherente $\mathscr G$,
$\mathscr F'$ a un sub-$\mathscr O_Y$-módulo $\mathscr G'$ de $\mathscr G$, y
basta tomar $\mathscr G''=\mathscr G/\mathscr G'$.

Para demostrar la equivalencia de b) y b$'$) en el caso general, observemos que,
cuando $Y$ es afín, sabemos que $\mathscr Q$ está determinado salvo isomorfismo
único; si entonces $U$ es un abierto afín del esquema afín $Y$, se deduce que
existe un isomorfismo funtorial
\[
  \mathscr T_{p-1}^{U}(\mathscr M|U)
  \xrightarrow{\sim}
  \mathscr{H}\!om_{\mathscr O_Y|U}
    (\mathscr Q|U,\mathscr M|U).
\]
En el caso general, para todo abierto afín $U$ de $Y$ existe un
$(\mathscr O_Y|U)$-módulo coherente $\mathscr Q_U$ y un isomorfismo funtorial
\[
  \mathscr T_{p-1}^{U}(\mathscr M|U)
  \xrightarrow{\sim}
  \mathscr{H}\!om_{\mathscr O_Y|U}
    (\mathscr Q_U,\mathscr M|U).
\]
La observación precedente muestra que, si $V$ es un abierto afín contenido en
$U$, se tiene $\mathscr Q_U|V=\mathscr Q_V$; de ahí la existencia y unicidad
del $\mathscr O_Y$-módulo $\mathscr Q$ que verifica \eqref{III.7.7.5.2}.

Queda por demostrar la equivalencia de a) y d). Es claro que d) es local sobre
$Y$, y acabamos de ver que lo mismo ocurre con a); además, d) es también local
sobre $Y'$. Ahora bien, cuando $Y=\operatorname{Spec}(A)$ e
$Y'=\operatorname{Spec}(A')$, vimos que $T_\bullet^{A'}$ es el funtor obtenido
\oldpage[III]{201}
a partir de $T_\bullet^A$ por extensión de escalares a $A'$, y es claro entonces
que a$'$) implica que \eqref{III.7.7.5.3} es un isomorfismo.
Recíprocamente, supongamos todavía que $Y=\operatorname{Spec}(A)$ es afín y sea
$A'$ la $A$-álgebra $A\oplus M$, donde $M$ es un $A$-módulo cualquiera y la
multiplicación en $A'$ viene dada por
\[
  (a_1,m_1)(a_2,m_2)
  =
  (a_1a_2,a_1m_2+a_2m_1).
\]
Entonces
\[
  T_p^{A'}(A')
  =
  T_p(A\oplus M)
  =
  T_p(A)\oplus T_p(M).
\]
La hipótesis de que \eqref{III.7.7.5.3} sea biyectivo implica que lo mismo
ocurre con la aplicación canónica
\[
  T_p(A)\otimes_A M\longrightarrow T_p(M).
\]
En otras palabras, d) implica a$''$), lo que termina la demostración.
\end{proof}

\begin{theorem}[7.7.6]
\label{III.7.7.6}
Sean $Y$ un preesquema localmente noetheriano, $f:X\to Y$ un morfismo propio y
$\mathscr F$ un $\mathscr O_X$-módulo coherente y plano sobre $Y$. Existe
entonces un $\mathscr O_Y$-módulo coherente $\mathscr Q$ (determinado salvo
isomorfismo único) y un isomorfismo de funtores en el
$\mathscr O_Y$-módulo cuasicoherente $\mathscr M$,
\[
  f_*(\mathscr F\otimes_{\mathscr O_Y}\mathscr M)
  \xrightarrow{\sim}
  \mathscr{H}\!om_{\mathscr O_Y}(\mathscr Q,\mathscr M).
  \tag{7.7.6.1}\label{III.7.7.6.1}
\]
De ello se deduce un isomorfismo de funtores
\[
  \Gamma(X,\mathscr F\otimes_{\mathscr O_Y}\mathscr M)
  \xrightarrow{\sim}
  \operatorname{Hom}_{\mathscr O_Y}(\mathscr Q,\mathscr M).
  \tag{7.7.6.2}\label{III.7.7.6.2}
\]
\end{theorem}

\begin{proof}
En efecto, como $\mathscr M\mapsto
\mathscr F\otimes_{\mathscr O_Y}\mathscr M$ es exacto
\sref[0]{0.6.7.4} y $f_*$ es exacto por la izquierda, el funtor
\[
  \mathscr M\longmapsto
  f_*(\mathscr F\otimes_{\mathscr O_Y}\mathscr M)
\]
es exacto por la izquierda. Basta entonces aplicar la equivalencia de b) y
b$'$) en \sref{III.7.7.5} para $p=1$.
\end{proof}

\begin{corollary}[7.7.7]
\label{III.7.7.7}
Sean $Y$ un preesquema localmente noetheriano, $f:X\to Y$ un morfismo propio,
$\mathscr F$ y $\mathscr F'$ dos $\mathscr O_X$-módulos coherentes y planos
sobre $Y$, y $u:\mathscr F\to\mathscr F'$ un homomorfismo. Consideremos los
dos funtores en el $\mathscr O_Y$-módulo cuasicoherente $\mathscr M$,
\begin{align*}
  \mathscr T(\mathscr M)
  &=
  \operatorname{Ker}
  \bigl(
    f_*(\mathscr F\otimes_{\mathscr O_Y}\mathscr M)
    \longrightarrow
    f_*(\mathscr F'\otimes_{\mathscr O_Y}\mathscr M)
  \bigr),\\
  T(\mathscr M)
  &=
  \Gamma(Y,\mathscr T(\mathscr M))\\
  &=
  \operatorname{Ker}
  \bigl(
    \Gamma(X,\mathscr F\otimes_{\mathscr O_Y}\mathscr M)
    \longrightarrow
    \Gamma(X,\mathscr F'\otimes_{\mathscr O_Y}\mathscr M)
  \bigr).
\end{align*}
Entonces existe un $\mathscr O_Y$-módulo coherente $\mathscr R$ (determinado
salvo isomorfismo único) y existen isomorfismos de funtores
\[
  \mathscr T(\mathscr M)
  \xrightarrow{\sim}
  \mathscr{H}\!om_{\mathscr O_Y}(\mathscr R,\mathscr M)
  \tag{7.7.7.1}\label{III.7.7.7.1}
\]
\[
  T(\mathscr M)
  \xrightarrow{\sim}
  \operatorname{Hom}_{\mathscr O_Y}(\mathscr R,\mathscr M).
  \tag{7.7.7.2}\label{III.7.7.7.2}
\]
\end{corollary}

\begin{proof}
Basta demostrar \eqref{III.7.7.7.2}; en efecto, esto probará
\eqref{III.7.7.7.1} en el caso en que $Y$ es afín, y de ahí se pasará al caso
general razonando como en la demostración de la equivalencia de b) y b$'$) en
\sref{III.7.7.5}, gracias a la unicidad salvo isomorfismo único de un
representante de un funtor representable \sref[0]{0.8.1.8}. De
\sref{III.7.7.6} resulta que existen dos $\mathscr O_Y$-módulos coherentes
$\mathscr Q$, $\mathscr Q'$ que definen isomorfismos funtoriales
\[
  \Gamma(X,\mathscr F\otimes_{\mathscr O_Y}\mathscr M)
  \xrightarrow{\sim}
  \operatorname{Hom}_{\mathscr O_Y}(\mathscr Q,\mathscr M),
  \qquad
  \Gamma(X,\mathscr F'\otimes_{\mathscr O_Y}\mathscr M)
  \xrightarrow{\sim}
  \operatorname{Hom}_{\mathscr O_Y}(\mathscr Q',\mathscr M).
\]
Ahora bien, $u:\mathscr F\to\mathscr F'$ define canónicamente un morfismo de
funtores
\[
  \Gamma(X,\mathscr F\otimes_{\mathscr O_Y}\mathscr M)
  \longrightarrow
  \Gamma(X,\mathscr F'\otimes_{\mathscr O_Y}\mathscr M).
\]

\oldpage[III]{202}
A este último le corresponde un homomorfismo único
$v:\mathscr Q'\to\mathscr Q$ de $\mathscr O_Y$-módulos tal que el diagrama
\[
\xymatrix@C=42pt@R=28pt{
  \Gamma(X,\mathscr F\otimes_{\mathscr O_Y}\mathscr M)
    \ar[r]\ar[d]_-{\sim}
  &
  \Gamma(X,\mathscr F'\otimes_{\mathscr O_Y}\mathscr M)
    \ar[d]^-{\sim}
  \\
  \operatorname{Hom}_{\mathscr O_Y}(\mathscr Q,\mathscr M)
    \ar[r]
  &
  \operatorname{Hom}_{\mathscr O_Y}(\mathscr Q',\mathscr M)
}
\]
sea conmutativo \sref[0]{0.8.1.4}. Como el funtor contravariante
$\mathscr N\mapsto
\operatorname{Hom}_{\mathscr O_Y}(\mathscr N,\mathscr M)$ es exacto por la
izquierda en la categoría de los $\mathscr O_Y$-módulos, basta tomar
$\mathscr R=\operatorname{Coker}(v)$ para obtener el isomorfismo
\eqref{III.7.7.7.2} buscado.
\end{proof}

\begin{corollary}[7.7.8]
\label{III.7.7.8}
Bajo las hipótesis de \sref{III.7.7.6} relativas a $X$, $Y$ y $f$, sean
$\mathscr F$ y $\mathscr G$ dos $\mathscr O_X$-módulos coherentes que verifican
las condiciones siguientes:
\begin{enumerate}
  \item[(i)] $\mathscr F$ es plano sobre $Y$;
  \item[(ii)] $\mathscr G$ es isomorfo al conúcleo de un homomorfismo
    $\mathscr E_1\to\mathscr E_0$ de $\mathscr O_X$-módulos localmente libres
    de tipo finito.
\end{enumerate}
Consideremos los dos funtores en el $\mathscr O_Y$-módulo cuasicoherente
$\mathscr M$,
\begin{align*}
  \mathscr T(\mathscr M)
  &=
  f_*
  \bigl(
    \mathscr{H}\!om_{\mathscr O_X}
      (\mathscr G,\mathscr F\otimes_{\mathscr O_Y}\mathscr M)
  \bigr),\\
  T(\mathscr M)
  &=
  \Gamma(Y,\mathscr T(\mathscr M))
  =
  \operatorname{Hom}_{\mathscr O_X}
    (\mathscr G,\mathscr F\otimes_{\mathscr O_Y}\mathscr M).
\end{align*}
Entonces existe un $\mathscr O_Y$-módulo coherente $\mathscr N$ (determinado
salvo isomorfismo único) y existen isomorfismos de funtores
\[
  \mathscr T(\mathscr M)
  \xrightarrow{\sim}
  \mathscr{H}\!om_{\mathscr O_Y}(\mathscr N,\mathscr M)
  \tag{7.7.8.1}\label{III.7.7.8.1}
\]
\[
  T(\mathscr M)
  \xrightarrow{\sim}
  \operatorname{Hom}_{\mathscr O_Y}(\mathscr N,\mathscr M).
  \tag{7.7.8.2}\label{III.7.7.8.2}
\]
\end{corollary}

\begin{proof}
En virtud del isomorfismo funtorial \sref[0]{0.5.4.2.1}, se tienen, de manera
funtorial en $\mathscr M$, los isomorfismos
\begin{align*}
  \mathscr{H}\!om_{\mathscr O_X}
    (\mathscr E_i,\mathscr F\otimes_{\mathscr O_Y}\mathscr M)
  &\xrightarrow{\sim}
  \check{\mathscr E}_i
    \otimes_{\mathscr O_X}
    (\mathscr F\otimes_{\mathscr O_Y}\mathscr M)\\
  &\xrightarrow{\sim}
  (\check{\mathscr E}_i\otimes_{\mathscr O_X}\mathscr F)
    \otimes_{\mathscr O_Y}\mathscr M\\
  &\xrightarrow{\sim}
  \mathscr{H}\!om_{\mathscr O_X}(\mathscr E_i,\mathscr F)
    \otimes_{\mathscr O_Y}\mathscr M
\end{align*}
para $i=0,1$. Pongamos
$\mathscr F_i=\mathscr{H}\!om_{\mathscr O_X}(\mathscr E_i,\mathscr F)$ para
$i=0,1$; estos son $\mathscr O_X$-módulos coherentes
\sref[0]{0.5.3.5} y planos sobre $Y$ \sref[0]{0.5.4.2}. Si
$v:\mathscr E_1\to\mathscr E_0$ denota el homomorfismo dado, sea
$u=\mathscr{H}\!om(v,1_{\mathscr F}):\mathscr F_0\to\mathscr F_1$. En virtud
de la exactitud por la izquierda del funtor
$\mathscr H\mapsto\mathscr{H}\!om_{\mathscr O_X}
(\mathscr H,\mathscr F\otimes_{\mathscr O_Y}\mathscr M)$, se tienen
isomorfismos funtoriales en $\mathscr M$,
\begin{align*}
  \mathscr{H}\!om_{\mathscr O_X}
    (\mathscr G,\mathscr F\otimes_{\mathscr O_Y}\mathscr M)
  &\xrightarrow{\sim}
  \operatorname{Ker}
  \bigl(
    \mathscr{H}\!om_{\mathscr O_X}
      (\mathscr E_0,\mathscr F\otimes_{\mathscr O_Y}\mathscr M)
    \longrightarrow
    \mathscr{H}\!om_{\mathscr O_X}
      (\mathscr E_1,\mathscr F\otimes_{\mathscr O_Y}\mathscr M)
  \bigr)\\
  &\xrightarrow{\sim}
  \operatorname{Ker}
  \bigl(
    \mathscr F_0\otimes_{\mathscr O_Y}\mathscr M
    \longrightarrow
    \mathscr F_1\otimes_{\mathscr O_Y}\mathscr M
  \bigr).
\end{align*}
Puesto que $f_*$ es exacto por la izquierda, se deduce un isomorfismo funtorial
\[
  f_*
  \bigl(
    \mathscr{H}\!om_{\mathscr O_X}
      (\mathscr G,\mathscr F\otimes_{\mathscr O_Y}\mathscr M)
  \bigr)
  \xrightarrow{\sim}
  \operatorname{Ker}
  \bigl(
    f_*(\mathscr F_0\otimes_{\mathscr O_Y}\mathscr M)
    \longrightarrow
    f_*(\mathscr F_1\otimes_{\mathscr O_Y}\mathscr M)
  \bigr),
\]
y basta entonces aplicar \sref{III.7.7.7}.
\end{proof}

\oldpage[III]{203}
\begin{remarks}[7.7.9]
\label{III.7.7.9}
\textnormal{(i)} En \sref{III.7.7.6}, \sref{III.7.7.7} y
\sref{III.7.7.8}, la formación de los $\mathscr O_Y$-módulos $\mathscr Q$,
$\mathscr R$ y $\mathscr N$ \emph{conmuta con los cambios de base}. Por ejemplo,
conservando las notaciones de \sref{III.7.7.2}, en el caso de
\sref{III.7.7.6} se tiene, para todo $\mathscr O_{Y'}$-módulo cuasicoherente
$\mathscr M'$, el isomorfismo
\[
  f'_*(\mathscr F\otimes_{\mathscr O_Y}\mathscr M')
  \xrightarrow{\sim}
  \mathscr{H}\!om_{\mathscr O_{Y'}}(g^*(\mathscr Q),\mathscr M'),
\]
pues, en virtud de la observación hecha en \sref{III.7.7.2}, todo se reduce a
ver que se tiene
\[
  \operatorname{Hom}_{\mathscr O_Y}(\mathscr Q,g_*(\mathscr M'))
  =
  \operatorname{Hom}_{\mathscr O_{Y'}}(g^*(\mathscr Q),\mathscr M'),
\]
lo cual no es sino \sref[0]{0.4.4.3.1}. Del mismo modo, cuando en
\sref{III.7.7.7} se sustituyen $Y$, $f$, $\mathscr M$, $\mathscr F$ y
$\mathscr F'$ por $Y'$, $f'$, $\mathscr M'$,
$\mathscr F\otimes_{\mathscr O_X}\mathscr O_{X'}$ y
$\mathscr F'\otimes_{\mathscr O_X}\mathscr O_{X'}$, respectivamente, hay que
sustituir $\mathscr R$ por $g^*(\mathscr R)$. Finalmente, en
\sref{III.7.7.8}, cuando se sustituyen $X$, $Y$, $f$, $\mathscr M$ y
$\mathscr F$ por $X'$, $Y'$, $f'$, $\mathscr M'$ y
$\mathscr F\otimes_{\mathscr O_X}\mathscr O_{X'}$, y $\mathscr G$ por
\[
  \mathscr G'=\mathscr G\otimes_{\mathscr O_X}\mathscr O_{X'},
\]
hay que sustituir $\mathscr N$ por $g^*(\mathscr N)$. Esto resulta de que se
tiene todavía una sucesión exacta
\[
  \mathscr E'_1\longrightarrow
  \mathscr E'_0\longrightarrow
  \mathscr G'\longrightarrow0
  \qquad\text{con}\qquad
  \mathscr E'_i=\mathscr E_i\otimes_{\mathscr O_X}\mathscr O_{X'},
\]
y de que
\[
  \check{\mathscr E}'_i
  \otimes_{\mathscr O_{X'}}
  (\mathscr F\otimes_{\mathscr O_Y}\mathscr M')
  =
  \check{\mathscr E}_i
  \otimes_{\mathscr O_X}
  (\mathscr F\otimes_{\mathscr O_Y}\mathscr M')
  \qquad(i=0,1).
\]

\medskip
\textnormal{(ii)} La condición (ii) del enunciado de \sref{III.7.7.8}
relativa a $\mathscr G$ se satisface siempre para cualquier
$\mathscr O_X$-módulo coherente $\mathscr G$ cuando existe un
$\mathscr O_X$-módulo invertible $Y$-amplio, por ejemplo cuando $Y$ es afín y
$f:X\to Y$ es un morfismo proyectivo. Basta observar, teniendo en cuenta
\sref[II]{II.5.5.1}, que existe un $\mathscr O_X$-módulo localmente libre de
tipo finito $\mathscr E_0$ tal que $\mathscr G$ es isomorfo a un cociente de
$\mathscr E_0$ \sref[II]{II.2.7.10}. Como $\mathscr E_0$ y $\mathscr G$ son
coherentes, también lo es el núcleo $\mathscr G_1$ de
$\mathscr E_0\to\mathscr G$; aplicando el mismo resultado a $\mathscr G_1$, se
obtiene una sucesión exacta
\[
  \mathscr E_1\longrightarrow
  \mathscr E_0\longrightarrow
  \mathscr G\longrightarrow0,
\]
donde $\mathscr E_0$ y $\mathscr E_1$ son localmente libres de tipo finito.

\medskip
\textnormal{(iii)} Demostraremos en el capítulo V que, en
\sref{III.7.7.8}, la hipótesis restrictiva (ii) es \emph{superflua}.
\end{remarks}

\begin{proposition}[7.7.10]
\label{III.7.7.10}
\textnormal{(Criterios locales para la propiedad de cambio de base).} Bajo las
condiciones generales de \sref{III.7.7.5}, sean $y$ un punto de $Y$ y $p$ un
entero. Las propiedades siguientes son equivalentes:
\begin{enumerate}
  \item[a)] El funtor $T_p^{\mathscr O_y}$ es exacto por la derecha.

  \item[b)] El homomorfismo canónico
    \[
      T_p^{\mathscr O_y}(\mathscr O_y)
      \longrightarrow
      T_p^{\mathscr O_y}(\kres(y))
    \]
    es sobreyectivo.

  \item[c)] Para todo entero $n$, el homomorfismo canónico
    \[
      T_p^{\mathscr O_y}
      (\mathscr O_y/\mathfrak m_y^{n+1})
      \longrightarrow
      T_p^{\mathscr O_y}(\kres(y))
    \]
    es sobreyectivo.
\end{enumerate}
Además, el conjunto de los $y\in Y$ que verifican estas condiciones es el
mayor abierto $U$ de $Y$ tal que $\mathscr T_p^U$ sea exacto por la derecha.
\end{proposition}

\begin{proof}
Habida cuenta de \sref{III.7.7.4}, la equivalencia de a), b) y c) resulta de
\sref{III.7.4.7} y \sref{III.7.5.2}. El hecho de que el conjunto $U$ donde
$T_p^{\mathscr O_y}$ es exacto por la derecha sea abierto es también
consecuencia de \sref{III.7.4.4}; recíprocamente, si $\mathscr T_p^V$ es exacto
por la derecha, lo mismo ocurre con $T_p^{\mathscr O_y}$ para todo $y\in V$,
por la condición c) de \sref{III.7.7.5} y \sref{III.7.6.1}.
\end{proof}

\begin{corollary}[7.7.11]
\label{III.7.7.11}
Si $\mathscr T_p$ es exacto por la derecha (respectivamente, por la izquierda),
entonces, para todo morfismo $g:Y'\to Y$, $\mathscr T_p^{Y'}$ es exacto por la
derecha (respectivamente, por la izquierda). El recíproco es cierto cuando el
morfismo $g$ es fielmente plano.
\end{corollary}

\oldpage[III]{204}
\begin{proof}
La primera aserción es consecuencia inmediata de \sref{III.7.6.1} y del hecho
de que la cuestión es local sobre $Y$ y $Y'$, por las condiciones c) y b) de
\sref{III.7.7.5}. Para demostrar la segunda aserción basta ver, en virtud de
\sref{III.7.7.10}, que para todo $y\in Y$, $T_p^{\mathscr O_y}$ es exacto por
la derecha (respectivamente, por la izquierda). Pero, por hipótesis, existe
$y'\in Y'$ tal que $g(y')=y$, y $\mathscr O_{y'}$ es un
$\mathscr O_y$-módulo fielmente plano; la conclusión resulta entonces de la
hipótesis y de \sref{III.7.6.1}.
\end{proof}

\begin{remarks}[7.7.12]
\label{III.7.7.12}
\textnormal{(i)} Bajo las hipótesis de \sref{III.7.7.4}, supongamos además que
$\mathscr P_\bullet$ es un complejo \emph{finito}. Resulta entonces de
\sref{III.7.7.1}, puesto que se puede tomar para $\mathfrak U$ un
recubrimiento abierto afín \emph{finito} de $X$, que el bicomplejo
$C^\bullet(\mathfrak U,
\mathscr P^\bullet\otimes_{\mathscr O_Y}\mathscr M)$ es también finito y, más
precisamente, que existe un conjunto finito $E$ de pares, \emph{independiente
de $\mathscr M$}, tal que
\[
  C^h(\mathfrak U,\mathscr P^k\otimes_{\mathscr O_Y}\mathscr M)=0
  \qquad\text{para todos los pares }(h,k)\notin E.
\]
Se concluye que existe $i_1$ tal que, para $i\geq i_1$, se tiene
$\mathscr T_i(\mathscr M)=0$ para \emph{todo} $\mathscr O_Y$-módulo
cuasicoherente $\mathscr M$. En particular, $\mathscr T_i$ es trivialmente un
funtor exacto en $\mathscr M$ para estos valores de $i$ y, por consiguiente,
por \sref{III.7.4.1}, $Z'_i(L_\bullet)$ es un $A$-módulo \emph{plano} de tipo
finito (y, por tanto, \emph{proyectivo} de tipo finito, puesto que $A$ es
noetheriano) para estos valores de $i$.

Consideremos entonces el complejo $L'_\bullet$ de $A$-módulos tal que
\[
  L'_i=L_i\quad(i<i_1),
  \qquad
  L'_{i_1}=Z'_{i_1}(L_\bullet),
  \qquad
  L'_i=0\quad(i>i_1),
\]
y sea $\mathscr L'_\bullet=\widetilde{L'_\bullet}$. Es claro que
\[
  \mathscr H_i(\mathscr L'_\bullet\otimes_{\mathscr O_Y}\mathscr M)
  =
  \mathscr H_i(\mathscr L_\bullet\otimes_{\mathscr O_Y}\mathscr M)
\]
para $i<i_1-1$ y también para $i\geq i_1$ (siendo entonces nulos ambos
miembros). Finalmente, como
\[
  \operatorname{Im}(Z'_{i_1}\otimes_A M)
  =
  \operatorname{Im}(L_{i_1}\otimes_A M)
\]
por definición, se tiene asimismo
\[
  \mathscr H_i(\mathscr L'_\bullet\otimes_{\mathscr O_Y}\mathscr M)
  =
  \mathscr H_i(\mathscr L_\bullet\otimes_{\mathscr O_Y}\mathscr M)
\]
para $i=i_1-1$. Se ve, pues, que se puede suponer en \sref{III.7.7.4} que
$\mathscr L_\bullet$ es también un complejo \emph{finito}, a condición de
exigir solamente que los $\mathscr L_i$ sean $\mathscr O_Y$-módulos
\emph{localmente libres}, asociados a $A$-módulos proyectivos de tipo finito.

Este razonamiento se aplica, en particular, al caso en que
$\mathscr P_\bullet$ se reduce a \emph{un solo} término $\mathscr F\ne0$, de
grado $0$ (en cuyo caso
$\mathscr T_n(\mathscr M)=\mathrm R^{-n}f_*
(\mathscr F\otimes_{\mathscr O_Y}\mathscr M)$). Se puede entonces suponer que
los $\mathscr L_i$ son nulos para $i>0$; en este caso se utilizarán
preferentemente las notaciones cohomológicas, escribiendo $\mathscr T^{-p}$ en
lugar de $\mathscr T_p$.

\medskip
\textnormal{(ii)} Cuando en el enunciado de \sref{III.7.7.5} ya no se supone
que los $\mathscr P_i$ sean planos sobre $Y$, las conclusiones siguen siendo
válidas a condición de poner esta vez
\[
  \mathscr T_p(\mathscr M)
  =
  \operatorname{Tor}^Y_p
    (f,1_Y;\mathscr P_\bullet,\mathscr M).
  \tag{7.7.12.1}\label{III.7.7.12.1}
\]
En efecto,
$\operatorname{Tor}^Y_n(f,1_Y;\mathscr P_\bullet,\mathscr O_Y)$ es entonces un
$\mathscr O_Y$-módulo \emph{coherente} en virtud de \sref{III.6.7.9}. La
demostración de \sref{III.6.10.5} se aplica sin cambio, teniendo en cuenta
\sref{III.6.10.1}, y muestra también que, cuando
$Y=\operatorname{Spec}(A)$ es afín, se tiene
\[
  \mathscr T_p(\mathscr M)
  =
  \mathscr H_p(\mathscr L_\bullet\otimes_{\mathscr O_Y}\mathscr M),
  \qquad
  \mathscr L_\bullet=\widetilde{L_\bullet},
\]
donde $L_\bullet$ es un complejo de $A$-módulos libres de tipo finito; esto
prueba la aserción.
\end{remarks}

\subsection{Aplicación a los morfismos propios: II. Criterios de planitud cohomológica}
\label{subsection:III.7.8}
\label{III.7.8}

\begin{definition}[7.8.1]
\label{III.7.8.1}
Sean $X$ e $Y$ dos preesquemas, sea $f:X\to Y$ un morfismo
cuasicompacto y separado, sea $\mathscr P_\bullet$ un complejo de
$\mathscr O_X$-módulos cuasicoherentes y planos sobre $Y$, y sea
$\mathscr T_\bullet$ el funtor homológico de $\mathscr O_Y$-módulos
cuasicoherentes definido por \eqref{III.7.7.1.1}. Sea $y$ un punto
de $Y$.

Se dice que $\mathscr P_\bullet$ es
\oldpage[III]{205}
\emph{homológicamente plano sobre
$Y$ en el punto $y$, en dimensión $p$} (o \emph{cohomológicamente
plano sobre $Y$ en el punto $y$, en dimensión $-p$}) si existe un
entorno abierto $U$ de $y$ en $Y$ tal que
\[
  \mathscr T_p^U=\mathscr T_U^{-p}
\]
es exacto.

Se dice que $\mathscr P_\bullet$ es \emph{homológicamente plano en
dimensión $p$ sobre $Y$} (o \emph{cohomológicamente plano en
dimensión $-p$ sobre $Y$}) si es homológicamente plano sobre $Y$ en
todo punto $y\in Y$, en dimensión $p$.

Cuando $\mathscr P_\bullet$ es homológicamente plano sobre $Y$
(respectivamente, sobre $Y$ en el punto $y$) en \emph{toda}
dimensión $p$, se dice simplemente que $\mathscr P_\bullet$ es
\emph{homológicamente plano sobre $Y$} (respectivamente, sobre $Y$
en el punto $y$), o \emph{cohomológicamente plano sobre $Y$}
(respectivamente, sobre $Y$ en el punto $y$).
\end{definition}

\begin{env}[7.8.2]
\label{III.7.8.2}
Por definición, la noción de planitud homológica sobre $Y$ es
\emph{local sobre $Y$}.

Si $Y$ es localmente noetheriano, o si $Y$ es un esquema, para que
$\mathscr P_\bullet$ sea homológicamente plano sobre $Y$ en
dimensión $p$, es necesario y suficiente que el funtor
$\mathscr T_p$ sea exacto. La demostración para $Y$ localmente
noetheriano se dio durante la demostración de \sref{III.7.7.5};
cuando $Y$ es un esquema, el razonamiento es el mismo, basado en
\sref[I]{I.9.4.2} aplicado a un abierto afín de un esquema
cuasicompacto.
\end{env}

\begin{proposition}[7.8.3]
\label{III.7.8.3}
Con las notaciones e hipótesis de \sref{III.7.8.1}, las condiciones
siguientes son equivalentes:
\begin{enumerate}
  \item[a)] $\mathscr P_\bullet$ es homológicamente plano sobre $Y$
    en el punto $y$, en dimensión $p$.
  \item[b)] Existe un entorno abierto $U$ de $y$ en $Y$ tal que
    $\mathscr T_p^U$ y $\mathscr T_{p+1}^U$ son exactos por la
    derecha.
  \item[c)] Existe un entorno abierto $U$ de $y$ en $Y$ tal que
    $\mathscr T_p^U$ y $\mathscr T_{p-1}^U$ son exactos por la
    izquierda.
  \item[d)] Existe un entorno abierto $U$ de $y$ en $Y$ tal que
    $\mathscr T_{p+1}^U$ es exacto por la derecha y
    $\mathscr T_{p-1}^U$ es exacto por la izquierda.
\end{enumerate}
\end{proposition}

\begin{proof}
Teniendo en cuenta la interpretación de $\mathscr T_p$ cuando $Y$
es afín, esto no es más que una traducción de una parte de
\sref{III.7.3.3}.
\end{proof}

\begin{proposition}[7.8.4]
\label{III.7.8.4}
Sean $Y$ un preesquema localmente noetheriano, $f:X\to Y$ un
morfismo propio, $\mathscr P_\bullet$ un complejo acotado
inferiormente de $\mathscr O_X$-módulos coherentes y planos sobre
$Y$, y $\mathscr T_\bullet$ el funtor definido por
\eqref{III.7.7.1.1}. Para todo $y\in Y$, las condiciones siguientes
son equivalentes:
\begin{enumerate}
  \item[a)] $\mathscr P_\bullet$ es homológicamente plano sobre $Y$
    en $y$, en dimensión $p$.
  \item[b)] El funtor $\mathscr T_p^{\mathscr O_y}$ es exacto.
  \item[c)] Existe un entero $n_0$ tal que, para $n\geq n_0$,
    \[
      \operatorname{length}
        \mathscr T_p^{\mathscr O_y}
          \bigl(\mathscr O_y/\mathfrak m_y^{\,n+1}\bigr)
      =
      \operatorname{length}
        \mathscr T_p^{\mathscr O_y}\bigl(\kres(y)\bigr)
      \cdot
      \operatorname{length}
        \bigl(\mathscr O_y/\mathfrak m_y^{\,n+1}\bigr).
      \tag{7.8.4.1}\label{III.7.8.4.1}
    \]
    Aquí todas las longitudes son longitudes de
    $\mathscr O_y$-módulos.
  \item[d)] Existe un entorno abierto $U$ de $y$ tal que
    \[
      \mathscr H^{-p}(f,\mathscr P^\bullet)|U
    \]
    es isomorfo a un $(\mathscr O_Y|U)$-módulo de la forma
    $(\mathscr O_Y|U)^m$ y, para todo $(\mathscr O_Y|U)$-módulo
    cuasicoherente $\mathscr M$, el homomorfismo canónico
    \[
      \mathscr H^{-p}(f,\mathscr P^\bullet)|U
        \otimes_{\mathscr O_Y|U}\mathscr M
      \longrightarrow
      \mathscr H^{-p}
        \bigl(f,(\mathscr P^\bullet|U)
          \otimes_{\mathscr O_Y|U}\mathscr M\bigr)
      \tag{7.8.4.2}\label{III.7.8.4.2}
    \]
    es biyectivo.
\end{enumerate}
Cuando se cumplen estas condiciones, se tiene además:
\begin{enumerate}
  \item[e)] Existe un entorno de $y$ en el cual la función
    $z\mapsto d_p(z)$, definida en \eqref{III.7.7.5.1}, es
    constante.
\end{enumerate}
Además, si $Y$ es reducido en $y$ \sref[0]{0.4.1.4}, la condición
e) es equivalente a las demás condiciones.
\end{proposition}

\oldpage[III]{206}
\begin{proof}
En efecto, la condición b) equivale a decir que
$\mathscr T_p^{\mathscr O_y}$ y $\mathscr T_{p+1}^{\mathscr O_y}$
son exactos por la derecha \sref{III.7.3.3}. La equivalencia de a)
y b) resulta entonces de \sref{III.7.7.10} y \sref{III.7.8.3}.

Como $\mathscr O_y/\mathfrak m_y^{\,n+1}$ es artiniano, y
\[
  \mathscr T_p^{\mathscr O_y}
    \bigl(\mathscr O_y/\mathfrak m_y^{\,n+1}\bigr)
  \quad\hbox{y}\quad
  \mathscr T_p^{\mathscr O_y}\bigl(\kres(y)\bigr)
\]
son $(\mathscr O_y/\mathfrak m_y^{\,n+1})$-módulos de tipo finito
\sref{III.7.7.4}, y por tanto de longitud finita, la equivalencia de
b) y c) resulta nuevamente de \sref{III.7.7.10} y
\sref{III.7.3.5.7}.

Que a) implica e), y que a) y e) son equivalentes cuando
$\mathscr O_y$ es reducido, resulta de \sref{III.7.6.9}.
Finalmente, a) implica que \eqref{III.7.8.4.2} es biyectivo, por la
definición \sref{III.7.8.1} y \sref{III.7.7.5}. También implica que
\[
  \bigl(\mathscr H^{-p}(f,\mathscr P^\bullet)\bigr)_y
\]
es un $\mathscr O_y$-módulo plano por la condición (f) de
\sref{III.7.3.3}, y por tanto un $\mathscr O_y$-módulo libre de
tipo finito \sref[0]{0.10.1.3}, pues es de tipo finito por
\sref{III.7.7.4}. Como $\mathscr H^{-p}(f,\mathscr P^\bullet)$ es
un $\mathscr O_Y$-módulo coherente \sref{III.7.7.4}, es localmente
libre en un entorno de $y$ \sref[0]{0.5.2.7}. Recíprocamente, d)
implica claramente a), por la definición del funtor
$\mathscr T_p^U$ \eqref{III.7.7.2.2}.
\end{proof}

\begin{proposition}[7.8.5]
\label{III.7.8.5}
Bajo las hipótesis de \sref{III.7.8.4}, las condiciones siguientes
son equivalentes:
\begin{enumerate}
  \item[a)] $\mathscr P_\bullet$ es homológicamente plano sobre $Y$
    en toda dimensión $i\leq p$.
  \item[b)] Para $i\leq p+1$, los funtores $\mathscr T_i$ son
    exactos por la derecha.
  \item[c)] Para $i\geq-p$, los $\mathscr O_Y$-módulos
    $\mathscr H^i(f,\mathscr P^\bullet)$ son localmente libres.
\end{enumerate}
\end{proposition}

\begin{proof}
La equivalencia de a) y b) es inmediata por \sref{III.7.8.3}, y a)
implica c) por \sref{III.7.8.4}.

Recíprocamente, supongamos que se cumple c). Recordemos también que
$\mathscr L_i=0$ para $i<i_0$ \sref{III.7.7.4}, y por tanto
$\mathscr T_i=0$ para $i\leq i_0$. Todo punto $y\in Y$ tiene, pues,
un entorno afín $U=\Spec(A)$ tal que $T_i^A(A)$ es un $A$-módulo
libre para $i\leq p$. De \sref{III.7.3.7} se deduce que
$\mathscr T_i^U=T_i^A$ es exacto para $i\leq p$.
\end{proof}

Aplicaremos principalmente los criterios de planitud cohomológica al
caso en que el complejo $\mathscr P_\bullet$ está concentrado en un
\emph{solo} $\mathscr O_X$-módulo coherente $\mathscr F$, plano
sobre $Y$, tomado como $\mathscr P_0$. Recordemos que en este caso
\[
  \mathscr T_p(\mathscr M)
  =R^{-p}f_*
    \bigl(\mathscr F\otimes_{\mathscr O_Y}\mathscr M\bigr).
\]

\begin{proposition}[7.8.6]
\label{III.7.8.6}
Sean $Y$ un preesquema localmente noetheriano, $f:X\to Y$ un
morfismo propio y plano, y $y$ un punto de $Y$. Denotemos por $X_y$
la fibra $f^{-1}(y)=X\otimes_Y\kres(y)$. Supongamos que
\[
  \Gamma(X_y,\mathscr O_{X_y})=R
\]
es una $\kres(y)$-álgebra separable (Bourbaki,
\emph{Algèbre}, cap.~VIII, \S~7, n\textsuperscript{o}~5), es decir,
un producto directo de un número finito de extensiones separables de
grado finito de $\kres(y)$. Entonces $\mathscr O_X$ es
cohomológicamente plano sobre $Y$ en el punto $y$, en dimensión $0$.
\end{proposition}

\begin{proof}
Por \sref{III.7.8.4}, podemos restringirnos al caso en que $Y$ es
el espectro del anillo local $A=\mathscr O_y$. La planitud de $f$
implica $\mathscr T_{-1}=0$, de modo que $T_0^A$ ya es exacto por la
izquierda y solo queda probar que es exacto por la derecha. Por la
condición c) de \sref{III.7.7.10}, podemos reducirnos incluso al
caso en que $A=\mathscr O_y$ es artiniano.

Sea $k'$ una extensión finita de $\kres(y)$ que descomponga $R$, de
modo que
\[
  R\otimes_{\kres(y)}k'
\]
es un producto directo de un número finito de cuerpos isomorfos a
$k'$. Existe un homomorfismo local de $A$ en un anillo local $A'$
tal que $A'$ es una $A$-álgebra libre finita y su cuerpo residual es
isomorfo a $k'$ \sref[0]{0.10.3.2}. Por \sref{III.7.6.1}, basta,
pues, probar que $T_0^{A'}$ es exacto por la derecha; en otras
palabras, podemos suponer que $R$ es un producto directo de $m$
cuerpos isomorfos a $\kres(y)$.

Utilizaremos el siguiente lema elemental.

\begin{lemma}[7.8.6.1]
\label{III.7.8.6.1}
Sea $Z$ un espacio localmente anillado. Para que $Z$ sea conexo,
\oldpage[III]{207}
es necesario y suficiente que el anillo $\Gamma(Z,\mathscr O_Z)$
no sea producto de dos anillos no nulos.
\end{lemma}

\begin{proof}
Es claro que, si $Z$ es la unión de dos abiertos no vacíos y
disjuntos $Z_1$ y $Z_2$, entonces $\Gamma(Z,\mathscr O_Z)$ es
isomorfo al producto de los dos anillos no nulos
$\Gamma(Z_1,\mathscr O_Z)$ y $\Gamma(Z_2,\mathscr O_Z)$.

Recíprocamente, decir que $\Gamma(Z,\mathscr O_Z)$ es tal producto
equivale a decir que contiene un idempotente $s$ distinto de $0$ y
de $1$. Para todo $z\in Z$, el germen $s_z$ es entonces un
idempotente del anillo local $\mathscr O_z$, y por tanto es igual a
$0$ o a $1$. El conjunto de los $z$ tales que $s_z=0$ es abierto.
Por otra parte, $s_z=1$ implica por definición que $s(z)\ne0$, y el
conjunto de los $z$ tales que $s_z=1$ es también abierto
\sref[0]{0.5.5.2}. De aquí se sigue la conclusión.
\end{proof}

Del lema se deduce que $X_y$ tiene exactamente $m$ componentes
conexas $X_i'$ y que
\[
  \Gamma(X_i',\mathscr O_{X_i'})=\kres(y)
\]
para todo $i$. Como $A$ es local y artiniano, su espectro consta de
un solo punto y, por consiguiente, $X$ y $X_y$ tienen el mismo
espacio subyacente. Así, $X$ tiene $m$ componentes conexas $X_i$
tales que
\[
  X_i'=X_i\otimes_Y\kres(y).
\]
Quedamos finalmente reducidos al caso $R=\kres(y)$. Por la
condición b) de \sref{III.7.7.10}, queda probar que el homomorfismo
canónico
\[
  \Gamma(X,\mathscr O_X)\longrightarrow
  \Gamma(X_y,\mathscr O_{X_y})
\]
es sobreyectivo. Esto es inmediato, pues el compuesto
\[
  \Gamma(Y,\mathscr O_Y)=A
  \longrightarrow\Gamma(X,\mathscr O_X)
  \longrightarrow\Gamma(X_y,\mathscr O_{X_y})=\kres(y)
\]
ya es sobreyectivo.
\end{proof}

\begin{corollary}[7.8.7]
\label{III.7.8.7}
Bajo las hipótesis de \sref{III.7.8.6}, existe un entorno abierto
$U$ de $y$ tal que:
\begin{enumerate}
  \item[(i)] $f_*(\mathscr O_X)|U$ es isomorfo a un
    $(\mathscr O_Y|U)$-módulo de la forma $(\mathscr O_Y|U)^m$;
  \item[(ii)] para todo $z\in U$, el homomorfismo canónico
    \[
      \bigl(f_*(\mathscr O_X)\bigr)_z
        \otimes_{\mathscr O_z}\kres(z)
      \longrightarrow
      \Gamma(X_z,\mathscr O_{X_z})
    \]
    es biyectivo.
\end{enumerate}
\end{corollary}

\begin{proof}
La aserción (i) resulta de \sref{III.7.8.6} y \sref{III.7.8.4}.
La aserción (ii) resulta de la exactitud de $\mathscr T_0^U$, para
una elección adecuada de $U$, y de \eqref{III.7.7.5.3}.
\end{proof}

\begin{corollary}[7.8.8]
\label{III.7.8.8}
Supongamos que se cumplen las hipótesis de \sref{III.7.8.6} y,
además, que
\[
  \Gamma(X_y,\mathscr O_{X_y})=\kres(y).
\]
Entonces existe un entorno abierto $U$ de $y$ tal que el
homomorfismo canónico
\[
  \mathscr O_Y|U\longrightarrow f_*(\mathscr O_X)|U
\]
es biyectivo.
\end{corollary}

\begin{proof}
En efecto, de la aserción (ii) de \sref{III.7.8.7} se deduce que el
entero $m$ que figura en la aserción (i) de ese corolario es
necesariamente igual a $1$.
\end{proof}

\begin{corollary}[7.8.9]
\label{III.7.8.9}
Bajo las hipótesis de \sref{III.7.8.6}, existen un entorno abierto
$U$ de $y$, un $\mathscr O_U$-módulo coherente $\mathscr Q$,
determinado salvo isomorfismo único, y un isomorfismo de funtores en
el $\mathscr O_U$-módulo cuasicoherente $\mathscr M$,
\[
  R^1f_*\bigl(f^*(\mathscr M)\bigr)
  \xrightarrow{\sim}
  \mathscr Hom_{\mathscr O_U}(\mathscr Q,\mathscr M).
  \tag{7.8.9.1}\label{III.7.8.9.1}
\]
\end{corollary}

\begin{proof}
La hipótesis implica que $\mathscr T_0^U$ es exacto para un $U$
adecuado. Basta, por tanto, aplicar la equivalencia de las
condiciones a) y b') de \sref{III.7.7.5}, con $p=0$ y tomando como
$\mathscr P_\bullet$ el complejo concentrado en grado $0$ cuyo
término es $\mathscr O_X$.
\end{proof}

\begin{remarks}[7.8.10]
\label{III.7.8.10}
\textnormal{(i)}
Bajo las hipótesis de \sref{III.7.8.6}, consideremos la
factorización de Stein de $f$ \sref{III.4.3.3},
\[
  \xymatrix{
    X \ar[r]^{f'} & Y' \ar[r]^{g} & Y .
  }
\]
\oldpage[III]{208}
Aquí
\[
  Y'=\Spec\bigl(f_*(\mathscr O_X)\bigr).
\]
El morfismo finito $g$ es tal que
\[
  g_*(\mathscr O_{Y'})=f_*(\mathscr O_X)
\]
es localmente libre cerca de $y$, y su fibra sobre $y$ es el
espectro de una álgebra separable sobre $\kres(y)$
\sref[II]{II.1.5.1}. Deduciremos en el capítulo~IV que existe un
entorno abierto $U$ de $y$ en $Y$ tal que, para la restricción
$g^{-1}(U)\to U$ de $g$, \emph{toda} fibra $g^{-1}(z)$, con
$z\in U$, es el espectro de una álgebra separable sobre $\kres(z)$.
Esto es lo que llamaremos un \emph{revestimiento étale} de $U$. Se
deducirá entonces de la aserción (ii) de \sref{III.7.8.7} que la
hipótesis sobre el punto $y$ en \sref{III.7.8.6} se cumple también
en todo punto de un entorno de $y$.

\medskip
\textnormal{(ii)}
Veremos en el capítulo~V que, incluso si $X$ es proyectivo sobre $Y$
(e incluso si además es ``simple'' sobre $Y$, propiedad que se
definirá en el capítulo~IV), el $\mathscr O_U$-módulo $\mathscr Q$
de \sref{III.7.8.9} no es necesariamente localmente libre. En otras
palabras, bajo estas hipótesis $\mathscr O_X$ no es necesariamente
cohomológicamente plano en \emph{dimensión $1$} sobre $Y$ en el
punto $y$. En el capítulo~V interpretaremos $\mathscr Q$ como el haz
de $1$-formas diferenciales del esquema de Picard de $X$ relativo a
$Y$, a lo largo de la sección unidad.
\end{remarks}

\subsection{Aplicación a los morfismos propios: III. Invariancia de la
característica de Euler-Poincaré y del polinomio de Hilbert}
\label{subsection:III.7.9}
\label{III.7.9}

\begin{env}[7.9.1]
\label{III.7.9.1}
Sean $A$ un anillo y $M$ un $A$-módulo proyectivo de tipo finito.
Recordemos (Bourbaki, \emph{Alg.~comm.}, cap.~II, \S5, n.\textsuperscript{o}~2)
que esto equivale a decir que el $\mathscr O_X$-módulo $\widetilde M$
asociado a $M$ sobre $X=\Spec(A)$ es localmente libre de tipo finito.
Para todo $\mathfrak p\in\Spec(A)$, se llama \emph{rango de $M$ en
$\mathfrak p$}, y se denota por $\operatorname{rank}_{\mathfrak p}(M)$,
al rango del $A_{\mathfrak p}$-módulo libre $M_{\mathfrak p}$ (o,
equivalentemente, al rango en $\mathfrak p$ del $\mathscr O_X$-módulo
localmente libre $\widetilde M$). Por tanto,
\[
  \operatorname{rank}_{\mathfrak p}M
  =
  \operatorname{rank}_{\mathfrak p}(M_{\mathfrak p})
  =
  \operatorname{rank}_{\kres(\mathfrak p)}
    \bigl(M\otimes_A\kres(\mathfrak p)\bigr).
  \tag{7.9.1.1}\label{III.7.9.1.1}
\]
\end{env}

\begin{proposition}[7.9.2]
\label{III.7.9.2}
Sea $P_\bullet$ un complejo finito de $A$-módulos proyectivos de
tipo finito y, para todo $A$-módulo $M$, pongamos
\[
  T_\bullet(M)=H_\bullet(P_\bullet\otimes_A M).
\]
Entonces, para todo $\mathfrak p\in\Spec(A)$, se tiene
\[
  \sum_i(-1)^i
    \operatorname{rank}_{\kres(\mathfrak p)}
      T_i\bigl(\kres(\mathfrak p)\bigr)
  =
  \sum_i(-1)^i\operatorname{rank}_{\mathfrak p}(P_i).
  \tag{7.9.2.1}\label{III.7.9.2.1}
\]
\end{proposition}

\begin{proof}
En efecto, por definición,
$T_i(\kres(\mathfrak p))
 =H_i(P_\bullet\otimes_A\kres(\mathfrak p))$.
Teniendo en cuenta \eqref{III.7.9.1.1}, la fórmula
\eqref{III.7.9.2.1} no es sino la invariancia de la característica
de Euler-Poincaré de un complejo finito de espacios vectoriales de
dimensión finita al pasar a la homología
\sref[0]{0.11.10.2}.
\end{proof}

\begin{corollary}[7.9.3]
\label{III.7.9.3}
La función
\[
  \mathfrak p\longmapsto
  \sum_i(-1)^i
    \operatorname{rank}_{\kres(\mathfrak p)}
      T_i\bigl(\kres(\mathfrak p)\bigr)
\]
es localmente constante en $\Spec(A)$.
\end{corollary}

\begin{theorem}[7.9.4]
\label{III.7.9.4}
Sean $Y$ un preesquema localmente noetheriano,
$f:X\to Y$ un morfismo propio y $\mathscr P_\bullet$ un complejo
finito de $\mathscr O_X$-módulos coherentes y planos sobre $Y$.
Si se pone
\[
  \mathscr T_\bullet(\mathscr M)
  =
  \mathscr H^{-\bullet}
    \bigl(f,\mathscr P^\bullet
      \otimes_{\mathscr O_Y}\mathscr M\bigr)
  \qquad\text{(cf.~\eqref{III.7.7.1.1})},
\]
entonces la función

\oldpage[III]{209}
\[
  y\longmapsto
  \sum_i(-1)^i
    \operatorname{rank}_{\kres(y)}
      \mathscr T_i\bigl(\kres(y)\bigr)
  \tag{7.9.4.1}\label{III.7.9.4.1}
\]
es localmente constante en $Y$.
\end{theorem}

\begin{proof}
Podemos limitarnos al caso en que $Y=\Spec(A)$ es afín y $A$ es
noetheriano. Como el complejo $\mathscr P_\bullet$ es finito,
\sref{III.7.7.12}{(i)} muestra que
\[
  \mathscr T_p(\mathscr M)
  =
  \mathscr H_p(\mathscr L_\bullet
    \otimes_{\mathscr O_Y}\mathscr M),
\]
donde $\mathscr L_\bullet=\widetilde{L}_\bullet$ y
$L_\bullet$ es un complejo finito de $A$-módulos proyectivos de
tipo finito. El teorema se deduce entonces de
\sref{III.7.9.3}.
\end{proof}

\begin{env}[7.9.5]
\label{III.7.9.5}
Bajo las hipótesis de \sref{III.7.9.4}, la función
\eqref{III.7.9.4.1} es \emph{constante} cuando $Y$ es
\emph{conexo}. Cuando $Y$ es conexo y no vacío, su único valor
(entero) se denota por
\[
  \operatorname{EP}(f,\mathscr P_\bullet),\qquad
  \operatorname{EP}(Y,\mathscr P_\bullet),
  \qquad\text{o simplemente}\qquad
  \operatorname{EP}(\mathscr P_\bullet)
\]
cuando no haya riesgo de confusión, y se dice que este entero es la
\emph{característica de Euler-Poincaré de $\mathscr P_\bullet$}
relativa a $f$ (o a $Y$). En general, se escribe asimismo
\[
  \operatorname{EP}(f,\mathscr P_\bullet;y),\qquad
  \operatorname{EP}(Y,\mathscr P_\bullet;y),\qquad
  \operatorname{EP}(\mathscr P_\bullet;y)
\]
para el segundo miembro de \eqref{III.7.9.4.1}.
\end{env}

\begin{env}[7.9.6]
\label{III.7.9.6}
Bajo las hipótesis de \sref{III.7.9.4} relativas a
$X$, $Y$ y $f$, sea
\[
  0\longrightarrow\mathscr P'_\bullet
    \xrightarrow{u}\mathscr P_\bullet
    \xrightarrow{v}\mathscr P''_\bullet
    \longrightarrow 0
\]
una sucesión exacta de complejos finitos de $\mathscr O_X$-módulos
coherentes y planos sobre $Y$, donde los homomorfismos $u$ y $v$
son de grados pares $2d$ y $2d'$, respectivamente.
Como $\mathscr T_\bullet$ es un funtor homológico
\sref{III.7.7.1}, se tiene una sucesión exacta de homología
\[
  \cdots\longrightarrow
  \mathscr T_i(\mathscr P'_\bullet,\kres(y))
  \longrightarrow
  \mathscr T_{i+2d}(\mathscr P_\bullet,\kres(y))
  \longrightarrow
  \mathscr T_{i+2d+2d'}(\mathscr P''_\bullet,\kres(y))
  \longrightarrow
  \mathscr T_{i-1}(\mathscr P'_\bullet,\kres(y))
  \longrightarrow\cdots ,
\]
que, además, solo tiene un número finito de términos. Al escribir
que la característica de Euler-Poincaré de este complejo es nula
\sref[0]{0.11.10.1}, se obtiene inmediatamente
\[
  \operatorname{EP}(\mathscr P_\bullet;y)
  =
  \operatorname{EP}(\mathscr P'_\bullet;y)
  +
  \operatorname{EP}(\mathscr P''_\bullet;y)
  \tag{7.9.6.1}\label{III.7.9.6.1}
\]
para todo $y\in Y$.

Supongamos, por ejemplo, que
$\mathscr P_\bullet=(\mathscr P_i)$ y que $\mathscr P_i=0$ para
$i<0$. Se tiene entonces la sucesión exacta de complejos
\[
\xymatrix@C=2.4em{
  \cdots \ar[r] & 0 \ar[r]\ar[d] & 0 \ar[r]\ar[d]
    & 0 \ar[r]\ar[d] & 0 \ar[r] & \cdots \\
  \cdots \ar[r] & 0 \ar[r]\ar[d] & 0 \ar[r]\ar[d]
    & \mathscr P_1 \ar[r]\ar[d] & \mathscr P_2 \ar[r] & \cdots \\
  \cdots \ar[r] & 0 \ar[r]\ar[d] & \mathscr P_0 \ar[r]\ar[d]
    & \mathscr P_1 \ar[r]\ar[d] & \mathscr P_2 \ar[r] & \cdots \\
  \cdots \ar[r] & 0 \ar[r]\ar[d] & \mathscr P_0 \ar[r]\ar[d]
    & 0 \ar[r]\ar[d] & 0 \ar[r] & \cdots \\
  \cdots \ar[r] & 0 \ar[r] & 0 \ar[r]
    & 0 \ar[r] & 0 \ar[r] & \cdots
}
\]
cuyas flechas verticales no nulas son los automorfismos identidad.
Aplicando \eqref{III.7.9.6.1} a esta sucesión exacta y razonando por
inducción sobre la longitud de $\mathscr P_\bullet$, se obtiene
\[
  \operatorname{EP}(\mathscr P_\bullet;y)
  =
  \sum_i(-1)^i\operatorname{EP}(\mathscr P_i;y).
  \tag{7.9.6.2}\label{III.7.9.6.2}
\]

\oldpage[III]{210}
Aquí, para todo $\mathscr O_X$-módulo coherente $\mathscr F$ plano
sobre $Y$, la notación $\operatorname{EP}(\mathscr F;y)$ (o
$\operatorname{EP}(f,\mathscr F;y)$, o
$\operatorname{EP}(Y,\mathscr F;y)$) significa
$\operatorname{EP}(\mathscr Q_\bullet;y)$ (con la notación relativa
correspondiente) para el complejo $\mathscr Q_\bullet$ cuyo único
término no nulo es $\mathscr F$ en grado $0$. Por tanto, el estudio
de las características de Euler-Poincaré de los complejos se reduce
al de los complejos con un solo término.
\end{env}

\begin{proposition}[7.9.7]
\label{III.7.9.7}
Bajo las hipótesis de \sref{III.7.9.4}, sean $Y'$ un preesquema
localmente noetheriano y $g:Y'\to Y$ un morfismo, pongamos
\[
  X'=X\times_Y Y',\qquad
  f'=f_{(Y')}:X'\longrightarrow Y',
\]
y sea
\[
  \mathscr P'_\bullet
  =
  \mathscr P_\bullet\otimes_{\mathscr O_Y}\mathscr O_{Y'}
\]
el complejo finito de $\mathscr O_{X'}$-módulos así obtenido.
Entonces $\mathscr P'_\bullet$ está formado por
$\mathscr O_{X'}$-módulos coherentes y planos sobre $Y'$, y, para
todo $y'\in Y'$, se tiene
\[
  \operatorname{EP}(\mathscr P'_\bullet;y')
  =
  \operatorname{EP}(\mathscr P_\bullet;g(y')).
  \tag{7.9.7.1}\label{III.7.9.7.1}
\]
\end{proposition}

\begin{proof}
Los $\mathscr O_{X'}$-módulos $\mathscr P'_i$, al ser imágenes
inversas de los $\mathscr P_i$ por la proyección $X'\to X$, son
coherentes. Son planos sobre $Y'$ en virtud de
\sref[0]{0.6.2.1} y \sref[I]{I.4.14.5}, pues la cuestión es local
en $X$, $Y$ e $Y'$. Por último, $f'$ es propio por
\sref[II]{II.5.4.2}, de modo que el primer miembro de
\eqref{III.7.9.7.1} está definido. La fórmula
\eqref{III.7.9.7.1} se deduce entonces de
\textup{(6.10.4.2)}, \sref{III.7.7.2} y
\sref{III.7.6.7}, reduciéndose, como siempre se puede hacer, al
caso en que $Y$ e $Y'$ son afines.
\end{proof}

\begin{proposition}[7.9.8]
\label{III.7.9.8}
Supongamos que se cumplen las hipótesis de \sref{III.7.9.4} y,
además, que existe un entero $i_0$ tal que
\[
  \mathscr T_i(\kres(y))=0
  \qquad\text{para $i\neq i_0$ y todo $y\in Y$.}
\]
Entonces
\[
  \mathscr T_{i_0}(\mathscr O_Y)
  =
  \mathscr H^{-i_0}(f,\mathscr P^\bullet)
\]
es un $\mathscr O_Y$-módulo localmente libre, cuyo rango en
$y\in Y$ es
\[
  (-1)^{i_0}\operatorname{EP}(f,\mathscr P_\bullet;y).
\]
\end{proposition}

\begin{proof}
Observemos primero que los funtores
$\mathscr T_i^{\mathscr O_y}$ cumplen las hipótesis de
\sref{III.7.4.4}, por lo que \sref{III.7.4.7} les es aplicable.
La hipótesis y \sref{III.7.5.3} implican que
$\mathscr T_i^{\mathscr O_y}$ es nulo para $i\neq i_0$.
En consecuencia, $\mathscr T_{i_0}$ es exacto por
\sref{III.7.3.3}, y \sref{III.7.8.4} muestra que
$\mathscr H^{-i_0}(f,\mathscr P^\bullet)$ es localmente libre.
Su rango en $y\in Y$ es
\[
  \operatorname{rank}_{\kres(y)}
    \mathscr T_{i_0}(\kres(y))
  =
  (-1)^{i_0}\operatorname{EP}(f,\mathscr P_\bullet;y),
\]
por definición, ya que $\mathscr T_i(\kres(y))=0$ para
$i\neq i_0$.
\end{proof}

\begin{corollary}[7.9.9]
\label{III.7.9.9}
Sean $Y$ un preesquema localmente noetheriano,
$f:X\to Y$ un morfismo propio y $\mathscr F$ un
$\mathscr O_X$-módulo coherente y plano sobre $Y$. Supongamos que
existe un entero $i_0$ tal que
\[
  H^i\bigl(f^{-1}(y),
    \mathscr F\otimes_{\mathscr O_Y}\kres(y)\bigr)=0
\]
para todo $i\neq i_0$ y todo $y\in Y$. Entonces
$R^{i_0}f_*(\mathscr F)$ es un $\mathscr O_Y$-módulo localmente
libre, cuyo rango en $y$ es
$(-1)^{i_0}\operatorname{EP}(f,\mathscr F;y)$.
\end{corollary}

En particular:

\begin{corollary}[7.9.10]
\label{III.7.9.10}
Bajo las hipótesis preliminares de \sref{III.7.9.9} relativas a
$X$, $Y$ y $\mathscr F$, supongamos que
$R^if_*(\mathscr F)=0$ para todo $i>0$. Entonces
$f_*(\mathscr F)$ es un $\mathscr O_Y$-módulo localmente libre,
cuyo rango en $y$ es $\operatorname{EP}(f,\mathscr F;y)$.
\end{corollary}

\begin{proof}
En virtud de \sref{III.7.9.9}, basta demostrar el lema siguiente.
\end{proof}

\begin{lemma}[7.9.10.1]
\label{III.7.9.10.1}
Bajo las hipótesis de \sref{III.7.9.10}, se tiene
\[
  H^i\bigl(f^{-1}(y),
    \mathscr F\otimes_{\mathscr O_Y}\kres(y)\bigr)=0
\]
para todo $i>0$ y todo $y\in Y$.
\end{lemma}

\begin{proof}
Podemos limitarnos al caso en que $Y=\Spec(A)$ es afín.
Con las notaciones de \sref{III.7.9.4}, y estando
$\mathscr P_\bullet$ reducido a su único término no nulo
$\mathscr F$ en grado $0$, la hipótesis da
$\mathscr T_p(\mathscr O_Y)=0$ para $p<0$. Se deduce de
\sref{III.7.3.7} que $\mathscr T_p$ es exacto para $p<0$, y el
lema resulta entonces de la equivalencia de las condiciones a) y
d) de \sref{III.7.7.5}.
\end{proof}

\oldpage[III]{211}
\begin{proposition}[7.9.11]
\label{III.7.9.11}
Bajo las hipótesis de \sref{III.7.9.4}, sea $\mathscr L$ un
$\mathscr O_X$-módulo inversible muy amplio sobre $Y$ y, para todo
$n\in\bb Z$, pongamos
\[
  \mathscr P_\bullet(n)
  =
  \mathscr P_\bullet
    \otimes_{\mathscr O_X}\mathscr L^{\otimes n}.
\]
Entonces, para todo $y\in Y$, la función
\[
  n\longmapsto
  \operatorname{EP}(f,\mathscr P_\bullet(n);y)
  \tag{7.9.11.1}\label{III.7.9.11.1}
\]
es un polinomio con coeficientes racionales, y es el mismo para
todos los puntos de una misma componente conexa de $Y$.
\end{proposition}

\begin{proof}
Es claro que $\mathscr P_\bullet(n)$ es un complejo de
$\mathscr O_X$-módulos planos sobre $Y$. En virtud de
\eqref{III.7.9.6.2}, podemos reducirnos al caso en que
$\mathscr P_\bullet$ tiene un único término no nulo
$\mathscr F\neq0$ en grado $0$. Como se trata de cuestiones
locales sobre $Y$, podemos suponer además que $Y$ es afín y que
$f$ es proyectivo \sref[II]{II.5.5.3}. Pongamos
\[
  X_y=f^{-1}(y),\qquad
  \mathscr L_y=\mathscr L\otimes_{\mathscr O_Y}\kres(y);
\]
entonces $\mathscr L_y$ es un $\mathscr O_{X_y}$-módulo muy amplio
\sref[II]{II.4.4.10}. Por \sref{III.7.7.2}, para el funtor
$\mathscr T_\bullet$ asociado a $\mathscr P_\bullet(n)$ se tiene
\[
  \mathscr T_i(\kres(y))
  =
  H^{-i}\bigl(
    X_y,\mathscr F_y\otimes\mathscr L_y^{\otimes n}\bigr),
  \qquad
  \mathscr F_y=\mathscr F\otimes_{\mathscr O_Y}\kres(y).
\]
Por consiguiente, $\operatorname{EP}(f,\mathscr F(n);y)$ no es
sino la característica de Euler-Poincaré
$\chi_{\kres(y)}(\mathscr F_y(n))$ definida en
\sref{III.2.5.1}. De \sref{III.2.5.3} se deduce que
\eqref{III.7.9.11.1} es un polinomio. Además, para cada $n$ fijo,
su valor es constante en cada componente conexa de $Y$ por
\sref{III.7.9.4}; esto termina la demostración.
\end{proof}

\begin{env}[7.9.12]
\label{III.7.9.12}
El polinomio \eqref{III.7.9.11.1}, con coeficientes racionales, se
denotará por
\[
  \operatorname{PH}(f,\mathscr P_\bullet;y)
  \qquad\text{o}\qquad
  \operatorname{PH}(\mathscr P_\bullet;y),
\]
y se llamará \emph{polinomio de Hilbert en $y$} relativo a
$\mathscr P_\bullet$, $f$ y $\mathscr L$ (o simplemente
\emph{polinomio de Hilbert en $y$} de $\mathscr P_\bullet$, o de
$f$, cuando no haya riesgo de confusión). Cuando $Y$ es conexo y
no vacío, se suprime la referencia a $y$ en la notación y la
terminología. Este invariante desempeñará un papel esencial en el
capítulo~V, en la teoría de los «móduli» de los haces cociente
coherentes de un haz coherente dado.

Con las notaciones de \sref{III.7.9.6} y
\sref{III.7.9.11}, se tiene
\[
  \operatorname{PH}(\mathscr P_\bullet;y)
  =
  \operatorname{PH}(\mathscr P'_\bullet;y)
  +
  \operatorname{PH}(\mathscr P''_\bullet;y),
  \tag{7.9.12.1}\label{III.7.9.12.1}
\]
y, en particular,
\[
  \operatorname{PH}(\mathscr P_\bullet;y)
  =
  \sum_i(-1)^i\operatorname{PH}(\mathscr P_i;y).
  \tag{7.9.12.2}\label{III.7.9.12.2}
\]
Estas fórmulas se deducen inmediatamente de
\eqref{III.7.9.6.1} y \eqref{III.7.9.6.2}. Asimismo, con las
notaciones e hipótesis de \sref{III.7.9.7}, se tiene
\[
  \operatorname{PH}(\mathscr P'_\bullet;y')
  =
  \operatorname{PH}(\mathscr P_\bullet;g(y')).
  \tag{7.9.12.3}\label{III.7.9.12.3}
\]
\end{env}

La fórmula \eqref{III.7.9.12.2} reduce el estudio de los
polinomios de Hilbert de un complejo al de los polinomios de
Hilbert de un solo $\mathscr O_X$-módulo plano sobre $Y$.
Estos últimos admiten la siguiente interpretación notable,
independiente de consideraciones homológicas.

\begin{corollary}[7.9.13]
\label{III.7.9.13}
Sean $Y$ un preesquema noetheriano, $f:X\to Y$ un morfismo propio,
$\mathscr L$ un $\mathscr O_X$-módulo inversible muy amplio sobre
$Y$ y $\mathscr F$ un $\mathscr O_X$-módulo coherente y plano
sobre $Y$. Existe un entero $n_0$ tal que, para $n\geq n_0$,
$f_*(\mathscr F(n))$ es un $\mathscr O_Y$-módulo localmente libre
cuyo rango en $y\in Y$ es
\[
  \operatorname{PH}(f,\mathscr F;y)(n).
\]
\end{corollary}

\begin{proof}
Como $f$ es proyectivo \sref[II]{II.5.5.3}, existe un entero $n_0$
tal que, para $n\geq n_0$,

\oldpage[III]{212}
\[
  R^if_*(\mathscr F(n))=0
  \qquad\text{para todo $i>0$}
\]
por \sref{III.2.2.1}. La conclusión se deduce entonces de
\sref{III.7.9.10}.
\end{proof}

El criterio de planitud siguiente será importante en el
capítulo~V, en la teoría de los «móduli» de los haces coherentes.

\begin{proposition}[7.9.14]
\label{III.7.9.14}
Sean $Y$ un preesquema noetheriano, $f:X\to Y$ un morfismo
proyectivo y $\mathscr L$ un $\mathscr O_X$-módulo inversible
amplio para $f$. Para todo $\mathscr O_X$-módulo $\mathscr F$ y
todo $n\in\bb Z$, pongamos
\[
  \mathscr F(n)=\mathscr F\otimes\mathscr L^{\otimes n}.
\]
Para que un $\mathscr O_X$-módulo coherente $\mathscr F$ sea plano
sobre $Y$, es necesario y suficiente que exista un entero $n_0$
tal que, para todo $n\geq n_0$, $f_*(\mathscr F(n))$ sea un
$\mathscr O_Y$-módulo localmente libre.
\end{proposition}

\begin{proof}
La necesidad de la condición se demuestra como en
\sref{III.7.9.13}; el resultado de \sref{III.2.2.1} se aplica a un
haz amplio porque $f$ es proyectivo.

Recíprocamente, podemos reducirnos al caso en que
$Y=\Spec(A)$ es afín. En virtud de la hipótesis y de
\sref{III.2.2.2}{(i)}, los $A$-módulos
$\Gamma(X,\mathscr F(n))$ son proyectivos y de tipo finito
(Bourbaki, \emph{Alg.~comm.}, cap.~II, \S5,
n.\textsuperscript{o}~2, teorema~1). Sea
\[
  S=\bigoplus_{n\geq0}\Gamma(X,\mathscr L^{\otimes n})
\]
el anillo graduado asociado. Se sabe que $X$ se identifica
canónicamente con $\Proj(S)$ por
\sref[II]{II.4.5.2}{(b)} y \sref[II]{II.5.4.4}. Sea
\[
  M=\bigoplus_{n\geq n_0}\Gamma(X,\mathscr F(n)).
\]
Sustituyendo, si es preciso, $\mathscr L$ por una potencia
$\mathscr L^{\otimes d}$, podemos suponer que $S$ está generado por
un número finito de elementos de grado $1$
\sref{III.2.3.5.1}; se deduce entonces de
\sref[II]{II.2.7.5} y \sref[II]{II.2.7.2} que $\mathscr F$ se
identifica con $\mathscr P\!roj_0(M)$.

Para todo elemento homogéneo $g\in S$ de grado positivo, se tiene,
por tanto,
\[
  \Gamma(X_g,\mathscr F)=M_{(g)}.
\]
Ahora bien, $M$, al ser una suma directa de $A$-módulos
proyectivos, es un $A$-módulo plano. Por consiguiente, $M_g$ es
plano por \sref[0]{0.6.3.2}, y también lo es $M_{(g)}$, que es un
factor directo de $M_g$ por \sref[0]{0.6.1.2}. Se concluye de
\sref[I]{I.4.14.5} que $\mathscr F$ es plano sobre $Y$ en todo
punto de $X_g$. Como los abiertos $X_g$ recubren $X$, queda
demostrada la proposición.
\end{proof}

\bigskip
\noindent\emph{(Continuará.)}

\clearpage
\oldpage[III]{213}

\section*{ÍNDICE DE NOTACIONES}

\begin{flushleft}
$\mathbf{Tor}_n^A(P_\bullet,Q_\bullet)$: 6.3.1.

\medskip
$\mathscr{T}\!or_n^{\mathscr O_S}(\mathscr P_\bullet,\mathscr Q_\bullet)$,
$\mathscr{T}\!or_n^S(\mathscr P_\bullet,\mathscr Q_\bullet)$: 6.4.1 y 6.5.1.

\medskip
$\mathscr{T}\!or_n^S(\mathscr F,\mathscr G)$: 6.5.1.

\medskip
$\mathscr{T}\!or_q^S(\mathscr P_\bullet^{(1)},\mathscr P_\bullet^{(2)},\ldots,
\mathscr P_\bullet^{(m)})$: 6.5.15.

\medskip
$\mathbf H^{-n}(\mathfrak U^{(i)},\mathscr P_\bullet^{(i)})$,
$\mathbf H^{-n}(X^{(i)},\mathscr P_\bullet^{(i)})$: 6.6.1.

\medskip
$\mathbf{Tor}_n^A(L_{\bullet\bullet}^{(1)},L_{\bullet\bullet}^{(2)})$,
$\mathbf{Tor}_n^S(\mathfrak U^{(1)},\mathfrak U^{(2)};
\mathscr P_\bullet^{(1)},\mathscr P_\bullet^{(2)})$: 6.6.2.

\medskip
$\mathscr{T}\!or_n^S(\mathfrak U^{(1)},\mathfrak U^{(2)};
\mathscr F^{(1)},\mathscr F^{(2)})$,
$\mathscr{T}\!or_n^S(\mathfrak U^{(1)},\mathfrak U^{(2)};
\mathscr P_\bullet^{(1)},\mathscr P_\bullet^{(2)})$: 6.6.2.

\medskip
$\mathscr{T}\!or_n^S(f_1,f_2;\mathscr P_\bullet^{(1)},
\mathscr P_\bullet^{(2)})$: 6.6.8, 6.7.1 y 6.7.3.

\medskip
$\mathscr{T}\!or_n^S((f_i)_{i\in I};(\mathscr P_\bullet^{(i)})_{i\in I})$,
$\mathscr{T}\!or_n^S(f_1,\ldots,f_m;\mathscr P_\bullet^{(1)},\ldots,
\mathscr P_\bullet^{(m)})$: 6.7.11.

\medskip
$\mathbf{Ab}$, $\mathbf{Ab}_A$: 7.1.1.

\medskip
$T_{(B)}$, $T^{(B)}$, $T\otimes_A B$
($T$ funtor de $\mathbf{Ab}_A$ en $\mathbf{Ab}$): 7.1.3.

\medskip
$t_M$: 7.2.2.

\medskip
$T_{\mathfrak p}$: 7.1.4.

\medskip
$\mathscr T_\bullet^{Y'}(\mathscr M')$,
$\mathscr T_\bullet^U(\mathscr M|U)$: 7.7.2.

\medskip
$\operatorname{EP}(f,\mathscr P_\bullet;y)$,
$\operatorname{EP}(Y,\mathscr P_\bullet;y)$,
$\operatorname{EP}(\mathscr P_\bullet;y)$: 7.9.5.

\medskip
$\operatorname{PH}(f,\mathscr P_\bullet;y)$,
$\operatorname{PH}(\mathscr P_\bullet;y)$: 7.9.12.
\end{flushleft}

\clearpage
\oldpage[III]{214}
\section*{ÍNDICE TERMINOLÓGICO}

\begin{flushleft}
Característica de Euler-Poincaré de un complejo de módulos: 7.9.5.

\medskip
Cohomológicamente plano en un punto $y\in Y$ en dimensión $p$,
cohomológicamente plano sobre $Y$ en dimensión $p$, cohomológicamente plano
sobre $Y$ en el punto $y$, cohomológicamente plano sobre $Y$ (complejo de
$\mathscr O_X$-módulos planos sobre $Y$): 7.8.1.

\medskip
Complejos homotópicamente equivalentes: 6.1.4.

\medskip
Espacio anillado de dimensión cohomológica $\leq n$: 6.5.5.

\medskip
Extensión de escalares en un funtor covariante aditivo: 7.1.3.

\medskip
Haz de anillos de dimensión cohomológica $\leq n$: 6.5.5.

\medskip
Fórmula de Künneth: 6.7.8.

\medskip
Homológicamente plano en un punto $y\in Y$ en dimensión $p$,
homológicamente plano sobre $Y$ en dimensión $p$, homológicamente plano sobre
$Y$ en el punto $y$, homológicamente plano sobre $Y$ (complejo de
$\mathscr O_X$-módulos planos sobre $Y$): 7.8.1.

\medskip
Equivalencia homotópica: 6.1.4.

\medskip
Hipertor de dos complejos de $A$-módulos: 6.3.1.

\medskip
Hipertor local de dos complejos de módulos: 6.4.1.

\medskip
Hipertor global de dos complejos de módulos: 6.6.2, 6.6.8 y 6.7.3.

\medskip
Polinomio de Hilbert relativo a un complejo de módulos: 7.9.12.

\medskip
Sucesiones espectrales de Künneth: 6.7.3.
\end{flushleft}

\clearpage
\oldpage[III]{215}
\section*{TABLA DE MATERIAS}

\begin{flushleft}
\textsc{Capítulo III.} --- \textbf{Estudio cohomológico de los haces coherentes
(continuación)} \dotfill 5

\medskip
\S~6. Funtores «Tor» locales y globales; fórmula de Künneth \dotfill 5

6.1. Introducción \dotfill 5

6.2. Hipercohomología de los complejos de módulos sobre un preesquema \dotfill 6

6.3. Hipertor de dos complejos de módulos \dotfill 9

6.4. Funtores hipertor locales de complejos de módulos cuasicoherentes: caso
de los esquemas afines \dotfill 14

6.5. Funtores hipertor locales de complejos de módulos cuasicoherentes: caso
general \dotfill 16

6.6. Funtores de hipertor global para complejos de módulos cuasicoherentes y
sucesiones espectrales de Künneth: caso de base afín \dotfill 21

6.7. Funtores de hipertor global de complejos de módulos cuasicoherentes y
sucesiones espectrales de Künneth: caso general \dotfill 25

6.8. Sucesiones espectrales de asociatividad de los hipertor globales \dotfill 32

6.9. Las sucesiones espectrales de cambio de base en los hipertor globales
\dotfill 34

6.10. Estructura local de ciertos funtores cohomológicos \dotfill 39

\medskip
\S~7. Estudio del cambio de base en los funtores homológicos covariantes de
módulos \dotfill 43

7.1. Funtores de $A$-módulos \dotfill 43

7.2. Caracterización del funtor producto tensorial \dotfill 44

7.3. Criterios de exactitud de los funtores homológicos de módulos \dotfill 48

7.4. Criterios de exactitud para los funtores
$\mathrm H_i(P\otimes_A M)$ \dotfill 53

7.5. Caso de los anillos locales noetherianos \dotfill 58

7.6. Descenso de las propiedades de exactitud. Teorema de semicontinuidad y
criterio de exactitud de Grauert \dotfill 60

7.7. Aplicación a los morfismos propios: I. La propiedad de cambio de base
\dotfill 65

7.8. Aplicación a los morfismos propios: II. Criterios de planitud
cohomológica \dotfill 72

7.9. Aplicación a los morfismos propios: III. Invariancia de la
característica de Euler-Poincaré y del polinomio de Hilbert \dotfill 76

\begin{flushright}
\emph{Recibido el 15 de diciembre de 1962.}
\end{flushright}
\end{flushleft}

% La página impresa 216 está en blanco.
\clearpage
\thispagestyle{empty}
\null
\clearpage

\oldpage[III]{217}
\section*{ERRATAS Y ADENDAS}
\begin{center}
(Lista 2)

\medskip
A) \emph{Errores tipográficos}
\end{center}

\noindent\textbf{$(0_{\mathrm I},\,3.2.1)$} En la línea 2 de la p. 27, sustituir
$X$ por $U$.

\noindent\textbf{$(0_{\mathrm I},\,3.2.6)$} En la primera línea contando desde
abajo de la p. 27 y en la línea 1 de la p. 28, sustituir (3 veces)
$\mathfrak S_\lambda$ por $\mathscr H_\lambda$.

\noindent\textbf{$(0_{\mathrm I},\,3.4.5)$} Añadir un paréntesis delante de
3.4.5).

\noindent\textbf{$(0_{\mathrm I},\,4.2.1)$} En la línea 10 de la p. 39,
sustituir $\psi_*(A)$ por $\psi_*(\mathscr A)$.

\noindent\textbf{$(0_{\mathrm I},\,5.1.3)$} En la línea 16 de la p. 45,
sustituir $\mathscr F|V$ por $\mathscr F|U$.

\noindent\textbf{$(0_{\mathrm I},\,5.3.9)$} En la línea 12 de la p. 48, antes
de «tal que», añadir: sobre un entorno abierto $U$ de $x$.

\noindent\textbf{$(0_{\mathrm I},\,7.6.9)$} En la tercera línea contando desde
abajo de la p. 73, sustituir $I_\lambda$ por $J_\lambda$.

\noindent\textbf{(I, 1.1.15)} En la línea 3 de la p. 83, sustituir
$\mathfrak a$ por $\mathfrak a\ne A$.

\noindent\textbf{(I, 1.4.1)} En la segunda línea contando desde abajo de la
p. 90, sustituir $\widetilde N$ por $N$.

\noindent\textbf{(I, 2.3.1)} En la tercera línea contando desde abajo de la
p. 101, sustituir (0, 4.1.6) por (0, 4.1.7).

\noindent\textbf{(I, 2.3.2)} En la línea 11 de la p. 101, sustituir $B$ por
$B_s$ y $C$ por $C_t$.

\noindent\textbf{(I, 2.5.5)} En la línea 19 de la p. 104, sustituir
$S$-morfismo por $S$-preesquema.

\noindent\textbf{(I, 3.3.7)} En la línea 17 de la p. 109, sustituir
$Y_{(S)}$ por $Y_{(S')}$.

\noindent\textbf{(I, 3.4.5)} En la línea 9 de la p. 113, sustituir
$\mathrm s$ por $s$.

\noindent\textbf{(I, 3.4.8)} En la línea 4 de la p. 114, sustituir $p(x)$ por
$f(x)$.

\noindent\textbf{(I, 3.5.1)} En la línea 3 de la p. 115, sustituir
$1_X\times g$ por $1_{X'}\times g$.

\noindent\textbf{(I, 3.7.2)} En la línea 8 de la p. 119, sustituir el primer
$X$ por $X'$.

\noindent\textbf{(I, 4.2.2)} En la línea 8 de la p. 123, en la flecha vertical
izquierda del diagrama, sustituir $\alpha_{\psi(y)}$ por $\varphi_{\psi(y)}$.

\noindent\textbf{(I, 4.4.3)} En la línea 16 de la p. 126, sustituir «isomorfa»
por «isomorfas».

\noindent\textbf{(I, 4.5.5)} En la línea 17 contando desde abajo de la p. 127,
sustituir (4.2.4) por (4.2.5).

\noindent\textbf{(I, 5.1.4)} En la línea 14 contando desde abajo de la p. 128,
sustituir (2.1.7) por (2.1.8).

\noindent\textbf{(I, 5.3.13)} En la línea 20 de la p. 134, sustituir (4.2.4)
por (4.2.5).

\noindent\textbf{(I, 6.4.2)} En la primera línea contando desde abajo de la
p. 147, sustituir recubrimiento por recubrimiento finito.

\noindent\textbf{(I, 9.1.13)} En la línea 15 de la p. 171, sustituir
$p^{-1}(\mathscr F)$ por $p^*(\mathscr F)$.

\noindent\textbf{(I, 9.5.11)} En la línea 7 de la p. 179, sustituir $Y$ por
$Y'$.

\noindent\textbf{(I, 9.6.5)} En la línea 17 contando desde abajo de la p. 180,
sustituir $(0,\,4.1.4)$ por $(0,\,4.1.3)$.

\noindent\textbf{(I, 10.12.2)} En la línea 14 contando desde abajo de la p.
206, sustituir $(X_n)_{(S_n)}$ por $(X_n)_{(S_m)}$.

\clearpage
\oldpage[III]{218}
\noindent\textbf{(I), Índice terminológico:} En la línea 6 contando desde
abajo de la p. 219, sustituir 0, 4.1.4 por 0, 4.1.5. En las líneas 16, 17 y 18
contando desde abajo de la p. 222, sustituir I, 2.1.7 por I, 2.1.8; en la línea
19 contando desde abajo de la p. 222, sustituir I, 2.1.8 por I, 2.1.7.

\noindent\textbf{(II, 1.5.2)} En la línea 17 de la p. 12, insertar una $f$
junto a la flecha vertical izquierda del diagrama.

\noindent\textbf{(II, 4.2.7)} En la línea 16 contando desde abajo de la p. 75,
sustituir (III, 2.1.14) por (III, 2.1.13).

\noindent\textbf{(II, 5.2.1)} En la línea 16 contando desde abajo de la p. 97,
sustituir $X-U$ por $U$. En la línea 8 contando desde abajo de la p. 97,
sustituir $X'$ por $X_f$.

\noindent\textbf{(II, 5.3.6)} En la línea 16 de la p. 100, sustituir (4.6.18)
por (4.6.17).

\noindent\textbf{(II, 6.2.7)} En la línea 9 contando desde abajo de la p. 116,
sustituir cap. V por cap. IV.

\noindent\textbf{(II, 6.6.5)} En la línea 15 contando desde abajo de la p.
132, sustituir cap. V por cap. IV.

\noindent\textbf{(II, 7.4.12)} En la línea 1 de la p. 151, sustituir cap. V
por cap. III.

\noindent\textbf{(II, 8.1.4)} En la línea 15 de la p. 154, sustituir (III,
2.3.8) por (III, 2.3.7).

\noindent\textbf{(II, 8.10.5)} En la línea 17 contando desde abajo de la p.
187, sustituir $g_{(j)x}$ por $g_{j(x)}$.

\noindent\textbf{$(0_{\mathrm{III}},\,10.2.7)$} En la línea 1 de la p. 20,
sustituir $u$-ádico por $\mathfrak m$-ádico. En la línea 13 contando desde abajo
de la p. 20, sustituir $k$ por $k_i$.

\noindent\textbf{$(0_{\mathrm{III}},\,11.1.1)$} En la línea 18 de la p. 23,
sustituir $\operatorname{Im}(d_r^{p+r,q-r+1})$ por
$\operatorname{Im}(d_r^{p-r,q+r-1})$.

\noindent\textbf{$(0_{\mathrm{III}},\,11.8.6)$} En la línea 14 contando desde
abajo de la p. 45, sustituir las flechas $\to$ por $\rightsquigarrow$ (4 veces).

\noindent\textbf{$(0_{\mathrm{III}},\,12.3.4)$} En la línea 4 de la p. 61,
sustituir $\rightsquigarrow$ por $\to$.

\noindent\textbf{$(0_{\mathrm{III}},\,13.2.4)$} En la segunda línea contando
desde abajo de la p. 67, sustituir $\varinjlim$ por $\varprojlim$ (2 veces).

\noindent\textbf{(III, 1.1.7)} En la línea 15 de la p. 85, sustituir $A^r$ por
$\bigwedge(A^r)$ (2 veces).

\noindent\textbf{(III, 1.2.4)} En la línea 7 de la p. 87, sustituir
$\Gamma(U,\mathscr F)$ por $\Gamma(U,\mathscr O_X)$ y sustituir $\mathscr G$
por $\mathscr F$. En la línea 8 de la p. 87, sustituir
$\mathrm H^p(\mathfrak U,\mathscr G)$ por $\mathrm H^p(\mathfrak U,\mathscr F)$.

\noindent\textbf{(III, 2.5.3.1)} En la línea 3 de la p. 111, sustituir $n>0$
por $n\geq0$. En la línea 8 de la p. 111, sustituir $-r\leq n\leq0$ por
$-r\leq n<0$.

\noindent\textbf{(III, 3.1.2)} En la línea 13 de la p. 116, sustituir
$\mathscr G$ por $\mathscr G\in\mathbf K'$.

\noindent\textbf{(III, 4.3.6)} En la línea 7 de la p. 133, sustituir
$K'=K\otimes_A\widehat A$ por $K'\supset K\otimes_A\widehat A$. En la línea 16
de la p. 133, sustituir $R(X)$ por $R(Y)$.

\noindent\textbf{(III, 4.4.12)} En la línea 10 de la p. 138, sustituir cap. V
por cap. IV.

\noindent\textbf{(III, 4.6.6)} En la línea 19 de la p. 142, sustituir «cap. V»
por «en un parágrafo posterior».

\begin{center}
B) \emph{Modificaciones del texto}\footnote{Para facilitar la búsqueda de las
referencias, las modificaciones pertenecientes a la lista de erratas insertada
en el capítulo N serán designadas en lo sucesivo por $\mathbf{Err}_{\mathrm N}$
seguido de un número.}
\end{center}

\noindent\textbf{$(\mathbf{Err}_{\mathrm{III}},\,1)$} En
$(0_{\mathrm I},\,5.1.3)$, en la línea 18 de la p. 45, sustituir «suma directa»
por «suma directa finita».

\noindent\textbf{$(\mathbf{Err}_{\mathrm{III}},\,2)$} En
$(0_{\mathrm I},\,5.4.3)$, en las líneas 8 a 4 contando desde abajo de la p. 49,
sustituir desde «En el caso general\ldots{}» por el texto siguiente:

En el caso general en que $X$ es un espacio anillado tal que $\mathscr O_x$ es
un anillo \emph{local} para
\oldpage[III]{219}
todo $x\in X$, si $\mathscr L$ es un $\mathscr O_X$-módulo \emph{de tipo finito}
tal que existe un $\mathscr O_X$-módulo $\mathscr F$ para el cual
$\mathscr L\otimes_{\mathscr O_X}\mathscr F$ es \emph{isomorfo a}
$\mathscr O_X$, el razonamiento anterior muestra que, para todo $x\in X$,
$\mathscr L_x$ es isomorfo a $\mathscr O_x$. De ello se deduce que $\mathscr L$
es entonces \emph{invertible}. En efecto, para todo $x\in X$, sea $U$ un
entorno abierto de $x$ tal que $\mathscr L|U$ esté engendrado por $n$ secciones
$s_i\ (1\leq i\leq n)$ sobre $U$ (5.2.1). Se puede suponer, por ejemplo, que
$s=s_1$ es tal que $s_x\ne0$ y, como $\mathscr L_x$ es isomorfo a
$\mathscr O_x$, existe para cada $i$ una sección $t_i$ de $\mathscr O_X$ sobre
un entorno abierto $V_i\subset U$ de $x$ tal que
$(t_i)_x s_x=(s_i)_x$. Por consiguiente, existe un entorno abierto
$V\subset\bigcap_{i=1}^nV_i$ de $x$ tal que $t_i s=s_i$ en $V$; dicho de otro
modo, $\mathscr L|V$ está engendrado por \emph{una sola} sección $s|V$. Además,
si $z$ es una sección, sobre un abierto $W\subset V$, del núcleo del
homomorfismo $\mathscr O_X|V\to\mathscr L|V$ definido por $s$ (5.1.1), $z_y$
anula $\mathscr L_y$ para todo $y\in W$, luego $z_y=0$ por hipótesis y, por
consiguiente, $z=0$, lo que termina de probar nuestra afirmación. Además, la
consideración del producto tensorial
$\mathscr L^{-1}\otimes\mathscr L\otimes\mathscr F$ muestra de inmediato que
$\mathscr F$ es isomorfo a $\mathscr L^{-1}$.

\noindent\textbf{$(\mathbf{Err}_{\mathrm{III}},\,3)$} En
$(0_{\mathrm I},\,7.2.4)$, es preciso imponer la condición de que las potencias
$\mathfrak J^n$ sean ideales \emph{cerrados} en el anillo admisible $A$ para
que la proposición sea correcta. La demostración dada es incorrecta, y el
enunciado rectificado resulta de Bourbaki, \emph{Top. gén.}, cap. III,
3.\textsuperscript{a} ed., \S~3, n.\textsuperscript{o}~5, cor. 1 de la prop. 9.
Asimismo, es preciso suponer que los $\mathfrak J^n$ son \emph{cerrados} en $A$
en los enunciados $(0_{\mathrm I},\,7.2.5)$ y $(0_{\mathrm I},\,7.2.6)$.

\noindent\textbf{$(\mathbf{Err}_{\mathrm{III}},\,4)$} Después de la
proposición (I, 2.2.5), añadir: con las notaciones de (I, 2.2.5), si
$\mathfrak J$ es un ideal de $A$, se denota por
$\widetilde{\mathfrak J\mathscr F}$ el sub-$\mathscr O_X$-módulo
$\mathfrak J\widetilde{\mathscr F}$ definido en $(0,\,4.3.5)$.

\noindent\textbf{$(\mathbf{Err}_{\mathrm{III}},\,5)$} En (I, 3.4.5), en la
primera línea contando desde abajo de la p. 112, sustituir «\emph{un cuerpo}
$K$» por «\emph{un cuerpo algebraicamente cerrado} $K$». En la línea 1 de la
p. 113, sustituir «\emph{extensión}» por «\emph{extensión algebraicamente
cerrada}». En las líneas 1 a 4 de la p. 113, sustituir desde «$K$ será
llamado\ldots{}» por el texto siguiente:

Para todo punto de $X$ con valores en un cuerpo $K$, $K$ se llamará el
\emph{cuerpo de valores} del punto correspondiente y, si $x$ es la localidad
de este punto, se dice también que este último está \emph{localizado en} $x$.
Se define así una aplicación $X(K)\to X$ que hace corresponder a un punto con
valores en $K$ su localidad.

En las líneas 7 y 16 de la p. 113, suprimir «geométrico»; en las líneas 10 y 13
de la p. 113, suprimir «geométricos». En las líneas 10 y 11 contando desde
abajo de la p. 115, suprimir «\emph{geométrico}».

\noindent\textbf{$(\mathbf{Err}_{\mathrm{III}},\,6)$} Al final de la
demostración de (I, 3.6.1), en la línea 12 contando desde abajo de la p. 117,
añadir: Si se pone
$X'=X\times_Y\operatorname{Spec}(\mathscr O_y/\mathfrak a_y)$, entonces, para
todo punto $x\in X'$, identificado por $p$ con un punto de $X$, se tiene
$\mathscr O_{X',x}=\mathscr O_{X,x}/\mathfrak a_y\mathscr O_{X,x}$. Puesto que
la cuestión es local en $X$ y $Y$, se puede suponer que
$X=\operatorname{Spec}(B)$, $Y=\operatorname{Spec}(A)$, y se tiene
$(B/\mathfrak a_yB)_x=B_x/\mathfrak a_yB_x$ por planitud $(0,\,1.3.2)$.

\noindent\textbf{$(\mathbf{Err}_{\mathrm{III}},\,7)$} En (I, 5.1.1), en la
tercera línea contando desde abajo de la p. 127, sustituir
«$\mathscr O_X$-módulo cuasicoherente» por «\emph{ideal cuasicoherente de}
$\mathscr B$».

\noindent\textbf{$(\mathbf{Err}_{\mathrm{III}},\,8)$} Después de (I, 5.1.10),
en la línea 16 de la p. 131, añadir:

\emph{Observación (5.1.11).} --- Una ligera adaptación del razonamiento hecho
en (5.1.9) muestra que la conclusión sigue siendo válida \emph{sin suponer a
priori que $X$ sea un preesquema}, sino

\oldpage[III]{220}
suponiendo solamente que $(X,\mathscr O_X)$ sea un espacio anillado en anillos
locales, que $\mathscr J$ sea un ideal de $\mathscr O_X$ tal que
$\mathscr J^n=0$, que el espacio anillado
$X_0=(X,\mathscr O_X/\mathscr J)$ sea un esquema afín y, por último, que los
ideales $\mathscr J^k/\mathscr J^{k+1}$ sean $\mathscr O_{X_0}$-módulos
cuasicoherentes. Basta, en efecto, utilizar (1.8.1) (cf.
$(\mathbf{Err}_{\mathrm{II}})$) en lugar de (2.2.4) en el razonamiento.

\noindent\textbf{$(\mathbf{Err}_{\mathrm{III}},\,9)$} En (I, 5.3.1), en la
línea 11 de la p. 132, insertar, después de «o $\Delta_X$», «o $\Delta_\varphi$
si $\varphi:X\to S$ es el morfismo estructural».

\noindent\textbf{$(\mathbf{Err}_{\mathrm{III}},\,10)$} En (I, 5.3.9), la
demostración es insuficiente, pues no prueba que $\Delta_X(X)$ sea localmente
cerrado en $X\times_SX$. Para obtener una demostración correcta, basta utilizar
(4.2.4, $a$)): para todo $x\in X$ y todo entorno afín $U$ de $x$ en $X$,
$U\times_SU$ es un entorno afín de $\Delta_X(x)$; teniendo en cuenta (5.3.16)
(cuya demostración solo utiliza la definición (5.3.1.1)), se reduce a demostrar
(5.3.9) cuando $S=\operatorname{Spec}(B)$ y $X=\operatorname{Spec}(A)$ son
esquemas afines; entonces resulta claro, en virtud de (5.3.1.1), que $\Delta_X$
corresponde al homomorfismo canónico $A\otimes_BA\to A$ que transforma
$x\otimes y$ en $xy$; como este homomorfismo es sobreyectivo, $\Delta_X$ es en
este caso una inmersión cerrada (4.2.3), lo que termina la demostración.

\noindent\textbf{$(\mathbf{Err}_{\mathrm{III}},\,11)$} En (I, 7.1.15) y
(I, 7.1.16), en las líneas 7 y 15 contando desde abajo de la p. 158, suprimir
«\emph{geométricos}».

\noindent\textbf{$(\mathbf{Err}_{\mathrm{III}},\,12)$} Sustituir (I, 7.4.7)
por: Se extienden (por abuso de lenguaje) las definiciones de (7.4.1) al caso
en que $X$ es un preesquema \emph{reducido} tal que todo punto admite un
entorno abierto que solo tiene un número \emph{finito} de componentes
irreducibles; se deduce entonces de (7.3.4) y (7.4.6) que, para un
$\mathscr O_X$-módulo cuasicoherente de tipo finito $\mathscr F$, decir que
$\mathscr F$ es un \emph{haz de torsión} equivale a decir que
$\operatorname{Supp}(\mathscr F)$ \emph{no contiene ninguna componente
irreducible de $X$}.

\noindent\textbf{$(\mathbf{Err}_{\mathrm{III}},\,13)$} En (I, 10.11.7), en
las líneas 17 a 23 de la p. 205, el final de la demostración desde «Queda por
probar\ldots{}» es innecesario en virtud de (10.11.6), pues los haces
considerados son coherentes por $(0,\,5.3.5)$.

\noindent\textbf{$(\mathbf{Err}_{\mathrm{III}},\,14)$} En (II, 1.7.8),
después de la línea 10 contando desde abajo de la p. 16, añadir: cuando
$\mathscr E=\mathscr O_S^n$, se escribe también $\mathbf V_S^n$ en lugar de
$\mathbf V(\mathscr E)$; si además $S=\operatorname{Spec}(A)$ es afín, se
escribe $\mathbf V_A^n$ en lugar de $\mathbf V_S^n$; se tiene
$\mathbf V_A^n=\operatorname{Spec}(A[T_1,\ldots,T_n])$, donde los $T_i$ son
indeterminadas.

\noindent\textbf{$(\mathbf{Err}_{\mathrm{III}},\,15)$} En (II, 1.7.10), en
la línea 10 de la p. 17, suprimir «\emph{geométricos}»; en las líneas 12 y
16--17 de la p. 17, suprimir «geométrico».

\noindent\textbf{$(\mathbf{Err}_{\mathrm{III}},\,16)$} En (II, 1.7.12),
añadir: Se pone asimismo:
\[
  \mathbf V(\mathscr O_S^n)=\mathbf V_S^n=S[T_1,\ldots,T_n].
\]

\noindent\textbf{$(\mathbf{Err}_{\mathrm{III}},\,17)$} En (II, 4.2.6), en la
línea 7 de la p. 75, suprimir «\emph{geométricos}»; en la línea 8 de la p. 75,
suprimir «\emph{geométrico}».

\noindent\textbf{$(\mathbf{Err}_{\mathrm{III}},\,18)$} En (II, 4.6.13), en la
línea 7 de la p. 92, es preciso especificar el razonamiento, pues, con las
notaciones de (4.4.10), hay que probar aquí que la inmersión $\Gamma_f$ es
\emph{cuasicompacta}, a fin de aplicar (i bis). En virtud de (4.6.4), para
probar que $\mathscr L$ es amplio para $f$, se puede limitar al caso
en que $Y$ es afín. Obsérvese, por otra parte, que en cada una de las hipótesis
de (v), $f$ es cuasicompacto (I, 6.6.4). Por otra parte, si $g$ es
\oldpage[III]{221}
separado, $\Gamma_f$ es una inmersión cerrada (I, 5.4.3), luego cuasicompacta
(I, 6.6.4). Si, por el contrario, $X$ es localmente noetheriano, como $Y$ es
afín y $f$ cuasicompacto, el espacio subyacente a $X$ es cuasicompacto, luego
noetheriano, y se puede aplicar de nuevo (I, 6.6.4) para probar que $\Gamma_f$
es cuasicompacta.

\noindent\textbf{$(\mathbf{Err}_{\mathrm{III}},\,19)$} El enunciado de (II,
5.2.2) es incorrecto, pues las condiciones $b)$, $c)$ y $c')$ no son locales
en $Y$ cuando no se sabe si, para un abierto $U$ de $Y$, un
$(\mathscr O_X|f^{-1}(U))$-módulo cuasicoherente es la restricción a
$f^{-1}(U)$ de un $\mathscr O_X$-módulo cuasicoherente. Por tanto, es preciso
sustituir en la línea 16 de la p. 98 «\emph{Sea $f:X\to Y$ un morfismo separado
y cuasicompacto}» por: «\emph{Sean $X$, $Y$ dos preesquemas tales que $X$ sea
un esquema o que el espacio subyacente a $X$ sea localmente noetheriano, y sea
$f:X\to Y$ un morfismo cuasicompacto.}» En la demostración, se observará que
las hipótesis implican que $f$ es \emph{separado} cuando se supone que $X$ es
un esquema, en virtud de (I, 5.5.5 y 5.5.8); por tanto, se puede aplicar
(I, 9.2.2) en ambos casos.

\noindent\textbf{$(\mathbf{Err}_{\mathrm{III}},\,20)$} En (II, 6.2.3),
añadir, después de la línea 18 contando desde abajo de la p. 115:

\emph{Se dice que un morfismo $f:X\to Y$ es cuasifinito en un punto $x\in X$
si existen un entorno abierto afín $V$ de $y=f(x)$ y un entorno abierto afín
$U$ de $x$ tales que $f(U)\subset V$ y que el morfismo $U\to V$, restricción de
$f$, sea cuasifinito. Se dice que un morfismo $f:X\to Y$ es localmente
cuasifinito si es cuasifinito en todo punto de $X$.}

\noindent\textbf{$(\mathbf{Err}_{\mathrm{III}},\,21)$} En (II, 6.4.3), la
demostración es incorrecta, pues los elementos de un sub-$A$-módulo de tipo
finito de $K$ no son necesariamente enteros sobre $A$. Sustituir las 9 últimas
líneas de la p. 121 por el texto siguiente: Como se puede suponer que $E$ es
sin torsión, luego fiel, $E$ es también un $A[u]$-módulo fiel ($A$ está
inmerso en el anillo de endomorfismos de $E$). Como $E$ es un $A$-módulo de
tipo finito, se deduce que $u$ es \emph{entero} sobre $A$ (Bourbaki,
\emph{Alg. comm.}, cap. V, \S~1, n.\textsuperscript{o}~1, lema 1). Se concluye
inmediatamente que los valores propios de $u\otimes1$ (en una clausura
algebraica de $K$) son elementos \emph{enteros} sobre $A$, y por tanto también
lo son los $\sigma_i(u)$.

\noindent\textbf{$(\mathbf{Err}_{\mathrm{III}},\,22)$} En
$(0_{\mathrm{III}},\,9.1.1)$, en la línea 18 contando desde abajo de la p. 12,
suprimir «abierta».

\noindent\textbf{$(\mathbf{Err}_{\mathrm{III}},\,23)$} En
$(0_{\mathrm{III}},\,11.7.3)$, en la línea 10 de la p. 43, después de
$\mathsf C''$, añadir: tal que, para todo objeto proyectivo $P$ de
$\mathsf C'$ (respectivamente, todo objeto proyectivo $P'$ de $\mathsf C''$),
el funtor $A'\rightsquigarrow T(P,A')$ (respectivamente,
$A\rightsquigarrow T(A,P')$) sea exacto en $\mathsf C'$ (respectivamente, en
$\mathsf C$).

\noindent\textbf{$(\mathbf{Err}_{\mathrm{III}},\,24)$} El enunciado de
$(0_{\mathrm{III}},\,13.7.7)$ es inexacto y debe modificarse como sigue: en la
línea 6 de la p. 78, suprimir «y
$(\mathrm R^{n+1}T(A_k))_{k\in\mathbf Z}$»; en la línea 7 de la p. 78, sustituir
«$\mathrm R'^nT(A)$ es un $S$-módulo de tipo finito» por
«$\mathrm R'^nT(A)$ y $\mathrm R'^{n+1}T(A)$ son $S$-módulos de tipo finito»;
en las líneas 12--13 de la p. 78, suprimir «y $p+q=n+1$». En la demostración,
en la línea 23 de la p. 78, suprimir «y para $n+1$».

\noindent\textbf{$(\mathbf{Err}_{\mathrm{III}},\,25)$} En (III, 1.4.15), en
la línea 6 contando desde abajo de la p. 92, sustituir «\emph{de tipo finito}»
por «\emph{cuasicompacto}».

\noindent\textbf{$(\mathbf{Err}_{\mathrm{III}},\,26)$} En (III, 2.2.4), en la
línea 1 de la p. 102, sustituir «\emph{que los soportes de $\mathscr F$ y de
$\mathscr H$ sean propios sobre $Y$}» por «\emph{que el soporte de
$\operatorname{Im}(\mathscr F\to\mathscr G)$ sea propio sobre $Y$}». En la
demostración, sustituir el texto de las líneas 10 a 16, desde «Dicho
esto\ldots{}», por:

\oldpage[III]{222}
Pongamos
$\mathscr G_1=\operatorname{Im}(\mathscr F\to\mathscr G)
=\operatorname{Ker}(\mathscr G\to\mathscr H)$, que es coherente. Como se tiene
la sucesión exacta
$0\to\mathscr G_1(n)\to\mathscr G(n)\to\mathscr H(n)$ y el funtor $f_*$ es
exacto por la izquierda, la sucesión
$0\to f_*(\mathscr G_1(n))\to f_*(\mathscr G(n))\to f_*(\mathscr H(n))$ es
exacta para todo $n$, y basta por tanto mostrar que, para $n$ suficientemente
grande, la sucesión $f_*(\mathscr F(n))\to f_*(\mathscr G_1(n))\to0$ es exacta.
Dicho de otro modo, se puede limitar al caso en que $\mathscr H=0$ y en que el
soporte de $\mathscr G$ es propio sobre $Y$, luego \emph{cerrado en $Z$}
(II, 5.4.10); por consiguiente, $\mathscr G'=i_*(\mathscr G)$ es un
$\mathscr O_Z$-módulo coherente tal que $\mathscr G'|X=\mathscr G$. Se sabe
(I, 9.4.3) que existe un $\mathscr O_Z$-módulo coherente $\mathscr F'$ tal que
$\mathscr F'|X=\mathscr F$. Además, como $X$ es abierto en $Z$ y
$\operatorname{Supp}(\mathscr G)\subset X$ es cerrado en $Z$, es inmediato que
se define un homomorfismo sobreyectivo $u':\mathscr F'\to\mathscr G'$ de haces
tal que $u'|X$ sea el homomorfismo sobreyectivo dado
$u:\mathscr F\to\mathscr G$, tomando $u'|U=0$ para todo abierto $U$ de $Z$ que
no interseque $\operatorname{Supp}(\mathscr G)$ y $u'|U=u|U$ para todo abierto
$U\subset X$. Dicho esto, se ve, como al comienzo de la demostración de
(2.2.2), que se puede limitar al caso en que $Y$ es afín, y se trata por tanto
de mostrar que, para $n$ suficientemente grande, el homomorfismo
$\Gamma(u):\Gamma(X,\mathscr F(n))\to\Gamma(X,\mathscr G(n))$ es sobreyectivo.
Ahora bien, se tiene el diagrama conmutativo
\[
\xymatrix@C=35mm@R=12mm{
\Gamma(Z,\mathscr F'(n)) \ar[r]^{\Gamma(u')} \ar[d]_v &
\Gamma(Z,\mathscr G'(n)) \ar[d]^w \\
\Gamma(X,\mathscr F(n)) \ar[r]_{\Gamma(u)} & \Gamma(X,\mathscr G(n))
}
\]
donde $v$ y $w$ son los homomorfismos de restricción. La definición de
$\mathscr G'$ muestra además que $w$ es \emph{biyectivo}; por otra parte, en
virtud de (2.2.3), $\Gamma(u')$ es \emph{sobreyectivo} para $n$ suficientemente
grande, luego lo mismo vale para $\Gamma(u)$.

\noindent\textbf{$(\mathbf{Err}_{\mathrm{III}},\,27)$} En (III, 2.2.5),
después de la línea 7 contando desde abajo de la p. 102, añadir:

(iii) Bajo las hipótesis de (2.2.4) relativas a $X$, $Y$, $f$ y $\mathscr L$,
es inmediato que, si $\mathscr H$ es un $\mathscr O_X$-módulo coherente cuyo
soporte es \emph{propio sobre $Y$}, un razonamiento análogo al de (2.2.4)
muestra que existe un entero $N$ tal que, para $n\geq N$, se tiene
$\mathrm R^if_*(\mathscr H(n))=0$ para todo $i>0$. Se deduce, por la sucesión
exacta de cohomología, que, si $u:\mathscr F\to\mathscr G$ es un homomorfismo
de $\mathscr O_X$-módulos coherentes tal que $\operatorname{Ker}(u)$ y
$\operatorname{Coker}(u)$ tienen soportes \emph{propios sobre $Y$}, entonces
existe $N$ tal que, para $n\geq N$, el homomorfismo correspondiente
$\mathrm R^if_*(\mathscr F(n))\to\mathrm R^if_*(\mathscr G(n))$ sea
\emph{biyectivo} para todo $i>0$.

\noindent\textbf{$(\mathbf{Err}_{\mathrm{III}},\,28)$} En (III, 4.4.9), en la
línea 17 contando desde abajo de la p. 137, sustituir «\emph{Sean $Y$ un
preesquema íntegro localmente noetheriano}» por «\emph{Sean $X$ e $Y$ dos
preesquemas íntegros localmente noetherianos}»; en la línea 16 contando desde
abajo de la p. 137, sustituir «\emph{de tipo finito}» por «\emph{localmente de
tipo finito}». La demostración permanece esencialmente inalterada, pues
$f^{-1}(y)$ es localmente de tipo finito sobre $k(y)$ y, por tanto, sigue
siendo discreto.

\noindent\textbf{$(\mathbf{Err}_{\mathrm{III}},\,29)$} En (I, 9.3.4), en la
línea 19 contando desde abajo de la p. 173, sustituir «\emph{noetheriano}» por
«\emph{cuasicompacto}»; en las líneas 18 y 19 contando desde abajo de la p.
173, sustituir «\emph{coherente}» por «\emph{cuasicoherente de tipo finito}».
En la demostración, se reduce al caso en que

\oldpage[III]{223}
$X=\operatorname{Spec}(A)$, $\mathscr F=\widetilde M$, donde $M$ es un
$A$-módulo de tipo finito, $\mathscr J=\widetilde{\mathfrak J}$, donde
$\mathfrak J$ es un ideal de tipo finito de $A$; el resto del razonamiento
permanece entonces inalterado.

\noindent\textbf{$(\mathbf{Err}_{\mathrm{III}},\,30)$} En (I, 9.3.5),
sustituir las líneas 1 a 7 contando desde abajo de la p. 173 por el texto
siguiente:

\emph{Proposición (9.3.5). --- Sean $X$ un preesquema y $\mathscr F$ un
$\mathscr O_X$-módulo cuasicoherente de tipo finito. Entonces existen un
subpreesquema cerrado $Y$ de $X$, cuyo espacio subyacente es igual a
$\operatorname{Supp}(\mathscr F)$, y un $\mathscr O_Y$-módulo cuasicoherente de
tipo finito $\mathscr G$ tales que, si $j:Y\to X$ es la inyección canónica,
$\mathscr F$ sea isomorfo a $j_*(\mathscr G)$.}

Bastará mostrar que el ideal $\mathscr J$ de $\mathscr O_X$,
\emph{anulador de $\mathscr F$}, es \emph{cuasicoherente}; se tomará entonces
como $Y$ el subpreesquema cerrado de $X$ definido por $\mathscr J$ (I, 4.1.2),
y, como $\mathscr J\mathscr F=0$, $\mathscr F$ es un
$(\mathscr O_X/\mathscr J)$-módulo y se resolverá la cuestión tomando
$\mathscr G=j^*(\mathscr F)$. Para ver que $\mathscr J$ es cuasicoherente, se
puede (puesto que la cuestión es local) limitar al caso en que
$X=\operatorname{Spec}(A)$, $\mathscr F=\widetilde M$, donde $M$ es un
$A$-módulo engendrado por un número finito de elementos
$x_i\ (1\leq i\leq r)$; el ideal $\mathscr J$ es entonces la intersección de
los anuladores de los $x_i$. Pero el anulador de $x_i$ es el núcleo del
homomorfismo $\mathscr O_X\to\mathscr F$ correspondiente al homomorfismo
$s\mapsto sx_i$ de $A$ en $M$; es, por tanto, un ideal cuasicoherente (I,
4.1.1), y toda intersección finita de tales ideales es asimismo un ideal
cuasicoherente (I, 1.3.10).

\end{document}
